\documentclass[11pt,a4paper]{article}
\usepackage[T1]{fontenc}
\usepackage[utf8]{inputenc}
\usepackage{lmodern}
\usepackage{amsmath,amssymb,amsthm,mathtools,mathrsfs}
\usepackage{microtype}
\usepackage[a4paper,margin=26mm,headheight=15pt]{geometry}
\usepackage{xcolor,array,booktabs,longtable,enumitem,fancyhdr,graphicx}
\usepackage[angle=-54.7,scale=1,color={gray!20},text={DRAFT UNDER REVIEW}]{draftwatermark}
\SetWatermarkFontSize{2.25cm}
\usepackage[pagebackref=true]{hyperref}
\renewcommand*{\backref}[1]{}
\renewcommand*{\backrefalt}[4]{\ifcase #1 \or\ \textit{(Cited on p.~#2.)}\else\ \textit{(Cited on pp.~#2.)}\fi}
\hypersetup{colorlinks=true,linkcolor=blue!55!black,citecolor=green!45!black,urlcolor=blue!60!black,
 pdftitle={Amplitude and fourth-moment obstructions for prime-corrected sublevel approximations - Research Note v27},
 pdfauthor={Victor Casana-Giner}}
\usepackage[capitalize,nameinlink]{cleveref}
\setlength{\parindent}{1.2em}
\setlength{\parskip}{2pt}
\emergencystretch=2em
\numberwithin{equation}{section}
\newtheorem{theorem}{Theorem}[section]
\newtheorem{proposition}[theorem]{Proposition}
\newtheorem{lemma}[theorem]{Lemma}
\newtheorem{corollary}[theorem]{Corollary}
\theoremstyle{definition}
\newtheorem{definition}[theorem]{Definition}
\theoremstyle{remark}
\newtheorem{remark}[theorem]{Remark}
\newcommand{\meas}{\operatorname{meas}}
\newcommand{\ind}{\mathbf 1}
\newcommand{\clip}{\operatorname{clip}}
\newcommand{\Cm}{\mathrm{Cm}}
\newcommand{\Cms}{\mathrm{Cms}}
\newcommand{\vol}{\operatorname{vol}}
\newcommand{\dd}{\,\mathrm dt}
\fancyhf{}
\fancyhead[L]{\small Prime-corrected sublevel approximation}
\fancyhead[R]{\small Research note v27 \textbullet\ 4 Nov 2026}
\fancyfoot[C]{\thepage}
\pagestyle{fancy}
\begin{document}
\begin{center}
{\LARGE\bfseries Amplitude and Fourth-Moment Obstructions for\\[3pt]
Prime-Corrected Sublevel Approximations\\[3pt]
to Riemann Zeta Zeros}\par
\vspace{10pt}
{\large V\'ictor Casa\~na-Giner}\par
\vspace{4pt}
{\small Research note v27 --- 4 November 2026}\par
\end{center}
\begin{abstract}
We study the author's Euler-prime sublevel points $\tau_k^{\mathrm C}$,
their version $\tau_k^{\mathrm{Cm}}$ clipped with the explicit bound of the accompanying Mathematics of Computation submission (MCOM), the distinct
weighted compensated family $\tau_k^{\mathrm{Cms}}(X;W)$, and the taucw points $\tau_k^{\mathrm w}$.
The effective MCOM bound $|S(t)|\le0.08382\log t+4.444$ for $t\ge1$
provides a uniform enclosure for the true counting discrepancy;
the author's sharper asymptotic bound
$\limsup_{t\to\infty}|S(t)|/\log t<0.0748596142711$ supplies a growing margin inside the
fixed MCOM truncation. The paper contains amplitude, fourth-moment,
sublevel, numerical, correlated-displacement and local-model results,
and separates exact identities from conclusions depending on
the estimates cited from the author's IJNT (International Journal of Number Theory) manuscript and from Fujii. It proves, given the MCOM bounds, a localized
crossing-budget inequality for the actual compensated ordinal points;
a transport estimate $O(T/(\log T)^3)$ between clipped and unclipped
ordinal sequences conditional on the cited fourth moment and on the MCOM bounds; a finite-step
prime-band inclusion; an exact Fourier-box formula with endpoint
terms and a parity-sensitive bound for repeated heights; and an exact
Euler--pole transform yielding a sufficient weighted-$L^2$ criterion
for RH, together with its MCOM-excess formulation. The transform
criterion does not assume convergence in measure of the prime phases.
A section on the exponential sums $\sum e^{i\omega\gamma_k}$,
$\sum e^{i\omega\mathsf x_k}$ over the zeros and over the sequences
$\gamma_k^{\mathrm{RvM}},\tau_k^{\mathrm C},\tau_k^{\mathrm{Cm}}$ and on their
centroids gives exact identities, a reduction of the Landau-type coefficient
of the corrected sequences to a sawtooth Fourier coefficient, the
Gaussian scale of the RvM displacements versus coherence of the
corrected ones, a band argument for the rearrangement term of $\tau^{\mathrm C}$ (these three conditional on cited inputs, as recorded in the status table), a harmonic criterion for the vanishing of the sawtooth
coefficient, and an exact reconstruction of the points, and of their
spacings, from the centroid at frequency~$1$. It states the formulas of the four centroids
($\gamma$, RvM, $\tau^{\mathrm C}$, $\tau^{\Cm}$) and of their differences, and what is known on the
convergence of each with and without RH: the RvM centroid converges, with a complete proof; the RH limit of the centroid of $\gamma$ is quoted; for the corrected points the question is reduced to the sawtooth limit, and convergence is open. It proves without RH the closed form of the real part of Cram\'er's function (the imaginary part being non-local), the distributional convergence of the four centroids on the window $(\log2,\log3)$ with the $1/\ell$ term, the exact defect between the $\gamma$-centroid and Cram\'er's function, and a spectral-rigidity theorem excluding any exact distributional transfer from $\gamma$ to the corrected points; it justifies without RH the Tauberian criterion for the centroid of $e^{i\gamma_k}$ (equivalence of the time-average and index-average centroids at non-resonant frequencies, under a mild hypothesis on multiple zeros, an exact index-based Laplace identity and its pseudofunction criterion for the four centroids and for their differences); it reduces the window criterion to one fixed test function with a family of modulations (an exact reformulation, with a non-negative scalar form) and proves that smooth averages of the sawtooth coefficient vanish automatically; it gives the exact moment expansion of the centroid difference, the band-occupation identity for the sublevel points and a counterexample showing that a bounded perturbation of the smooth count does not force separation; and it proves, without RH, the continuation of Cram\'er's function across the window with explicit principal parts at the prime powers, deducing that a prime-built sequence whose centroid has that limit on the window equals the zeros exactly when RH holds, together with tests on the zero file.
A self-contained section treats the taucw points, the sublevel points of $N_0$ plus the Selberg-weighted prime phase with cutoff $X=t^\theta$, whose definition and computation contain no zero of $\zeta$: the compensator is $O(1/t)$ for $\theta\le1$; the difference between $N_0+P_X^W$ and $N$ is, up to $O(1/t)$, the sum of an on-line part, which has a fourth moment $\le C(\theta_0)\,T$ for $\theta\ge\theta_0>0$ without hypothesis on the zeros (given the increment estimate \eqref{eq:FujiiB} of Fujii quoted in \cref{sec:definitions} and Selberg's fourth moment of $S$), and an off-line part in which each zero enters through the explicit factor $X^{2y}$; and the taucw points, centroids and exponential sums for the first $3000$ indices are compared with those of the zeros.
The same section proves, without any hypothesis on the zeros, the existence, order and coarse localization of the taucw points; an exact transport inequality and an $L^1$ contraction of sublevel distributions; the derivative of the phase in the cutoff exponent, the variance of the phase for $\theta<\frac12$ together with the exact variance of an increment in the cutoff, and signed transport up to $\theta=1$; the fourth moment of the prime-power part, the mean of the long-prime part and its Chebyshev representation; the fixed-frequency Fourier coefficients of the phase and a spectral formula for the points; a quadratic estimate for the local reordering defect; Plancherel and band-limited identities for counting discrepancies; and the exact relation between the prime-defined counting surrogate and $S$ (\crefrange{sec:taucw-exist}{sec:taucw-S}). It computes in closed form the contribution of a pair of zeros off the critical line to the phase, with its spectral density, its exact $L^2$ mass and lower bounds for it; shows that, up to $O(\log t)$ and uniformly in the cutoff, the prime sum with its compensator equals this contribution; proves that the absence of zeros off the critical line is equivalent to the boundedness of its weighted $L^2$ norm along any unbounded sequence of cutoffs; and gives the cost, in $L^1$ displacement of the strictly increasing taucw points, of an imaginary part shared by two zeros (\cref{sec:taucw-detect}). Finally it extends the second moment of the phase to the support $n\le t(\log t)^{3-\varepsilon}$, computes the diagonal constant $(5\log2-3)/(2\pi^2)$ of the increment from $\theta=\frac12$ to $\theta=1$, writes the Selberg weight as a difference of two Fej\'er-type autocorrelations whose Fourier transform is the kernel of the explicit formula, bounds the energy of the off-line part for $\theta<\frac12$ by $O(T\log\log T)$, shows that zeros off the line of bounded height change the taucw error at height $t$ by $O(t^{\theta-2}\log^{-2}t)$, reduces the weighted norm to the heights $t\le(X^2(1+\log\log3X))^{1/3}$ and writes it as a positive Gram form, gives the arithmetic form of the harmonics of the phase, and records the threshold $0.1608326\ldots$ for the Hilbert--Schmidt correction that the rank--trace inequality of Alp\"oge and Furman would need at twice their bandwidth (\cref{sec:taucw-beyond}).
It is not claimed that the required $L^2$ bound or arithmetic
noncollapse has been proved, or that the weighted zero-explicit formula has been
proved for weights other than Selberg's.
The imported weighted kernel is explicitly distinguished from the
author's original prime construction. The paper does not establish RH,
simplicity or an improved Selberg--Fujii fourth-moment theorem.
\end{abstract}
\noindent\textbf{Keywords:} Riemann zeta function; Riemann--von Mangoldt formula; prime powers; Selberg approximation; sublevel measures; fourth moments; argument of zeta.\\
\noindent\textbf{MSC 2020:} 11M26, 11M06, 11N05.
\tableofcontents

\section*{List of variables and notation}
\addcontentsline{toc}{section}{List of variables and notation}
Symbols are listed roughly in order of introduction. Where the same letter is used with different meanings, the scope of each meaning is indicated; the context (subscripts, arguments) always distinguishes them.
\begingroup\footnotesize
\setlength{\LTpre}{4pt}\setlength{\LTpost}{6pt}
\begin{longtable}{@{}p{0.22\linewidth}p{0.56\linewidth}p{0.16\linewidth}@{}}
\toprule
Symbol & Meaning & Introduced\\
\midrule
\endfirsthead
\toprule
Symbol & Meaning & Introduced\\
\midrule
\endhead
$\rho_k=\beta_k+i\gamma_k$ & nontrivial zeros with $\gamma_k>0$, with multiplicity, $\gamma_k\le\gamma_{k+1}$ & \S\ref{sec:definitions}\\
$N(t)$ & number of zeros with $0<\Im\rho\le t$ (right-continuous) & \S\ref{sec:definitions}\\
$N_0(t)$ & Riemann--von Mangoldt main term \eqref{eq:N0} & \eqref{eq:N0}\\
$S(t)$ & $\frac1\pi\arg\zeta(\frac12+it)$; at zero heights defined by \eqref{eq:RVM} & \eqref{eq:RVM}\\
$R(t)$ & $\vartheta(t)/\pi+1-N_0(t)$, $|R(t)|\le0.53/(\pi t)$ for $t\ge1$ & \eqref{eq:RVM}\\
$\vartheta(t)$ & Riemann--Siegel theta function & \S\ref{sec:definitions}\\
$N_\vartheta(t)$ & $\vartheta(t)/\pi+1$ & \cref{prop:forward-v3}\\
$\clip_a(x)$ & symmetric truncation to $[-a,a]$ & \eqref{eq:clip}\\
$u_k(F)$ & sublevel measure $\meas\{F<k-\frac12\}$ & \eqref{eq:sublevel}\\
$\mathsf k,\ \alpha_{\rm F},\ H,\ Y,\ h$ & parameters of the Fujii input (moment order, range exponent, interval length, prime cutoff, shift) & \eqref{eq:FujiiA}--\eqref{eq:FujiiB}\\
$p_\nu$, $m_j$ & $\nu$-th prime; prime budget $\lfloor2^{j/2}\rfloor$ on the dyadic block $[2^j,2^{j+1}]$ & \eqref{eq:Phi-m}, \eqref{eq:interpolation}\\
$\Phi_m$, $\Phi^\circ$ & Euler phase of the first $m$ primes; its dyadic interpolation & \eqref{eq:Phi-m}, \eqref{eq:interpolation}\\
$h(t)$ & original truncation $\log\log\max(t,e^e)$ (in \S\ref{sec:kernel-v4} the letter $h$ also denotes $Lg$; $h$ is also the shift of \eqref{eq:FujiiB} and, in $\mathsf G_k^{(h)}$, a lag) & \eqref{eq:old}\\
$F_h$, $\tau_k^{\mathrm C}$ & original IJNT function and its sublevel points & \eqref{eq:old}\\
$A(t)$ & generic truncation bound in \cref{lem:countconstraint} & \cref{lem:countconstraint}\\
$q$ & in \cref{lem:countconstraint}: $N(t)$; elsewhere $q(v)$ is the weight density \eqref{eq:weights} and $q_F$ the integer discrepancy; a prime in the proof of \cref{thm:smallX} & --\\
$r$ & index displacement (\cref{thm:omega}); prime-power exponent (\eqref{eq:Phi-m}); radii $r_1,r_2$ (\S\ref{sec:filter-v4}) & --\\
$H(t)$ & MCOM truncation $0.08382\log\max(t,e)+5.5$ (not the interval length $H$ of the Fujii input) & \eqref{eq:MCOMH}\\
$A_{\mathrm M},B_{\mathrm M}$ & $0.08382$ and $5.5$ & \cref{prop:rvmbound-v5}\\
$C_{\mathrm H}$ & the author's asymptotic MCOM coefficient $0.0748596142710907\ldots$ (the bound \eqref{eq:asymptMCOM} is stated with the rounded-up value $0.0748596142711$) & \eqref{eq:asymptMCOM}\\
$G(t)$ & $N(t)-N_0(t)$ (in \S\ref{sec:filter-v4}, $G$, $G_L$ also denote generic functions) & \eqref{eq:MCOMH}\\
$F^{\Cm}$, $\tau_k^{\Cm}$ & taucm: MCOM-truncated function and its sublevel points & \eqref{eq:taucm}\\
$q_F(t)$ & integer discrepancy $|\lfloor F+\frac12\rfloor-N|$ & \cref{thm:taucm}\\
$F_Q$, $u_k(Q)$, $Q$ & structural saturation model and its sublevel measures; oscillation parameter & \eqref{eq:sat-model}\\
$\varphi_\pm$ & $N_0\pm H$ & \cref{thm:structural}\\
$\mathcal W$, $W$, $q$ & weight class, weight, density $q=-W'$ on $[1,2]$ & \eqref{eq:weights}\\
$W_{1/4}$, $q_{1/4}$ & weight with $q(v)=v/2+1/4$ (the index labels the weight) & \eqref{eq:weights}, \eqref{eq:q14-v4}\\
$\Lambda$ & von Mangoldt function & \S\ref{sec:weights}\\
$X$, $L$ & prime cutoff; $L=\log X$ & \eqref{eq:PhiW}\\
$\Phi_X^W$, $\Pi_X^W$, $P_X^W$ & weighted Euler phase, pole compensator, $P_X^W=\Phi_X^W-\Pi_X^W$ & \eqref{eq:PhiW}--\eqref{eq:PW}\\
$F_X^W$, $\tau_k^{\Cms}(X;W)$ & taucms: truncated weighted function and its sublevel points & \eqref{eq:Cms}\\
$\mathrm{DBSRH}(W)$ & $\lim_{X\to\infty}\tau_k^{\Cms}(X;W)=\gamma_k$ for each $k$ & \eqref{eq:DBSRH}\\
$D$, $E_T$, $\mathcal K_T^{\mathrm{RvM}}$ & local meanings in \cref{thm:v7-transport}: $D=(|\Phi^\circ|-H)_+$, $E_T=\{t\in[T/2,3T]:D(t)>0\}$, and $\mathcal K_T^{\mathrm{RvM}}=\{k:T\le\gamma_k^{\mathrm{RvM}}\le2T\}$ (the letter $D$ also has the meanings of the neighbouring rows) & \cref{thm:v7-transport}\\
$a_p$, $D(t)$, $D_T$ & coefficients $1-W(\log p/\log X)$, the difference sum \eqref{eq:diffprime}, $D_T=\log(\log T/\log X)$ & \cref{thm:smallX}\\
$A_X$, $B_X$, $b_n$, $w_p$, $V_X$, $\delta$ & prime and prime-square sums, coefficients of $A_X^2$, weights, variance, $\delta=\min(\varepsilon/2,1/200)$ & \cref{thm:moderate}\\
$B_W$, $\mathcal E_{T,X}$ & bound for $q$ of the two fixed weights; saturated set & \cref{thm:moderate-log}, \cref{prop:sat-measure}\\
$\mathcal Z_L$ & smoothed zero sum & \eqref{eq:ZL}\\
$\mathcal L_q(L,z)$ & Laplace integral, $z=d+i(\gamma-t)$, $d=\beta-\frac12$ & \eqref{eq:laplace}\\
$\mathcal A_k^\pm$ & signed sublevel contributions & \eqref{eq:Aplus}, \eqref{eq:Aminus}\\
$\bar u_k$ & $(u_k+u_{k+1})/2$ & \S\ref{sec:cells}\\
$\gamma_k^{\mathrm{RvM}}$, $v_k$ & solution of $N_0=k-\frac12$ with $\gamma_k^{\mathrm{RvM}}>2\pi$; $v_k=\gamma_k^{\mathrm{RvM}}$ in proofs & \S\ref{sec:spacing-v3}\\
$\Delta_k(X;W)$, $s_k$ & consecutive spacing of taucms; $s_k=\Delta_k$ & \eqref{eq:gap-v3}, \eqref{eq:center-v5}\\
$a_k,b_k$ & barriers $N_0\pm H=k-\frac12$ (not the moment bounds $a_T,b_T$ nor the band exponents $a,b$ of \eqref{eq:v20-band}) & \cref{thm:block-sep-v3}\\
$I_k^\pm$, $B_k^\pm$ & occupation intervals and saturated measures & \cref{prop:occupation-v3}\\
$e_k(F)$ & evaluation defect $k-\frac12-F(\gamma_k)$ & \eqref{eq:ek-v3}\\
$\Psi_j$, $j$ & $\int_{z-1/2}^{z+1/2}v^{2j}dv$; moment order (in \eqref{eq:interpolation} $j$ is the dyadic index) & \eqref{eq:Psi-v3}\\
$M$ & number of indices in repeated-height groups & \eqref{eq:duplicate-count-v3}\\
$\mathscr F_X$ & generic family in the forward criterion & \cref{prop:forward-v3}\\
$\theta$ & prime-support exponent, $X=t^\theta$ (not $\vartheta$) & \S\ref{sec:v9-data}\\
$g_k$ & in \S\ref{sec:v9-data}: normalized nearest-neighbour separation; from \S\ref{sec:near-v4} on: $\gamma_{k+1}-\gamma_k$ & \eqref{eq:gfit-v3}, \eqref{eq:poserror-v4}\\
$\delta_k$ & $\frac12-|e_k(F)|$ & \S\ref{sec:v9-data}\\
$\eta$, $\eta_T(F)$ & mean of $e_k(F)^2$ on a block (unrelated to $\eta_k$ below) & \S\ref{sec:v9-data}, \S\ref{sec:jone-v3}\\
$M_T$, $D_F$, $\Phi_c$ & $N(2T)-N(T)$; $N-F$; truncated phase of the numerical script & \eqref{eq:jone-AB-v3}\\
$A_T,B_T$ & $\int D_F^2N_0'$ and $\int D_F^2\Phi_c'$ on $[T,2T]$ & \eqref{eq:jone-AB-v3}\\
$a_T,b_T,\kappa_T,r_T$ & second moment, fourth moment, derivative ratio, endpoint bound & \eqref{eq:jone-input-v3}\\
$\mathcal E_\partial(F;T)$ & endpoint contribution in \eqref{eq:jone-identity-v3} & \eqref{eq:jone-identity-v3}\\
$d_k(X;W)$ & true-zero position error $\tau_k^{\Cms}-\gamma_k$ & \eqref{eq:poserror-v4}\\
$\mathcal M_k(X;W)$ & symmetric margin & \eqref{eq:margin-v4}\\
$\mathscr L_{k,X}$ & local Lipschitz constant of $F_X^W$ on $[a_k,b_{k+1}]$ & \eqref{eq:Mbound-v4}\\
$\star$, $\mathsf x_k^\star$, $\sigma_k^\star$ & family $\star\in\{\mathrm{RvM},\mathrm C,\Cm\}$; its points $\gamma_k^{\mathrm{RvM}},\tau_k^{\mathrm C},\tau_k^{\Cm}$; displacement $\gamma_k-\mathsf x_k^\star$ & \cref{def:v9-objects}\\
$\mathfrak J_T$ & index block $\{k:T<\gamma_k\le2T\}$, of size $M_T$ & \cref{def:v9-objects}\\
$\omega,\ \omega'$ & real frequencies of the exponential sums & \eqref{eq:v9-Theta}\\
$\Theta_T^\star(\omega,\omega')$, $\Xi_T^\star(\omega)$ & joint transform of position and displacement; $\Theta(\omega,\omega)-\Theta(\omega,0)$ & \eqref{eq:v9-Theta}\\
$\mathrm{Par}_\ell,\ \mathrm{Cen}_\ell$, $\mathbb S_\ell$, $\nabla^2$ & partial sum, centroid, double partial sum $\ell\,\mathrm{Cen}_\ell$, second difference & \eqref{eq:v9-Cen}, \cref{lem:v9-exact}\\
$\mathsf y_k$ & generic real sequence in $\mathrm{Par},\mathrm{Cen}$ & \eqref{eq:v9-Cen}\\
$\chi(\omega)$ & Landau coefficient $-\Lambda(e^\omega)e^{-\omega/2}$ & \eqref{eq:v9-chi}\\
$N^{\mathsf x}$, $F^*$ & counting function of a sequence; increasing rearrangement of $F$ & \cref{lem:v9-rearr}\\
$z_F$, $\varrho_F$, $\mathrm{saw}$ & slip $N-\lfloor F+\frac12\rfloor$; rearrangement term; $y-\lfloor y+\frac12\rfloor$ & \eqref{eq:v9-slip}, \eqref{eq:v9-sawD}\\
$\varsigma_T(\omega;F)$ & sawtooth coefficient of $F$ on $[T,2T]$ & \eqref{eq:v9-varsigma}\\
$\mu_F(c)$, $E_F$, $B_F$ & $\meas\{F<c\}$; $\{t:F(t)+\frac12\in\mathbb Z\}$; $\{s:F^*(s)+\frac12\in\mathbb Z\}$ & \cref{lem:v9-rearr}\\
$F^\circ$ (in \S\ref{sec:v9-centroid}) & unclipped phase $N_0+\Phi^\circ$ & \cref{lem:v9-sawcmp}\\
$\mathsf Z_T$ & endpoint term of the Stieltjes identity & \cref{thm:v9-reduction}\\
$\varpi_T$, $\omega_T(\psi)$, $\psi$ & Gaussian scale $\sqrt{2\log\log T}/\log T$; frequency $\psi/\varpi_T$; its parameter & \S\ref{sec:v9-coherence}\\
$\mathfrak b^\star(\omega)$ & block coefficient of $\Xi_T^\star$ & \cref{thm:v9-cen}\\
$\mathsf f_k$, $\mathbb Y(U)$ & $e^{i\omega\gamma_k}-e^{i\omega\mathsf x_k^\star}$; its sum over $\gamma_k\le U$ & \cref{thm:v9-cen}\\
$n$ & running index of partial sums (local) & \eqref{eq:v9-Cen}\\
$\mathsf r_\omega$ & model value of $\varsigma_T$ in the numerical paragraph & \S\ref{sec:v9-numbers}\\
$\mathsf U_k,\ \mathsf V_k,\ \mathsf I_k,\ \mathfrak w_k$ & $\{F<k-\frac12\}$; $[0,\tau_k^{\mathrm C})$; band $[r_k^-,r_k^+)$ containing their symmetric difference; its length & \cref{prop:v9-rearr}\\
$\mathfrak h,\ \mathfrak n,\ \mathrm{Fej}_{\mathfrak n},\ \mathfrak d$ & harmonic index; Fej\'er order and mean of $\mathrm{saw}$; non-concentration width & \cref{prop:v9-harm}\\
$\Delta_k^{\mathsf x},\ \mathfrak c_{\mathsf x},\ \omega_\sharp,\ \mathfrak E_k$ & spacing $\mathsf x_{k+1}-\mathsf x_k$; MCOM coefficient of the family; frequency cutoff; collision energy & \eqref{eq:v9-energy}, \cref{prop:v9-recon}\\
$\mathsf D_\ell^\star(\omega)$ & centroid of $\gamma$ minus centroid of $\mathsf x^\star$ & \eqref{eq:v11-diff}\\
$\mathfrak m(\omega)$, $K^0$ & limit of the RvM centroid; limit of $\mathsf D_\ell^{\mathrm{RvM}}(1)$ & \eqref{eq:v11-mid}, \cref{prop:v11-gamma}\\
$\mathscr N_\infty^{(\mathsf x)}(f)$, $\mathscr O_{\mathsf x}(f)$, $\mathscr U_{\mathsf x}(w)$, $\mathscr V_{\mathrm C}(w)$, $\mathsf o_k$, $\mathsf n_\gamma$ & distributional limit of a centroid and its $1/\ell$ coefficient; holomorphic difference of two exponential sums; same for $\tau^{\mathrm C}$ against RvM; $|\beta_k-\frac12|$; counting measure of the $\gamma_k$ & \cref{prop:v13-dist}, \cref{prop:v13-rig}, \cref{prop:v13-def}\\
$\mathscr C(w),\ \mathfrak p(w),\ \mathfrak R$ & $e^{-w/2}\mathfrak V(w)$; its real part $\frac12(e^{w/2}-e^{-w/2}/(e^{2w}-1))$; region $\operatorname{Re}w>0,|\operatorname{Im}w|<\pi$ & \cref{thm:v17-cont}\\
$\mathsf H,\ \mathsf H^\natural,\ \mathfrak u,\ \mathfrak g,\ \mathfrak s_n,\ \mathscr C^\sharp,\ \mathfrak q$ & generic holomorphic function; $\overline{\mathsf H(\bar w)}$, its boundary value, real-symmetric function; principal part at $\log n$; $\mathscr C-\mathfrak s_2-\mathfrak s_3$, $\operatorname{Im}\mathscr C^\sharp$ on the real line & \cref{lem:v17-refl}, \cref{thm:v17-cont}, \cref{cor:v17-local}\\
$\mathsf M_\ell^{(j)},\ \mathsf s_n,\ \mathsf m,\ \mathsf R_{\ell,\mathsf m}$ & moments of the displacement $\sigma_k^{\mathrm C}$ against $e^{i\omega\tau_k^{\mathrm C}}$; their partial sums; Taylor order; Taylor remainder & \cref{def:v16-moments}, \cref{prop:v16-moment}\\
$\mathsf E,\ \mathsf w,\ \mathsf v,\ \mathfrak o,\ \mathsf K_{\mathfrak o},\ t_\natural,\ t_\lambda$ & steepness of $F_{\mathsf E}$ and $\mathsf w=1/\mathsf E$; rescaled variable; level offset $\pm\frac12$; the equation $F_{\mathsf E}=k+\mathfrak o$; root of $N_0=k$; root of $F_{\mathsf E}=\lambda$ & \cref{prop:v16-counter}\\
$J_0$, $\mathsf a$, $\mathscr I_f,\mathscr Y_f,\mathscr Y_f^\star,\mathscr Z_f^\star,\mathscr W_f$ & subinterval of $J$ where the test function is positive; modulation; Cram\'er side, defect, defect for $\mathsf x^\star$, difference $\gamma$ against $\mathsf x^\star$, weighted norm & \cref{thm:v15-reduction}, \cref{lem:v15-tail}\\
$\mathbb B_u,\mathbb A_u,\mathbb M_u,\mathbb G_u$, $\mathsf L$, $\mathsf W_{\mathrm{rep}}$ & partial sum, its integral and mean, Laplace series over $\gamma_k$; $\log\frac T{2\pi}$; excess multiplicity & \cref{thm:v14-transfer}, \cref{thm:v14-taub}\\
$\mathbb P_{\mathsf x},\mathbb Q_{\mathsf x}$, $\mathscr P^{\mathrm t},\mathscr P^{\mathrm o}$, $\mathscr X$ & index-based Laplace series, pseudofunction test functions, transform of a difference of centroids & \cref{prop:v14-index}\\
$\mathfrak V(z)$, $\varsigma_\infty$ & Cram\'er's $\sum_{\operatorname{Im}\rho>0}e^{\rho z}$ (continued); $\lim_T\varsigma_T(\omega;F_h)$ & \cref{prop:v11-gamma}, \cref{prop:v11-tau}\\
$K(z)$, $\mu_m$, $w(y)$ & two-source kernel, moments of $q_{1/4}$, kernel weight $\int_{\max(1,y)}^2q_{1/4}$ & \eqref{eq:q14-v4}, \eqref{eq:moments-v4}, \cref{thm:pair-injective-v4}\\
$C(z,h)$, $\mathcal R_K$ & pair profile $K(z-h/2)+K(z+h/2)$ and its remainder; $h=Lg$ & \cref{thm:resolution-v4}\\
$\mathcal C_{c,g,L}$, $c$ & complete local pair profile, centre & \eqref{eq:pair-complete-v4}\\
$F_L^\circ$ & untruncated $N_0+P_{e^L}^W$ & \S\ref{sec:filter-v4}\\
$\mathscr D_{L;r}$, $\mathscr Q_{L;r_1,r_2}$, $\mathscr B_{r_1,r_2}$ & symmetric difference, spatial filter, filter symbol & \eqref{eq:Dr-v4}, \eqref{eq:Qdef-v4}\\
$B$, $\mathcal R_B$ & smooth background and its Taylor remainder & \eqref{eq:Qbackground-v4}\\
$\dot P_L$ & $\partial_LP_{e^L}^W$ at fixed $t$ & \eqref{eq:dLprime-v4}\\
$\mathsf A$, $\mathbf v$, $\mathsf P$, $\mathbf y$, $\mathbf b$, $\mathbf r$, $\mathbf r_4$, $\widehat{h^2}$ & sampling matrix, $K''$ vector, projection, data, nuisance, residual, quartic remainder, estimator & \cref{thm:identifiability-v4}\\
$h_k^{\mathrm{RvM}}$, $\varepsilon_k^{\mathrm{RvM}}$ & RvM spacing; RvM displacement $\tau_k^{\Cms}-\gamma_k^{\mathrm{RvM}}$ & \eqref{eq:hrvm-v5}, \eqref{eq:epsrvm-v5}\\
$\eta_k$, $X_k$ & coherence fraction and threshold cutoff & \eqref{eq:coherence-eta-v5}\\
$\xi_{k,\pm}$ & mean-value points & \cref{prop:rvmbound-v5}\\
$J_k$, $p$, $K_k$ & band $[a_k,b_{k+1}]$, Lebesgue exponent, uniform derivative bound & \cref{prop:sobolev-v5}\\
$c_k(X;W)$ & pair centre $\frac12(\tau_k^{\Cms}+\tau_{k+1}^{\Cms})$ & \eqref{eq:center-v5}\\
$\mathcal B_T$, $\Phi_{a,b}$, $Y_k$ & prime band $T^a<p\le T^b$; its phase; its increment at consecutive RvM points & \eqref{eq:v20-band}, \eqref{eq:v20-Y}\\
$\mathfrak r_{a,b}$, $\mathfrak v_{a,b}$, $\mathfrak z_{a,b}$ & band limits of correlation, variance and covariance of $Y_k$ & \eqref{eq:v20-main}, \eqref{eq:v20-var-cov}\\
$\mathsf G_k^{(h)}$ & RvM gap of lag $h$, $\gamma_{k+h}^{\mathrm{RvM}}-\gamma_k^{\mathrm{RvM}}$ & \cref{lem:v20-spacing}\\
$\mathfrak r^{\mathrm{diag}}_{W,\theta}$ & diagonal (heuristic) prediction for the correlation of consecutive increments of the whole phase & \eqref{eq:v21-diag}\\
$\varepsilon$ & fixed small exponent in the cutoff ranges (unrelated to $\varepsilon_k^{\mathrm{RvM}}$) & \cref{thm:moderate}\\
$E_X(t)$ & excess $(|P_X^W(t)|-H(t))_+$ of the \emph{untruncated} weighted prime phase & \eqref{eq:v7-excess}\\
$u_k^\circ$ & auxiliary untruncated sublevel points (not the author's original taucs) & \eqref{eq:v7-untruncated}\\
$\mathcal T$ and $k_s$ & auxiliary Fourier--Laplace transform of the weighted Euler phase & \eqref{eq:v7-ks}\\
$\Psi_g$ & finite-interval box convolution of a test function & \cref{thm:v7-box}\\
$\mathrm J(t)$, $r_0(t)$ & Dirichlet-type integral and the limit correction $r_0=-1+\pi^{-1}\arctan(2t)+\pi^{-1}\mathrm J(t)$ & \cref{lem:v8-smooth}\\
$E(s)$ & $s^{-2}-2(s-\frac12)^{-2}+\int_1^2(\log u)u^{-s}du$, the transform of $r_0$ & \cref{lem:v8-smooth}\\
$P_\infty$, $F_\infty$ & limit in measure of $P_{X_j}^W$; $N_0+P_\infty$ & \cref{prop:v8-first}\\
$t_\sharp$ & $\meas\{0<t<\gamma_1:R(t)+r_0(t)\le-\frac12\}$, the first-index defect & \cref{prop:v8-first}\\
$\tau_k^\flat(X;W)$ & $2\pi+\meas\{t\ge2\pi:F_X^W<k-\frac12\}$, sublevel points measured from $2\pi$ & \cref{prop:v8-first}\\
$\mathcal K_X(s)$ & $\sum_k k_s(\tau_k^\flat(X;W))$ & \cref{prop:v8-criterion}\\
$F^{\mathrm w}_\theta$, $\tau_k^{\mathrm w}$ & taucw phase $N_0+\Phi^{\mathrm w}_\theta$ with cutoff $X=t^\theta$ and its sublevel points; $F^{\mathrm w,H}_\theta$ with truncation at $H$ & \eqref{eq:taucw-F}, \eqref{eq:taucwH}\\
$\mathfrak G_L$, $\mathfrak K_L$, $\mathfrak I_L$, $\mathfrak H_L$ & kernel of the smoothed explicit formula; its kernel on the imaginary axis, $\mathfrak K_L(\delta)=-\pi^{-1}\Re\mathfrak G_L(i\delta)$; the integral of $\mathfrak G_L$ along the horizontal ray and that of $\mathfrak K_L$ along $\mathbb R$ & \eqref{eq:kernel-G}\\
$\mathcal O_L(t)$, $\mathcal Y_L(t)$ & on-line (imaginary parts of the zeros only) and off-line parts of $N_0+P_X^W-N$, $L=\log X$ & \eqref{eq:taucw-on-off}\\
$\tau_k^{\mathrm w}(\theta)$, $N^{\mathrm w}_\theta$, $N^{\mathrm w}$, $S^{\mathrm w}$, $\lambda_k$ & taucw point of cutoff $X=t^\theta$ ($\tau_k^{\mathrm w}=\tau_k^{\mathrm w}(1)$); counting function of the points ($N^{\mathrm w}=N^{\mathrm w}_1$); prime-defined surrogate $N^{\mathrm w}-N_0-R$; level $k-\frac12$ & \eqref{eq:taucw-count}\\
$\mathbf w_{\theta,t}(n)$, $\mathfrak A_\theta$, $\mathsf T_k^\pm$ & moving weight $W(\log n/(\theta\log t))$; increasing majorant of $|\Phi^{\mathrm w}_\theta|$; barriers $(N_0\pm\mathfrak A_\theta)(\mathsf T_k^\pm)=\lambda_k$ & \cref{prop:taucw-ord}\\
$\vol_{f,I}$, $\mathcal U$, $q_F$ & sublevel distribution $\meas\{t\in I:f(t)<v\}$; finite set of indices; integer discrepancy of \cref{thm:taucm} & \eqref{eq:taucw-vol}, \cref{thm:taucw-L1}\\
$\Phi^{\mathrm{fr}}_{\theta,T}$, $\mathfrak L$, $\mathfrak k_n$, $\mathbf d_n$ & phase with weights frozen at height $T$; length $(2T)^{2\theta}$ of the prime-power sum; coefficient $\Lambda(n)/(\pi\sqrt n\log n)$; $s$-derivative of the weight & \eqref{eq:taucw-frozen}, \cref{lem:taucw-freeze}\\
$\iota_n$, $\mathsf X$, $\mathscr K_n$ & coefficients of a frozen sine polynomial; its exponential polynomial; difference of two moving weights & \cref{thm:taucw-var}, \cref{thm:taucw-signed}\\
$\mathfrak S(\theta_-,\theta_+)$ & variance constant of the increment in the cutoff & \eqref{eq:taucw-S}\\
$\mathfrak a_T$, $\Phi^{\mathrm{sm}}_T$, $\Phi^{\mathrm{pw}}_T$, $\Phi^{\mathrm{lg}}_T$ & the $\lfloor\sqrt T\rfloor$-th prime; small-prime, prime-power and long-prime parts of $\Phi^{\mathrm w}_1$ on $[T,2T]$ & \eqref{eq:taucw-split}\\
$\mathfrak Y$, $\mathsf F$, $\mathsf Y$, $\upsilon_m$ & threshold $(T/2)^{1/4}$; exponential polynomial $\sum_{p\le\mathfrak Y}p^{-1}p^{2it}$; square-prime part of $\Phi^{\mathrm{pw}}_T$; coefficients of $\mathsf F^2$ & \cref{thm:taucw-pow}\\
$\mathscr J_T$, $\mathscr G_T$, $\mathfrak f_t$, $\mathfrak j_t$ & continuous and Chebyshev-fluctuation integrals of the long-prime part; integrand amplitude and $x\mathfrak f_t'(x)$ & \eqref{eq:taucw-cont}, \eqref{eq:taucw-G}\\
$\mathfrak T$, $\mathfrak e$ & Chebyshev function $\sum_{p\le x}\log p$ and $\mathfrak T(x)-x$ & \cref{cor:taucw-cheb}\\
$\Upsilon$ & $e^\omega$ in \cref{thm:taucw-coef}, as in (L) of \S\ref{sec:v9-sawtooth}; $\Lambda(\Upsilon)=0$ unless a prime power & \cref{thm:taucw-coef}\\
$\varrho_{f;I}$, $\mathsf O_\lambda$, $\mathfrak M_\lambda$, $\mathsf J_\lambda$, $\mathsf B_\lambda$, $\mathfrak N_\lambda$, $\mathsf Q_\lambda$, $\mathsf e_\pm$ & local reordering defect of a bounded function on $I$; sublevel set $\{f<\lambda\}$ and the end of its left-packed copy $\mathsf J_\lambda=[a,\mathfrak M_\lambda)$; sublevel set $\{g<\lambda\}$ and its end $\mathfrak N_\lambda$; weighted symmetric-difference mass; measures of the two differences & \eqref{eq:taucw-local}, \cref{thm:taucw-quad}\\
$\mathsf d$, $\mathsf b$, $\mathfrak l$, $\operatorname{Lip}$ & Lipschitz test function; smooth cutoff; lower bound of $g'$; Lipschitz constant & \cref{thm:taucw-quad}, \cref{cor:taucw-lowfreq}\\
$\mathscr M_\ell$, $\mathsf C_\ell$, $\mathfrak x_j$, $\mathfrak y_j$, $\mathsf j_1,\mathsf j_2$ & counting discrepancy and spectral difference of two ordered lists of length $\ell$; the lists; their counting functions in $j$ & \eqref{eq:taucw-M}, \cref{thm:taucw-plan}\\
$\mathfrak F_{\mathfrak B}$, $\mathfrak B$ & Fej\'er kernel and its bandwidth & \cref{thm:taucw-fejer}\\
$N_{[\ell]}$, $N^{\mathrm w}_{[\ell]}$, $\mathsf g^{\mathrm w}$, $\mathsf h^{\mathrm w}$ & counting functions of the first $\ell$ points of $\gamma$ and of $\tau^{\mathrm w}$; $\limsup|N^{\mathrm w}-N_0|/\log t$ and $\limsup|N-N^{\mathrm w}|/\log t$ & \cref{prop:taucw-S}, \eqref{eq:taucw-gh}\\
$\mathcal J_L(y;v)$, $\mathscr R_{L,y}(u)$ & profile of the pair $\{\rho,1-\bar\rho\}$, $y=\beta-\frac12$, as a function of $v=\gamma-t$; its spectral density on $0<u<2L$ & \eqref{eq:taucw-J}\\
$\mathcal N$, $\mathcal J_{\mathcal N}$, $\mathcal H$ & finite set of pairs $(y_{\mathbf p},\gamma_{\mathbf p})$; sum of their profiles; weighted exponential sum $\sum_{\mathbf p}\mathscr R_{L,y_{\mathbf p}}(u)e^{iu\gamma_{\mathbf p}}$ & \cref{prop:taucw-energy}\\
$\mathcal Q_\ell$ & indices $j<\ell$ with a shared imaginary part $\gamma_j=\gamma_{j+1}$ and $j-\frac12>F^{\mathrm w}_\theta(2\pi)$ & \cref{prop:taucw-shared}\\
$A^{\mathrm w}_X$, $c^{\mathrm w}_n$ & Dirichlet polynomial $\pi\sum_{n\le X^2}\mathfrak k_nW(\frac{\log n}{\log X})n^{it}$ and its coefficients $\pi\mathfrak k_nW(\frac{\log n}{\log X})$ & \cref{lem:taucw-mvsel}\\
$\mathrm B_1$ & Meissel--Mertens constant, $\sum_{p\le x}p^{-1}=\log\log x+\mathrm B_1+o(1)$ & \cref{thm:v20-anticorr}, \cref{sec:taucw-beyond}\\
$\mathsf z$, $X(t)$, $t_0(\mathsf z)$, $\Phi^{\mathrm w}_{\langle\mathsf z\rangle}$, $\Phi^{\mathrm{fr}}_{\langle\mathsf z\rangle,T}$ & exponent $0\le\mathsf z<\frac32$; cutoff $t^{1/2}(\log t)^{\mathsf z}$ for $t\ge t_0(\mathsf z)$; phase with that cutoff; its version with weights frozen at $T$ & \eqref{eq:taucw-beyond-def}, \cref{prop:taucw-beyond}\\
$\Phi^{\mathrm w}_{[1/2,1]}$, $\mathscr E^{\mathrm{off}}_T$, $W_{[1/2,1]}$ & increment $\Phi^{\mathrm w}_1-\Phi^{\mathrm w}_{1/2}$ (with $\mathscr K_n$ at $(\theta_-,\theta_+)=(\frac12,1)$); off-diagonal part of its second moment; $W(v)-W(2v)$ & \eqref{eq:taucw-inc1-def}, \eqref{eq:taucw-inc1-off}, \cref{prop:taucw-inc1}\\
$W^{\mathrm{ev}}$, $\mathbb T_L$ & even extension of the Selberg weight; triangle $(L-|y|)_+$ & \cref{lem:taucw-fejer}\\
$L(t)$, $L_\theta$ & $\theta\log t$; $\theta\log T$ & \eqref{eq:taucw-L4}, \cref{prop:taucw-Yenergy}\\
$C_{\mathfrak K}$ & $\int|\mathfrak K_1|$ & \cref{lem:taucw-kernel}\\
$\mathsf s_\circ$ & spacing threshold & \eqref{eq:taucw-shared-count}\\
$\gamma^\diamond$, $\mathscr Z^\diamond$ & height bound; zeros off the line with $|\gamma|\le\gamma^\diamond$ & \cref{prop:taucw-bh}\\
$\mathsf t_X$, $C_{\mathrm{tail}}$ & truncation height $(X^2(1+\log\log3X))^{1/3}$; constant of the tail bound & \cref{lem:taucw-tail}\\
$\nu_X$, $\mathfrak C$, $\mathbb J$ & signed measure representing $P^W_X$; cosine transform of $\frac12(1+t)^{-3}$; Gram kernel & \cref{lem:taucw-gram}\\
$\mathcal P_\theta(s;t)$, $\ell^{\mathrm w}_\theta(t)$ & finite weighted Euler product; its logarithm at $s=\frac12+it$ & \eqref{eq:taucw-Euler}\\
$\psi_{\mathrm{ind}}$, $\widetilde{\mathbf G}_2$, $N(T,2T)$, $N^{\mathrm{sc}}(T,2T)$, $c_{\mathrm{HS}}$, $C_{\mathrm{AF}}$, $C_{\mathrm{MT}}$ & indicator window; matrix of \cite{AlpogeFurman2026} at $L=2\mathsf L(T)$; zeros and simple critical zeros with $T<\gamma\le2T$; Hilbert--Schmidt correction; $\frac32-\frac1{\sqrt2}\cot\frac1{\sqrt2}$; $2-C_{\mathrm{AF}}$ & \cref{prop:taucw-AF2}, \cref{rem:taucw-AF2}\\
$\mathscr A(t)$ & pole term $\Im E_1((-\frac12+it)\log2)$, independent of $X$ & \cref{lem:taucw-expl}\\
$C_0$, $C$, $\mathcal E_L$ & constants and bounded remainder in the decomposition of $P^W_X$ & \cref{thm:taucw-P}\\
\bottomrule
\end{longtable}
\endgroup

\section{Definitions and scope}\label{sec:definitions}
Let $\rho_k=\beta_k+i\gamma_k$ denote the nontrivial zeros of $\zeta(s)$ with $\gamma_k>0$, counted with multiplicity and indexed so that $\gamma_k\le\gamma_{k+1}$. No assumption $\beta_k=\tfrac12$ is made. Let $N(t)$ count these zeros with $0<\Im\rho\le t$, including multiplicity, and put
\begin{equation}\label{eq:N0}
N_0(t)=\frac{t}{2\pi}\log\frac{t}{2\pi e}+\frac78\qquad(t>0).
\end{equation}
The unconditional Riemann--von Mangoldt identity, away from the imaginary parts of zeros, is
\begin{equation}\label{eq:RVM}
N(t)=N_0(t)+S(t)+R(t),\qquad
S(t)=\frac1\pi\arg\zeta(\tfrac12+it),\qquad R(t)=O(t^{-1}).
\end{equation}
Here $R(t)=\vartheta(t)/\pi+1-N_0(t)$, with $\vartheta$ the standard Riemann--Siegel theta function. At the imaginary parts of zeros, $N$ is right-continuous ($N(t)$ counts the zeros with $0<\Im\rho\le t$), and $S(t)$ is \emph{defined} there by \eqref{eq:RVM}, that is $S(t):=N(t)-N_0(t)-R(t)$; thus \eqref{eq:RVM} holds for every $t>0$ and $S$ is right-continuous. Bounds for $|S|$ are quoted from the literature only at heights that are not imaginary parts of zeros, where $S$ has no convention-dependence; they are extended to a zero height $\gamma$ by letting $t\downarrow\gamma$ through non-heights, which is legitimate because $S$ is right-continuous by the definition just given and the bounding functions are continuous. No pointwise statement of a cited theorem at a jump is assumed. The distinction between $N(t)$ and the count of zeros on the critical line is essential throughout.

For $a\ge0$ and $x\in\mathbb R$ define the symmetric truncation
\begin{equation}\label{eq:clip}
\clip_a(x)=\max\{-a,\min\{a,x\}\}.
\end{equation}
For a continuous function $F:(0,\infty)\to\mathbb R$ whose sublevel sets below the levels $k-\frac12$ have finite measure, put
\begin{equation}\label{eq:sublevel}
u_k(F)=\meas\{t>0:F(t)<k-\tfrac12\},\qquad k\ge1.
\end{equation}
No monotonicity assumption on $F$ is made. The central problem is to determine how accurately $u_k(F)$ approximates $\gamma_k$, including the possible occurrence of repeated heights $\gamma_k=\gamma_{k+1}$.

\paragraph{Status convention.}
An \emph{unconditional} assertion means no RH or simplicity assumption. References to the IJNT fourth-moment and gap estimates are expressly identified as inputs of the cited manuscript, rather than independently reverified here. Statements about the full weighted Euler phase as $X\to\infty$ are \emph{not} inferred from bounds proved only for $X\le T^{1/4}$. Section~\ref{sec:open} lists the analytic steps that remain open. Sections~\ref{sec:near-v4}--\ref{sec:observed-v4} further isolate the local two-source model from the actual prime phase, including its residual; \cref{sec:correlated-v5} gives the RvM-displacement comparison.

\paragraph{The Fujii input.}
All uses of Fujii's theorem below refer to the statement in \cite{CasanaIJNT} (Theorem ``Fujii, for the modulus $1$, with the corrections of \cite[\S5]{Fujii1981}'', label \texttt{thm:known\_fujii}): Fujii's Main Theorem \cite[p.~225]{Fujii1974}, together with Selberg's approximation as given in \cite[\S3, pp.~232--233]{Fujii1974} and \cite[Lemma~5]{Fujii1974}. Let $\mathsf k\ge1$ be an integer, let $0<\alpha_{\rm F}\le\frac12$ be fixed, and let $T^{1/2+\alpha_{\rm F}}\le H\le T$.
\begin{enumerate}[label=(\alph*),leftmargin=*]
\item For $T^{\alpha_{\rm F}/(20\mathsf k)}\le Y\le H^{1/\mathsf k}$,
\begin{equation}\label{eq:FujiiA}
\int_T^{T+H}\Bigl|S(t)-\frac1\pi\operatorname{Im}\sum_{p<Y}p^{-1/2-it}\Bigr|^{2\mathsf k}dt\ll_{\mathsf k,\alpha_{\rm F}}H.
\end{equation}
\item For $0<h\le1$ and $T\ge T_0(\mathsf k,\alpha_{\rm F})$,
\begin{equation}\label{eq:FujiiB}
\int_T^{T+H}\bigl(S(t+h)-S(t)\bigr)^{2\mathsf k}dt\ll_{\mathsf k,\alpha_{\rm F}}H\bigl(\log(2+h\log T)\bigr)^{\mathsf k}.
\end{equation}
\end{enumerate}
Only $\mathsf k\le2$ and $H=T$ are used below; as explained in \cite{CasanaIJNT}, the corrections of \cite[\S5]{Fujii1981} do not affect (a) or (b) in this range.

\section{The original taucs and their amplitude obstruction}\label{sec:original}
Write $p_1<p_2<\dots$ for the rational primes and set
\begin{equation}\label{eq:Phi-m}
\Phi_m(t)=-\frac1\pi\sum_{\nu\le m}\arg(1-p_\nu^{-1/2-it})
=-\frac1\pi\sum_{\nu\le m}\sum_{r\ge1}
 \frac{\sin(rt\log p_\nu)}{r p_\nu^{r/2}},\qquad \Phi_0=0.
\end{equation}
Here $\arg(1-p^{-1/2-it})$ denotes the continuous branch determined by the absolutely convergent logarithmic series $\log(1-z)=-\sum_{r\ge1}z^r/r$, $|z|<1$; thus $\arg(1-p^{-1/2-it})=\sum_{r\ge1}\sin(rt\log p)/(rp^{r/2})$, and the second equality in \eqref{eq:Phi-m} fixes the branch instead of being an identity modulo $2\pi$.
For $2^j\le t\le2^{j+1}$ define the dyadically interpolated phase
\begin{equation}\label{eq:interpolation}
m_j=\lfloor2^{j/2}\rfloor,\qquad
\Phi^\circ(t)=\left(2-\frac{t}{2^j}\right)\Phi_{m_j}(t)
 +\left(\frac{t}{2^j}-1\right)\Phi_{m_{j+1}}(t).
\end{equation}
The original construction in \cite{CasanaIJNT} is
\begin{equation}\label{eq:old}
h(t)=\log\log\max(t,e^e),\qquad
F_h(t)=N_0(t)+\clip_{h(t)}(\Phi^\circ(t)),\qquad
\tau_k^{\mathrm C}=u_k(F_h).
\end{equation}
The IJNT manuscript establishes its sublevel construction and, for the indices in the asymptotic range, strict increase and that
\begin{equation}\label{eq:oldMS}
\sum_{T<\gamma_k\le2T}|\tau_k^{\mathrm C}-\gamma_k|^2\ll\frac{T}{\log T}.
\end{equation}
It does not prove pointwise containment between consecutive values of $\gamma_k$.

\begin{lemma}[Necessary counting-function bound]\label{lem:countconstraint}
Let $A(t)\ge0$ be continuous. Suppose $N_0(t)\pm A(t)$ are eventually strictly increasing,
$N_0(t)-A(t)\to\infty$, and
$\sup_{0<s\le t}(N_0(s)+A(s))=N_0(t)+A(t)$ for all sufficiently large $t$.
If a continuous $F$ satisfies $|F(t)-N_0(t)|\le A(t)$ and
\begin{equation}\label{eq:neighbor}
\gamma_{k-1}\le u_k(F)<\gamma_{k+1}\qquad\text{for every sufficiently large }k,
\end{equation}
then at every sufficiently large $t$ which is not the imaginary part of a zero,
\begin{equation}\label{eq:countnecessary}
|N(t)-N_0(t)|\le A(t)+\tfrac32.
\end{equation}
\end{lemma}
\begin{proof}
Choose $T_0$ and $k_0$ such that the eventual monotonicity of $N_0\pm A$, the running-maximum identity and the containment \eqref{eq:neighbor} hold for $t\ge T_0$ and $k\ge k_0$. Let $t\ge T_0$ be so large that $q:=N(t)$ satisfies $q-1\ge k_0$, and let $t$ not be a zero height, so that $\gamma_q<t<\gamma_{q+1}$. If
$q-\frac32>N_0(t)+A(t)$, the upper bound for $F(s)$ holds below $q-\frac32$ on $(0,t]$, whence $u_{q-1}(F)\ge t$. This contradicts $u_{q-1}(F)<\gamma_q<t$. If $q+\frac32<N_0(t)-A(t)$, the lower bound for $F(s)$ holds above $q+\frac32$ for every $s\ge t$ by eventual monotonicity, whence $u_{q+2}(F)\le t$, contradicting $u_{q+2}(F)\ge\gamma_{q+1}>t$.
\end{proof}

\begin{theorem}[Unconditional failure of eventual neighboring-zero containment]\label{thm:omega}
For the original $\tau_k^{\mathrm C}$ of \eqref{eq:old}, the inequalities
\begin{equation}\label{eq:impossible}
\gamma_{k-1}\le\tau_k^{\mathrm C}<\gamma_{k+1}
\end{equation}
cannot hold for all sufficiently large $k$. For every fixed integer $r\ge1$ there are arbitrarily large $k$ with $\tau_k^{\mathrm C}\ge\gamma_{k+r}$ and arbitrarily large $k$ with $\tau_k^{\mathrm C}\le\gamma_{k-r}$. The same statements hold for $u_k(F)$ with any continuous $F$ satisfying $|F-N_0|\le h$ on $(0,\infty)$ (finiteness of $u_k(F)$ is automatic because $N_0-h\to\infty$).
\end{theorem}
\begin{proof}
The truncation $h$ satisfies the conditions of \cref{lem:countconstraint}. The unconditional Tsang $\Omega_\pm$ theorem \cite{Tsang1986,Dobner2024} gives
\begin{equation}\label{eq:omegaS}
S(t)=\Omega_\pm\!\left(\left(\frac{\log t}{\log\log t}\right)^{1/3}\right)
\end{equation}
along arbitrarily large heights away from zeros. Since
$(\log t/\log\log t)^{1/3}/\log\log t\to\infty$,
\eqref{eq:countnecessary} with $A=h$ contradicts \eqref{eq:RVM}.
For the stronger assertion, at a positive large value of $S(t)$ put $q=N(t)$ and use each index $k$ with $N_0(t)+h(t)<k-\tfrac12$ and $k+r\le q$; then
$u_k(F)\ge t>\gamma_{k+r}$. The number of such indices tends to infinity along a positive $\Omega$ sequence. The negative sequence gives the analogous assertion on the other side, using $N_0-h$.
\end{proof}

\begin{corollary}[Necessary truncation scale]\label{cor:amp}
Under the monotonicity hypotheses of \cref{lem:countconstraint}, no continuous $F$ with $|F-N_0|\le A$ satisfies the eventual containment \eqref{eq:neighbor} if
\begin{equation}\label{eq:amp-scale}
A(t)=o\!\left((\log t/\log\log t)^{1/3}\right).
\end{equation}
\end{corollary}
\begin{proof}Combine \cref{lem:countconstraint} with \eqref{eq:omegaS}.\end{proof}

\section{Modified taucm from the explicit MCOM bound}\label{sec:taucm}
The accompanying Mathematics of Computation submission \cite{CasanaMCOM} obtains, without RH,
\begin{equation}\label{eq:MCOM}
|S(t)|\le0.08382\log t+4.444\qquad(t\ge1),
\end{equation}
with the compatible convention at zero heights. Its other bound is
$|N(t)-\frac{t}{2\pi}\log\frac{t}{2\pi e}|\le0.08382\log t+5.326$ for $t\ge1$ (not used below).
Combined with the explicit estimate $|R(t)|\le0.53/(\pi t)$ for $t\ge1$ and the elementary range $0<t<1$, \eqref{eq:MCOM} gives
\begin{equation}\label{eq:MCOMH}
|G(t)|\le H(t),\qquad
G(t):=N(t)-N_0(t),\qquad
H(t):=0.08382\log\max(t,e)+5.5
\end{equation}
for almost every $t>0$ (and under the stated conventions also at the jumps). Indeed $4.444+0.53/\pi<5.5$; for $0<t<1$ one has $N(t)=0$, so $G(t)=-N_0(t)$, and $N_0(t)=\frac7{8}+\frac{t}{2\pi}\log\frac t{2\pi e}$ with $-\frac{1}{2\pi}\bigl(e^{-1}+1+\log2\pi\bigr)<\frac t{2\pi}\log\frac t{2\pi e}\le0$ (because $t|\log t|\le e^{-1}$ and $t<1$), whence $0<N_0(t)\le\frac78<5.5\le H(t)$.

\begin{definition}[Improved taucm]\label{def:taucm}
Retain the original $\Phi^\circ$ and its prime budget $m_j=\lfloor2^{j/2}\rfloor$, and set
\begin{equation}\label{eq:taucm}
F^{\Cm}(t)=N_0(t)+\clip_{H(t)}(\Phi^\circ(t)),\qquad
\tau_k^{\Cm}=u_k(F^{\Cm}).
\end{equation}
\end{definition}

\begin{proposition}[Contractivity of covered truncation]\label{prop:contract}
If $|y|\le a$, then for every real $x$,
\begin{equation}\label{eq:contract}
|\clip_a(x)-y|\le|x-y|.
\end{equation}
Consequently, taking $x=\Phi^\circ(t)$, $y=G(t)$, $a=H(t)$ and using \eqref{eq:MCOMH}, for almost all $t>0$,
\begin{equation}\label{eq:contract-F}
|F^{\Cm}(t)-N(t)|\le|N_0(t)+\Phi^\circ(t)-N(t)|.
\end{equation}
\end{proposition}
\begin{proof}
Truncation to $[-a,a]$ is the nearest-point projection onto a closed interval. Apply its elementary distance inequality with $x=\Phi^\circ(t)$ and $y=G(t)$.
\end{proof}

\begin{theorem}[Fourth moment and mean-square localization for taucm]\label{thm:taucm}
Assume the explicit bound \eqref{eq:MCOMH} (from \eqref{eq:MCOM}, \cite{CasanaMCOM}) and the following inputs, established without RH in the cited IJNT manuscript \cite{CasanaIJNT}: its fourth-moment estimate for the \emph{untruncated} phase (Lemma ``Fourth moment of the counting error'', label \texttt{lem:fourth}),
\begin{equation}\label{eq:IJNTinput}
\int_U^{2U}|N_0(t)+\Phi^\circ(t)-N(t)|^4\dd\ll U
\quad(U=2^j\text{ sufficiently large}),
\end{equation}
and its third-moment estimate for the successive differences $\gamma_{k+1}-\gamma_k$ (Lemma ``Third moment of the gaps'', label \texttt{lem:gaps}, used in the form $\sum_{\gamma_j\in[T/2,3T]}(\gamma_{j+1}-\gamma_j)^3\ll T(\log T)^{-2}$) and deterministic sublevel-transfer theorem (Theorem ``Sublevel transfer'', label \texttt{thm:transfer}, with constant $2\cdot24^{1/3}$).
Then, without RH and for all sufficiently large $T$,
\begin{align}
\int_T^{2T}|F^{\Cm}(t)-N(t)|^4\dd&\ll T,\label{eq:cm-fourth}\\
\sum_{T<\gamma_k\le2T}|\tau_k^{\Cm}-\gamma_k|^2
&\ll\frac{T}{\log T}.\label{eq:cm-mean}
\end{align}
The implied constants and the threshold for $T$ are absolute, and the prime budget is unchanged.
\end{theorem}
\begin{proof}
By \cref{prop:contract}, \eqref{eq:IJNTinput} passes directly to $F^{\Cm}$. If $2^j\le T<2^{j+1}$, then $[T,2T]\subset[2^j,2^{j+2}]=[2^j,2^{j+1}]\cup[2^{j+1},2^{j+2}]$, and applying \eqref{eq:IJNTinput} on these two dyadic blocks proves \eqref{eq:cm-fourth}. For large $T$, \eqref{eq:MCOMH} and $H(t)=O(\log(2+t))$ ensure that all sublevel sets for indices with $T<\gamma_k\le2T$ contain $(0,T/2]$ and are contained in $(0,3T)$, exactly as in the corresponding localization step of IJNT: the separations in $N_0$ across these endpoints have order $T\log T$, whereas both the correction and the difference $N-N_0$ are $O(\log T)$. Explicitly, let $T<\gamma_k\le2T$ with $T$ large, so that $N(T)\le k-1$ and $N(2T)\ge k$. Fix $T_0$ such that $N_0\pm H$ are increasing on $[T_0,\infty)$ and put $C_{T_0}=\sup_{0<t\le T_0}(N_0(t)+H(t))$; take $T$ so large that $T/2\ge T_0$, $N_0(3T)-H(3T)>N_0(2T)+H(2T)$ and $N_0(T)-H(T)>\max\bigl(C_{T_0},\,N_0(T/2)+H(T/2)\bigr)$. If $0<t\le T/2$, then $F^{\Cm}(t)\le N_0(t)+H(t)\le\max\bigl(C_{T_0},N_0(T/2)+H(T/2)\bigr)<N_0(T)-H(T)\le N(T)\le k-1$ by \eqref{eq:MCOMH}, which covers the initial range $t\le T_0$ as well. If $t\ge3T$, then $F^{\Cm}(t)\ge N_0(t)-H(t)\ge N_0(3T)-H(3T)>N_0(2T)+H(2T)\ge N(2T)\ge k$. The deterministic sublevel-transfer inequality of IJNT, with integer discrepancy
$q_F(t)=|\lfloor F(t)+\frac12\rfloor-N(t)|$, gives
\begin{equation}
\sum_{T<\gamma_k\le2T}|\tau_k^{\Cm}-\gamma_k|^2
\ll \left(\sum_{\gamma_j\in[T/2,3T]}(\gamma_{j+1}-\gamma_j)^3\right)^{1/2}
\left(\int_{T/2}^{3T}q_{F^{\Cm}}(t)^4\dd\right)^{1/2}.
\end{equation}
The first factor is $O(T^{1/2}/\log T)$ by the cited gap input; the second is $O(T^{1/2})$ because $q_F\le|F-N|+\frac12$ and \eqref{eq:cm-fourth} applies on a bounded collection of dyadic intervals. Their product proves \eqref{eq:cm-mean}.
\end{proof}

\begin{remark}[Projection does not order the two truncations]
\label{rem:no-h-cm-comparison}
The contractivity estimate \eqref{eq:contract-F} compares the improved function
with the \emph{untruncated} phase $N_0+\Phi^\circ$; it does not yield
$|F^{\Cm}-N|\le |F_h-N|$. A larger truncation interval need not improve
pointwise distance to a target lying outside the smaller interval.
Hence the input required above is specifically the untruncated fourth moment
\eqref{eq:IJNTinput}, not a fourth-moment estimate deduced solely for $F_h$.
The elementary bound $|F^{\Cm}-N|\le 2H$ holds whenever
$|N-N_0|\le H$.
\end{remark}

\begin{remark}[Asymptotic refinement from the MCOM result]\label{rem:MCOM-asympt}
The same MCOM submission also states, without RH, the stronger asymptotic inequality
\begin{equation}\label{eq:asymptMCOM}
\limsup_{t\to\infty}\frac{|S(t)|}{\log t}<0.0748596142711.
\end{equation}
Consequently, for every fixed $c>0.0748596142711$ there exists $T_c$ such that
$|N(t)-N_0(t)|\le c\log t$ for all $t\ge T_c$ away from zero heights. One may therefore replace $H$ by a positive continuous truncation bound that agrees with $c\log t$ above $T_c$, and that dominates $|N-N_0|$ below $T_c$. The same projection argument applies. Unlike \eqref{eq:MCOMH}, this refinement does not provide an effective value of $T_c$ and cannot be used as an explicit estimate at a specified height.
\end{remark}

\begin{remark}[Exact scope]\label{rem:taucm-scope}
The modification removes the specific impossibility mechanism in \cref{thm:omega}: its truncation is large enough to contain $N-N_0$. It does not prove that every $\tau_k^{\Cm}$ lies between the imaginary parts of neighboring zeros, nor does it establish RH or simplicity. The argument transfers established IJNT estimates to a different truncation; it does not derive a uniform result for arbitrary prime cutoffs $X$.
\end{remark}

\section{The structural saturation obstruction}\label{sec:saturation}
The following is the precise mathematical content of the objection that ordering, truncation and the general sublevel-transfer inequality cannot by themselves prohibit saturation. It is not a theorem asserting that the \emph{actual} weighted Euler phases saturate.

\begin{theorem}[Structural saturation model]\label{thm:structural}
Let $H$ be given by \eqref{eq:MCOMH}. For $Q>0$ define
\begin{equation}\label{eq:sat-model}
F_Q(t)=N_0(t)+H(t)\tanh\bigl(Q\sin(Qt)\bigr),
\qquad u_k(Q)=u_k(F_Q).
\end{equation}
For every $Q>0$, $F_Q$ is continuous, $|F_Q-N_0|\le H$, and its sublevel measures are finite. For all sufficiently large $Q$ (depending on $k$ where $k$ is involved), these sublevel measures are moreover positive and strictly increasing in $k$, and every deterministic sublevel-transfer inequality applicable to continuous $F$ applies to $F_Q$. For each fixed sufficiently large $k$,
\begin{equation}\label{eq:sat-limit}
\lim_{Q\to\infty}u_k(Q)
=\frac12\meas\{N_0+H<k-\tfrac12\}
 +\frac12\meas\{N_0-H<k-\tfrac12\}.
\end{equation}
Writing $\varphi_\pm=N_0\pm H$, which are strictly increasing near the relevant levels for large $k$, one also has
\begin{equation}\label{eq:sat-spacing}
\lim_{Q\to\infty}(u_{k+1}(Q)-u_k(Q))
=\frac12\bigl[\varphi_+^{-1}(k+\tfrac12)-\varphi_+^{-1}(k-\tfrac12)
+\varphi_-^{-1}(k+\tfrac12)-\varphi_-^{-1}(k-\tfrac12)\bigr].
\end{equation}
These properties do not characterize the prime-corrected functions $F_X^W$ and do not provide information on $\beta_k$.
\end{theorem}
\begin{proof}
For every $Q>0$ the bounds and continuity follow immediately from \eqref{eq:sat-model}. Since $N_0-H\to+\infty$, all relevant sublevel measures are finite. For $Q$ large, a point near $3\pi/(2Q)$ has $F_Q<1/2$, whereas $F_Q(t)\to+\infty$ as $t\to\infty$. Continuity gives an open interval on which $F_Q<\frac12$, so $u_k(Q)>0$ for every $k\ge1$. Moreover, by the intermediate-value theorem $F_Q$ takes a value in the open interval $(k-\frac12,k+\frac12)$ at some point; the inverse image of this open interval is a nonempty open subset of $(0,\infty)$, hence of positive Lebesgue measure, and it is contained in $\{k-\frac12\le F_Q<k+\frac12\}$. Therefore $u_{k+1}(Q)-u_k(Q)>0$, and $0<u_k(Q)<u_{k+1}(Q)$.

Fix a compact interval $I$ and $c\in\mathbb R$ such that the level sets $\{N_0\pm H=c\}\cap I$ have measure zero. Put $\eta_Q=Q^{-1/2}$ and $I_Q^\pm=\{t\in I:\pm\sin(Qt)\ge\eta_Q\}$. Then $\meas(I\setminus(I_Q^+\cup I_Q^-))\to0$, and on $I_Q^\pm$ one has $|F_Q-(N_0\pm H)|\le H(1-\tanh\sqrt Q)\to0$ uniformly. Since $\meas(\{|N_0\pm H-c|\le\epsilon\}\cap I)\to0$ as $\epsilon\to0$, it follows that $\meas(I_Q^\pm\cap\{F_Q<c\})-\meas(I_Q^\pm\cap\{N_0\pm H<c\})\to0$. Finally, for every measurable $A\subset I$, $\meas(A\cap I_Q^\pm)\to\frac12\meas(A)$: the indicator of $\{\pm\sin(Qt)\ge\eta_Q\}$ differs from that of $\{\pm\sin(Qt)>0\}$ on a set of measure $o(1)$, and the latter is $\frac12$ plus a $2\pi/Q$-periodic bounded function $g(Qt)$ of mean zero, $g=\ind_{\{\sin>0\}}-\frac12$. By the following averaging lemma, $\int_Ig(Qt)\ind_A(t)\,dt\to0$. \emph{Lemma.} If $g$ is bounded, $2\pi$-periodic with mean zero and $f\in L^1(I)$, then $\int_If(t)g(Qt)\,dt\to0$ as $Q\to\infty$. \emph{Proof.} Given $\epsilon>0$ choose a step function $f_\epsilon=\sum_jc_j\ind_{J_j}$ (finitely many intervals) with $\|f-f_\epsilon\|_{L^1}<\epsilon$. Over any interval $J$ the integral $\int_Jg(Qt)\,dt$ is the sum of an integral over full periods, which vanishes, and an integral over a remainder of length at most $2\pi/Q$, so $|\int_Jg(Qt)\,dt|\le2\pi\|g\|_\infty/Q$. Hence $\limsup_Q|\int_Ifg(Q\cdot)|\le\epsilon\|g\|_\infty$ for every $\epsilon$. The level boundaries of $N_0\pm H$ are finite in any compact interval, for sufficiently large $k$, because the functions eventually increase. Since $N_0-H\to\infty$, the sublevel sets lie in a common compact interval for each fixed $k$. This proves \eqref{eq:sat-limit}. Taking differences yields \eqref{eq:sat-spacing}; for large $k$ the two sublevel measures equal the corresponding inverse-function values.
\end{proof}

\begin{theorem}[Fourth-moment obstruction to saturation]\label{thm:sat-fourth}
Let $G=N-N_0$. For every fixed $T\ge2$, the functions in \eqref{eq:sat-model} satisfy
\begin{equation}\label{eq:sat-L4}
\lim_{Q\to\infty}\int_T^{2T}|F_Q(t)-N(t)|^4\dd
=\int_T^{2T}\bigl(H^4+6H^2G^2+G^4\bigr)\dd
\ge\int_T^{2T}H(t)^4\dd.
\end{equation}
Consequently no estimate $\int_T^{2T}|F_Q-N|^4\dd\le CT$ can hold for all sufficiently large $T$ and all $Q>0$ with an absolute $C$, because
\begin{equation}\label{eq:sat-scale}
\int_T^{2T}H(t)^4\dd\asymp T(\log T)^4.
\end{equation}
This does not contradict the fixed-budget IJNT estimate: $F_Q$ is an abstract test construction, not its Euler phase.
\end{theorem}
\begin{proof}
For fixed $T$, $G$ and $H$ are bounded on $[T,2T]$. Weak equidistribution of the two signs, established in the proof of \cref{thm:structural} and applied on each interval between consecutive jumps of $N$, gives the first equality; these jumps form a finite set. The integrand's average is
$\frac12(H-G)^4+\frac12(-H-G)^4=H^4+6H^2G^2+G^4$.
To pass to the limit, note $|F_Q-N|^4=|H\tanh(Q\sin Qt)-G|^4\le(H+|G|)^4$ for $\sin(Qt)$-driven values of the saturated model, which is bounded on the fixed interval $[T,2T]$ and so provides an integrable dominant independent of $Q$; on the set where $|\sin Qt|\ge Q^{-1/2}$ the integrand differs from $(\pm H-G)^4$ by $o(1)$ uniformly, and the set where $|\sin Qt|<Q^{-1/2}$ has measure tending to $0$. The last estimate follows from $H(t)=0.08382\log t+O(1)$ throughout $[T,2T]$.
\end{proof}

\begin{corollary}[A conditional exclusion principle]\label{cor:exclusion}
Suppose a family of \emph{actual} weighted prime phases $F_X^W$ satisfies, along some $X_n\to\infty$ and for every fixed $T\ge2$, the limit relation \eqref{eq:sat-L4} with $F_Q$ replaced by $F_{X_n}^W$ (the saturated two-sign distribution of \cref{thm:sat-fourth}). Then the global estimate
\begin{equation}\label{eq:desired}
\sup_{X\ge T^\varepsilon}\int_T^{2T}|F_X^W-N|^4\dd\ll_\varepsilon T
\end{equation}
for fixed $\varepsilon>0$, uniformly in $T$, is impossible. Equivalently, a proof of \eqref{eq:desired} would exclude any scenario for the prime phase for which this saturation assertion is \emph{independently established}.
\end{corollary}
\begin{proof}
For fixed $T$ all sufficiently large $X_n$ satisfy $X_n\ge T^\varepsilon$. Take $n\to\infty$ in the asserted uniform estimate and use \eqref{eq:sat-L4}; then let $T\to\infty$. The resulting $T(\log T)^4\ll_\varepsilon T$ is impossible.
\end{proof}

\begin{remark}[What the objection does not prove]\label{rem:claude}
The preceding theorems show that the three elementary properties (increasing sublevel points, a bounded correction, and a universal transfer inequality) do not rule out saturation. They do \emph{not} prove that RH is false, that an arithmetic family necessarily saturates when RH is false, or that no further theorem using the specific prime coefficients can lead to RH. Establishing an off-critical-line saturation statement requires a separate explicit-formula and non-cancellation argument; its present status in \cite{CasanaDBSRH} is a proof sketch under additional hypotheses.
\end{remark}

\section{Smoothed prime powers and taucm(s)}\label{sec:weights}
Let $\mathcal W$ denote the class of weights
\begin{equation}\label{eq:weights}
W(u)=\int_1^2q(v)\ind_{\{u\le v\}}\,dv,
\quad q\ge0,\quad q\in BV[1,2],\quad \int_1^2q(v)\,dv=1.
\end{equation}
Thus $0\le W\le1$, $W=1$ for $0\le u\le1$, $W=0$ for $u\ge2$, and $W$ is nonincreasing. The Selberg weight corresponds to $q(v)=1$; the tested weight $W_{1/4}$ corresponds to $q(v)=v/2+1/4$. Let $\Lambda$ denote von Mangoldt's function. For $X\ge2$, set
\begin{align}
\Phi_X^W(t)&=\frac1\pi\operatorname{Im}\sum_{n\ge2}
 W\!\left(\frac{\log n}{\log X}\right)
\frac{\Lambda(n)}{n^{1/2+it}\log n},\label{eq:PhiW}\\
\Pi_X^W(t)&=\frac1\pi\operatorname{Im}\int_2^\infty
 W\!\left(\frac{\log u}{\log X}\right)
\frac{u^{-1/2-it}}{\log u}\,du,\label{eq:PiW}\\
P_X^W(t)&=\Phi_X^W(t)-\Pi_X^W(t),\label{eq:PW}\\
F_X^W(t)&=N_0(t)+\clip_{H(t)}(P_X^W(t)),\qquad
\tau_k^{\Cms}(X;W)=u_k(F_X^W).\label{eq:Cms}
\end{align}
The sum and the integral are finite because $W$ is supported on $[0,2]$. The compensator $\Pi_X^W$ corresponds to the pole contribution. These definitions do not use $\beta_k$.

\begin{definition}[DBSRH (Deep Blue Sea Riemann Hypothesis) for a weight]\label{def:DBSRH}
For a fixed $W\in\mathcal W$, the convergence statement is
\begin{equation}\label{eq:DBSRH}
\mathrm{DBSRH}(W):\qquad
\lim_{X\to\infty}\tau_k^{\Cms}(X;W)=\gamma_k
\quad\text{for each fixed }k\ge1.
\end{equation}
No equivalence between \eqref{eq:DBSRH} and RH is assumed or proved here.
\end{definition}

\begin{lemma}[Pole compensator]\label{lem:pole}
Uniformly over $W\in\mathcal W$, $X\ge2$, and $t\ge1$,
\begin{equation}\label{eq:pole-bound}
|\Pi_X^W(t)|\ll\frac{1+X}{t}.
\end{equation}
\end{lemma}
\begin{proof}
With $L=\log X$ and $v=\log u$, the integral in \eqref{eq:PiW} has amplitude $a(v)=W(v/L)e^{v/2}/v$, supported on $[\log2,2L]$. The total variation of $a$ plus its endpoint values is $O(1+X)$, uniformly in $W$: $W$ is nonincreasing with total variation at most one, and $e^{v/2}/v$ has bounded variation $O(1+X)$ on this interval. Integration by parts against $e^{-itv}$ proves the result.
\end{proof}

\section{A rigorous restriction on small prime cutoffs}\label{sec:subpoly}
The next theorem demonstrates that a fourth-moment estimate uniform over \emph{all} diverging cutoffs $X=X(T)$ is false, even without any hypothesis on RH.

\begin{theorem}[Subpolynomial cutoff obstruction]\label{thm:smallX}
Subject to the Fujii input \eqref{eq:FujiiA}, let $X=X(T)\ge2$ satisfy $\log X=o(\log T)$ (as $T\to\infty$, or along a sequence $T_n\to\infty$, in which case the conclusion holds along that sequence). For every choice $W=W_T\in\mathcal W$ (allowed to depend on $T$),
\begin{equation}\label{eq:smallX}
\int_T^{2T}|F_X^W(t)-N(t)|^4\dd
\ge\left(\frac1{4\pi^4}+o(1)\right)
 T\left(\log\frac{\log T}{\log X}\right)^2.
\end{equation}
The $o(1)$ is uniform over the indicated weights and, along the stated cutoff sequence, tends to zero. In particular, the left-hand side divided by $T$ diverges.
\end{theorem}
\begin{proof}
Put $Y=T^{1/200}$. By \eqref{eq:FujiiA} with $\mathsf k=2$, $\alpha_{\rm F}=\frac15$ and Fujii's interval length equal to $T$ (so that $Y=T^{\alpha_{\rm F}/(20\mathsf k)}$), and H\"older's inequality, its $L^2$ consequence is
\begin{equation}\label{eq:FujiiL2}
\left\|S(t)-\frac1\pi\operatorname{Im}\sum_{p<Y}p^{-1/2-it}\right\|_{L^2[T,2T]}
\ll\sqrt T.
\end{equation}
For large $T$, $X^2<Y$. The difference between the short prime sum and the prime terms of $\Phi_X^W$ is the imaginary part divided by $\pi$ of
\begin{equation}\label{eq:diffprime}
D(t)=\sum_{p<Y}\frac{a_p}{\sqrt p}p^{-it},\qquad
a_p=1-W\!\left(\frac{\log p}{\log X}\right),
\end{equation}
where $a_p=0$ for $p\le X$, $0\le a_p\le1$ for $X<p\le X^2$, and $a_p=1$ for $X^2<p<Y$.
Mertens' theorem \cite{Mertens1874} yields, uniformly in $W$,
\begin{equation}\label{eq:diffvar}
\sum_{p<Y}\frac{a_p^2}{p}
=\log\frac{\log T}{\log X}+O(1).
\end{equation}
Indeed the interval $(X,X^2]$ contributes $O(1)$, and the interval $(X^2,Y)$ gives the claimed logarithm up to an absolute constant.

Montgomery--Vaughan \cite{MontgomeryVaughan} gives
\begin{equation}
\int_T^{2T}|D(t)|^2\dd
=T\sum_{p<Y}\frac{a_p^2}{p}+O(Y).
\end{equation}
Also $\int_T^{2T}D(t)^2\dd=O(Y)$ by direct integration of $(pq)^{-it}$, since $\log(pq)\ge\log4$. Hence
\begin{equation}\label{eq:imD}
\int_T^{2T}\left|\frac{\operatorname{Im}D(t)}\pi\right|^2\dd
=\frac{T}{2\pi^2}\left(\log\frac{\log T}{\log X}+O(1)\right)+O(Y).
\end{equation}
The sum over prime powers $p^r$, $r\ge2$, in \eqref{eq:PhiW} has $L^2[T,2T]$ norm $O(\sqrt T)$ uniformly for $X=T^{o(1)}$: for $r=2$ use Montgomery--Vaughan on
$\tfrac12\sum_{p\le X}W(2\log p/\log X)p^{-1-2it}$, whose diagonal $\sum_p p^{-2}$ converges and whose remainder is $O(\pi(X))=o(T)$; for $r\ge3$, since $0\le W\le1$, $\bigl|\sum_p\sum_{r\ge3}W(r\log p/\log X)\,p^{-rit}/(rp^{r/2})\bigr|\le\sum_p\sum_{r\ge3}1/(rp^{r/2})<\infty$ with a constant independent of $X$, $W$ and $t$. By \cref{lem:pole}, $\|\Pi_X^W\|_2\ll X/\sqrt T=o(\sqrt T)$; also $\|R\|_2\ll T^{-1/2}$. Combining these estimates with \eqref{eq:FujiiL2} and \eqref{eq:imD}, and writing $D_T=\log(\log T/\log X)\to\infty$, gives
\begin{equation}\label{eq:preclipL2}
\|N_0+P_X^W-N\|_2
\ge \left(\frac1{\sqrt2\pi}+o(1)\right)\sqrt{TD_T}.
\end{equation}

We now check that truncation does not alter the principal size. Because $X^4=T^{o(1)}$, a second application of Montgomery--Vaughan to the square of the prime polynomial yields
\begin{equation}\label{eq:powerfour}
\int_T^{2T}|P_X^W(t)|^4\dd\ll T\bigl(1+\log\log(3X)\bigr)^2.
\end{equation}
Indeed, the square of the prime part $\sum_{p\le X^2}W(\log p/\log X)p^{-1/2-it}$ is $\sum_{n\le X^4}c_nn^{-it}$ with $c_{pq}=2w_pw_q/\sqrt{pq}$ for $p<q$ and $c_{p^2}=w_p^2/p$ (here $w_p=W(\log p/\log X)$; the two ordered factorisations of $pq$, $p\ne q$, are already counted), so that $\sum_n|c_n|^2\le2\bigl(\sum_{p\le X^2}p^{-1}\bigr)^2\ll(1+\log\log(3X))^2$ and $\sum_nn|c_n|^2\le2\pi(X^2)^2\le X^4=T^{o(1)}$, so Montgomery--Vaughan gives the bound for this part. The contributions of the prime squares and higher powers are handled as in \eqref{eq:B4} (Montgomery--Vaughan applied to the square of the prime-square sum, with remainder $O(\pi(X)^2)=T^{o(1)}$) and by the absolute bound for $r\ge3$; for the pole term the $L^2$ bound is not enough, and one uses the pointwise bound of \cref{lem:pole}: $\int_T^{2T}|\Pi_X^W|^4\,dt\ll T\bigl((1+X)/T\bigr)^4=(1+X)^4/T^3=o(T)$ for $X=T^{o(1)}$. For every $a>0$, $|x-\clip_a(x)|^2\le |x|^4/a^2$; therefore $H(t)\gg\log T$ on $[T,2T]$ implies
\begin{equation}
\int_T^{2T}|P_X^W-\clip_H(P_X^W)|^2\dd
\ll T\frac{(1+\log\log(3X))^2}{(\log T)^2}=o(T).
\end{equation}
Thus \eqref{eq:preclipL2} remains valid with $N_0+P_X^W$ replaced by $F_X^W$. Finally, H\"older's inequality gives
\begin{equation}
\int_T^{2T}|F_X^W-N|^4\dd
\ge T^{-1}\left(\int_T^{2T}|F_X^W-N|^2\dd\right)^2,
\end{equation}
which proves \eqref{eq:smallX}.
\end{proof}

\begin{remark}[Meaning of the lower restriction]\label{rem:necessary}
A universal $O(T)$ fourth-moment estimate valid for \emph{every} diverging cutoff is therefore impossible. The theorem forces $\liminf_{T\to\infty}\log X(T)/\log T>0$ along any family satisfying a uniform $O(T)$ bound. It does not assert that $X\ge T^\varepsilon$ is sufficient when there is no upper bound on $X$.
\end{remark}

\section{Uniform fourth moment in a polynomial cutoff range}\label{sec:moderate}
\begin{theorem}[Uniform weighted fourth moment]\label{thm:moderate}
Fix $0<\varepsilon<\frac14$. Subject to the Fujii input \eqref{eq:FujiiA} and to \eqref{eq:MCOMH}, there is $C_\varepsilon>0$, independent of $T$, $X$, and $W$, such that for $T$ sufficiently large,
\begin{equation}\label{eq:moderate}
\boxed{\displaystyle
\sup_{\substack{W\in\mathcal W\\ T^\varepsilon\le X\le T^{1/4}}}
\int_T^{2T}|F_X^W(t)-N(t)|^4\dd\le C_\varepsilon T.}
\end{equation}
No RH assumption is used.
\end{theorem}
\begin{proof}
Fix $\delta=\min(\varepsilon/2,1/200)>0$ and set $Y=T^\delta$. By \eqref{eq:FujiiA} with $\mathsf k=2$, $\alpha_{\rm F}=40\delta\le\frac15$ and Fujii's interval length equal to $T$ (so that $Y=T^{\alpha_{\rm F}/(20\mathsf k)}\le T^{1/2}$),
\begin{equation}\label{eq:Fujii4}
\int_T^{2T}\left|S(t)-\frac1\pi\operatorname{Im}\sum_{p<Y}p^{-1/2-it}\right|^4\dd\ll_\varepsilon T.
\end{equation}
Because $Y<X$, the difference between the prime terms of $\Phi_X^W$ and the sum in \eqref{eq:Fujii4} is
\begin{equation}
A_X(t)=\sum_{Y\le p\le X^2}\frac{w_p}{\sqrt p}p^{-it},
\qquad w_p=W(\log p/\log X)\in[0,1].
\end{equation}
Mertens' theorem and $T^\varepsilon\le X\le T^{1/4}$ give
\begin{equation}\label{eq:variance}
V_X:=\sum_{Y\le p\le X^2}\frac{w_p^2}p
\ll\log\log(X^2)-\log\log Y+O(1)\ll_\varepsilon1.
\end{equation}
Write $A_X^2=\sum_{n\le X^4}b_n n^{-it}$. If $d(n)$ counts ordered decompositions $n=pq$ with primes $Y\le p,q\le X^2$, then $d(n)\le2$ and
\begin{equation}
\sum_n|b_n|^2\le2V_X^2\ll_\varepsilon1,
\quad
\sum_n n|b_n|^2\le2\pi(X^2)^2\ll X^4\le T.
\end{equation}
Consequently the Montgomery--Vaughan mean-value theorem applied to $A_X^2$ yields
\begin{equation}\label{eq:A4}
\int_T^{2T}|A_X(t)|^4\dd\ll_\varepsilon T.
\end{equation}

The prime-square terms are, up to the coefficient $1/(2\pi)$,
\begin{equation}
B_X(t)=\sum_{p\le X}\frac{W(2\log p/\log X)}p\,p^{-2it}.
\end{equation}
Applying the same argument to $B_X^2$, now with frequencies $(pq)^2$, gives
\begin{equation}\label{eq:B4}
\int_T^{2T}|B_X(t)|^4\dd\ll T;
\end{equation}
indeed $\sum_p p^{-2}<\infty$ controls the diagonal and the Montgomery--Vaughan remainder is $O(\pi(X)^2)=O(X^2)$. The sum of the prime-power terms with exponent $r\ge3$ is bounded absolutely independently of $X$. By \cref{lem:pole}, $|\Pi_X^W(t)|\ll X/T\ll T^{-3/4}$ on $[T,2T]$, and $R(t)=O(T^{-1})$ by \eqref{eq:RVM}. Combining these estimates with \eqref{eq:Fujii4} by Minkowski,
\begin{equation}\label{eq:untruncatedmoderate}
\int_T^{2T}|N_0(t)+P_X^W(t)-N(t)|^4\dd\ll_\varepsilon T.
\end{equation}
Finally, $|G(t)|\le H(t)$ by \eqref{eq:MCOMH}. The contractivity in \cref{prop:contract}, with $x=P_X^W(t)$ and $y=G(t)$, transfers \eqref{eq:untruncatedmoderate} directly to \eqref{eq:moderate}.
\end{proof}

\begin{remark}[The precise upper restriction]\label{rem:quarter}
The exponent $1/4$ comes from the support $n\le X^2$ in the weighted prime polynomial and the support $n\le X^4$ after squaring it. Montgomery--Vaughan controls the remainder by $O(X^4)$. This is a limitation of this \emph{particular fourth-moment proof}; it is not a theorem that a larger $X$ is impossible or incompatible with RH. In particular, \cref{thm:moderate} does not establish \eqref{eq:desired}.
\end{remark}

\begin{theorem}[Logarithmic improvement for the two fixed weights]\label{thm:moderate-log}
Fix $0<\varepsilon<1/4$ and choose either the Selberg weight $q(v)=1$
or $W_{1/4}$ with $q(v)=v/2+1/4$. Under the same Fujii input as
\cref{thm:moderate}, there exist $C_{\varepsilon,W}$ and $T_0(\varepsilon,W)$ such that, for $T\ge T_0$,
\begin{equation}\label{eq:moderate-log}
\sup_{T^\varepsilon\le X\le T^{1/4}(\log T)^{3/2}}
\int_T^{2T}|F_X^W(t)-N(t)|^4\dd\le C_{\varepsilon,W}T.
\end{equation}
This is an extension in the \emph{weighted support parameter} $X$,
not a claim for every $W\in\mathcal W$.
\end{theorem}
\begin{proof}
For the indicated weights, $0\le q(v)\le B_W$ and
$W(v)\le B_W(2-v)$ for $1\le v\le2$. Thus
\begin{equation}\label{eq:wprime2}
\sum_{p\le X^2}W(\log p/\log X)^2
\ll_W\frac{X}{\log X}+\frac{X^2}{(\log X)^3}
\ll_W\frac{X^2}{(\log X)^3}.
\end{equation}
Indeed, for $p\le X$, use the prime-counting bound;
for $X<p\le X^2$, split into strips
$X^2e^{-(r+1)}<p\le X^2e^{-r}$ ($r\ge0$),
use $W\le B_W(r+1)/\log X$ and
$\pi(y)\ll y/\log y$ when $y\ge X$,
then sum the absolutely convergent series
$\sum_{r\ge0}e^{-r}(r+1)^2$.
For $A_X$ as in the proof of \cref{thm:moderate}, the bound
$\sum_n|b_n|^2\ll_\varepsilon1$ is unchanged, while
\begin{equation}
\sum_n n|b_n|^2
\le2\left(\sum_{Y\le p\le X^2} w_p^2\right)^2
\ll_W\frac{X^4}{(\log X)^6}\ll_\varepsilon T
\end{equation}
throughout the enlarged range: $X\le T^{1/4}(\log T)^{3/2}$ gives $X^4\le T(\log T)^6$, and $X\ge T^\varepsilon$ gives $\log X\ge\varepsilon\log T$, so $X^4/(\log X)^6\le\varepsilon^{-6}T$. The treatment of prime squares,
higher prime powers and the pole compensator remains valid:
$X^2\ll T$ and $X/T\ll1$. Apply Montgomery--Vaughan and the
projection contraction as before.\end{proof}

\begin{proposition}[Small proportion of saturated values]\label{prop:sat-measure}
Under the assumptions of \cref{thm:moderate-log} and for $T\ge T_0(\varepsilon,W)$, let
$\mathcal E_{T,X}=\{t\in[T,2T]:|P_X^W(t)|\ge H(t)\}$.
Then
\begin{equation}\label{eq:sat-measure}
\sup_{T^\varepsilon\le X\le T^{1/4}(\log T)^{3/2}}
\meas(\mathcal E_{T,X})\ll_{\varepsilon,W}\frac{T}{(\log T)^4}.
\end{equation}
\end{proposition}
\begin{proof}
Split $\mathcal E_{T,X}$ according to $|G(t)|\le H(t)/2$ and its
complement, where $G=N-N_0$. In the first case $|F_X^W-N|\ge H/2$;
Markov's inequality and \eqref{eq:moderate-log} give
$O_{\varepsilon,W}(T/(\log T)^4)$. In the second case,
$|S(t)|\gg\log T$ for large $T$ because $R(t)=O(T^{-1})$.
The eighth-moment estimate of Selberg \cite{Selberg1946,Titchmarsh} gives
$O(T(\log\log T)^4/(\log T)^8)$ for the second set.
\end{proof}
\begin{remark}
The last proposition concerns a jointly growing $X$ and $T$ in a specified
range. It provides no information about $X\to\infty$ with $T$ held fixed.
\end{remark}

\section{A zero-sum fourth-moment estimate without an upper cutoff}\label{sec:zero-sum}
A distinct estimate gives supporting evidence for considering $X\ge T^\varepsilon$ without an upper restriction, but is not an estimate for the full arithmetic phase.

\begin{theorem}[Unconditional smoothed zero-sum bound]\label{thm:ZL}
For $L>0$ let
\begin{equation}\label{eq:ZL}
\mathcal Z_L(t)=\sum_{\gamma_k>0}\frac1{(1+L|t-\gamma_k|)^2},
\end{equation}
where the imaginary parts are counted with multiplicity. Subject to the Fujii increment input \eqref{eq:FujiiB}, for each fixed $\varepsilon>0$ and all $T\ge T_0(\varepsilon)$,
\begin{equation}\label{eq:ZL4}
\boxed{\displaystyle
\sup_{L\ge\varepsilon\log T}\int_T^{2T}\mathcal Z_L(t)^4\dd
\ll_\varepsilon T.}
\end{equation}
No RH hypothesis is required.
\end{theorem}
\begin{proof}
For $v\in\mathbb R$,
\begin{equation}\label{eq:kernelintegral}
(1+L|v|)^{-2}=\int_{|v|}^\infty\frac{2L}{(1+Lu)^3}\,du.
\end{equation}
Define $N(y)=0$ for $y\le0$. By Tonelli and \eqref{eq:kernelintegral},
\begin{equation}\label{eq:Zrepresentation}
\mathcal Z_L(t)=\int_0^\infty\frac{2L}{(1+Lu)^3}
\bigl(N(t+u)-N(t-u)\bigr)\,du
\end{equation}
for almost every $t$. On $0<u\le1$, \eqref{eq:FujiiB} with $\mathsf k=2$, $\alpha_{\rm F}=\frac12$ and Fujii's interval length equal to $T$, together with \eqref{eq:RVM}, implies
\begin{equation}\label{eq:inc4}
\|N(t+u)-N(t-u)\|_{L^4[T,2T]}
\ll T^{1/4}(1+u\log T).
\end{equation}
To see this, split the difference into the two increments of length $u$, shift the integration variable for one of them, and use the bound $\int |S(t+u)-S(t)|^4\dd\ll T(\log(2+u\log T))^2$; the main term changes by $O(u\log T)$, and $\sqrt{\log(2+v)}\ll1+v$.
For $u>1$, the Riemann--von Mangoldt formula and $S(y)=O(\log(2+y))$ yield the elementary pointwise bound
\begin{equation}\label{eq:inc-tail}
N(t+u)-N(t-u)\ll(u+1)\log(2T+u),\qquad T\le t\le2T.
\end{equation}
Apply Minkowski to \eqref{eq:Zrepresentation}, using \eqref{eq:inc4} on $(0,1]$ and \eqref{eq:inc-tail} on $(1,\infty)$:
\begin{equation}
\|\mathcal Z_L\|_4\ll T^{1/4}
\left(1+\frac{\log T}{L}\right)
+T^{1/4}\,O\!\left(\frac{1+\log T}{L^2}\right).
\end{equation}
Here $\int_0^\infty 2L(1+Lu)^{-3}\,du=1$ and
$\int_0^\infty 2Lu(1+Lu)^{-3}\,du=1/L$; the tail integral is $O((1+\log T)/L^2)$ for $L\ge1$.
When $L\ge\varepsilon\log T$ the right-hand side is $O_\varepsilon(T^{1/4})$, proving \eqref{eq:ZL4}.
\end{proof}

\begin{remark}[No hidden assumption about $\beta_k$]\label{rem:ZL}
The sum in \eqref{eq:ZL} uses only $\gamma_k$ and their multiplicities. Thus \eqref{eq:ZL4} is compatible both with RH and with the existence of zeros outside $\Re s=1/2$. To transfer it to the actual weighted Euler phase under RH, one must \emph{first prove} an explicit formula $N_0+P_X^W-N=\sum_\rho K_W(\rho,t;X)+(\text{remainder})$ with a kernel $K_W$ (for the Selberg weight, $K_W(\rho,t;X)=\pi^{-1}\Im\,\mathfrak I_L(\rho-\tfrac12-it)$ with $L=\log X$) satisfying
$|K_W(\frac12+i\gamma,t;X)|\ll_W(1+|t-\gamma|\log X)^{-2}$ and a suitably controlled remainder. That transfer is not contained in \cref{thm:ZL}. For the Selberg weight, a formula of this type with the kernel bounds of \cref{lem:taucw-kernel}(b) is \cref{lem:taucw-expl}.
\end{remark}

\section{Off-critical zeros, weight smoothness, and the unresolved range}\label{sec:open}
For $W$ as in \eqref{eq:weights}, the leading contribution of a hypothetical zero with $\rho=\frac12+d+i\gamma$ ($d>0$), in the proposed explicit-formula analysis of \cite{CasanaDBSRH}, involves a Laplace integral of the form
\begin{equation}\label{eq:laplace}
\mathcal L_q(L,z)=\int_1^2q(v)e^{vLz}\,dv,
\qquad z=d+i(\gamma-t),\quad L=\log X.
\end{equation}
For a sufficiently regular $q$ with $q(2)\ne0$, integration by parts gives, at fixed $z\ne0$ with $\Re z>0$,
\begin{equation}\label{eq:Lap-asympt}
\mathcal L_q(L,z)=\frac{q(2)e^{2Lz}}{Lz}
+O_z\!\left(\frac{e^{2Ld}}{L^2}\right)
+O_z\!\left(\frac{e^{Ld}}L\right),
\end{equation}
where the displayed remainder assumes, for instance, that $q$ has a bounded derivative. If $q$ vanishes to finite order at the endpoint, repeated integration by parts introduces additional powers of $1/L$ but does not suppress the factor $e^{2Ld}$. The Selberg weight and $W_{1/4}$ have $q(2)>0$. These statements explain a \emph{possible} saturation mechanism, not a proved lower bound for the total zero sum: distinct off-critical contributions can cancel.

\paragraph{The precise remaining statement.}
For a fixed $\varepsilon>0$ and a fixed admissible weight $W$, the candidate estimate is
\begin{equation}\label{eq:fullgoal}
\boxed{\displaystyle
\sup_{X\ge T^\varepsilon}
\frac1T\int_T^{2T}|F_X^W(t)-N(t)|^4\dd\le C_{\varepsilon,W}
\qquad(T\ge T_0).}
\end{equation}
The lower restriction is consistent with the unconditional obstruction in \cref{thm:smallX}; there is deliberately no upper restriction. However \eqref{eq:fullgoal} has \emph{not} been proved. It cannot follow just from the deterministic truncation inequality, because the latter bounds the truncated error by the \emph{untruncated} arithmetic error, whose fourth moment is controlled here for the cutoffs $X\le T^{1/4}$ (\cref{thm:moderate}), and $X\le T^{1/4}(\log T)^{3/2}$ for the two fixed weights (\cref{thm:moderate-log}); for the Selberg weight and every $X\ge T^\varepsilon$ only the on-line part of $N_0+P_X^W-N$ is controlled (\cref{thm:taucw-online}).

\paragraph{What remains to be supplied.}
\begin{enumerate}[leftmargin=*,itemsep=4pt]
\item \textbf{A fully justified weighted explicit formula.} The formula relating $P_X^W-S$ to contributions of zeros, the pole and any remaining terms must be established with uniformity in $T$ and in arbitrarily large $X$, including the treatment of jumps and of $t$ near $\gamma_k$. The corresponding proposition in \cite{CasanaDBSRH} is presently a sketch. For the Selberg weight a formula of this kind, for $t$ not the imaginary part of a zero and uniform in $X\ge2$, is \cref{lem:taucw-expl}; the other weights of \eqref{eq:weights} are not covered.
\item \textbf{An arithmetic fourth-moment estimate beyond $X\le T^{1/4}$ (resp.\ $T^{1/4}(\log T)^{3/2}$ for the two fixed weights).} The bound \eqref{eq:fullgoal} requires cancellation among prime powers beyond what the direct Montgomery--Vaughan argument supplies. A fourth-moment estimate on long intervals need not be replaced by a uniform mean square on windows of length one: the latter is sufficient in some approaches but is not logically necessary.
\item \textbf{Non-cancellation for off-critical zeros.} The abstract model in \cref{thm:structural} cannot establish the behavior of the actual $P_X^W$. A genuine implication towards RH needs a proved saturation or alternative obstruction when $\beta_k>1/2$, with no unjustified hypothesis that a finite set of zeros achieves an isolated maximal real part.
\item \textbf{Separate cell containment from convergence.} A sequence of strictly increasing $\tau_k^{\Cms}(X;W)$ can converge as $X\to\infty$ to a nondecreasing sequence with equal consecutive values. Thus \eqref{eq:DBSRH} alone does not directly establish RH or simplicity by an argument based solely on strict increase at finite $X$.
\end{enumerate}

\section{Exact sublevel identity and the stronger cell condition}\label{sec:cells}
The following identity is valid for every continuous $F$ with finite sublevel measures (so that $u_k(F)$ of \eqref{eq:sublevel} is defined), irrespective of RH:
\begin{equation}\label{eq:area}
u_k(F)-\gamma_k
=\int_0^\infty\bigl(\ind_{\{F(t)<k-\frac12\}}-
\ind_{\{N(t)<k-\frac12\}}\bigr)\dd.
\end{equation}
Indeed $\{N<k-\frac12\}$ equals $(0,\gamma_k)$ up to a set of measure zero. For $F=F_X^W$, write the right-hand side as
\begin{align}
\mathcal A_k^+(X)&=\meas\{t:F_X^W(t)<k-\tfrac12\le N(t)\},\label{eq:Aplus}\\
\mathcal A_k^-(X)&=\meas\{t:N(t)<k-\tfrac12\le F_X^W(t)\},\label{eq:Aminus}
\end{align}
so that $\tau_k^{\Cms}(X;W)-\gamma_k=\mathcal A_k^+(X)-\mathcal A_k^-(X)$.
A global $L^4$ bound concerns the \emph{size} of the arithmetic discrepancy, not uniform cancellation between the two signed sublevel contributions for every $k$.

For any strictly increasing sequence $u_k$, put $\bar u_k=(u_k+u_{k+1})/2$.
\begin{proposition}[A sufficient cell condition for RH and simplicity]\label{prop:RH}
If $\gamma_1<\bar u_1$ and
\begin{equation}\label{eq:cell}
\bar u_{k-1}<\gamma_k<\bar u_k
\qquad(k\ge2),
\end{equation}
then every nontrivial zero has real part $\frac12$ and is simple.
\end{proposition}
\begin{proof}
The interval $(-\infty,\bar u_1)$, which contains $\gamma_1$, and the intervals in \eqref{eq:cell} are pairwise disjoint. Consequently consecutive $\gamma_k$ cannot be equal. A zero of multiplicity greater than one would repeat its imaginary part. A zero $\rho=\beta+i\gamma$ with $\beta\ne\frac12$ would have the distinct reflected zero $1-\beta+i\gamma$, again repeating the imaginary part. Both cases are excluded.
\end{proof}
\begin{remark}
The weaker condition $\gamma_{k-1}\le u_k<\gamma_{k+1}$ does not imply RH or simplicity; it can accommodate equal consecutive $\gamma_k$. The proposition is a sufficient criterion, \emph{not} a claim that the required cell containment has been proved for the proposed taucm(s).
\end{remark}

\section{Finite computations recorded in the associated research files}\label{sec:data}
For context, the author's research notes \cite{CasanaNotes} for the branch
\texttt{kc-tauC-2026-11-02} of the repository \texttt{vcasana1/prueba}
record the following root-mean-square discrepancies, normalized by the mean spacing between successive $\gamma_k$, over the first $2{,}001{,}052$ zeros. These are \emph{reported measurements}, not computations independently repeated or certified for this note; they do not imply a uniform-in-$k$ result.
\begin{center}
\begin{tabular}{@{}p{0.68\linewidth}r@{}}
\toprule
Approximation scheme & Reported RMS \\
\midrule
Main term $N_0$ alone & $0.332$\\
Original IJNT $\tau_k^{\mathrm C}$ & $0.0557$\\
Sharp prime cutoff at $p\le t^{1.4}$ & $0.01456$\\
Selberg weight with $X=t^{0.7}$ (support to $X^2=t^{1.4}$) & $0.01456$\\
$W_{1/4}$ with $X=t^{0.7}$ (same support) & $\mathbf{0.01376}$\\
\bottomrule
\end{tabular}
\end{center}
The notes also report that the $W_{1/4}$ weight reduces RMS by approximately $5.5\%$ relative to the sharp cutoff at the same support, but performs less well for the two exceptionally close pairs indexed $4765$--$4766$ and $6709$--$6710$. These measurements concern earlier prime-phase approximations; no new numerical estimate for $\tau_k^{\Cm}$ under the MCOM truncation bound is claimed. Importantly, $X=t^{0.7}$ lies \emph{outside} the rigorously uniform range $X\le T^{1/4}$ proved in \cref{thm:moderate}. This distinction prevents a numerical experiment from being mistaken for an extension of that theorem.


\section{Separation from the Riemann--von Mangoldt points}\label{sec:spacing-v3}
Let $\gamma_k^{\mathrm{RvM}}$ be the unique solutions of
$N_0(\gamma_k^{\mathrm{RvM}})=k-\tfrac12$ with
$\gamma_k^{\mathrm{RvM}}>2\pi$. Since $N_0'(t)=(2\pi)^{-1}\log(t/2\pi)$,
\begin{equation}\label{eq:rvm-separation-v3}
\gamma_{k+1}^{\mathrm{RvM}}-\gamma_k^{\mathrm{RvM}}
=\frac{1}{N_0'(\xi_k)}>0
\end{equation}
for some $\xi_k$ between the two points, and this difference is
$\sim2\pi/\log\gamma_k^{\mathrm{RvM}}$ as $k\to\infty$.
For the weighted constructions, put
\begin{equation}\label{eq:gap-v3}
\Delta_k(X;W)=\tau_{k+1}^{\Cms}(X;W)-\tau_k^{\Cms}(X;W)
=\meas\{t>0:k-\tfrac12\le F_X^W(t)<k+\tfrac12\}.
\end{equation}
The identity is immediate from the definition of sublevel measure.
A continuous function attaining both levels takes, by the intermediate value theorem, the values of $(k-\frac12,k+\frac12)$ on a nonempty open set, so $\Delta_k(X;W)>0$;
for sufficiently large $k$ this follows from $F_X^W=N_0+O(H)$ and
$N_0-H\to\infty$, uniformly in $X$. For a small index, the existence
of a point below the lower level must be checked: continuity and the
upper bound alone do not guarantee strict positivity for every $X$.
Even if all finite-cutoff spacings are positive,
$\lim_{X\to\infty}\Delta_k(X;W)$ may be zero.

\begin{theorem}[Uniform separation of sufficiently distant indices]
\label{thm:block-sep-v3}
For $k$ sufficiently large, let $a_k,b_k$ be the unique solutions in
the eventual increasing region of $N_0(t)+H(t)=k-\tfrac12$ and
$N_0(t)-H(t)=k-\tfrac12$, respectively. Then
\begin{equation}\label{eq:barriers-v3}
a_k\le\tau_k^{\Cms}(X;W)\le b_k\qquad(X\ge2).
\end{equation}
For every positive integer $m$ satisfying
$m>2H(b_k)$,
\begin{equation}\label{eq:blocksep-v3}
\inf_{X\ge2}\bigl[\tau_{k+m}^{\Cms}(X;W)-
\tau_k^{\Cms}(X;W)\bigr]\ge a_{k+m}-b_k>0.
\end{equation}
The same conclusion applies to taucm.
\end{theorem}
\begin{proof}
The inequalities $N_0-H\le F_X^W\le N_0+H$ and the eventual
monotonicity of both barriers show that the sublevel measure is at least
$a_k$ and at most $b_k$, once the values of $N_0\pm H$ on the initial range where they are not monotone cannot
reach the levels in question. Further,
$N_0(b_k)+H(b_k)=k-\tfrac12+2H(b_k)<k+m-\tfrac12$.
Strict increase of $N_0+H$ implies $b_k<a_{k+m}$; subtract
\eqref{eq:barriers-v3} at the two indices.\end{proof}
\begin{remark}[Asymptotic MCOM improvement]
For large $k$ the sufficient index separation is
$m>(0.16764+o(1))\log\gamma_k^{\mathrm{RvM}}+O(1)$.
Using a non-explicit large-height truncation bound obtained from the
MCOM asymptotic coefficient, the coefficient can be replaced by
$2C_{\mathrm H}+o(1)=0.1497192285421\ldots+o(1)$, where
$C_{\mathrm H}=0.0748596142710907\ldots$.
Neither statement covers adjacent indices $m=1$ for all large $k$.
\end{remark}

\begin{proposition}[Two local arithmetic occupation tests]
\label{prop:occupation-v3}
For $k$ sufficiently large put
\begin{align}
I_k^+&=\{t>0:k-\tfrac12<N_0(t)+H(t)<k+\tfrac12\},\\
I_k^-&=\{t>0:k-\tfrac12<N_0(t)-H(t)<k+\tfrac12\},\\
B_k^+(X;W)&=\meas\{t\in I_k^+:P_X^W(t)\ge H(t)\},\\
B_k^-(X;W)&=\meas\{t\in I_k^-:P_X^W(t)\le-H(t)\}.
\end{align}
Then $I_k^+\cap I_k^-=\varnothing$ for large $k$ and
\begin{equation}\label{eq:occup-v3}
\boxed{\quad\Delta_k(X;W)\ge B_k^+(X;W)+B_k^-(X;W).\quad}
\end{equation}
Consequently, $\liminf_{X\to\infty}(B_k^++B_k^-)>0$
is a sufficient arithmetic condition for noncollapse at index $k$, that is, for $\liminf_{X\to\infty}\Delta_k(X;W)>0$ (cf.\ \eqref{eq:noncollapse-v5}).
\end{proposition}
\begin{proof}
On $I_k^+$, the condition $P_X^W\ge H$ makes
$F_X^W=N_0+H\in(k-\tfrac12,k+\tfrac12)$; the negative case is
analogous. Both sets therefore contribute to \eqref{eq:gap-v3}.
Since $2H>1$ in the large-height region, the two intervals are disjoint.
\end{proof}

\begin{remark}[Relevance and exact open issue]
The occupation inequality uses only $N_0$, the MCOM truncation and
Euler--Selberg phases. It does not establish that either $B_k^+$ or
$B_k^-$ stays positive as $X\to\infty$. The structural saturation
model of \cref{thm:structural} \emph{does} have nonzero limiting spacings;
a single increasingly sharp passage from $N_0-H$ to $N_0+H$ can, in
contrast, make one spacing vanish. Therefore the separation of
$\gamma_k^{\mathrm{RvM}}$ cannot alone be transferred to the $\tau_k^{\Cms}$.
If DBSRH$(W)$ and $\limsup_{X\to\infty}\Delta_k(X;W)>0$
hold for every $k$, then $\gamma_{k+1}>\gamma_k$ for all $k$ and the
functional equation implies RH and simplicity. This implication is
conditional on \emph{both} hypotheses, not a proof of either.
\end{remark}

\section{The exact index defect at zero heights}\label{sec:uk-v3}
To prevent a notation collision with the sublevel measure $u_k(F)$
from \eqref{eq:sublevel}, denote the \emph{evaluation defect} by
\begin{equation}\label{eq:ek-v3}
e_k(F)=k-\tfrac12-F(\gamma_k).
\end{equation}
For integer $j\ge1$ define
\begin{equation}\label{eq:Psi-v3}
\Psi_j(z)=\int_{z-1/2}^{z+1/2}v^{2j}\,dv
=\frac{(z+1/2)^{2j+1}-(z-1/2)^{2j+1}}{2j+1}.
\end{equation}
In particular,
\begin{equation}\label{eq:psi1-v3}
\Psi_1(z)=z^2+\frac1{12}.
\end{equation}
The IJNT exact-defect identity relates an integral of $(N-F)^{2j}F'$
to the sum of $\Psi_j(e_k(F))$, with explicit terms at the two
integration endpoints. When applying that identity, the endpoint
convention and the factor $F'(t)$ must not be omitted, especially if
$F$ is nonmonotone. The exact form, with its endpoint terms, is
reproduced in \eqref{eq:IJNTidentity} at the end of this section.

\begin{proposition}[Repeated imaginary parts force an evaluation defect]
\label{prop:duplicate-v3}
Suppose $F$ is a single-valued real function at all the $\gamma_k$.
If $\gamma_{k+1}=\gamma_k$, then
\begin{equation}\label{eq:dup-diff-v3}
e_{k+1}(F)-e_k(F)=1,
\qquad e_k(F)^2+e_{k+1}(F)^2\ge\frac12.
\end{equation}
More generally, among any finite group of $m\ge2$ equal imaginary
parts, $\sum e_k(F)^2\ge m(m^2-1)/12$.
In any consecutive index window, let $M$ count the indices lying in
fully contained repeated-height groups. Then
\begin{equation}\label{eq:duplicate-count-v3}
M\le4\sum_{k\text{ in window}}e_k(F)^2.
\end{equation}
In particular, if $|e_k(F)|<\tfrac12$ for each index of the window,
no two \emph{consecutive indices inside the window} can have the same
imaginary part.
\end{proposition}
\begin{proof}
At a shared imaginary part, the value of $F$ is unchanged, so
$e_{k+1}=e_k+1$. The minimum of $z^2+(z+1)^2$ is $1/2$.
For $m$ successive values $z,z+1,\ldots,z+m-1$, the sum of squares is
minimized at $z=-(m-1)/2$, giving $m(m^2-1)/12$.
Since $m\le4m(m^2-1)/12$ for $m\ge2$, summing over distinct
repeated-height groups proves \eqref{eq:duplicate-count-v3}.
The strict cell statement follows immediately from $e_{k+1}-e_k=1$.
\end{proof}

\begin{remark}[Finite-window and high-moment interpretation]
For a group of two coincident imaginary parts,
$\Psi_j(e_k)+\Psi_j(e_{k+1})\ge2/(2j+1)$, whereas the two undisplaced
terms, each equal to $\Psi_j(0)=4^{-j}/(2j+1)$, contribute $2\cdot4^{-j}/(2j+1)$. Their excess is at least
$2(1-4^{-j})/(2j+1)$. This argument rigorously applies to
\emph{disjoint} coincident pairs. Counting overlapping pairs with the
same coefficient requires a separate combinatorial proof.
For a finite window, strict cell inequalities $|e_k(F)|<1/2$ imply
absence of internal coincidences. They do not imply the inequalities
for all zero heights, nor do high-moment estimates at a fixed prime
budget settle the all-height problem.
\end{remark}

\begin{proposition}[Forward evaluation: the proved implication]
\label{prop:forward-v3}
Let $\mathscr F_X(t)$ be any family of real-valued functions, and suppose
\begin{equation}\label{eq:forward-v3}
\lim_{X\to\infty}\mathscr F_X(\gamma_k)=k-\tfrac12
\qquad\text{for every }k\ge1.
\end{equation}
Then RH holds and every nontrivial zero is simple.
\end{proposition}
\begin{proof}
If two consecutive zero indices share the same imaginary part,
$\mathscr F_X$ has the same value at both indices for every $X$,
contradicting the two different required limits. Thus all imaginary
parts are distinct, which excludes multiple zeros. If $\rho=\beta+i\gamma$ were a nontrivial zero with $\beta\ne\frac12$, the functional equation together with conjugation would give the distinct zero $1-\bar\rho=1-\beta+i\gamma$ with the same imaginary part, so two consecutive indices would have equal imaginary parts, which has just been excluded. Hence RH holds and all zeros are simple.
\end{proof}
\begin{remark}
For the proposed family
$\mathscr F_X=N_\vartheta+P_X^W$, with
$N_\vartheta(t)=\vartheta(t)/\pi+1$, the converse implication
(RH and simplicity $\Rightarrow$ \eqref{eq:forward-v3}) requires a
rigorous pointwise weighted explicit formula evaluated at the zeros.
It is not proved here and must not be replaced by an unverified
midpoint-kernel argument. The same caution applies to the suggested
uncorrected evaluations at a hypothetical off-critical zero.
\end{remark}

\paragraph{The exact defect identity with its endpoint terms.}
For indices $1\le k_1<k_2$ and a real absolutely continuous $F$ on a
compact interval containing $[\gamma_{k_1},\gamma_{k_2}]$, the IJNT
manuscript \cite{CasanaIJNT} (Theorem ``Exact defect identity'', label
\texttt{thm:identity}) gives, for every integer $j\ge1$ and with
$e_k=e_k(F)$,
\begin{equation}\label{eq:IJNTidentity}
\int_{\gamma_{k_1}}^{\gamma_{k_2}}\bigl(N(t)-F(t)\bigr)^{2j}F'(t)\,dt
=\sum_{k=k_1+1}^{k_2-1}\Psi_j(e_k)
+\frac{(\frac12+e_{k_1})^{2j+1}+(\frac12-e_{k_2})^{2j+1}}{2j+1}.
\end{equation}
No monotonicity of $F$ is required, and repeated imaginary parts are
admitted: a degenerate gap $\gamma_k=\gamma_{k+1}$ contributes zero to
both sides.

\section{Numerical evidence from the 2,001,052-zero dataset}\label{sec:v9-data}
The following results are \emph{reported numerical computations},
corroborated here against the text outputs in the repository directory
\texttt{KC\_tauC\_2026-11-02/v8\_identidad\_en\_ordenadas}.
The binary arrays were not independently recomputed and the
NUFFT/\texttt{float64} calculations have not been certified with
interval arithmetic. In particular ``passes the cell test numerically''
does not mean a new rigorous verification of RH.
The computation uses the evaluation defects \eqref{eq:ek-v3}
for three specific prime-phase rules, with
$X=t^\theta$ for the weighted rules. The numerical variable
$\theta$ is a prime-support exponent; it is not the
Riemann--Siegel function $\vartheta$.

\begin{center}\scriptsize
\begin{tabular}{@{}p{0.30\linewidth}rrrr@{}}
\toprule
Rule & RMS $e_k$ & $\eta$ on $(2^{19},2^{20}]$ & $1-4\eta$ & $|e_k|\ge1/2$\\
\midrule
IJNT original & $0.0910$ & $0.00840$ & $0.9664$ & $220$\\
$W_{1/4}$, $\theta=0.5$ & $0.0569$ & $0.00330$ & $0.9868$ & $3$\\
$W_{1/4}$, $\theta=0.7$ & $\mathbf{0.0381}$ & $0.00149$ & $\mathbf{0.9941}$ & $0$\\
\bottomrule
\end{tabular}
\end{center}
Here $\eta$ is the average of $e_k(F)^2$ in the indicated dyadic
height block; $1-4\eta$ is the numerical value of the lower bound
that \cref{prop:duplicate-v3} would supply for the fraction of
indices whose imaginary part $\gamma_k$ is not shared with another index, if the $e_k$ values were certified.
These values are \emph{not} unconditional asymptotic proportions.

\begin{center}\small
\begin{tabular}{@{}p{0.27\linewidth}p{0.22\linewidth}p{0.19\linewidth}p{0.17\linewidth}@{}}
\toprule
Rule & Passing / 2001 & median/max $j^*$ & minimum $\delta_k$\\
\midrule
IJNT original & $1790$ (89.46\%) & $7/23$ & negative (in 211 windows)\\
$W_{1/4}$, $\theta=0.5$ & $1998$ (99.85\%) & $5/10$ & negative (in 3 windows)\\
$W_{1/4}$, $\theta=0.7$ & $2001$ (100\%) & $4/7$ & $0.0117$\\
\bottomrule
\end{tabular}
\end{center}
Windows comprise 1000 consecutive indices each. The numerical script
checks $\max|e_k|<1/2$ before computing a sufficient high-moment
threshold $j^*$ for each window. The final 52 indices of the dataset
are not part of these 2001 complete windows; in the $\theta=0.7$
run the maximum $|e_k|$ across the \emph{full} dataset remains below $1/2$.
The label ``not passing'' does not assert that the high-moment
inequality fails for every $j$, only that this sufficient implementation
rejects the window when an evaluation lies outside the strict cell.

Let $g_k$ denote the nearest-neighbor separation normalized by the
local mean separation $2\pi/\log(\gamma_k/2\pi)$, and put
$\delta_k=1/2-|e_k(F)|$. For the 510 reported points with $g_k<0.05$,
least-squares fits are
\begin{equation}\label{eq:gfit-v3}
\delta_k\approx
\begin{cases}
1.49\,g_k-0.0002&\text{original IJNT},\\
1.79\,g_k-0.0003& W_{1/4},\ \theta=0.5,\\
2.50\,g_k+0.0001& W_{1/4},\ \theta=0.7.
\end{cases}
\end{equation}
The slopes exhibit the numerical ratio $2.50/1.79\approx1.40$,
matching $0.7/0.5=1.4$ to the reported precision; this is evidence for,
not a proof of, linear dependence on $\theta$. The average
$\int_1^2 v(v/2+1/4)\,dv=37/24\approx1.542$ gives the predicted slopes $2\theta\cdot\frac{37}{24}\approx1.54$ and $2.16$ for $\theta=0.5,0.7$, whereas the fitted slopes $1.79$ and $2.50$ are about $16\%$ larger; the heuristic leading dependence $2\theta\int_1^2 vq_{1/4}(v)\,dv$ (stated without derivation) thus accounts for the proportionality in $\theta$ but not for the full size, and the
neglected $O(1)$ terms and the other zeros have not been rigorously controlled.
For the closest reported pair ($k=1{,}115{,}578$ and $k+1=1{,}115{,}579$;
$g_k\approx0.005$), the two defects for $\theta=0.7$ are
$-0.4883$ and $+0.4846$.
The 1\% lower-fit intercepts are approximately
$-0.1395,-0.0642,-0.0246$ for the three rules.

The extrapolation $\theta_{\rm needed}(n)\approx0.21 n^{0.09}$
uses a GUE-inspired minimum-gap law together with two measured
choices of $\theta$; it is a heuristic only, not a theorem, and it
cannot justify applying mean-value estimates beyond their proved
length ranges. In particular, the weighted cutoff $X=t^{0.5}$ has
prime-power support $n\le t$, well beyond the $X\le T^{1/4}$
condition proved for \cref{thm:moderate}, and even the logarithmically
enlarged range of \cref{thm:moderate-log}. The original IJNT rule
retains a different prime budget: approximately $\sqrt t$ \emph{primes},
not weighted prime powers supported up to $t$.


\section{The \texorpdfstring{$j=1$}{j=1} conditional proportion budget: finite computations}\label{sec:jone-v3}
The repository outputs \cite{CasanaV9} record the exact-defect
identity at $j=1$ in a form which distinguishes two contributions that
must \emph{not} be conflated. Here $F=N_0+\Phi_c$, where $\Phi_c$ is the
specific truncated weighted phase used in the numerical script and
$D_F(t)=N(t)-F(t)$. Write
\begin{equation}\label{eq:jone-AB-v3}
\begin{split}
M_T&=N(2T)-N(T),\\
A_T(F)&=\int_T^{2T}D_F(t)^2 N_0'(t)\,dt,\\
B_T(F)&=\int_T^{2T}D_F(t)^2 \Phi_c'(t)\,dt.
\end{split}
\end{equation}
The derivative $\Phi_c'$ is the almost-everywhere derivative, including
the derivative of the moving truncation bound on its active portions.
In particular, \emph{it is not permissible to omit $B_T$}.
Define
$\eta_T(F)=M_T^{-1}\sum_{T<\gamma_k\le2T}e_k(F)^2$.
The exact IJNT identity gives, with its precise endpoint contribution
$\mathcal E_{\partial}(F;T)$,
\begin{equation}\label{eq:jone-identity-v3}
\eta_T(F)=\frac{A_T(F)+B_T(F)}{M_T}-\frac1{12}
+\mathcal E_{\partial}(F;T).
\end{equation}
Here $\mathcal E_{\partial}$, including its sign, is \emph{defined} by
\eqref{eq:jone-identity-v3}; the boundary terms in
\eqref{eq:IJNTidentity} are required if one claims a numerical upper bound for it.
The numerical comparison between the two sides is a consistency test,
not an independent analytic control of the endpoint contribution.

\begin{center}\small
\begin{tabular}{@{}rrrrr@{}}
\toprule
Dyadic $T$ & $A_T/M_T$ & $B_T/M_T$ & $A_T/M_T+B_T/M_T-1/12$
& measured $\eta_T$ \\
\midrule
$2^{10}$ & $0.04647$ & $+0.03883$ & $0.00196$ & $0.00200$\\
$2^{15}$ & $0.05064$ & $+0.03546$ & $0.00277$ & $0.00281$\\
$2^{19}$ & $0.05253$ & $+0.03407$ & $0.00327$ & $0.00330$\\
\bottomrule
\end{tabular}
\end{center}
The table is copied from \texttt{decomp\_w14\_th0.5.txt}. The phase
in these computations uses $W_{1/4}$ at support exponent $\theta=1/2$
and the original $\log\log t$ truncation bound: these numbers are
\emph{not} measurements for $H(t)$ of MCOM, nor does this table
establish a bound as $T\to\infty$. The displayed discrepancies are
at most $5\cdot10^{-5}$ at the reported precision.

\begin{theorem}[Conditional lower bound for the simple critical-line fraction]\label{thm:conditional-proportion-v3}
Let $F=N_0+\Phi_c$ be locally absolutely continuous, so that \eqref{eq:IJNTidentity} applies on $[\gamma_{k_1},\gamma_{k_2}]$ and the identity \eqref{eq:jone-identity-v3}, which defines $\mathcal E_\partial$, holds on dyadic intervals.
Let $T$ be sufficiently large (the threshold being absolute) and suppose that, for this $T$,
\begin{equation}\label{eq:jone-input-v3}
\begin{aligned}
T^{-1}\int_T^{2T}|D_F(t)|^2\,dt&\le a_T,&
T^{-1}\int_T^{2T}|D_F(t)|^4\,dt&\le b_T,\\
\frac{\bigl(T^{-1}\int_T^{2T}|\Phi_c'(t)|^2dt\bigr)^{1/2}}
{T^{-1}\int_T^{2T}N_0'(t)dt}&\le\kappa_T,&
|\mathcal E_{\partial}(F;T)|&\le r_T.
\end{aligned}
\end{equation}
Then the fraction of indices $k$ with $T<\gamma_k\le2T$
for which $\rho_k$ is simple and $\beta_k=\tfrac12$ is at least
\begin{equation}\label{eq:jone-fraction-v3}
\boxed{\begin{gathered}
1-4\max\Bigl(0,\,(1+O(1/\log T))\,[a_T+\kappa_T\sqrt{b_T}]
-\tfrac1{12}+r_T\Bigr)\\
-O\!\left(\frac{\log T}{M_T}\right)
\end{gathered}}
\end{equation}
The constants implicit in $O(1/\log T)$ and in $O(\log T/M_T)$ are absolute. This is a
\emph{conditional estimate on the three explicitly specified moments},
not an unconditional improvement of any published percentage.
\end{theorem}
\begin{proof}
By \eqref{eq:duplicate-count-v3}, at most $4\sum e_k(F)^2$
indices in fully contained groups of equal imaginary parts can be
non-simple or associated with an off-critical zero. The functional
equation pairs any off-critical zero with a distinct zero at the same
imaginary part; thus every excluded index belongs to such a group.
Only the at most two groups intersecting the endpoints of the window
need separate treatment. Their total multiplicities are $O(\log T)$,
using the unconditional $S(t)=O(\log t)$ and the zero-counting formula.
This gives a simple-critical-line fraction at least
$1-4\eta_T(F)-O(\log T/M_T)$.
Moreover
\begin{equation}
\frac{A_T(F)}{M_T}\le
\frac{T\sup_{[T,2T]}N_0'}{M_T}a_T
=(1+O(1/\log T))a_T.
\end{equation}
By Cauchy--Schwarz,
\begin{equation}
\frac{|B_T(F)|}{M_T}
\le \frac{T}{M_T}\sqrt{b_T}\left(\frac1T
\int_T^{2T}|\Phi_c'(t)|^2dt\right)^{1/2}
\le(1+O(1/T))\kappa_T\sqrt{b_T}.
\end{equation}
Insert these estimates and the endpoint bound into
\eqref{eq:jone-identity-v3}. Since $\eta_T(F)\ge0$, the resulting upper
bound for $\eta_T(F)$ may be replaced by its positive part.
\end{proof}

\begin{remark}[Numerical budget and its logical status]\label{rem:jone-budget-v3}
For the measured block $T=2^{19}$, the outputs
\texttt{budget\_w14\_th0.5\_j19.txt} and
\texttt{mv\_dphi\_th0.5\_j19.txt} report
\begin{equation}
 a_{\rm obs}=0.05253,\qquad b_{\rm obs}=0.00813,
\qquad\kappa_{\rm model}=0.8241.
\end{equation}
Substitution into the leading expression in
\eqref{eq:jone-fraction-v3} gives approximately $82.6\%$.
This value is a \emph{numerical conditional budget}; it is neither a
new verified lower bound at infinite height nor a rigorous finite-height
certificate. In particular, the Montgomery--Vaughan--Preissmann
calculation reported in \texttt{mv\_dphi\_th0.5\_j19.txt} treats a fixed-coefficient approximation to $\Phi_c'$.
The actual dyadic interpolation, the $t$-dependence of the weights,
the pole compensator (if included), and derivatives of the active
truncation introduce additional terms that must be bounded before the model
constant $0.8241$ may be used as a verified bound for $\kappa_T$.
If in addition one could establish $b_T\le3a_T^2$, $r_T=o(1)$,
and $\kappa_T\le0.8241+o(1)$ uniformly as $T\to\infty$,
then the lower bound \eqref{eq:jone-fraction-v3} would exceed $0.417$ for all large $T$ if
$\limsup_{T\to\infty}a_T<0.09437\ldots$, and would be positive for all large $T$
if $\limsup_{T\to\infty}a_T<0.13732\ldots$. These thresholds are
\emph{algebraic consequences of the conditional assumptions}.
The observed $a_{\rm obs}=0.05253$ is not a proof of either bound.
\end{remark}

\paragraph{Next analytic target.}
For the actual $W_{1/4}$ family with $\theta=1/2$, derive explicit,
height-uniform estimates for $a_T,b_T,\kappa_T$ and the endpoint term
$r_T$ to be inserted into \eqref{eq:jone-input-v3}, so that \eqref{eq:jone-fraction-v3} becomes a nontrivial bound. The estimate of the
fourth moment with an unspecified implied constant is insufficient.
The calculation suggests that $B_T$ is about $40\%$ of the sum
$A_T+B_T$ on the measured high block, so dropping it would be a
substantial mathematical error. This approach seeks a new result about
the author's \emph{prime-corrected construction} via its exact defect
identity, not a stand-alone classical characterization of RH.

%======================================================================
\section{Neighboring pairs: exact margins and finite-cutoff separation}
\label{sec:near-v4}
The difference between a strict spacing at each finite cutoff and a
nonzero limiting spacing is essential. This section studies that
question for the arithmetic sublevel points themselves. The symbol
$e_k(F)$ in \eqref{eq:ek-v3} remains the \emph{evaluation defect};
here the position error is denoted by $d_k(X;W)$ to avoid a collision.
For $k$ in the eventually increasing region of $N_0$ set
\begin{equation}\label{eq:poserror-v4}
 d_k(X;W):=\tau_k^{\Cms}(X;W)-\gamma_k,\qquad
 g_k:=\gamma_{k+1}-\gamma_k\ge0.
\end{equation}
\begin{proposition}[Exact symmetric-margin identity]
\label{prop:triangle-v4}
Define the symmetric margin
\begin{equation}\label{eq:margin-v4}
 \mathcal M_k(X;W):=\Delta_k(X;W)-|d_k(X;W)|-|d_{k+1}(X;W)|.
\end{equation}
Writing $x_+:=\max(x,0)$, one has exactly
\begin{equation}\label{eq:margin-exact-v4}
 \boxed{\ \mathcal M_k(X;W)
 =g_k-2\bigl[(d_k(X;W))_++(-d_{k+1}(X;W))_+\bigr]\le g_k.\ }
\end{equation}
Equality holds precisely when the two sublevel points enclose the
corresponding imaginary parts, i.e.
$\tau_k^{\Cms}(X;W)\le\gamma_k$ and
$\tau_{k+1}^{\Cms}(X;W)\ge\gamma_{k+1}$.
Consequently $\mathcal M_k(X;W)>0$ implies $g_k>0$.
Here $g_k$ is the unnormalized gap of \eqref{eq:poserror-v4}, not the
normalized nearest-neighbour separation of \cref{sec:v9-data}.
\end{proposition}
\begin{proof}
By definition $\Delta_k=g_k+d_{k+1}-d_k$.
Use $x-|x|=-2(-x)_+$ and $-x-|x|=-2x_+$.
The criterion for equality follows because both terms subtracted
from $g_k$ are nonnegative.
\end{proof}
\begin{remark}[What a numerical margin does not certify]
The sublevel points are prime-defined, whereas \eqref{eq:margin-v4}
uses the true $\gamma_k$. Thus positivity of a margin computed from
verified zeros is a useful finite test, not by itself a prime-only
inequality valid at every height. A repeated imaginary part gives
$g_k=0$ and hence $\mathcal M_k\le0$, even if
$\Delta_k(X;W)>0$ at every finite $X$.
\end{remark}

\begin{proposition}[Quantitative separation at a fixed prime cutoff]
\label{prop:Lipschitz-v4}
Fix $W\in\mathcal W$, $X\ge2$ and $k$ sufficiently large for the two barriers in
\cref{thm:block-sep-v3} to be defined. Set
\begin{equation}\label{eq:Mbound-v4}
 \mathscr L_{k,X}:=\operatorname*{ess\,sup}_{t\in[a_k,b_{k+1}]}
 |(F_X^W)'(t)|.
\end{equation}
Then $0<\mathscr L_{k,X}<\infty$ and
\begin{equation}\label{eq:Lipgap-v4}
 \boxed{\ \Delta_k(X;W)\ge\mathscr L_{k,X}^{-1}>0.\ }
\end{equation}
This bound is finite-cutoff only: the hypotheses do not give an
upper bound on $\mathscr L_{k,X}$ independent of $X$.
\end{proposition}
\begin{proof}
By the barriers, $F_X^W$ attains a value at or below
$k-\tfrac12$ in the initial part of $[a_k,b_{k+1}]$ and a value
at or above $k+\tfrac12$ in its final part. Let $r$ be a point of the
initial part with $F_X^W(r)\le k-\tfrac12$, put
$s_1=\inf\{t>r:F_X^W(t)\ge k+\tfrac12\}$ and
$s_0=\sup\{t\in[r,s_1]:F_X^W(t)\le k-\tfrac12\}$. By continuity and
these extremal choices, $F_X^W(s_0)=k-\tfrac12$,
$F_X^W(s_1)=k+\tfrac12$, and $k-\tfrac12<F_X^W<k+\tfrac12$ on
$(s_0,s_1)$. The Lipschitz bound gives
$1\le \mathscr L_{k,X}(s_1-s_0)\le\mathscr L_{k,X}\Delta_k$.
In particular $\mathscr L_{k,X}>0$: a locally absolutely continuous
function whose derivative vanishes almost everywhere on $[s_0,s_1]$
would be constant there.
On a compact interval the prime phase, compensator and $H$ are
locally Lipschitz for fixed $X$; so is the truncation, hence
$\mathscr L_{k,X}<\infty$.
\end{proof}

%======================================================================
\section{The \texorpdfstring{$W_{1/4}$}{W\textunderscore 1/4} two-source kernel: an exact resolution law}
\label{sec:kernel-v4}
The following kernel analysis is a \emph{local two-source model};
it is not an unconditional decomposition of the actual Euler phase
into individual contributions of zeta zeros. Define
\begin{equation}\label{eq:q14-v4}
 q_{1/4}(v)=\frac v2+\frac14\quad (1\le v\le2),\qquad
 K(z)=\frac12+\frac1\pi\int_1^2q_{1/4}(v)
 \operatorname{Si}(vz)\,dv,
\end{equation}
where $\operatorname{Si}(x)=\int_0^x(\sin y)/y\,dy$ is continuous at
$x=0$. The index $1/4$ labels the weight, not the real part of any
zero. For $m\ge0$, write
\begin{equation}\label{eq:moments-v4}
 \mu_m:=\int_1^2v^m q_{1/4}(v)\,dv.
\end{equation}
In particular $\mu_1=37/24$, $\mu_2=59/24$,
$\mu_3=323/80$ and $\mu_4=34/5$.

\begin{lemma}[Bounded derivatives of the two-source kernel]
\label{lem:Kder-v4}
The function $K$ extends to an entire function of $z$. On the real
axis,
\begin{equation}\label{eq:Kder-v4}
\begin{split}
 K'(z)&=\frac1\pi\int_1^2q_{1/4}(v)\int_0^v\cos(zy)\,dy\,dv,\\
 K''(z)&=-\frac1\pi\int_1^2q_{1/4}(v)\int_0^v y\sin(zy)\,dy\,dv,\\
 K'''(z)&=-\frac1\pi\int_1^2q_{1/4}(v)\int_0^v y^2\cos(zy)\,dy\,dv,\\
 K^{(4)}(z)&=\frac1\pi\int_1^2q_{1/4}(v)
 \int_0^v y^3\sin(zy)\,dy\,dv.
\end{split}
\end{equation}
Thus $K(-z)=1-K(z)$ and
\begin{equation}\label{eq:Kbounded-v4}
 \sup_{z\in\mathbb R}|K''(z)|\le\frac{\mu_2}{2\pi},\qquad
 \sup_{z\in\mathbb R}|K^{(4)}(z)|\le\frac{\mu_4}{4\pi}.
\end{equation}
Moreover $K'(z),K''(z)\to0$ as $|z|\to\infty$ through real values,
and neither is identically zero.
\end{lemma}
\begin{proof}
Since $\operatorname{Si}(x)=\sum_{n\ge0}(-1)^nx^{2n+1}/((2n+1)(2n+1)!)$
is entire and the $v$-integral is over a compact set, $K$ is entire.
The identity
$\operatorname{Si}(vz)=\int_0^v(\sin(zy))/y\,dy$, whose integrand has a
removable singularity at $y=0$, permits
termwise differentiation on compact subsets of the complex plane.
All estimates follow by integrating $y$ and $y^3$ against the
nonnegative weight. The decay of $K'$ and $K''$ is the
Riemann--Lebesgue lemma for the integrable, compactly supported
kernels obtained after interchanging the $v,y$ integrals, namely
$K'(z)=\frac1\pi\int_0^2w(y)\cos(zy)\,dy$ and
$K''(z)=-\frac1\pi\int_0^2yw(y)\sin(zy)\,dy$, with $w$ as in the proof
of \cref{thm:pair-injective-v4}. Finally
$K'(0)=\mu_1/\pi>0$, while
$K'''(0)=-\mu_3/(3\pi)<0$.
\end{proof}

\begin{theorem}[Global quadratic resolution estimate]
\label{thm:resolution-v4}
For real $z$ and $h$, put
$C(z,h):=K(z-h/2)+K(z+h/2)$.
Then
\begin{equation}\label{eq:pair-taylor-v4}
 \boxed{\quad C(z,h)=2K(z)+\frac{h^2}{4}K''(z)
 +\mathcal R_K(z,h),\qquad
 |\mathcal R_K(z,h)|\le\frac{\mu_4}{768\pi}|h|^4
 =\frac{17}{1920\pi}|h|^4.\quad}
\end{equation}
In addition,
\begin{equation}\label{eq:pair-lipschitz-v4}
 \sup_{z\in\mathbb R}|C(z,h)-C(z,0)|
 \le\frac{\mu_2}{8\pi}h^2.
\end{equation}
Consequently, for two sources at separation $g$ and local scale $L=\log X$ (so that $h=Lg$), a gap $g$ small
compared with $1/L$ changes the full two-source profile only at
quadratic order $L^2g^2$. This is a theorem \emph{about $C$}, not an
estimate of the contribution of all zeta zeros.
\end{theorem}
\begin{proof}
Apply Taylor's theorem with fourth-order remainder to
$K(z+h/2)+K(z-h/2)$. Odd powers cancel. The remainder is
at most $2(|h|/2)^4\sup|K^{(4)}|/4!$, proving
\eqref{eq:pair-taylor-v4}; the second symmetric-difference
inequality and \eqref{eq:Kbounded-v4} prove
\eqref{eq:pair-lipschitz-v4}.
\end{proof}

\begin{theorem}[Injectivity of the complete local pair profile]
\label{thm:pair-injective-v4}
Fix $L>0$ and define, for any center $c\in\mathbb R$ and
separation $g\ge0$,
\begin{equation}\label{eq:pair-complete-v4}
 \mathcal C_{c,g,L}(t):=
 K\!\left(L(t-c-g/2)\right)+
 K\!\left(L(t-c+g/2)\right).
\end{equation}
If $\mathcal C_{c_1,g_1,L}$ and $\mathcal C_{c_2,g_2,L}$ coincide on
any nonempty open real interval, then $c_1=c_2$ and $g_1=g_2$.
Thus, in the exact two-source model, the \emph{complete profile}
identifies even an arbitrarily small positive gap. This is an
injectivity statement, not a stable recovery estimate in the
presence of unknown contributions.
\end{theorem}
\begin{proof}
Define the compactly supported nonnegative function
\begin{equation}
 w(y)=\int_{\max(1,y)}^2q_{1/4}(v)\,dv
 \quad(0\le y\le2).
\end{equation}
For $0<y<1$, $w(y)=1$. Interchanging the integrals in
\eqref{eq:Kder-v4} gives
\begin{equation}
 K'(z)=\frac1\pi\int_0^2w(y)\cos(zy)\,dy.
\end{equation}
Consequently,
\begin{equation}\label{eq:pair-transform-v4}
 \mathcal C'_{c,g,L}(t)=\frac{2L}{\pi}
 \int_0^2w(y)\cos\bigl(L(t-c)y\bigr)
 \cos\bigl(Lgy/2\bigr)\,dy.
\end{equation}
The profiles are entire in $t$, so agreement on an open real
interval implies agreement everywhere, and then their derivatives agree.
Writing $\cos(L(t-c)y)$ as the average of $e^{\pm iL(t-c)y}$ and
substituting $y\mapsto-y$ in the second term,
\begin{equation}
 \mathcal C'_{c,g,L}(t)=\frac L\pi\int_{-2}^2f_{c,g}(y)e^{iLty}\,dy,
 \qquad f_{c,g}(y)=w(|y|)\cos(Lg|y|/2)e^{-iLcy},
\end{equation}
with $f_{c,g}$ integrable and supported in $[-2,2]$. Uniqueness of the
Fourier transform of $L^1$ functions gives $f_{c_1,g_1}=f_{c_2,g_2}$
almost everywhere; both are continuous on $0<y<1$, where $w=1$, so
\begin{equation}
 e^{-iLc_1y}\cos(Lg_1y/2)
 =e^{-iLc_2y}\cos(Lg_2y/2),\qquad 0<y<1.
\end{equation}
By analyticity this also holds at $y=0$. Differentiate once at
$y=0$ to get $c_1=c_2$, and twice to get $g_1^2=g_2^2$.
Since both separations are nonnegative, $g_1=g_2$.
\end{proof}

\begin{remark}[Why $\theta=0.7$ is effective, without claiming optimality]
If $X=t^\theta$ for the numerical weighted construction, the
support of the prime-power phase extends to $n\le t^{2\theta}$.
The choice $\theta=0.7$ therefore reaches $t^{1.4}$, whereas
$\theta=0.5$ reaches $t$. This increases the local scale $L$ and hence the scale
$L^2g^2/4$ of the quadratic term in \eqref{eq:pair-taylor-v4}; the term
itself still vanishes at the zeros of $K''$.
The weight has $q_{1/4}(2)=5/4$ and $\mu_1=37/24$; its upper
endpoint receives more emphasis than in Selberg's constant density
$q\equiv1$. The observed superiority of $\theta=0.7$ is empirical;
neither its optimality nor an unconditional arithmetic separation
theorem at that cutoff follows from this kernel calculation.
\end{remark}

%======================================================================
\section{A prime-only spatial filter and its cutoff derivative}
\label{sec:filter-v4}
For the rest of this section, $X=e^L$ with $L>0$ is a \emph{fixed}
cutoff and the phase is \emph{untruncated}:
$F_L^\circ(t):=N_0(t)+P_{e^L}^W(t)$. This differs from
$F_{e^L}^W(t)$ in \eqref{eq:Cms}; replacing it by the clipped
phase in a derivative identity would be an error.
Here ``prime-only'' means built from prime powers through $\Lambda$,
the pole compensator and $N_0$, without any information on the zeros.
For a function $G$ defined near $c>0$ and $r>0$ with $c>r/L$ set
\begin{equation}\label{eq:Dr-v4}
 \mathscr D_{L;r}[G](c):=
 \frac{G(c+r/L)-G(c-r/L)}{2r}.
\end{equation}
For $0<r_1<r_2$ (and $c>r_2/L$) set
\begin{equation}\label{eq:Qdef-v4}
 \mathscr Q_{L;r_1,r_2}[G](c)
 :=\mathscr D_{L;r_1}[G](c)-\mathscr D_{L;r_2}[G](c),\qquad
 \mathscr B_{r_1,r_2}(v):=\frac{\sin(r_1v)}{r_1}
 -\frac{\sin(r_2v)}{r_2}.
\end{equation}
\begin{proposition}[Exact Euler formula for the spatial filter]
\label{prop:Qexact-v4}
For any admissible weight $W$ of \eqref{eq:weights}, the spatial
filter of $F_L^\circ$ satisfies
\begin{equation}\label{eq:Qprime-v4}
\begin{split}
\mathscr Q_{L;r_1,r_2}[F_L^\circ](c)
={}&\mathscr Q_{L;r_1,r_2}[N_0](c)\\
&-\frac1\pi\sum_{n\le e^{2L}}
 W\!\left(\frac{\log n}{L}\right)
 \frac{\Lambda(n)}{\sqrt n\log n}\cos(c\log n)
 \mathscr B_{r_1,r_2}\!\left(\frac{\log n}{L}\right)\\
&-\mathscr Q_{L;r_1,r_2}[\Pi_{e^L}^W](c).
\end{split}
\end{equation}
Thus the right side is determined by the primes and the
Riemann--von Mangoldt main term; there is no appearance of
$\beta_k$ or of the actual $\gamma_k$.
If a background $B$ is five times continuously differentiable on
$[c-r_2/L,c+r_2/L]$, then
\begin{equation}\label{eq:Qbackground-v4}
\begin{split}
\mathscr Q_{L;r_1,r_2}[B](c)
={}&\frac{r_1^2-r_2^2}{6L^3}B'''(c)+\mathcal R_B,\\
|\mathcal R_B|&\le\frac{r_1^4+r_2^4}{120L^5}
 \sup_{|t-c|\le r_2/L}|B^{(5)}(t)|.
\end{split}
\end{equation}
In particular the constant and affine parts of $B$ are eliminated
\emph{exactly}, not merely approximately.
\end{proposition}
\begin{proof}
The finite Euler sum has imaginary part
\begin{equation}
-\frac1\pi\sum_n\frac{\Lambda(n)}{\sqrt n\log n}
 W\!\left(\frac{\log n}{L}\right)\sin(t\log n).
\end{equation}
Use the sine subtraction identity at $c\pm r/L$, then subtract
the two radii. The compensator contributes the last term in
\eqref{eq:Qprime-v4}. Taylor's theorem for the symmetric first
difference gives
$\mathscr D_{L;r}[B](c)=B'(c)/L+
 r^2B'''(c)/(6L^3)+\mathcal R$ with
$|\mathcal R|\le r^4\sup|B^{(5)}|/(120L^5)$.
\end{proof}

\begin{proposition}[Cutoff derivative and the moving-point correction]
\label{prop:dL-v4}
Suppose $q$ is continuous on $[1,2]$, with $W$ given by
\eqref{eq:weights}. For $L>0$ away from the discrete values
$L=\log n$ and $L=\tfrac12\log n$ (that is, $e^L=n$ or $e^{2L}=n$) of prime powers $n\ge2$,
write $\dot P_L(t):=(\partial/\partial L)
P_{e^L}^W(t)$ with $t$ \emph{held fixed}. Then
\begin{equation}\label{eq:dLprime-v4}
 \boxed{\quad
 L^2\dot P_L(t)=\frac1\pi\operatorname{Im}
 \left[\sum_{e^L<n<e^{2L}}
 q\!\left(\frac{\log n}{L}\right)
 \frac{\Lambda(n)}{n^{1/2+it}}
 -\int_{e^L}^{e^{2L}}
 q\!\left(\frac{\log u}{L}\right)u^{-1/2-it}\,du
 \right].\quad}
\end{equation}
The derivative is taken piecewise at the exceptional values and
exists almost everywhere in $L$. For any family $G_L(t)$
differentiable in both variables, with $G_L'$ the derivative in $t$ and
$\partial_LG_L$ the derivative in $L$, the \emph{total} derivative of the moving spatial filter is
\begin{equation}\label{eq:commutator-v4}
\begin{split}
\frac{d}{dL}\mathscr Q_{L;r_1,r_2}[G_L](c)
={}&\mathscr Q_{L;r_1,r_2}[\partial_LG_L](c)\\
&-\frac{1}{2L^2}\Bigl(G_L'(c+r_1/L)+G_L'(c-r_1/L)\\
&\hspace{32mm}-G_L'(c+r_2/L)-G_L'(c-r_2/L)\Bigr).
\end{split}
\end{equation}
In particular, differentiating at fixed $t$ and then taking the
spatial filter is \emph{not} interchangeable with differentiating
the filter as a function of $L$.
\end{proposition}
\begin{proof}
Inside the transition region $1<\log n/L<2$, one has
\begin{equation}
\partial_L W\!\left(\frac{\log n}{L}\right)
=\frac{\log n}{L^2}\,q\!\left(\frac{\log n}{L}\right);
\end{equation}
outside this region the derivative is zero. For the pole compensator write
$\Pi_{e^L}^W(t)=\frac1\pi\operatorname{Im}\int_{\log2}^{\infty}W(v/L)e^{v/2-itv}\,\frac{dv}v$;
for $L$ in a compact neighbourhood of a non-exceptional value the
integrand is supported in a fixed compact $v$-interval and its
$L$-derivative is bounded there, so differentiation under the integral
sign (dominated convergence) gives the integral term. This proves \eqref{eq:dLprime-v4}.
For \eqref{eq:commutator-v4}, differentiate
$G_L(c\pm r/L)$ by the chain rule, retaining the moving-point
contributions from $\mp r/L^2$, and subtract the two radii.
\end{proof}

%======================================================================
\section{Finite-sample identifiability and the remaining arithmetic error}
\label{sec:identify-v4}
The kernel calculation can be promoted to a precise \emph{conditional}
local inverse statement. It does not require GUE. It also does not
supply, on its own, a decomposition of $F_L^\circ$ when RH is not
assumed.
\begin{theorem}[Nondegenerate quadratic local inverse]
\label{thm:identifiability-v4}
Fix any nonempty open interval $J\subset\mathbb R$. The six entire
functions
\begin{equation}\label{eq:six-v4}
 1,\quad z,\quad z^2,\quad z^3,\quad K'(z),\quad K''(z)
\end{equation}
are linearly independent. Consequently there exist six distinct
points $z_1,\dots,z_6\in J$ such that the evaluation matrix of
\eqref{eq:six-v4} is invertible.
Let $\mathsf A$ be the $6\times5$ matrix obtained by evaluating
the first five functions at these points, and let
$\mathbf v=(K''(z_i))_{i=1}^6$. Let $\mathsf P$ be the orthogonal
projection onto the orthogonal complement of the range of
$\mathsf A$, so that $\mathsf P\mathbf v\ne0$.
If measured prime-phase values, after subtraction of the
coincident-pair model $2K$, admit the representation
\begin{equation}\label{eq:vector-v4}
 \mathbf y=\mathsf A\mathbf b+\frac{h^2}{4}\mathbf v
   +\mathbf r+\mathbf r_4,
 \qquad\|\mathbf r_4\|_2\le
 \frac{\mu_4}{768\pi}\sqrt6\,|h|^4,
\end{equation}
where $\mathbf b\in\mathbb R^5$ is an unknown nuisance vector,
$\mathbf r_4$ is the vector of fourth-order Taylor remainders of \eqref{eq:pair-taylor-v4} at the points $z_i$,
$\mathbf r$ is an unspecified arithmetic residual, and $h=Lg\ge0$,
then the estimator, observable from the sampled vector $\mathbf y$
(and prime-observable once $\mathbf y$, the sample points, the centre and
$L$ are constructed from prime-defined data),
\begin{equation}\label{eq:estimator-v4}
 \widehat{h^2}:=4\frac{\langle\mathsf P\mathbf y,
 \mathsf P\mathbf v\rangle}{\|\mathsf P\mathbf v\|_2^2}
\end{equation}
satisfies the rigorous conditional stability estimate
\begin{equation}\label{eq:estimate-error-v4}
 \boxed{\quad
 |\widehat{h^2}-h^2|\le
 \frac{4}{\|\mathsf P\mathbf v\|_2}
 \left(\|\mathsf P\mathbf r\|_2+
 \frac{\mu_4\sqrt6}{768\pi}|h|^4\right).\quad}
\end{equation}
A proved bound on the right (valid for the actual $h$) smaller than a \emph{positive measured}
$\widehat{h^2}$ would certify $h>0$ \emph{provided} the
representation \eqref{eq:vector-v4} is valid for the actual phase.
\end{theorem}
\begin{proof}
Suppose a linear combination of \eqref{eq:six-v4} vanishes on $J$.
By analyticity it vanishes on all of $\mathbb R$. Its odd part gives
$az+bz^3+cK''(z)=0$; since $K''(z)\to0$ at infinity,
$a=b=0$, then $c=0$ because $K''\not\equiv0$.
Its even part gives $d+ez^2+fK'(z)=0$, hence
$d=e=f=0$ by decay and $K'(0)>0$.
Linear independence of analytic functions implies that evaluations
at some six distinct points of any open interval span $\mathbb R^6$;
otherwise the evaluation vectors would lie in a common proper
subspace, giving a forbidden linear relation on $J$.
Thus $\mathsf P\mathbf v\ne0$. Apply $\mathsf P$ to
\eqref{eq:vector-v4}, take the scalar product with
$\mathsf P\mathbf v$ and use Cauchy--Schwarz. The bound on
$\mathbf r_4$ is the remainder estimate of \cref{thm:resolution-v4}.
\end{proof}

\begin{remark}[Center uncertainty and the limit of the result]
A small shift of the assumed center changes $2K(z)$ by a term
proportional to $K'(z)$ at first order, which explains the fifth
column of $\mathsf A$. Higher powers of the center error must be
included in $\mathbf r$ or estimated separately. A smooth local
background is not \emph{automatically} a cubic polynomial, even when
its derivatives exist; its interpolation residual is also included
in $\mathbf r$. Most importantly, the off-critical-zero contribution
to the full Euler phase cannot be declared part of a controlled
background without proof. Thus \eqref{eq:estimate-error-v4} is
an exact inverse stability theorem, \emph{not} a proof that the
actual prime residual satisfies the necessary bound. Explicitly, if a
proved bound $\mathcal B$ for the right side of
\eqref{eq:estimate-error-v4}, valid for the actual $h$, satisfies $\mathcal B<\widehat{h^2}$, then
$h^2\ge\widehat{h^2}-\mathcal B>0$. The estimator is an algebraic
quantity and may be negative when a residual is present.
\end{remark}

\paragraph{Research reduction for the actual weighted constructions.}
For every fixed $k$ and cutoffs $X$ for which $h=Lg$ is of moderate size (for large $h$ the quartic term of \eqref{eq:estimate-error-v4} dominates the quadratic signal), the
new quantitative task is to establish a representation of the
form \eqref{eq:vector-v4} in a neighborhood selected using only
the prime-defined sublevel points, together with a bound for
$\|\mathsf P\mathbf r\|_2$ that is \emph{smaller than the quadratic
signal} $\frac14\|\mathsf P\mathbf v\|_2h^2$ at the true separation. Such a bound must be
proved without imposing that the relevant zeros lie on the critical
line; if this cannot be achieved, the local inverse gives an effective
finite-height diagnostic but not RH. The required bound is independent of whether
any hypothetical off-line zero has a maximal real part.

\section{Local spacing observations: scope}
\label{sec:observed-v4}
The following data are included solely to identify the regime
motivating \cref{sec:kernel-v4,sec:filter-v4,sec:identify-v4}. They are
not numerically certified.
Unfolded spacing here means applying $N_0$ to the corresponding
positive heights, where $N_0$ is increasing. In the index range
$10^6\le k<2\cdot10^6$, the nearest-neighbor study reported:
\begin{center}\footnotesize
\begin{tabular}{@{}p{0.31\linewidth}p{0.15\linewidth}p{0.25\linewidth}p{0.19\linewidth}@{}}
\toprule
Prime rule & Minimum gap & Median tauc gap for 122 tight pairs & Median midpoint error\\
\midrule
IJNT original & $0.235$ & $0.354$ & $0.0092$\\
$W_{1/4}$, $\theta=0.5$ & $0.222$ & $0.293$ & $0.0023$\\
$W_{1/4}$, $\theta=0.7$ & $0.173$ & $0.210$ & $0.0008$\\
Sharp support, $\theta=1.4$ & $0.144$ & $0.163$ & $0.0023$\\
\bottomrule
\end{tabular}
\end{center}
These numbers concern the numerical \emph{original clipped phases}
with the small $\log\log t$ truncation in the supplied scripts; they
must not be silently reinterpreted as calculations for the
MCOM-truncated $F_X^W$ with the pole compensator.
For gaps below $0.05$ unfolded units, the reported Spearman
correlation between the two tauc gaps and the real gaps was close
to zero for all tested rules, whereas the midpoint error was much
smaller than individual point errors. A separate margin study
reports the equality $\mathcal M_k=g_k$ for the tight pairs in the weighted
$\theta=0.5,0.7$ constructions (with the phases of that study, which need not be the MCOM-truncated $F_X^W$ of \eqref{eq:margin-v4}); by
\eqref{eq:margin-exact-v4} this equality holds precisely when the two sublevel points enclose the
actual pair. These observations motivate, but do not establish, the
local quadratic model. In particular they do not imply that an
off-line pair can be discriminated by a fixed prime cutoff.




%======================================================================
\section{Correlated displacement from the Riemann--von Mangoldt points:
       what proximity does and does not prove}\label{sec:correlated-v5}
This section gives the precise response to a natural question: since the
Riemann--von Mangoldt points are strictly separated and the prime-defined
sublevel points are strictly separated at each finite cutoff (for every sufficiently large index; see the discussion after \eqref{eq:gap-v3}), can their
\emph{proximity} prevent two consecutive sublevel points from colliding as
$X\to\infty$?  The answer is a conditional yes, provided the \emph{difference
of consecutive displacements} is controlled.  Neither finite-cutoff
strictness nor an amplitude bound alone proves the required control.
The conclusions here concern the author's own constructions and do not
introduce a separate classical criterion for RH.

\subsection{Exact decomposition into a positive main spacing and a
           correlated arithmetic displacement}
Let $\gamma_k^{\mathrm{RvM}}$ be the unique solution greater than $2\pi$ of
$N_0(\gamma_k^{\mathrm{RvM}})=k-\tfrac12$, and put
\begin{equation}\label{eq:hrvm-v5}
 h_k^{\mathrm{RvM}}:=\gamma_{k+1}^{\mathrm{RvM}}-
 \gamma_k^{\mathrm{RvM}}>0.
\end{equation}
The positivity holds for every $k\ge1$ since $N_0$ is strictly
increasing on $(2\pi,\infty)$ and its range on that interval contains
every $k-\tfrac12$: indeed $N_0(2\pi)=-1+\frac78=-\frac18$, so the
range is $(-\frac18,\infty)$.  For the weighted, MCOM-truncated sublevel points
\eqref{eq:Cms}, define the \emph{RvM displacement}, distinct from the
true-zero position error $d_k(X;W)$ in \eqref{eq:poserror-v4}, by
\begin{equation}\label{eq:epsrvm-v5}
 \varepsilon_k^{\mathrm{RvM}}(X;W):=
 \tau_k^{\Cms}(X;W)-\gamma_k^{\mathrm{RvM}}.
\end{equation}
It is determined by the primes, the pole compensator and $N_0$,
without using any $\gamma_k$ or $\beta_k$.
\begin{proposition}[Exact correlated-displacement identity]
\label{prop:correlated-v5}
For every $k\ge1$ and every cutoff $X$ for which the defining
sublevel measures are finite,
\begin{equation}\label{eq:correlated-v5}
\boxed{\quad
 \Delta_k(X;W)=h_k^{\mathrm{RvM}}
 +\varepsilon_{k+1}^{\mathrm{RvM}}(X;W)
 -\varepsilon_k^{\mathrm{RvM}}(X;W).
\quad}
\end{equation}
In particular, for a fixed index $k$ the following two properties are
equivalent:
\begin{align}
 &\liminf_{X\to\infty}\Delta_k(X;W)>0,
 \label{eq:noncollapse-v5}\\
 &\limsup_{X\to\infty}
 \bigl[\varepsilon_k^{\mathrm{RvM}}(X;W)-
       \varepsilon_{k+1}^{\mathrm{RvM}}(X;W)\bigr]
 <h_k^{\mathrm{RvM}}.
 \label{eq:coherence-v5}
\end{align}
A sufficient quantified form is: for some $\eta_k\in(0,1)$ and
$X_k\ge2$,
\begin{equation}\label{eq:coherence-eta-v5}
 \varepsilon_k^{\mathrm{RvM}}(X;W)-
 \varepsilon_{k+1}^{\mathrm{RvM}}(X;W)
 \le(1-\eta_k)h_k^{\mathrm{RvM}}
 \qquad(X\ge X_k).
\end{equation}
Under \eqref{eq:coherence-eta-v5},
$\Delta_k(X;W)\ge\eta_k h_k^{\mathrm{RvM}}>0$ for $X\ge X_k$.
There is no requirement that $\eta_k$ be uniform in $k$.
\end{proposition}
\begin{proof}
Subtract \eqref{eq:epsrvm-v5} at the consecutive indices and
use \eqref{eq:hrvm-v5}.  Taking upper and lower limits yields
\begin{equation}
 \liminf_{X\to\infty}\Delta_k(X;W)
 =h_k^{\mathrm{RvM}}-
 \limsup_{X\to\infty}
 (\varepsilon_k^{\mathrm{RvM}}-
  \varepsilon_{k+1}^{\mathrm{RvM}}),
\end{equation}
with the usual extended-real interpretation if a limit is infinite.
The quantified conclusion follows by direct substitution.
\end{proof}
\begin{remark}[Exact identity, not an independent arithmetic bound]
Condition \eqref{eq:coherence-v5} is an exact reformulation of
noncollapse, not a proof of noncollapse.  Its value is to identify
\emph{which arithmetic difference} one must estimate.  Controlling
$|\varepsilon_k^{\mathrm{RvM}}|$ and
$|\varepsilon_{k+1}^{\mathrm{RvM}}|$ separately is sufficient if
\begin{equation}
 |\varepsilon_k^{\mathrm{RvM}}|+
 |\varepsilon_{k+1}^{\mathrm{RvM}}|
 \le(1-\eta_k)h_k^{\mathrm{RvM}},
\end{equation}
but is stronger than necessary.  A large common displacement cancels
exactly in \eqref{eq:correlated-v5}.  The quantities
$\varepsilon_k^{\mathrm{RvM}}$ must not be confused either with the
evaluation defects $e_k(F)$ of \eqref{eq:ek-v3} or with the
true-zero errors $d_k(X;W)$ of \eqref{eq:poserror-v4}.
\end{remark}

\subsection{What the existing MCOM truncation actually yields}
The large-index barriers of \cref{thm:block-sep-v3} provide an
\emph{individual} proximity estimate.  The asymptotic MCOM coefficient
$C_{\mathrm H}$ must not be used for the \emph{fixed} truncation $H$;
the coefficient appropriate to $H$ is computed here.
\begin{proposition}[Uniform RvM proximity from the fixed truncation]
\label{prop:rvmbound-v5}
Let
\begin{equation}
 H(t)=A_{\mathrm M}\log\max(t,e)+B_{\mathrm M},\qquad
 A_{\mathrm M}=0.08382,\quad B_{\mathrm M}=5.5,
\end{equation}
and let $a_k,b_k$ be the large-index barriers of
\cref{thm:block-sep-v3}.  For $k$ sufficiently large, uniformly in $X\ge2$,
\begin{equation}\label{eq:rvmbound-v5}
 -\frac{H(a_k)}{N_0'(\xi_{k,-})}
 \le\varepsilon_k^{\mathrm{RvM}}(X;W)
 \le\frac{H(b_k)}{N_0'(\xi_{k,+})},
\end{equation}
where $\xi_{k,-}\in(a_k,\gamma_k^{\mathrm{RvM}})$ and
$\xi_{k,+}\in(\gamma_k^{\mathrm{RvM}},b_k)$ are supplied by the mean-value
theorem.  Consequently,
\begin{equation}\label{eq:rvmasympt-v5}
 \boxed{\quad
 \sup_{X\ge2}|\varepsilon_k^{\mathrm{RvM}}(X;W)|
 \le 2\pi A_{\mathrm M}+o(1)
 \qquad(k\to\infty).
 \quad}
\end{equation}
On the other hand,
\begin{equation}\label{eq:rvmsmallgap-v5}
 h_k^{\mathrm{RvM}}\sim
 \frac{2\pi}{\log\gamma_k^{\mathrm{RvM}}}\to0.
\end{equation}
Thus \eqref{eq:rvmasympt-v5} does not imply either
\eqref{eq:coherence-v5} or \eqref{eq:noncollapse-v5}.
\end{proposition}
\begin{proof}
Set $v_k:=\gamma_k^{\mathrm{RvM}}$.  Since
\(N_0(a_k)+H(a_k)=N_0(v_k)=N_0(b_k)-H(b_k)\),
the mean-value theorem gives
\begin{equation}
 v_k-a_k=\frac{H(a_k)}{N_0'(\xi_{k,-})},\qquad
 b_k-v_k=\frac{H(b_k)}{N_0'(\xi_{k,+})}.
\end{equation}
The barriers $a_k\le\tau_k^{\Cms}(X;W)\le b_k$ yield
\eqref{eq:rvmbound-v5}.  Moreover $a_k/v_k\to1$ and $b_k/v_k\to1$.
For fixed $\delta\in(0,1)$, $N_0(v_k)-N_0((1-\delta)v_k)\gg_\delta
v_k\log v_k$, whereas $H((1-\delta)v_k)=O(\log v_k)$; hence
$N_0+H<k-\frac12$ at $(1-\delta)v_k$, and $a_k>(1-\delta)v_k$ for large
$k$ because $N_0+H$ is increasing there. Similarly
$N_0((1+\delta)v_k)-N_0(v_k)\gg_\delta v_k\log v_k>H((1+\delta)v_k)$
gives $b_k<(1+\delta)v_k$. Since $\delta$ is arbitrary, the two ratios
tend to $1$.
On this common asymptotic scale,
$H(t)=A_{\mathrm M}\log v_k+O(1)+o(\log v_k)$ and
$N_0'(t)=(\log v_k)/(2\pi)+O(1)$, so the two extremal
differences are $2\pi A_{\mathrm M}+o(1)$.
Finally, the mean-value theorem between $v_k$ and $v_{k+1}$
gives \eqref{eq:rvmsmallgap-v5}.
\end{proof}
\begin{remark}[The asymptotic MCOM coefficient requires a different $H$]
Suppose the author's asymptotic bound \eqref{eq:asymptMCOM} for $S(t)$ holds, with coefficient
$C_{\mathrm H}<A_{\mathrm M}$.  For any fixed $c>C_{\mathrm H}$,
one may \emph{define a different sublevel family} whose truncation
bound is eventually $\widetilde H_c(t)=c\log t$ and whose
small-height extension is chosen sufficiently large to contain
$|N(t)-N_0(t)|$.  For that \emph{different} family, the argument of
\cref{prop:rvmbound-v5} yields
$\sup_X|\widetilde\varepsilon_k^{\mathrm{RvM}}|
\le2\pi c+o(1)$ as $k\to\infty$; one may then take
$c\downarrow C_{\mathrm H}$ at the level of asymptotic coefficients.
This does \emph{not} replace $A_{\mathrm M}$ by
$C_{\mathrm H}$ in \eqref{eq:rvmasympt-v5} for the fixed
$H(t)=0.08382\log\max(t,e)+5.5$ of the present weighted taucms.
Nor does the improved individual bound alone imply noncollapse.
\end{remark}

\subsection{What DBSRH would transfer, and what mean-square proximity cannot}
\begin{corollary}[Conditional transfer of RvM separation]
\label{cor:rh-rvm-v5}
Assume DBSRH$(W)$ as formulated in \eqref{eq:DBSRH} and, for every
$k\ge1$, condition \eqref{eq:coherence-v5}.  Then all positive
imaginary parts of the zeros of $\zeta$ are pairwise distinct,
so RH holds and every nontrivial zero is simple.
Conversely, \emph{under DBSRH$(W)$}, condition
\eqref{eq:coherence-v5} for all $k$ is equivalent to RH
\emph{together with} simplicity.
\end{corollary}
\begin{proof}
By DBSRH$(W)$, at each fixed index
$\tau_k^{\Cms}(X;W)\to\gamma_k$.  Therefore
\begin{equation}
 \Delta_k(X;W)\longrightarrow\gamma_{k+1}-\gamma_k.
\end{equation}
\Cref{prop:correlated-v5} makes the limit strictly positive for every
$k$, hence no imaginary parts repeat.  If a zero
$\rho=\beta+i\gamma$ has $\beta\ne\tfrac12$, the distinct zero
$1-\overline\rho=1-\beta+i\gamma$ repeats the same positive
imaginary part.  A multiple zero does so as well.  Conjugation treats
negative imaginary parts.  This proves RH and simplicity.  Conversely,
if RH and simplicity hold, all consecutive $\gamma_k$ are strictly
increasing, and the displayed limit is positive.  This implies
\eqref{eq:coherence-v5} by \cref{prop:correlated-v5}.
Thus, under DBSRH$(W)$: \eqref{eq:coherence-v5} for every $k$ $\iff$
$\gamma_{k+1}>\gamma_k$ for every $k$ $\iff$ RH and simplicity. This is
not an equivalence between DBSRH$(W)$ and RH.
\end{proof}
\begin{remark}[Exact signature of a repeated imaginary part]
Under DBSRH$(W)$, a repeated imaginary part at indices $k,k+1$
would force
\begin{equation}\label{eq:dup-rvm-v5}
 \lim_{X\to\infty}
 \bigl[\varepsilon_k^{\mathrm{RvM}}(X;W)-
       \varepsilon_{k+1}^{\mathrm{RvM}}(X;W)\bigr]
 =h_k^{\mathrm{RvM}}.
\end{equation}
Thus the entire positive RvM spacing would be consumed by a
\emph{difference} of arithmetic displacements.  This is a precise
target for a prime-based non-collapse estimate; it is not, without
such an estimate, a contradiction.
Mean-square bounds such as \eqref{eq:oldMS} or their taucm analogue
control an \emph{average over indices in height windows}; they do not
exclude a single exceptional pair at a fixed $k$, even if the
average RMS is comparable to the RvM mean spacing.
\end{remark}

\subsection{A derivative-based sufficient condition, and why it is too strong}
The local Lipschitz inequality of \cref{prop:Lipschitz-v4}
illustrates one way to obtain a cutoff-independent spacing, but
its required uniform derivative control is not available for the
Euler phase.  The following sharpened criterion also identifies an
\emph{obstruction to using that particular proof strategy}.
\begin{proposition}[Sobolev sufficient condition on a fixed index band]
\label{prop:sobolev-v5}
Fix a sufficiently large $k$ and let
$J_k=[a_k,b_{k+1}]$, with the barriers from
\cref{thm:block-sep-v3}.  If, for some $p>1$, some $X_k\ge2$ and
$K_k\in(0,\infty)$,
\begin{equation}\label{eq:uniform-derivative-v5}
 \sup_{X\ge X_k}\bigl\|(F_X^W)'\bigr\|_{L^p(J_k)}\le K_k,
\end{equation}
then, for every $X\ge X_k$,
\begin{equation}\label{eq:sobolev-spacing-v5}
 \boxed{\quad\Delta_k(X;W)\ge
 K_k^{-p/(p-1)}>0.\quad}
\end{equation}
\end{proposition}
\begin{proof}
For each cutoff the barrier values yield a passage within $J_k$
from the level $k-\tfrac12$ to $k+\tfrac12$.  Choose $s_0<s_1$ in
$J_k$ by the extremal construction in the proof of
\cref{prop:Lipschitz-v4}, so that the passage occurs without leaving
the open level band on $(s_0,s_1)$.  It is part of the
set counted by $\Delta_k$, and absolute continuity and H\"older's
inequality give
\begin{equation}
 1\le\int_{s_0}^{s_1}|(F_X^W)'(t)|\,dt
 \le K_k(s_1-s_0)^{1-1/p}
 \le K_k\Delta_k(X;W)^{1-1/p}.
\end{equation}
Rearranging proves the claimed bound.
\end{proof}
\begin{remark}[Why uniform derivatives are not the proposed solution]
The hypothesis \eqref{eq:uniform-derivative-v5} is much stronger
than the occupation condition of \cref{prop:occupation-v3} and is
\emph{not} a consequence of the mean-value bounds of this note.
Moreover, if an arithmetic family converges pointwise away from
the zero heights to a function with jumps on $J_k$, a uniform
$W^{1,p}(J_k)$ bound with $p>1$, together with the fixed MCOM
amplitude bounds, is impossible: the one-dimensional Morrey
inequality would make the family equicontinuous, and the uniform bound
$|F_X^W|\le|N_0|+H$ on $J_k$ then makes it compact in $C(J_k)$ by
Arzel\`a--Ascoli, whereas a continuous subsequential limit cannot agree
on both sides of a jump with the discontinuous pointwise limit.
Thus controlling all derivatives uniformly in the cutoff is not the
plausible path.  A more local prime-only occupation or correlated
\emph{increment} estimate is required; the obstruction is to this
strong regularity strategy, not to the taucm(s) programme itself.
\end{remark}

\subsection{Why the numerically accurate pair centers do not close the proof}
For clarity, define
\begin{equation}\label{eq:center-v5}
 c_k(X;W):=\tfrac12\bigl(\tau_k^{\Cms}(X;W)
 +\tau_{k+1}^{\Cms}(X;W)\bigr),\qquad
 s_k(X;W):=\Delta_k(X;W).
\end{equation}
Knowledge of $c_k$ and $s_k$ is equivalent to knowledge of the
ordered pair, but an accurate center does not bound $s_k$ from below:
$c_k\pm s_k/2$ may approach the same limit with $s_k\downarrow0$.
The accurate numerical centers in \cref{sec:observed-v4} are therefore
valuable for choosing sampling windows for the untruncated Euler
phase, not a proof of separation.  The quadratic inverse theorem
\cref{thm:identifiability-v4} and the prime-only occupation inequality
\eqref{eq:occup-v3} are two possible ways to attack the missing
arithmetic estimate.  Neither currently provides the uniform error
control for all indices or the fixed-index limit $X\to\infty$.

\paragraph{Precise research target.}
For every fixed $k$, find a statement proved directly for the
weighted prime phase $P_X^W$ and the MCOM-truncated sublevel measures
that yields one of the equivalent inequalities
\eqref{eq:noncollapse-v5}--\eqref{eq:coherence-v5}, preferably through
positive occupation of the level band or a controlled correlation of
consecutive arithmetic displacements.  Such a theorem must allow the
phase derivative to grow with $X$, respect the full weight support
$n\le X^2$, and not assume information about $\beta_k$ or the
absence of repeated $\gamma_k$.  Together with an independently
proved DBSRH$(W)$ it would give RH and simplicity; without DBSRH,
it does not by itself locate the genuine imaginary parts.

\subsection{A band theorem: the increments of a prime band at the RvM points are negatively correlated}\label{sec:v20-band}
The research target stated at the end of the preceding subsection asks for a controlled correlation of consecutive arithmetic displacements. This subsection proves a related statement for a \emph{band} of the weighted prime phase, evaluated at the RvM points. It uses no hypothesis on the zeros (neither RH nor simplicity) and no information on the complete phase. It does \emph{not} give the correlation of the complete spacings $\Delta_k(X;W)$; the precise scope is recorded in \cref{rem:v20-scope}.

\paragraph{Setting and local notation.}
In this subsection $N_0'(t)=\frac1{2\pi}\log\frac t{2\pi}$, $h\in\{0,1,2\}$ is a lag, $\omega$ is a real frequency, and $\mathsf G_k^{(h)}:=\gamma^{\mathrm{RvM}}_{k+h}-\gamma^{\mathrm{RvM}}_k$ (thus $\mathsf G_k^{(1)}=h_k^{\mathrm{RvM}}$ in \eqref{eq:hrvm-v5}). The symbol $\|y\|$ denotes the distance from $y\in\mathbb R$ to the nearest integer. Fix exponents $\tfrac14<a<b<\tfrac12$ and, for $T\ge T_0(a,b)$, put
\begin{equation}\label{eq:v20-band}
 \mathcal B_T:=\{p\ \text{prime}:\ T^a<p\le T^b\},\qquad
 \Phi_{a,b}(t;T):=-\frac1\pi\sum_{p\in\mathcal B_T}\frac{\sin(t\log p)}{\sqrt p},
\end{equation}
\begin{equation}\label{eq:v20-Y}
 \mathcal K_T:=\{k:\ T\le\gamma^{\mathrm{RvM}}_k\le2T\},\qquad
 Y_k:=\Phi_{a,b}(\gamma^{\mathrm{RvM}}_{k+1};T)-\Phi_{a,b}(\gamma^{\mathrm{RvM}}_k;T).
\end{equation}
Write $\langle x_k\rangle_T:=(\#\mathcal K_T)^{-1}\sum_{k\in\mathcal K_T}x_k$, and let $\operatorname{Corr}_T(x_k,y_k)$ be the Pearson coefficient over $\mathcal K_T$,
\begin{equation}
 \operatorname{Corr}_T(x_k,y_k)=\frac{\langle x_ky_k\rangle_T-\langle x_k\rangle_T\langle y_k\rangle_T}
 {\sqrt{(\langle x_k^2\rangle_T-\langle x_k\rangle_T^2)(\langle y_k^2\rangle_T-\langle y_k\rangle_T^2)}}.
\end{equation}
The terms of \eqref{eq:v20-band} are the terms $n=p$ of \eqref{eq:PhiW}, because $\Lambda(p)/\log p=1$, and $W(\log p/\log X)=1$ for $p\le X$. Hence, for every $W\in\mathcal W$ and every $X\ge T^b$, the function $\Phi_{a,b}(\cdot\,;T)$ is an exact sub-sum of $\Phi_X^W$; the compensator $\Pi_X^W$ plays no role in it.

\begin{theorem}[Anticorrelation of the increments of a prime band]\label{thm:v20-anticorr}
Let $\tfrac14<a<b<\tfrac12$ and $T\ge T_0(a,b)$. Then, with no hypothesis on the zeta function,
\begin{equation}\label{eq:v20-main}
 \operatorname{Corr}_T(Y_k,Y_{k+1})=\mathfrak r_{a,b}+O_{a,b}\Bigl(\frac1{\log T}\Bigr),\qquad
 \mathfrak r_{a,b}:=\frac{\displaystyle\int_a^b\cos(2\pi v)\,\bigl[1-\cos(2\pi v)\bigr]\frac{dv}{v}}
 {\displaystyle\int_a^b\bigl[1-\cos(2\pi v)\bigr]\frac{dv}{v}}<0 .
\end{equation}
More precisely, with $\mathfrak v_{a,b}:=\pi^{-2}\int_a^b(1-\cos 2\pi v)\,dv/v>0$ and $\mathfrak z_{a,b}:=\pi^{-2}\int_a^b\cos(2\pi v)(1-\cos 2\pi v)\,dv/v<0$,
\begin{equation}\label{eq:v20-var-cov}
 \langle Y_k\rangle_T=O\bigl(T^{b/2-1}\bigr),\qquad
 \langle Y_k^2\rangle_T=\mathfrak v_{a,b}+O\Bigl(\frac1{\log T}\Bigr),\qquad
 \langle Y_kY_{k+1}\rangle_T=\mathfrak z_{a,b}+O\Bigl(\frac1{\log T}\Bigr).
\end{equation}
\end{theorem}

\begin{remark}[Value for one band]\label{rem:v20-value}
For $(a,b)=(0.30,0.45)$ the right-hand side of \eqref{eq:v20-main} is $\mathfrak r_{a,b}=-0.6811226\ldots$. This value is quoted, not recomputed in this note, and it is not used in any proof.
\end{remark}

\begin{lemma}[Spacings of the RvM points]\label{lem:v20-spacing}
Let $T\ge T_0$, $k\in\mathcal K_T$, $h\in\{0,1,2\}$, and for $p\in\mathcal B_T$ put $v_p:=\log p/\log T\in(a,b]$. Then, uniformly,
\begin{enumerate}[label=(\roman*),leftmargin=*]
\item $\mathsf G_k^{(h)}=\dfrac{2\pi h}{\log T}+O_h\Bigl(\dfrac1{\log^2T}\Bigr)$ and $\mathsf G_k^{(h)}\log p=2\pi h\,v_p+O_h\Bigl(\dfrac1{\log T}\Bigr)$;
\item $k\mapsto\mathsf G_k^{(h)}$ is nonincreasing and $\mathsf G_{k_1}^{(h)}\log p-\mathsf G_{k_2}^{(h)}\log p=O_h(1/\log T)$ for $k_1,k_2\in\mathcal K_T$ and $p\in\mathcal B_T$;
\item $\#\mathcal K_T=\dfrac T{2\pi}\log T\,\bigl(1+O(1/\log T)\bigr)$.
\end{enumerate}
\end{lemma}
\begin{proof}
Since $N_0(\gamma^{\mathrm{RvM}}_{k+h})-N_0(\gamma^{\mathrm{RvM}}_k)=h$, the mean value theorem gives $\xi_k\in(\gamma^{\mathrm{RvM}}_k,\gamma^{\mathrm{RvM}}_{k+h})$ with $\mathsf G_k^{(h)}=h/N_0'(\xi_k)=2\pi h/\log(\xi_k/2\pi)$. For $T\ge T_0$ the spacings are $<1$, so $\xi_k\in[T,2T+2]$ for every $k\in\mathcal K_T$ and $\log(\xi_k/2\pi)=\log T+O(1)$. This is the first formula of (i); the second follows by multiplication by $\log p=v_p\log T$ and $v_p\le b$. The map $g(x):=N_0^{-1}(x-\frac12)$ satisfies $g'=1/N_0'(g)>0$ and $g''<0$ because $N_0'$ is increasing; hence $g$ is concave and $\mathsf G_k^{(h)}=g(k+h)-g(k)$ is nonincreasing in $k$. Moreover, for every $k\in\mathcal K_T$,
\begin{equation}
 \frac{2\pi h}{\log((2T+2)/2\pi)}\le\mathsf G_k^{(h)}\le\frac{2\pi h}{\log(T/2\pi)},
\end{equation}
so for $k_1,k_2\in\mathcal K_T$
\begin{equation}
 \bigl|\mathsf G_{k_1}^{(h)}-\mathsf G_{k_2}^{(h)}\bigr|\le2\pi h\,\frac{\log(2+2/T)}{\log(T/2\pi)\,\log((2T+2)/2\pi)}\ll_h\frac1{\log^2T},
\end{equation}
and multiplication by $\log p\le b\log T$ gives (ii). Finally $\#\mathcal K_T=N_0(2T)-N_0(T)+O(1)=\frac T{2\pi}\bigl(\log\frac T{2\pi e}+2\log2\bigr)+O(1)$, which is (iii).
\end{proof}

\begin{lemma}[Exponential sums over the RvM points]\label{lem:v20-lattice}
There is an absolute constant $C_{\mathrm{KL}}$ with the following property. Let $T\ge T_0$, $j\in\{0,1,2\}$, $\omega\in\mathbb R\setminus\{0\}$ and $\varkappa\in(0,\frac12]$ be such that $\|\omega/\log(t/2\pi)\|\ge\varkappa$ for all $t\in[T,2T+2]$. Let $(\mathfrak t_k)_{k\in\mathcal K_T}$ be complex numbers of modulus one with total variation $\operatorname{TV}(\mathfrak t):=\sum_k|\mathfrak t_{k+1}-\mathfrak t_k|$. Then
\begin{equation}\label{eq:v20-lattice}
 \Bigl|\sum_{k\in\mathcal K_T}\mathfrak t_k\,e^{i\omega\gamma^{\mathrm{RvM}}_{k+j}}\Bigr|\le\frac{C_{\mathrm{KL}}\,(1+\operatorname{TV}(\mathfrak t))}{\varkappa}.
\end{equation}
\end{lemma}
\begin{proof}
Put $f(x):=\omega\,g(x+j)/2\pi$ with $g$ as in the proof of \cref{lem:v20-spacing}, so that $e^{i\omega\gamma^{\mathrm{RvM}}_{k+j}}=e(f(k))$ with $e(y)=e^{2\pi iy}$. Then $f'(x)=\omega/\log(g(x+j)/2\pi)$ is continuous and monotonic. Let $\mathcal K_T=\{k_1,\dots,k_2\}$. For $x\in[k_1,k_2]$ one has $g(x+j)\in[\gamma^{\mathrm{RvM}}_{k_1},\gamma^{\mathrm{RvM}}_{k_2+j}]\subset[T,2T+2]$ when $T\ge T_0$, hence $\|f'(x)\|\ge\varkappa$. The Kusmin--Landau inequality \cite[Thm.~2.1, p.~7]{GrahamKolesnik1991} (if $f$ is continuously differentiable, $f'$ is monotonic and $\|f'\|\ge\lambda>0$ on an interval $I$, then $\sum_{n\in I}e(f(n))\ll\lambda^{-1}$ with an absolute constant) applied on $I=[k_1,m]$ gives $|S_m|\le C_{\mathrm{KL}}/\varkappa$ for the partial sums $S_m:=\sum_{k_1\le k\le m}e(f(k))$, for every $m\in[k_1,k_2]$. Summation by parts,
\begin{equation}
 \sum_{k=k_1}^{k_2}c_ke(f(k))=c_{k_2}S_{k_2}-\sum_{m=k_1}^{k_2-1}(c_{m+1}-c_m)S_m,
\end{equation}
gives \eqref{eq:v20-lattice}.
\end{proof}

\begin{proof}[Proof of \cref{thm:v20-anticorr}]
Throughout, $T\ge T_0(a,b)$ and implied constants may depend on $a,b$. For $t\in[T,2T+2]$ one has $\log T-\log2\pi\le\log(t/2\pi)<\log T$; hence for every $\omega>0$
\begin{equation}\label{eq:v20-ratio}
 \frac{\omega}{\log T}\ \le\ \frac{\omega}{\log(t/2\pi)}\ \le\ \frac{\omega}{\log T}\Bigl(1+\frac C{\log T}\Bigr)\qquad(T\le t\le2T+2).
\end{equation}
Three ranges of $\omega$ are used in \cref{lem:v20-lattice}.

\emph{(A) $\omega=\log p$, $p\in\mathcal B_T$.} By \eqref{eq:v20-ratio} the ratio lies in $[a,\,b(1+C/\log T)]\subset(0,\frac12)$, so $\|\cdot\|\ge a$.

\emph{(B) $\omega=\log p+\log q$ with $p,q\in\mathcal B_T$ (including $p=q$).} The ratio lies in $[2a,\,2b(1+C/\log T)]\subset(\frac12,1-\varkappa_b]$ with $\varkappa_b:=\frac12(1-2b)$, so $\|\cdot\|\ge\varkappa_b$. \emph{The hypothesis $b<\frac12$ is used here, in (A), in (C) and in \eqref{eq:v20-offdiag}; $a>\frac14$ is used for the sign of $\mathfrak r_{a,b}$ (Step~4) and makes the ranges in (B) and (C) explicit.}

\emph{(C) $\omega=\log(q/p)$, $p<q$ in $\mathcal B_T$.} Here $0<\omega\le(b-a)\log T$ and $b-a<\frac14$, so the ratio lies in $(0,\frac12)$ and $\|\cdot\|\ge\omega/\log T$. Moreover $\omega=\log(q/p)\ge1-p/q=(q-p)/q\ge1/q\ge T^{-b}$.

\emph{Step 1: means.} By (A) with $\mathfrak t_k\equiv1$, $|\sum_{k\in\mathcal K_T}e^{i\log p\,\gamma^{\mathrm{RvM}}_{k+j}}|\le C_{\mathrm{KL}}/a$ for $j=0,1,2$. Since $\sum_{n\le x}n^{-1/2}\le2\sqrt x$,
\begin{equation}
 \Bigl|\sum_{k\in\mathcal K_T}\Phi_{a,b}(\gamma^{\mathrm{RvM}}_{k+j};T)\Bigr|\le\frac1\pi\sum_{p\in\mathcal B_T}\frac{C_{\mathrm{KL}}}{a\sqrt p}\ll T^{b/2},
\end{equation}
so $\langle\Phi_{a,b}(\gamma^{\mathrm{RvM}}_{k+j};T)\rangle_T=O(T^{b/2-1})$ by \cref{lem:v20-spacing}(iii).

\emph{Step 2: second moments.} Put $\Gamma_h(T):=\langle\Phi_{a,b}(\gamma^{\mathrm{RvM}}_k;T)\,\Phi_{a,b}(\gamma^{\mathrm{RvM}}_{k+h};T)\rangle_T$. Then
\begin{equation}\label{eq:v20-gamma-expand}
 \Gamma_h(T)=\frac1{\pi^2}\sum_{p,q\in\mathcal B_T}\frac1{\sqrt{pq}}\,\bigl\langle\sin(\gamma^{\mathrm{RvM}}_k\log p)\,\sin(\gamma^{\mathrm{RvM}}_{k+h}\log q)\bigr\rangle_T ,
\end{equation}
and $\sin x\sin y=\frac12[\cos(x-y)-\cos(x+y)]$ with $\gamma^{\mathrm{RvM}}_{k+h}=\gamma^{\mathrm{RvM}}_k+\mathsf G_k^{(h)}$ gives
\begin{multline}
 \sin(\gamma^{\mathrm{RvM}}_k\log p)\sin(\gamma^{\mathrm{RvM}}_{k+h}\log q)=\tfrac12\cos\bigl((\log p-\log q)\gamma^{\mathrm{RvM}}_k-\mathsf G_k^{(h)}\log q\bigr)\\
 -\tfrac12\cos\bigl((\log p+\log q)\gamma^{\mathrm{RvM}}_k+\mathsf G_k^{(h)}\log q\bigr).
\end{multline}
By \cref{lem:v20-spacing}(ii) and $|e^{i\phi}-e^{i\phi'}|\le|\phi-\phi'|$, the unimodular factors $\mathfrak t_k=e^{\mp i\mathsf G_k^{(h)}\log q}$ have $\operatorname{TV}(\mathfrak t)\ll1/\log T$, so \cref{lem:v20-lattice} applies to each of the two cosines.

\emph{Diagonal terms $p=q$.} The second cosine has $\omega=2\log p$, range (B), hence contributes $O(1/\#\mathcal K_T)$ for each $p$. The first is $\cos(\mathsf G_k^{(h)}\log p)=\cos(2\pi h\,v_p)+O(1/\log T)$ by \cref{lem:v20-spacing}(i), uniformly in $k\in\mathcal K_T$. Since $\sum_{p\in\mathcal B_T}1/p=O(1)$ and $\#\mathcal K_T\gg T$,
\begin{equation}\label{eq:v20-diag}
 \frac1{2\pi^2}\sum_{p\in\mathcal B_T}\frac{\cos(2\pi h\,v_p)}p+O\Bigl(\frac1{\log T}\Bigr).
\end{equation}
By Mertens' theorem \cite{Mertens1874} in the form $\mathcal S_1(x):=\sum_{p\le x}1/p=\log\log x+\mathrm B_1+E(x)$, $|E(x)|\le5/\log x$ \cite[Thm.~1.10, p.~18]{Tenenbaum2015}, and Stieltjes integration by parts with the $C^1$ function $G(v)=\cos(2\pi hv)$ on $[a,b]$,
\begin{equation}
 \sum_{p\in\mathcal B_T}\frac{G(v_p)}p=\int_{(a,b]}G(v)\,d\mathcal S_1(T^v)=\int_a^bG(v)\frac{dv}v+O_h\Bigl(\frac1{\log T}\Bigr),
\end{equation}
because $\mathcal S_1(T^v)=\log v+\log\log T+\mathrm B_1+O(1/(v\log T))$. Hence the diagonal contributes $\frac1{2\pi^2}\int_a^b\cos(2\pi hv)\,dv/v+O(1/\log T)$.

\emph{Off-diagonal terms $p\ne q$.} The second cosine is in range (B); each pair $(p,q)$ contributes $O\bigl(1/(\sqrt{pq}\,\#\mathcal K_T)\bigr)$, so all pairs together contribute $\ll T^b/\#\mathcal K_T$, which is absorbed in \eqref{eq:v20-offdiag}. The first has $\omega=\log(p/q)$, range (C) (after complex conjugation if $p>q$), and \cref{lem:v20-lattice} gives a bound $\ll T^b\log T$ for the sum over $k$. Therefore the off-diagonal part of \eqref{eq:v20-gamma-expand} is
\begin{equation}\label{eq:v20-offdiag}
 \ll\frac{T^b\log T}{\#\mathcal K_T}\Bigl(\sum_{p\le T^b}\frac1{\sqrt p}\Bigr)^2\ll\frac{T^b\log T\cdot T^b}{T\log T}=T^{2b-1}\ll\frac1{\log T},
\end{equation}
which is a further use of $b<\frac12$. Altogether
\begin{equation}\label{eq:v20-moments}
 \Gamma_h(T)=\frac1{2\pi^2}\int_a^b\cos(2\pi hv)\,\frac{dv}v+O\Bigl(\frac1{\log T}\Bigr)\qquad(h=0,1,2).
\end{equation}

\emph{Step 3: increments.} Since $|\Phi_{a,b}|\le\frac1\pi\sum_{p\le T^b}p^{-1/2}\le2T^{b/2}/\pi$, shifting the summation range by one or two indices changes any of the averages below by $O(T^b/\#\mathcal K_T)=O(T^{b-1})$. Hence $\langle Y_k\rangle_T=O(T^{b/2-1})$ by Step 1, and
\begin{align}
 \langle Y_k^2\rangle_T&=2\Gamma_0(T)-2\Gamma_1(T)+O(T^{b-1}),\\
 \langle Y_kY_{k+1}\rangle_T&=2\Gamma_1(T)-\Gamma_0(T)-\Gamma_2(T)+O(T^{b-1}).
\end{align}
Insert \eqref{eq:v20-moments}. The first limit is $\pi^{-2}\int_a^b(1-\cos 2\pi v)\,dv/v=\mathfrak v_{a,b}$. The second is $(2\pi^2)^{-1}\allowbreak\int_a^b[2\cos2\pi v-1-\cos4\pi v]\,dv/v$, and since $\cos4\pi v=2\cos^22\pi v-1$, the integrand equals $2\cos2\pi v\,(1-\cos2\pi v)$; the limit is $\mathfrak z_{a,b}$. This proves \eqref{eq:v20-var-cov}.

\emph{Step 4: the coefficient.} The empirical means are $O(T^{b/2-1})$, so the empirical covariance and variances differ from $\langle Y_kY_{k+1}\rangle_T$ and $\langle Y_k^2\rangle_T$, $\langle Y_{k+1}^2\rangle_T$ by $O(T^{b-2})$. As $\mathfrak v_{a,b}>0$ ($1-\cos2\pi v>0$ almost everywhere) and depends only on $a,b$, the quotient is $\mathfrak z_{a,b}/\mathfrak v_{a,b}+O(1/\log T)=\mathfrak r_{a,b}+O(1/\log T)$. Finally, for $v\in(a,b)\subset(\frac14,\frac12)$ one has $\cos2\pi v<0$ and $1-\cos2\pi v>0$, so the numerator integrand is negative and $\mathfrak r_{a,b}<0$. The hypothesis $a>\frac14$ is used for this sign and in ranges (B) and (C).
\end{proof}

\begin{remark}[Scope of \cref{thm:v20-anticorr}]\label{rem:v20-scope}
\begin{enumerate}[label=(\roman*),leftmargin=*]
\item The theorem is about the band \eqref{eq:v20-band}, an exact sub-sum of $\Phi_X^W$ for $X\ge T^b$, at the points $\gamma_k^{\mathrm{RvM}}$. It uses no zero, no RH and no property of the complete phase. It is not a statement about the sublevel points $\tau_k^{\Cms}(X;W)$.
\item By \eqref{eq:correlated-v5}, noncollapse of the complete spacings is decided by $\varepsilon_{k+1}^{\mathrm{RvM}}-\varepsilon_k^{\mathrm{RvM}}$, where $\varepsilon_k^{\mathrm{RvM}}=\tau_k^{\Cms}-\gamma_k^{\mathrm{RvM}}$. The increments $Y_k$ are the contribution of one band to the difference of the phase at consecutive RvM points; they are not that difference. The phase increment at the RvM points and the displacements $\varepsilon_k^{\mathrm{RvM}}$ are related through the level equation, which is nonlinear in the phase; replacing one by the other is a linearization whose error is governed by the size of $(P_X^W)'/N_0'$, and no such replacement is made or claimed here. Hence the theorem does not give the limit of $\operatorname{Corr}(\Delta_k,\Delta_{k+1})$ for the complete points.
\item The theorem does not determine how the contributions of different bands combine when the complete phase is assembled: the primes with $p\le T^{1/4}$, those with $T^{1/2}\le p\le X$, and the weighted range $1<\log n/\log X\le2$ are not covered. In particular the correlation of the complete spacings is not determined by \eqref{eq:v20-main}.
\item The condition $b<\frac12$ enters through ranges (A) and (B) (in (B) the quantity $\log(pq)/\log(t/2\pi)$ stays at distance $\ge\varkappa_b$ from the integer $1$), through (C), and through the exponent in \eqref{eq:v20-offdiag}. For $b\ge\frac12$ the bound of \cref{lem:v20-lattice} is lost for these frequencies and the argument does not extend.
\end{enumerate}
\end{remark}

\begin{remark}[Two different facts, and a diagonal heuristic for combining bands]\label{rem:v21-two-facts}
\begin{enumerate}[label=(\roman*),leftmargin=*]
\item \emph{Numerical observation for the complete points.} The author's computation gives, for the complete points, $\operatorname{Corr}(s_k,s_{k+1})=-0.3835683$, with $s_k=\Delta_k(X;W)$ as in \eqref{eq:center-v5}. This value is quoted, not recomputed here; the cutoff $X$, the weight and the window of indices are not recorded in this note. The value $-0.6811226\ldots$ of \cref{rem:v20-value} is the limit proved in \cref{thm:v20-anticorr} for the increments of one band at the RvM points. The two numbers refer to different objects, and neither determines the other.
\item \emph{Diagonal heuristic (not a theorem).} Write $\log X=\theta\log T$ and $w(v):=W(v/\theta)$ for the weight at $n=T^v$. Keep only the prime terms and, in the second moments of the increments of $\Phi_X^W$ at the RvM points, only the diagonal terms $n=m$, as in \eqref{eq:v20-diag}; use Mertens' theorem as there. The computation of the proof of \cref{thm:v20-anticorr} then gives, for the correlation of consecutive increments of the whole phase,
\begin{equation}\label{eq:v21-diag}
 \mathfrak r^{\mathrm{diag}}_{W,\theta}:=\frac{\displaystyle\int_0^{2\theta}w(v)^2\cos(2\pi v)\,\bigl[1-\cos(2\pi v)\bigr]\frac{dv}v}
 {\displaystyle\int_0^{2\theta}w(v)^2\,\bigl[1-\cos(2\pi v)\bigr]\frac{dv}v}.
\end{equation}
For $w^2=\ind_{(a,b]}$ this is $\mathfrak r_{a,b}$ of \eqref{eq:v20-main}. For $w=\ind_{(0,1]}$ it equals $\tfrac12\operatorname{Cin}(4\pi)/\operatorname{Cin}(2\pi)-1$, where $\operatorname{Cin}(x)=\int_0^x(1-\cos u)\,du/u$; it has not been evaluated here. In a linear model, \eqref{eq:correlated-v5} with $\varepsilon_k^{\mathrm{RvM}}\approx-P_X^W(\gamma_k^{\mathrm{RvM}})/N_0'(\gamma_k^{\mathrm{RvM}})$ and the near-constancy of $h_k^{\mathrm{RvM}}$ would make $\mathfrak r^{\mathrm{diag}}_{W,\theta}$ the limit of $\operatorname{Corr}(\Delta_k,\Delta_{k+1})$. Nothing in this item is proved. For $v>\tfrac12$ the pairs $(n,m)$ with $nm$ near $t/2\pi$ have a stationary point in the sums of \cref{lem:v20-lattice} (the barrier at $b=\frac12$ of \cref{rem:v20-scope}(iv)), and the omitted terms are not shown to be negligible; the prime powers and the compensator $\Pi_X^W$ are also left out; and the linearization is not justified, by (iii).
\item \emph{The linearization is not small.} The same diagonal computation for the derivative, averaged over $t\in[T,2T]$, gives, to leading order in $\log T$,
\begin{equation}\label{eq:v21-slope}
 \Bigl\langle\bigl(\Phi_X^{W\prime}/N_0'\bigr)^2\Bigr\rangle\approx2\int_0^{2\theta}w(v)^2\,v\,dv .
\end{equation}
For one band $T^a<p\le T^b$ the right-hand side is $b^2-a^2$, a root-mean-square of $0.335\ldots$ for $(a,b)=(0.30,0.45)$. For $\theta=1$ and the Selberg weight ($q\equiv1$, so $w=1$ on $(0,1]$ and $w(v)=2-v$ on $(1,2]$) it is $2\bigl(\tfrac12+\tfrac5{12}\bigr)=\tfrac{11}6$, a root-mean-square of $1.35\ldots$. Since this is of order one, the first-order replacement in (ii) is not automatic, and $\mathfrak r^{\mathrm{diag}}_{W,\theta}$ need not equal the limit of $\operatorname{Corr}(\Delta_k,\Delta_{k+1})$; in particular nothing here predicts $-0.3835683$. The status of \eqref{eq:v21-slope} is the same as that of (ii).
\end{enumerate}
\end{remark}

% ======================================================================
% V7 ADDITIONS.  All results below distinguish proved statements,
% consequences depending on the cited IJNT input, and open reductions.
% ======================================================================
\section{Arithmetic scope and source attribution}\label{sec:v7-scope}
This section integrates the deductions developed below
with the author's public working notes at commit
\texttt{38aac7c} \cite{CasanaV10Notes,CasanaV10Transport}.
The sources are not identical in status. The earlier taucs
$\tau_k^{\mathrm C}$ and taucm $\tau_k^{\Cm}$ use the author's
interpolated \emph{first-prime Euler phase} $\Phi^\circ$; the taucms
$\tau_k^{\Cms}(X;W)$ use a \emph{different, compensated, weighted}
prime-power phase. In particular the weight $W_{1/4}$ is a chosen
external weight and is \emph{not} presented as an invention
of the author. The elementary sublevel, projection, Fourier,
complex-analysis and moment inequalities used below are mathematical
techniques, not separately claimed discoveries. The possible novelty
lies in the combination with the author's ordinal objects and the
specific MCOM coefficient. None of the present deductions is a
proof of RH.

Three different assertions must remain separate:
\begin{enumerate}[label=(\roman*),leftmargin=*]
\item strict positivity of $\tau_{k+1}^{\Cms}(X;W)-\tau_k^{\Cms}(X;W)$
for each finite $X$;
\item the noncollapse condition $\liminf_{X\to\infty}
[\tau_{k+1}^{\Cms}(X;W)-\tau_k^{\Cms}(X;W)]>0$ for a fixed index;
\item convergence of the \emph{same} sublevel points to the
corresponding imaginary parts of zeros (DBSRH).
\end{enumerate}
For every sufficiently large index the first follows from the
continuity and coercive MCOM barriers; for the initial indices
it requires checking that the corresponding two levels are
attained. The
second plus the third, for all indices, would imply RH and simplicity;
none follows from a bound for $S(t)$ alone.

\section{The genuine MCOM margin, and the distinction from moment inputs}
\label{sec:v7-mcom}
Write $A_{\mathrm M}=0.08382$, $C_{\mathrm H}=0.0748596142710907\ldots$
and keep the \emph{fixed, effective} truncation $H$ of
\eqref{eq:MCOMH}; the asymptotic coefficient is not silently substituted
into that definition. For $t$ outside zero heights write
$G(t)=N(t)-N_0(t)$, with the right-continuous convention at the heights.

\begin{proposition}[MCOM interior margin]\label{prop:v7-margin}
Assume the two inequalities of the author's MCOM manuscript,
\eqref{eq:MCOM} and \eqref{eq:asymptMCOM}. For every fixed
$0<\eta<A_{\mathrm M}-C_{\mathrm H}$ there exists $T_\eta$ such that
\begin{equation}\label{eq:v7-margin}
 H(t)-|G(t)|\ge\eta\log t\qquad(t\ge T_\eta)
\end{equation}
away from zero heights and, by right limits, at the heights.
\end{proposition}
\begin{proof}
Choose $c$ strictly between $C_{\mathrm H}$ and
$A_{\mathrm M}-\eta$. The asserted limsup implies
$|S(t)|\le c\log t$ eventually. Since $G=S+R$ and $R(t)=O(t^{-1})$,
choose a slightly larger intermediate coefficient, still below
$A_{\mathrm M}-\eta$, to absorb $R(t)$. The $5.5$ in $H$ is
nonnegative, and \eqref{eq:v7-margin} follows.
\end{proof}

The margin is a mathematically concrete use of the \emph{two}
MCOM estimates. It is not, by itself, a new Selberg--Fujii moment
bound. The inequality $|G|\le H$ also supplies the exact projection
identity, for any real prime phase $P$,
\begin{equation}\label{eq:v7-exact-projection}
 |P-G|-|\clip_H(P)-G|=(|P|-H)_+.
\end{equation}
The right side is zero on nonsaturated points and equals the distance
removed by the truncation on saturated points. The explicit MCOM
bound makes the target $G$ lie inside the interval $[-H,H]$; the
asymptotic refinement gives the stronger margin
\eqref{eq:v7-margin}. Neither property implies pointwise
convergence of the prime phases.

\begin{proposition}[Crossing budget localized by the MCOM barriers]
\label{prop:v7-crossing}
Let $F_X^W$ be the actual taucms function from \eqref{eq:Cms},
$P_X^W$ its \emph{untruncated} phase, and
$\lambda_k=k-\tfrac12$. For every $k$ sufficiently large, let
$a_k<b_k$ be the MCOM barriers specified in
\cref{thm:block-sep-v3} by
$N_0(a_k)+H(a_k)=\lambda_k$ and
$N_0(b_k)-H(b_k)=\lambda_k$.
Then, for every $X\ge2$ and $W\in\mathcal W$,
\begin{equation}\label{eq:v7-crossing}
 |\tau_k^{\Cms}(X;W)-\gamma_k|
 \le \meas\bigl\{t\in[a_k,b_k]:
 |P_X^W(t)-G(t)|\ge\tfrac12\bigr\}.
\end{equation}
The same statement holds for $\tau_k^{\Cm}$ with $P_X^W$
replaced by $\Phi^\circ$.
\end{proposition}
\begin{proof}
The sublevel indicator of $N(t)$ at $\lambda_k$ equals
$\mathbf1_{\{t<\gamma_k\}}$ a.e., even for repeated imaginary parts,
because the zeros are counted with multiplicity. For any integer $n$
and real $x$, differing indicators of $x<k-\frac12$ and
$n<k-\frac12$ imply $|x-n|\ge\frac12$.
The MCOM bound $|G|\le H$ and projection contractivity yield
$|F_X^W-N|\le|P_X^W-G|$. Finally, the MCOM barrier inequalities
place both sublevel boundaries in $[a_k,b_k]$ up to null sets;
outside this interval the two indicators agree. Integration
proves the claim. The constants here are \emph{not} claimed to
be sharp enough to force the right side to vanish.
\end{proof}

\begin{remark}[What the $1/2$ crossing does and does not do]
The asymptotic MCOM coefficient yields
$\limsup_{k\to\infty}|\gamma_k-\gamma_k^{\mathrm{RvM}}|
\le2\pi C_{\mathrm H}<1/2$ under the MCOM assertion.
This strict numerical threshold is not supplied by the earlier
coefficients discussed in the source notes. It is a proximity
statement in \emph{physical units of the imaginary part}; it is not a bound by
half the RvM spacing, which is $\asymp1/\log k$ and tends to zero.
The equality of the limsup with $2\pi$ times the asymptotic coefficient, quoted in the exploratory
notes, is \emph{not} used: the necessary consequence here is the
inequality $\le$. Neither this threshold nor
\eqref{eq:v7-crossing} excludes repeated imaginary parts on its own.
\end{remark}

\section{Quantitative ordinal transport under the cited fourth moment}
\label{sec:v7-transport}
The following is an audited formulation of the transport statement
first written in the conversation and saved in
\cite{CasanaV10Transport}. It does not improve the fourth moment
of Selberg or Fujii. It improves the \emph{consequence for the
ordinal points} relative to what the moment statement alone says.

\begin{definition}[Auxiliary untruncated sublevels]
\label{def:v7-untruncated}
Define, solely for comparison,
\begin{equation}\label{eq:v7-untruncated}
 u_k^\circ:=\meas\{t>0:N_0(t)+\Phi^\circ(t)<k-\tfrac12\}.
\end{equation}
These are neither the original taucs nor the taucm(s). Let
$\mathcal K_T^{\mathrm{RvM}}=\{k:T\le\gamma_k^{\mathrm{RvM}}\le2T\}$.
\end{definition}

\begin{theorem}[Transport of an MCOM-clipped Euler phase]
\label{thm:v7-transport}
Assume the untruncated fourth-moment input \eqref{eq:IJNTinput}
from the author's IJNT manuscript and both MCOM estimates
\eqref{eq:MCOM} and \eqref{eq:asymptMCOM}. Then, as $T\to\infty$,
\begin{align}
 \sum_{k\in\mathcal K_T^{\mathrm{RvM}}}|\tau_k^{\Cm}-u_k^\circ|
 &\ll\frac{T}{(\log T)^3},\label{eq:v7-transport}\\
 \frac1{\#\mathcal K_T^{\mathrm{RvM}}}\sum_{k\in\mathcal K_T^{\mathrm{RvM}}}
 |\tau_k^{\Cm}-u_k^\circ|
 &\ll\frac1{(\log T)^4}.\label{eq:v7-transport-avg}
\end{align}
For fixed $\epsilon>0$, the fraction of indices $k\in\mathcal K_T^{\mathrm{RvM}}$ satisfying
$|\tau_k^{\Cm}-u_k^\circ|>
\epsilon(\gamma_{k+1}^{\mathrm{RvM}}-
\gamma_k^{\mathrm{RvM}})$ is
$O_\epsilon((\log T)^{-3})$.
The analogous total $L^1$ bound for the difference of consecutive
spacings holds with both indices in $\mathcal K_T^{\mathrm{RvM}}$.
\end{theorem}
\begin{proof}
Let $P(t)=\Phi^\circ(t)$ and
$D(t)=(|P(t)|-H(t))_+$.
By \cref{prop:v7-margin}, on $[T/2,3T]$, for $T$ sufficiently
large, the set $E_T=\{t\in[T/2,3T]:D(t)>0\}$ satisfies
$|P-G|\ge D+c\log T$, with $c>0$ fixed. Thus
\begin{equation}\label{eq:v7-local-D}
 \meas(E_T)\ll\frac{1}{(\log T)^4}
 \int_{T/2}^{3T}|P-G|^4dt\ll\frac{T}{(\log T)^4},
 \quad
 \int_{T/2}^{3T}D\,dt\ll\frac{T}{(\log T)^3}.
\end{equation}
The last inequality follows from
$D\le |P-G|^4/(c\log T)^3$ on $E_T$.
The cited IJNT fourth moment applies on the requisite finite
collection of dyadic blocks.

To justify localization independently of clipping, expand the finite
Euler logarithms and sum the elementary bounds:
$|\Phi^\circ(t)|\ll (1+t)^{1/4}\sqrt{\log(3+t)}$.
Consequently $N_0+P$ and $N_0+\clip_H P$ are both below every
level with $k\in\mathcal K_T^{\mathrm{RvM}}$ when $t\le T/2$ and above it when
$t\ge3T$, for all large $T$. The initial compact interval is
absorbed by the same asymptotic comparison.

For fixed $t$ the sublevel indicators for these two functions
can disagree only at levels $k-\frac12$ within an interval of
length $D(t)$. If $D(t)>0$, that interval contains at most
$D(t)+1$ such half-integer levels. Therefore
\begin{align}
\sum_{k\in\mathcal K_T^{\mathrm{RvM}}}|\tau_k^{\Cm}-u_k^\circ|
&\le\int_{T/2}^{3T}
\bigl[D(t)+\mathbf1_{E_T}(t)\bigr]dt
\ll T(\log T)^{-3}.
\end{align}
The cardinality of $\mathcal K_T^{\mathrm{RvM}}$ is $\asymp T\log T$.
The pointwise triangle inequality for two consecutive
spacings proves their $L^1$ comparison; Markov's inequality and
$\gamma_{k+1}^{\mathrm{RvM}}-
\gamma_k^{\mathrm{RvM}}\asymp(\log T)^{-1}$
give the stated exceptional proportion.
\end{proof}

\begin{corollary}[Uniform transport over the proved weighted cutoff range]
\label{cor:v7-weighted-transport}
Assume the uniform fourth-moment bound \eqref{eq:untruncatedmoderate} for the untruncated phase, proved in
\cref{sec:moderate} under the Fujii input \eqref{eq:FujiiA}, and the
MCOM estimates. For a fixed $0<\varepsilon<1/4$ define the
\emph{auxiliary} untruncated points
\begin{equation}
 u_k^\circ(X;W)=\meas\{t>0:N_0(t)+P_X^W(t)<k-\tfrac12\}.
\end{equation}
For $T\ge T_0(\varepsilon)$, uniformly for $W\in\mathcal W$ and
$T^\varepsilon\le X\le T^{1/4}$,
\begin{equation}\label{eq:v7-weighted-transport}
 \sum_{k\in\mathcal K_T^{\mathrm{RvM}}}
 |\tau_k^{\Cms}(X;W)-u_k^\circ(X;W)|
 \ll_\varepsilon\frac{T}{(\log T)^3}.
\end{equation}
\end{corollary}
\begin{proof}
The uniform untruncated bound \eqref{eq:untruncatedmoderate} replaces the fixed $\Phi^\circ$
input in \eqref{eq:v7-local-D}. The proof of \cref{thm:moderate} uses $X$ only through $X\ge T^\varepsilon$ and $X^4\ll T$, so \eqref{eq:untruncatedmoderate} holds, with a larger constant, on each of the blocks $[T/2,T]$, $[T,2T]$, $[2T,4T]$, which cover $[T/2,3T]$. The finite weighted Euler sum
has the absolute bound $O(X)$, while the pole compensator has the
bound $O(X/t)$ for $t\ge1$; at $0<t<1$ it has the trivial bound
$O(X)$ as well. For $X\le T^{1/4}$ these bounds are
$o(T\log T)$, enough to localize the relevant sublevel
indicators to $[T/2,3T]$. The preceding level-counting argument
then gives the bound uniformly.
\end{proof}

\begin{remark}[No improvement of Selberg--Fujii itself]
Neither \eqref{eq:v7-transport} nor
\eqref{eq:v7-weighted-transport} improves a fourth-moment
exponent, constant or allowed prime cutoff in Selberg--Fujii.
Further, the \emph{rate} $T/(\log T)^3$ follows from a positive
logarithmic margin and the cited fourth moment; once those inputs
are assumed, that elementary step does not uniquely require the
specific decimal $0.08382$. The particular MCOM estimates are
what justify the margin and the covered projection for the
\emph{specified} taucm construction. The auxiliary points
$u_k^\circ$ are \emph{not} the true imaginary parts, and these results
do not imply noncollapse for a fixed $k$ as $X\to\infty$.
\end{remark}

\begin{remark}[The asymptotic estimate \eqref{eq:asymptMCOM} is not needed for the rate]
\label{rem:v8-transport-noasympt}
The conclusions of \cref{thm:v7-transport} and
\cref{cor:v7-weighted-transport} hold with the explicit bound
\eqref{eq:MCOM} alone, provided Selberg's unconditional moment
estimate $\int_{T/2}^{3T}|S(t)|^{6}\,dt\ll T(\log\log T)^3$
\cite{Selberg1946} is added. Split $[T/2,3T]$ into
$B_1=\{|G|\le H/2\}$ and $B_2=\{|G|>H/2\}$. On $E_T\cap B_1$ one has
$|P-G|\ge D+H/2\ge D+\tfrac12A_{\mathrm M}\log(T/2)$, and the Markov
step of the proof applies verbatim with $c=\tfrac12A_{\mathrm M}$.
On $B_2$, $|S|\ge H/2-|R|\ge\tfrac12A_{\mathrm M}\log T+2$ for large
$T$, so $\meas(B_2)\ll T(\log\log T)^3(\log T)^{-6}$. Since $|G|\le H$,
$D\le|P-G|$ pointwise, and H\"older's inequality with the fourth
moment gives
$\int_{B_2}D\,dt\le\bigl(\int|P-G|^4\bigr)^{1/4}\meas(B_2)^{3/4}
\ll T(\log\log T)^{9/4}(\log T)^{-9/2}\ll T(\log T)^{-3}$,
while $\meas(E_T\cap B_2)\le\meas(B_2)\ll T(\log T)^{-4}$.
Thus the only property of the truncation used by the transport rate
is $H\asymp\log t$ together with the covering $|G|\le H$; the
asymptotic coefficient $C_{\mathrm H}$ removes the factor
$(\log\log T)^{O(1)}$ from the exceptional set but does not change
the order $T(\log T)^{-3}$.
\end{remark}

\section{A prime-by-prime inequality and its exact limitation}
\label{sec:v7-prime-step}
For the original Euler phase \eqref{eq:Phi-m}, adding a prime
$p$ changes the finite phase by
\begin{equation}\label{eq:v7-single-prime}
 \phi_p(t)=-\frac1\pi\arg(1-p^{-1/2-it}),\qquad
 |\phi_p(t)|\le\frac{\arcsin(p^{-1/2})}{\pi}=:\alpha_p.
\end{equation}
The bound is elementary geometry of the circle centered at $1$
of radius $p^{-1/2}$. It is independent of MCOM; MCOM enters only
when these phases are used with the fixed, covered truncation $H$.
If an \emph{auxiliary} phase at a fixed number of primes is called
$B_m(t)=N_0(t)+\clip_{H(t)}\Phi_m(t)$, then
$|B_{m+1}-B_m|\le\alpha_{p_{m+1}}$. Consequently
\begin{equation}\label{eq:v7-persistence-one-step}
\begin{split}
&\meas\{k-\tfrac12\le B_{m+1}<k+\tfrac12\}\\
&\quad\ge \meas\{k-\tfrac12+\alpha_{p_{m+1}}
\le B_m<k+\tfrac12-\alpha_{p_{m+1}}\}.
\end{split}
\end{equation}
This is a finite-step interior-band inclusion, not persistence
for an infinite sequence of prime cutoffs. The bounds
$\alpha_p\sim1/(\pi\sqrt p)$ are not summable over primes.
The structural saturation model of \cref{sec:saturation} also
shows why pointwise clipping and diminishing step sizes cannot
alone rule out a band collapsing in a limit. The compensated
weighted family $P_X^W$ is \emph{not} the same as the discrete
family $\Phi_m$, and \eqref{eq:v7-persistence-one-step} is not
claimed as a theorem about its increments.

\section{A zero-free Euler transform: a conditional RH criterion}
\label{sec:v7-transform}
This alternative route refines the transform idea in
\cite{CasanaV10Notes}. It uses the \emph{actual compensated
weighted prime phase} $P_X^W$, not a new approximation or a
relabeling of $S(t)$. The transform itself is elementary
Fourier--Laplace analysis; only its application to this ordinal
program is under discussion. This is an \emph{alternative
analytic route}, not an established prime-band persistence
result, and not a proof of RH.

Fix $W\in\mathcal W$ once and for all. For $\Re s>1/2$ put
\begin{equation}\label{eq:v7-ks}
 k_s(t)=\frac1{(s-\frac12-it)^2}+
        \frac1{(s-\frac12+it)^2},\qquad
 \mathcal T[f](s)=-\int_0^\infty f(t)\partial_t k_s(t)\,dt.
\end{equation}
For every compact $K\Subset\{\Re s>1/2\}$,
$|\partial_t k_s(t)|\le C_K(1+t)^{-3}$,
so the transform is holomorphic whenever
$\int_0^\infty|f(t)|(1+t)^{-3}dt<\infty$.

\begin{lemma}[Exact finite-cutoff Euler transform]
\label{lem:v7-euler-transform}
For $X\ge2$ and $\Re s>1/2$,
\begin{equation}\label{eq:v7-euler-transform}
\mathcal T[P_X^W](s)
=-\sum_{n\ge2} W\!\left(\frac{\log n}{\log X}\right)
\frac{\Lambda(n)\log n}{n^s}
+\int_2^\infty W\!\left(\frac{\log u}{\log X}\right)
\frac{\log u}{u^s}\,du.
\end{equation}
\end{lemma}
\begin{proof}
The sum and integral defining $P_X^W$ have compact support
in their $n,u$ variables for fixed $X$. For $a>0$ and
$\Re(s-\frac12)>0$, Fourier--Laplace inversion gives
$\int_0^\infty k_s(t)\cos(at)dt=\pi a e^{-(s-\frac12)a}$.
Apply integration by parts to each term of $P_X^W$; every
sine term vanishes at $t=0$ and its derivative contributes
$-a\cos(at)/\pi$. Set $a=\log n$ or $\log u$.
\end{proof}

\begin{theorem}[Weighted $L^2$ normality implies RH]
\label{thm:v7-normality}
Fix a weight $W\in\mathcal W$. Suppose that for some unbounded
sequence of cuts $X_j\to\infty$,
\begin{equation}\label{eq:v7-normality-assumption}
\sup_j \int_0^\infty
\frac{|P_{X_j}^W(t)|^2}{(1+t)^3}\,dt<\infty.
\end{equation}
Then all nontrivial zeros of $\zeta$ lie on $\Re s=\tfrac12$.
No Cauchy-in-measure hypothesis is needed.
\end{theorem}
\begin{proof}
Cauchy--Schwarz and the bound following \eqref{eq:v7-ks}
show that the holomorphic family
$\{\mathcal T[P_{X_j}^W]\}_j$ is locally bounded on
$\Re s>1/2$. It is therefore normal by Montel's theorem.
For $\Re s>1$, dominated convergence in
\eqref{eq:v7-euler-transform}, using $0\le W\le1$ and
$W(v)=1$ for $v\le1$, gives the pointwise limit
\begin{equation}\label{eq:v7-meromorphic-limit}
 Q(s):=-\left(\frac{\zeta'}\zeta\right)'(s)
 +\int_2^\infty\frac{\log u}{u^s}\,du,
 \qquad\Re s>1.
\end{equation}
The integral equals
$2^{1-s}\bigl(\log2/(s-1)+(s-1)^{-2}\bigr)$;
its double pole at $s=1$ cancels the double pole of the
first summand. Thus $Q$ admits a meromorphic continuation
to $\Re s>1/2$, with a nonremovable double pole at any
zero $\rho$ in that half-plane. Any locally uniform
subsequential limit provided by Montel is holomorphic on the
whole half-plane and agrees with $Q$ for $\Re s>1$.
Uniqueness of analytic continuation therefore forbids
such a $\rho$. Hence $\zeta(s)\neq0$ for $\Re s>1/2$ (no trivial zero lies there). If a nontrivial zero $\rho$ had $\Re\rho<1/2$, then $1-\overline\rho$, a nontrivial zero by the functional equation and conjugation symmetry, would satisfy $\Re(1-\overline\rho)>1/2$, a contradiction. Hence every nontrivial zero has real part $1/2$.
\end{proof}

\begin{corollary}[MCOM excess criterion: an exact sufficient route]
\label{cor:v7-excess}
For a fixed $W\in\mathcal W$ and the fixed effective MCOM truncation $H(t)$,
let the \emph{untruncated excess} be
\begin{equation}\label{eq:v7-excess}
 E_X(t)=\bigl(|P_X^W(t)|-H(t)\bigr)_+.
\end{equation}
If for some $X_j\to\infty$
\begin{equation}\label{eq:v7-excess-criterion}
 \sup_j\int_0^\infty\frac{E_{X_j}(t)^2}{(1+t)^3}\,dt<\infty,
\end{equation}
then RH holds. Contrapositively, if RH is false, then for
every fixed $W\in\mathcal W$,
\begin{equation}\label{eq:v7-excess-diverges}
 \lim_{X\to\infty}\int_0^\infty
 \frac{\bigl(|P_X^W(t)|-H(t)\bigr)_+^2}{(1+t)^3}\,dt=+\infty.
\end{equation}
\end{corollary}
\begin{proof}
Pointwise, $|P_X^W|\le H+E_X$. As
$\int_0^\infty H(t)^2(1+t)^{-3}dt<\infty$, the triangle
inequality in weighted $L^2$ converts
\eqref{eq:v7-excess-criterion} to
\eqref{eq:v7-normality-assumption}. For the contrapositive,
if the limit did not diverge, there would be an unbounded
sequence $X_j$ with bounded excess norm, contradicting
the theorem.
\end{proof}

\begin{remark}[How MCOM enters and what remains open]
The last proof uses $H(t)$ as an explicit \emph{scale}, but
uses only $H\in L^2((1+t)^{-3}dt)$; its logical implication
to RH is therefore not a new theorem that relies numerically
on the coefficient $0.08382$. To exploit the \emph{sharpness}
of MCOM one must actually establish
\eqref{eq:v7-excess-criterion} from the arithmetic coefficients
or an MCOM-sensitive inequality. This has \emph{not} been done.
The normal-family technique is classical, and finite Euler products are not new in this context \cite{Gonek2012}; the finite Euler
identity \eqref{eq:v7-euler-transform} and the precise
weighted-excess condition are recorded to identify a
sharply formulated open task. There is no demonstrated
improvement over Selberg--Fujii, For the Selberg weight the converse ``RH
implies \eqref{eq:v7-excess-criterion}'' holds by \cref{thm:taucw-equiv}, since $E_X\le|P^W_X|$; for the other weights of \eqref{eq:weights} it is not established.
\end{remark}

\section{Identification of the limit phase, and the first index}
\label{sec:v8-limit}
The criterion of \cref{thm:v7-normality} uses no convergence of the
prime phases. If convergence in measure is added, the limit phase is
identified exactly, and the first sublevel point of \eqref{eq:Cms} is
seen \emph{not} to converge to $\gamma_1$. The external inputs of this
section are Malmst\'en's integral for $\log\Gamma$
\cite[\S12.31]{WhittakerWatson}, the boundary formula for Cauchy
integrals of $L^2$ functions \cite[Thm.~93]{TitchmarshFI}, the Hadamard
product of $\xi$, and the explicit bound $|R(t)|\le0.53/(\pi t)$
already used in \eqref{eq:MCOMH}. Throughout, $N_\vartheta=\vartheta/\pi+1$,
so that $N=N_\vartheta+S$ for every $t>0$ by \eqref{eq:RVM}, $b:=s-\frac12$,
and $\psi=\Gamma'/\Gamma$.

\begin{lemma}[Transforms of the smooth terms]\label{lem:v8-smooth}
Let $\Re s>\frac12$.
\begin{enumerate}[label=(\roman*),leftmargin=*]
\item For $a>0$, $\mathcal T[\sin(a\,\cdot)](s)=\pi a^2e^{-ab}$;
 moreover $\mathcal T[1](s)=2/b^2$ and $\mathcal T[t](s)=0$.
\item $\mathcal T[N_\vartheta](s)=2/b^2-\tfrac14\psi'(s/2)$.
\item Put
\begin{equation}\label{eq:v8-r0}
 \mathrm J(t)=\int_0^{\log2}e^{v/2}\frac{\sin(tv)}{v}\,dv,\qquad
 r_0(t)=-1+\frac1\pi\arctan(2t)+\frac1\pi \mathrm J(t).
\end{equation}
Then $r_0$ is real-analytic, $r_0(0)=-1$, $\mathrm J(t)>0$ and
$|r_0(t)|\le1.46/t$ for $t>0$, and
\begin{equation}\label{eq:v8-E}
 \mathcal T[r_0](s)=E(s):=\frac1{s^2}-\frac2{b^2}
 +\int_1^2(\log u)\,u^{-s}\,du .
\end{equation}
\item For $\Re s>1$ the limit $Q$ of \eqref{eq:v7-meromorphic-limit}
 satisfies $Q(s)=-(\xi'/\xi)'(s)-E(s)-\mathcal T[N_\vartheta](s)$.
\end{enumerate}
\end{lemma}
\begin{proof}
(i) Integrating by parts,
$\mathcal T[\sin(a\,\cdot)](s)=a\int_0^\infty\cos(at)k_s(t)\,dt
=\frac a2\int_{\mathbb R}e^{iat}k_s(t)\,dt$, since $k_s$ is even.
For $a>0$ and $\Re b>0$, residues give
$\int_{\mathbb R}e^{iat}(b+it)^{-2}dt=2\pi ae^{-ab}$ and
$\int_{\mathbb R}e^{iat}(b-it)^{-2}dt=0$. Next,
$\mathcal T[1]=k_s(0)-k_s(\infty)=2/b^2$, and
$\mathcal T[t]=\int_0^\infty k_s\,dt=\frac12\int_{\mathbb R}k_s\,dt=0$
because each of $(b\pm it)^{-2}$ integrates to zero over $\mathbb R$.

(ii) Malmst\'en's formula gives, for $\Re z>0$,
$\log\Gamma(z)=\int_0^\infty\bigl[(z-1)e^{-u}
-\frac{e^{-u}-e^{-zu}}{1-e^{-u}}\bigr]\frac{du}u$. With
$z=\frac14+\frac{it}2$,
$\Im\log\Gamma(z)=\int_0^\infty\bigl[\frac t2e^{-u}
-\frac{e^{-u/4}\sin(tu/2)}{1-e^{-u}}\bigr]\frac{du}u$,
whose integrand is $O(t+ut^3)$ for $tu\le1$ and $O(t/u+u^{-2})$ for $tu>1$ (when $0<u\le1$), so that its integral against $(1+t)^{-3}\,dt$ is $O(1+\log(1/u))$, which is integrable on $(0,1)$; hence Fubini's theorem
allows $\mathcal T$ to be applied under the integral. By (i),
the term linear in $t$ contributes nothing and
$\mathcal T[\sin(\tfrac u2\,\cdot)]=\pi\tfrac{u^2}4e^{-bu/2}$, so
$\mathcal T[\Im\log\Gamma(\tfrac14+\tfrac{i\cdot}2)]
=-\frac\pi4\int_0^\infty\frac{u\,e^{-su/2}}{1-e^{-u}}\,du
=-\frac\pi4\psi'(s/2)$. Since
$\vartheta(t)=\Im\log\Gamma(\frac14+\frac{it}2)-\frac t2\log\pi$ and
$\mathcal T[t]=0$, (ii) follows.

(iii) From $\arctan(2t)=\int_0^\infty e^{-v/2}\sin(tv)\,dv/v$ and (i),
$\mathcal T[\pi^{-1}\arctan(2\,\cdot)]=\int_0^\infty ve^{-sv}dv=s^{-2}$
and $\mathcal T[\pi^{-1}\mathrm J]=\int_0^{\log2}ve^{(1-s)v}dv
=\int_1^2(\log u)u^{-s}du$; with $\mathcal T[-1]=-2/b^2$ this is
\eqref{eq:v8-E}. Analyticity and $r_0(0)=-1$ are clear. For the
bounds write $\mathrm J(t)=\operatorname{Si}(t\log2)+\int_0^{\log2}\varphi(v)\sin(tv)\,dv$
with $\varphi(v)=(e^{v/2}-1)/v$ increasing from $\frac12$ to $0.598$.
Since $\varphi$ is positive and increasing, integration by parts gives $|\int_0^{\log2}\varphi\sin(tv)\,dv|\le\bigl(\varphi(0)+\varphi(\log2)+\int_0^{\log2}\varphi'\bigr)/t=2\varphi(\log2)/t\le1.2/t$.
The positivity of $\mathrm J$ is proved in three ranges. For $0<t\le2$, the inequalities $\sin x\ge x-x^3/6$ for $x\ge0$ and $1\le e^{v/2}\le\sqrt2$ on $[0,\log2]$ give $\mathrm J(t)\ge t\log2-\frac{\sqrt2(\log2)^3}{18}t^3\ge\bigl(\log2-\frac{4\sqrt2(\log2)^3}{18}\bigr)t>0.58\,t$. For $2\le t\le\pi/\log2$, the argument $t\log2$ lies in $[2\log2,\pi]$, where $\operatorname{Si}$ is increasing, and $\operatorname{Si}(x)\ge x-x^3/18$ gives $\operatorname{Si}(t\log2)\ge2\log2-(2\log2)^3/18>1.23$; since $1.2/t\le0.6$, $\mathrm J(t)>0.63$. For $t\ge\pi/\log2$, $|\operatorname{Si}(x)-\frac\pi2|\le2/x$ gives $\mathrm J(t)\ge\frac\pi2-\bigl(\frac2{\log2}+1.2\bigr)\frac1t\ge\frac\pi2-\frac2\pi-\frac{1.2\log2}\pi>0.66$. Hence $\mathrm J(t)>0$ for all $t>0$.
Finally $r_0=\pi^{-1}\bigl[\mathrm J-\frac\pi2-\arctan\frac1{2t}\bigr]$ and
$|\operatorname{Si}(x)-\frac\pi2|\le2/x$ give
$|r_0(t)|\le\pi^{-1}(2/\log2+1.2+\tfrac12)/t\le1.46/t$.

(iv) From $\xi(s)=\frac12s(s-1)\pi^{-s/2}\Gamma(s/2)\zeta(s)$,
$(\xi'/\xi)'=(\zeta'/\zeta)'-s^{-2}-(s-1)^{-2}+\frac14\psi'(s/2)$.
Using (ii), (iii) and $\int_1^\infty(\log u)u^{-s}du=(s-1)^{-2}$,
\begin{equation}
-(\xi'/\xi)'-E-\mathcal T[N_\vartheta]
=-(\zeta'/\zeta)'+\int_1^\infty(\log u)u^{-s}du-\int_1^2(\log u)u^{-s}du
=Q(s).\qedhere
\end{equation}
\end{proof}

\begin{lemma}[Uniqueness for the transform]\label{lem:v8-unique}
Let $g$ be measurable on $(0,\infty)$ with
$\int_0^\infty|g(t)|^2(1+t)^{-3}dt<\infty$ and $\mathcal T[g](s)=0$
for all $\Re s>\frac12$. Then $g=0$ almost everywhere.
\end{lemma}
\begin{proof}
Let $g_o$ be the odd extension of $g$ to $\mathbb R$. Since
$\partial_tk_s$ is odd, $\mathcal T[g](s)=-\frac12\int_{\mathbb R}g_o\,\partial_tk_s\,dt$,
and $\partial_tk_s(t)=2(t+ib)^{-3}+2(t-ib)^{-3}$. Put
$C(z)=\int_{\mathbb R}g_o(t)(t-z)^{-3}dt$ for $z\notin\mathbb R$
(absolutely convergent and holomorphic). Oddness gives $C(-z)=C(z)$,
so $\mathcal T[g](s)=-2C(ib)$, and the hypothesis says $C\equiv0$ on
$\{\Im z>0\}$, hence on $\{\Im z<0\}$ as well. Let
$C_1(z)=\int_{\mathbb R}g_o(t)\bigl[(t-z)^{-1}-(t-i)^{-1}-(z-i)(t-i)^{-2}\bigr]dt$;
the bracket is $O(|t|^{-3})$, $C_1$ is holomorphic off $\mathbb R$ and
$C_1''=2C\equiv0$, so $C_1$ is a polynomial $p_\pm$ of degree at most one on each
half-plane. The bracket equals $(z-i)^2(t-z)^{-1}(t-i)^{-2}$, so
$C_1(z)=(z-i)^2\int_{\mathbb R}w(t)(t-z)^{-1}dt$ with
$w(t)=g_o(t)(t-i)^{-2}\in L^2(\mathbb R)$. By the boundary formula for
Cauchy integrals of $L^2$ functions \cite[Thm.~93]{TitchmarshFI},
$2\pi i\,w(x)=\lim_{y\to0^+}[C_1(x+iy)-C_1(x-iy)](x-i)^{-2}
=(p_+(x)-p_-(x))(x-i)^{-2}$ for almost every $x$. Hence $g_o$ agrees
a.e. with a polynomial of degree at most one; being odd, $g_o(t)=ct$,
and $\int_0^\infty c^2t^2(1+t)^{-3}dt<\infty$ forces $c=0$.
\end{proof}

\begin{proposition}[Limit phase and the first index]\label{prop:v8-first}
Fix $W\in\mathcal W$ and assume the explicit bound \eqref{eq:MCOM}. Suppose that a sequence $X_j\to\infty$
satisfies \eqref{eq:v7-normality-assumption} and that
$P_{X_j}^W$ converges in measure on every compact subset of
$(0,\infty)$ to a function $P_\infty$. Then:
\begin{enumerate}[label=(\alph*),leftmargin=*]
\item RH holds and $P_\infty=S-r_0$ almost everywhere on $(0,\infty)$,
 with $r_0$ from \eqref{eq:v8-r0};
\item $|S(t)-r_0(t)|<H(t)$ for every $t>0$ (so $|P_\infty|<H$ almost everywhere): the MCOM truncation is
 inactive in the limit, and $F_\infty:=N_0+\clip_H(P_\infty)=N-R-r_0$;
\item $\lim_{j\to\infty}\tau_k^{\Cms}(X_j;W)=\gamma_k$ for every $k\ge2$;
\item $\lim_{j\to\infty}\tau_1^{\Cms}(X_j;W)=\gamma_1-t_\sharp$, where
\begin{equation}\label{eq:v8-tsharp}
 t_\sharp:=\meas\{0<t<\gamma_1:\ R(t)+r_0(t)\le-\tfrac12\}>0 .
\end{equation}
\end{enumerate}
Consequently $\mathrm{DBSRH}(W)$ in the form \eqref{eq:DBSRH} is
incompatible with these hypotheses. With the sublevel points
measured from $2\pi$, the minimum of $N_0$,
\begin{equation}\label{eq:v8-flat}
 \tau_k^\flat(X;W):=2\pi+\meas\{t\ge2\pi:\ F_X^W(t)<k-\tfrac12\},
\end{equation}
one has $\lim_j\tau_k^\flat(X_j;W)=\gamma_k$ for every $k\ge1$.
\end{proposition}
\begin{proof}
(a) RH is \cref{thm:v7-normality}. By Fatou's lemma $P_\infty$
satisfies the weighted $L^2$ bound; the bound gives uniform
integrability on compact sets and, by Cauchy--Schwarz, uniformly small
contributions of $(0,\delta)\cup(B,\infty)$ to
$\int|P_{X_j}^W|(1+t)^{-3}dt$, so $\mathcal T[P_{X_j}^W](s)\to\mathcal T[P_\infty](s)$
for every $\Re s>\frac12$. Hence $\mathcal T[P_\infty]$ is the
holomorphic continuation of $Q$ obtained in the proof of
\cref{thm:v7-normality}. Integrating by parts,
$\mathcal T[N](s)=\sum_kk_s(\gamma_k)=\sum_\rho(s-\rho)^{-2}$ under RH,
which equals $-(\xi'/\xi)'(s)$ by the Hadamard product. By
\cref{lem:v8-smooth}(iv), on $\Re s>\frac12$ (where both sides are
holomorphic under RH),
$\mathcal T[S-r_0]=\mathcal T[N]-\mathcal T[N_\vartheta]-\mathcal T[r_0]
=-(\xi'/\xi)'-\mathcal T[N_\vartheta]-E=Q=\mathcal T[P_\infty]$.
\Cref{lem:v8-unique} applied to $g=P_\infty-S+r_0$, which satisfies
the weighted $L^2$ bound since $|S|\ll\log(2+t)$ and $r_0$ is bounded,
proves (a).

(b) For $0<t<\gamma_1$ we have $N=0$, so $S=-1-\vartheta/\pi\in[-1,0.13]$
because $\vartheta\ge-3.531$ on $(0,\infty)$; also
$|J(t)|\le t\int_0^{\log2}e^{v/2}dv\le0.83\,t$, so for $t<2$,
$|r_0|\le1.53$ and $|S-r_0|\le2.6<5.5\le H(t)$. For $t\ge2$,
\eqref{eq:MCOM} (extended to zero heights by right limits) and
$|r_0|\le0.73$ give $|S-r_0|\le0.08382\log t+5.18<H(t)$.
Thus $\clip_H(P_\infty)=P_\infty$ and $F_\infty=N_0+S-r_0=N-R-r_0$.

(c), (d) Since $F_X^W\ge N_0-H$, every sublevel set at level
$\lambda_k=k-\frac12$ lies in $(0,B_k]$ with
$B_k:=\sup\{N_0-H<\lambda_k\}$ independent of $X$. The truncation is
$1$-Lipschitz, so $F_{X_j}^W\to F_\infty$ in measure on compact subsets
of $(0,B_k]$. On each interval between consecutive distinct zero heights,
$F_\infty-\lambda_k$ is real-analytic and not constant (as
$R+r_0\to0$ at infinity and $R+r_0\to\frac18-1$ as $t\to0^+$), so
$\meas\{F_\infty=\lambda_k\}=0$. Therefore, for each $\delta>0$,
$\meas(\{F_{X_j}^W<\lambda_k\}\cap[\delta,B_k])\to\meas(\{F_\infty<\lambda_k\}\cap[\delta,B_k])$,
while the part in $(0,\delta)$ has measure at most $\delta$; letting
$\delta\to0$, $\lim_j\tau_k^{\Cms}(X_j;W)=\meas\{F_\infty<\lambda_k\}$.
We compute this set. For $t\ge\gamma_1$, $|R|\le0.53/(\pi\gamma_1)<0.012$
and $|r_0|\le1.46/\gamma_1<0.104$, so $|R+r_0|<\frac12$ and, $N$ being
an integer, $N-R-r_0<k-\frac12$ holds exactly when $N\le k-1$, that is
for $t<\gamma_k$. For $0<t<\gamma_1$, $F_\infty=-R-r_0$; since $J>0$,
$-r_0\le1-\pi^{-1}\arctan(2t)<1$, and $R>-\frac12$ there (for
$t\ge1$ by $|R|\le0.53/\pi$; for $t<1$ by $\vartheta(t)\ge\vartheta(1)>-1.8$
and $N_0(t)\le\frac78$, which give $R(t)>-0.45$), so $F_\infty<\frac32$.
Hence $\{F_\infty<\lambda_k\}=(0,\gamma_k)$ up to a null set for every
$k\ge2$, proving (c), while
$\{F_\infty<\frac12\}=\{t<\gamma_1:\ R+r_0>-\frac12\}$, whose measure is
$\gamma_1-t_\sharp$. Since $R(t)+r_0(t)\to-\frac78<-\frac12$ as
$t\to0^+$ and both functions are continuous on $(0,\gamma_1)$, the set in
\eqref{eq:v8-tsharp} contains an interval $(0,\epsilon)$, so $t_\sharp>0$.
For \eqref{eq:v8-flat}: on $[2\pi,\gamma_1)$, $|R|\le0.53/(2\pi^2)<0.03$
and $-r_0\le|r_0|\le1.46/(2\pi)<0.24$, so $F_\infty<\frac12$
there, giving $\lim_j\tau_1^\flat=2\pi+(\gamma_1-2\pi)=\gamma_1$; for
$k\ge2$ the interval $(0,2\pi)$ lies inside $\{F_\infty<\lambda_k\}$, so
$\tau_k^\flat$ and $\tau_k^{\Cms}$ have the same limit.
\end{proof}

\begin{remark}[Scope]
Apart from \eqref{eq:MCOM}, the two hypotheses of \cref{prop:v8-first} concern only the compensated
weighted prime phase. The proposition does not prove RH; it shows that,
in the only scenario in which the transform criterion applies, the
limit of the prime phase is $S-r_0$ and not $S$, that the correction
$r_0$ is invisible at every zero height above $\gamma_1$ (where
$|R+r_0|<\frac12$), and that it displaces the first sublevel point by a
positive amount $t_\sharp$. The endpoint observation $P_X^W(0)=0$ versus
$N_0(0^+)=\frac78$ in \cref{sec:v7-complex} is the finite-cutoff trace of
the same fact. The conditional statements of
\cref{cor:rh-rvm-v5} and \cref{sec:v7-open} are stated for $\tau_1^{\Cms}$; if DBSRH is used there in the scenario of \cref{prop:v8-first}, its index $k=1$ is to be read with $\tau_1^\flat$ of \eqref{eq:v8-flat}.
\end{remark}

\begin{proposition}[Transform criterion for the normalized sublevel points]
\label{prop:v8-criterion}
Fix $W\in\mathcal W$ and put, for $\Re s>\frac12$,
\begin{equation}\label{eq:v8-KX}
 \mathcal K_X(s):=\sum_{k\ge1}k_s\bigl(\tau_k^\flat(X;W)\bigr),
\end{equation}
which converges absolutely and locally uniformly, uniformly in $X\ge2$,
because $\tau_k^\flat\ge a_k$ for $k\ge k_0$, where $k_0$ is the index from which the barrier $a_k>2\pi$ of \cref{thm:block-sep-v3} is defined,
and $\sum_ka_k^{-2}<\infty$, while $2\pi\le\tau_k^\flat\le B_k$ for
$k<k_0$, with $B_k:=\sup\{t>0:N_0(t)-H(t)<k-\tfrac12\}$. Let $\Sigma\subset\{\Re s>\frac12\}$ have an accumulation point
in that half-plane. Then
\begin{equation}\label{eq:v8-criterion}
 \begin{aligned}
 &\lim_{X\to\infty}\mathcal K_X(s)=-\Bigl(\frac{\xi'}{\xi}\Bigr)'(s)
 \ \ \text{for every }s\in\Sigma\\
 &\qquad\Longleftrightarrow\quad
 \mathrm{RH}\ \text{and}\ \lim_{X\to\infty}\tau_k^\flat(X;W)=\gamma_k
 \ \text{for every }k\ge1 .
 \end{aligned}
\end{equation}
For $\Sigma\subset\{\Re s>1\}$ the right side of the limit is the
convergent prime expression
$-(\zeta'/\zeta)'(s)+s^{-2}+(s-1)^{-2}-\frac14\psi'(s/2)$, so the
criterion compares the sublevel points with Euler, $\Gamma$ and the
pole only; for $s=\frac12+d+iT$ the real and imaginary parts of
$k_s(\tau)$ are the kernels of order $j=2$ used in \cite{CasanaMCOM}.
\end{proposition}
\begin{proof}
($\Leftarrow$) Under RH, $-(\xi'/\xi)'(s)=\sum_\rho(s-\rho)^{-2}=\sum_kk_s(\gamma_k)$,
and dominated convergence with the bound
$\sup_{s\in K}|k_s(\tau)|\le C_K\min(1,\tau^{-2})$ for each compact $K\subset\{\Re s>\frac12\}$ (indeed $\Re(s-\frac12)\ge\delta_K>0$ gives $|s-\frac12\mp i\tau|\ge\delta_K$, and $|s-\frac12\mp i\tau|\ge\tau/2$ for $\tau\ge2\sup_K|s|+1$), summable against $\tau_k^\flat\ge a_k$, $\sum_ka_k^{-2}<\infty$, gives the limit, locally uniformly.
($\Rightarrow$) Let $X_n\to\infty$. As $\tau_k^\flat(X_n)\in[2\pi,B_k]$
is nondecreasing in $k$, a diagonal subsequence satisfies
$\tau_k^\flat(X_{n_i})\to\lambda^\infty_k$ for every $k$, with
$2\pi\le\lambda^\infty_k\le\lambda^\infty_{k+1}$ and $\lambda^\infty_k\ge a_k$ for $k\ge k_0$.
By dominated convergence $\mathcal K_{X_{n_i}}\to\mathcal K_\infty:=\sum_kk_s(\lambda^\infty_k)$
locally uniformly on $\Re s>\frac12$; the same series converges locally
uniformly on $\mathbb C$ away from the points $\frac12\pm i\lambda^\infty_k$, so
$\mathcal K_\infty$ is meromorphic on $\mathbb C$ with double poles
exactly at $\frac12\pm i\lambda^\infty_k$, the leading coefficient at
$\frac12+i\mu$ being the number of $k$ with $\lambda^\infty_k=\mu$. All poles of $\mathcal K_\infty$ lie on the line $\Re s=\frac12$, so $\mathcal K_\infty$ is holomorphic in the connected half-plane $\Re s>\frac12$. The difference $f=\mathcal K_\infty-\bigl(-(\xi'/\xi)'\bigr)$ is meromorphic there and vanishes on $\Sigma$; an accumulation point of $\Sigma$ cannot be a pole of $f$, since $|f|\to\infty$ at a pole, so $f$ vanishes near it and, by the identity theorem on the connected domain, $f\equiv0$ there. Thus
$\mathcal K_\infty=-(\xi'/\xi)'$ on $\Re s>\frac12$; as
$\mathcal K_\infty$ is holomorphic there, $\xi$ has no zero in that
half-plane, which is RH. Then $-(\xi'/\xi)'=\sum_kk_s(\gamma_k)$, and
comparing the double poles and their coefficients of the two meromorphic
functions (every term $k_s(\lambda^\infty_k)$ has leading coefficient $+1$ at its pole, so there is no cancellation between terms and a height occurring $m$ times gives a double pole with coefficient exactly $m$) shows that the multisets $\{\lambda^\infty_k\}$ and $\{\gamma_k\}$
coincide; both being enumerated in nondecreasing order,
$\lambda^\infty_k=\gamma_k$ for every $k$. Since this holds along a subsequence
of every sequence $X_n\to\infty$, the full limits exist.
\end{proof}

\begin{remark}
\Cref{prop:v8-criterion} is the precise content of the
Laguerre--P\'olya proposal in \cite{CasanaV10Notes}: the finite-cutoff
points are real by construction, so the only way the limit identity can
fail is through a defect in the \emph{positions} of the $\lambda_k$,
never through a non-real limit. A repeated imaginary part is
\emph{not} excluded by the right side of \eqref{eq:v8-criterion}: if
$\gamma_k=\gamma_{k+1}$, the right side of \eqref{eq:v8-criterion} holds with $\lambda_k=\lambda_{k+1}=\gamma_k$, and so does the left side. Simplicity therefore has to come from the arithmetic
noncollapse of \cref{prop:correlated-v5}, exactly as in
\cref{cor:rh-rvm-v5}.
\end{remark}

\section{A box-transform and a parity-sensitive multiplicity test}
\label{sec:v7-parity}
This section formalizes the third route in
\cite{CasanaV10Notes} \emph{without} importing a zero-free
claim or assuming RH. The identity is deterministic and holds
for an arbitrary locally absolutely continuous real function
$F$, including the actual taucm and taucms functions where
their stated local regularity is available.

\begin{theorem}[Finite-interval box identity, with the endpoint term]
\label{thm:v7-box}
Let $0<a<b$ not be imaginary parts of zeros. Suppose $F$ is
absolutely continuous on $[a,b]$. For a continuous test
function $g$ with a primitive $A_g$ ($A_g'=g$), put
$e_k(F)=k-\frac12-F(\gamma_k)$ and
$\Psi_g(z)=\int_{z-1/2}^{z+1/2}g(v)\,dv$. Then
\begin{equation}\label{eq:v7-box}
\begin{split}
\sum_{a<\gamma_k<b}\Psi_g(e_k(F))
={}&\int_a^b g(N(t)-F(t))F'(t)\,dt\\
 &+A_g\bigl(N(b)-F(b)\bigr)
   -A_g\bigl(N(a)-F(a)\bigr).
\end{split}
\end{equation}
The sum counts imaginary parts of zeros with multiplicity.
\end{theorem}
\begin{proof}
On every interval between consecutive distinct zero heights,
$N$ is constant, and with $x=N-F$ we have
$g(N-F)F'=-\frac d{dt} A_g(N-F)$ a.e. Sum these
identities over the finitely many pieces in $[a,b]$.
At an imaginary part with multiplicity $m$, the missing telescopic
contribution is
$A_g(x^-+m)-A_g(x^-)$, which equals
$\sum_{i=0}^{m-1}\Psi_g(e_{k+i}(F))$ because the boxes
$[e_{k+i}-1/2,e_{k+i}+1/2]$ tessellate $[x^-,x^-+m]$.
The remaining terms are exactly the two endpoint values
shown in \eqref{eq:v7-box}.
\end{proof}

For $\omega\ne0$, taking $g(v)=e^{i\omega v}$ yields the exact
finite-interval Fourier form
\begin{equation}\label{eq:v7-fourier-box}
\begin{split}
 \frac{2\sin(\omega/2)}{\omega}
 \sum_{a<\gamma_k<b}e^{i\omega e_k(F)}
={}&\int_a^b e^{i\omega(N-F)}F'(t)\,dt\\
 &+\frac{e^{i\omega(N(b)-F(b))}
          -e^{i\omega(N(a)-F(a))}}{i\omega}.
\end{split}
\end{equation}
No ``boundary terms are negligible'' statement is assumed.

\begin{proposition}[Parity bound for repeated heights]
\label{prop:v7-parity}
Let $\mathcal I$ be a union of complete repeated-height
groups in the index list, and let $r(\mathcal I)$ be the number
of distinct imaginary parts represented. For any real $c$
and any single-valued real function $F$,
\begin{equation}\label{eq:v7-parity}
\#\mathcal I-r(\mathcal I)
\le\sum_{k\in\mathcal I}
  \bigl[1-\cos\pi(e_k(F)-c)\bigr].
\end{equation}
\end{proposition}
\begin{proof}
For a group of multiplicity $m$ we have
$e_{k+i}=e_k+i$, $0\le i<m$; its right side is
$m-\cos(\pi(e_k-c))\sum_{i=0}^{m-1}(-1)^i\ge m-1$.
Sum over disjoint complete groups.
\end{proof}

The familiar relation $(-1)^{N(t)}=-\operatorname{sgn}Z(t)$
outside zeros, with the Hardy function $Z(t)=e^{i\vartheta(t)}\zeta(\frac12+it)$ (so that $Z(0)=\zeta(\frac12)<0$; the relation holds for the $N$ of \S\ref{sec:definitions}, since a zero off the line and its mirror $1-\bar\rho$ share the same height), illustrates the information carried by the parity
observable; it is \emph{not} a new detection algorithm for
off-line zeros. The box identity alone cannot make the
right side of \eqref{eq:v7-parity} vanish, and, since no sign or size
condition is imposed on $F'$, the oscillatory integral in
\eqref{eq:v7-fourier-box} cannot be discarded. This is therefore a separate exact
potential tool, not an unconditional simplicity theorem.

\section{Complex signed-pair kernel and proposed zero formula:
       status and limitations}\label{sec:v7-complex}
In the external notes, a conditional complex continuation of the
two-source kernel from \cref{sec:kernel-v4} was proposed to
separate a small real spacing from two mirror-symmetric
zeros with the same imaginary part. It is important not to import
this weighted model as if it were the author's \emph{original}
taucs or as if it supplied the full zero-explicit formula.

For the already defined entire kernel $K$ and any complex $h$,
$C(z,h):=K(z-h/2)+K(z+h/2)$ is entire and even in $h$.
The Taylor expansion and its global remainder for real $h$
are established in \cref{sec:kernel-v4}. Algebraically a
real pair has $h^2\ge0$, a repeated imaginary part has $h^2=0$,
and a formal transverse pair at $\gamma\mp id$ corresponds
to $h^2=-4L^2d^2$. The analytic continuation of the
\emph{model} is rigorous; the identification of its complex
arguments with terms in the \emph{actual weighted prime
explicit formula} remains unproved. In particular, no
infinite zero sum $\sum_\rho K(L(t-u_\rho))$, with $u_\rho:=\gamma-i(\beta-\frac12)$ for $\rho=\beta+i\gamma$ (so that $t-u_\rho=iz$ with $z$ as in \eqref{eq:laplace}), is asserted
here as an unconditional identity. The contour deformation,
branch constants, uniform convergence and contributions of
the noncentral zeros must all be controlled first.

\begin{proposition}[Signed central slope for an auxiliary weight]
\label{prop:v7-signed-slope}
Let $q\ge0$ be any nonzero, integrable weight on $[1,2]$,
normalized by $\int_1^2q(v)\,dv=1$, and let
$K_q(z)=\frac12+\pi^{-1}\int_1^2q(v)\operatorname{Si}(vz)\,dv$.
For the \emph{model} $C_q(z,h)=K_q(z-h/2)+K_q(z+h/2)$,
its derivative at the pair center is
\begin{equation}
 \mathscr S_q(h)=\frac2\pi\int_1^2q(v)
 \frac{\sin(vh/2)}{h/2}\,dv
 \quad(\mathscr S_q(0)=\tfrac2\pi\int_1^2vq(v)\,dv).
\end{equation}
For every real $h\ne0$, $\mathscr S_q(h)<\mathscr S_q(0)$;
for every real $d>0$,
$\mathscr S_q(2id)>\mathscr S_q(0)$.
These comparisons classify the \emph{model} parameter with a
sign; they are not a criterion about actual off-line zeta
zeros without a validated zero-explicit formula.
\end{proposition}
\begin{proof}
For real $x\ne0$, $\sin x/x<1$; for real $x\ne0$,
$\sinh x/x>1$. Substitute into the integral, with $q\ge0$.
\end{proof}
The symmetry $K_q(-z)=1-K_q(z)$ makes the derivative of
$C_q$ at the center stationary under a first-order shift of
the sampling center, a useful \emph{model} fact; the
uncontrolled influence of neighboring contributions remains.
The special weight $q(v)=v/2+1/4$ used in earlier numerical
sections is one \emph{borrowed choice} and is not needed for
this auxiliary proposition.

The inverse estimator from \cref{sec:identify-v4} measures a
signed algebraic quadratic parameter only if its sampling
points and center are independently specified and the
projected arithmetic residual is rigorously bounded.
No such prime-only bound is currently known. The small-gap
numerical experiments cannot supply one. For these reasons
the signed-pair model is recorded as an \emph{auxiliary
exploratory route}, not as the main proposed RH mechanism.

A separate Laguerre--P\'olya proposal in the public notes
forms an even entire product from the real $\tau_k^{\Cms}$.
Reality of finite-cutoff sublevel points and normality of such
products cannot by themselves exclude a multiple real zero
in a limit, so this is not an independent RH route. Any
product-normality assertion also requires its stated uniform
tail bounds and control of the initial factors.

The public notes also flag a lower-endpoint defect
for the naive ``DBSRH for every $k\ge1$'' statement: the
Euler phase satisfies $P_X^W(0)=0$, whereas the fixed main
term extends to $N_0(0^+)=7/8>1/2$. \Cref{prop:v8-first} makes
this precise: under the hypotheses of the transform criterion
together with convergence in measure, $\tau_1^{\Cms}$ converges to
$\gamma_1-t_\sharp$ with $t_\sharp>0$, while the points
$\tau_k^\flat=2\pi+\meas\{t\ge2\pi:F_X^W(t)<k-\frac12\}$ of
\eqref{eq:v8-flat} converge to $\gamma_k$ for every $k\ge1$.
Finite-cutoff strictness of $\tau_k^\flat$ in the initial indices is
not established here. We therefore retain the original definitions
for existing theorems and do not silently replace
$\mathrm{DBSRH}(W)$ by its corrected variant.

\section{Exponential sums over the zeros and over the approximating sequences; centroids}\label{sec:v9-centroid}
This section studies the sums $\sum e^{i\omega\gamma_k}$, $\sum e^{i\omega x_k}$ for $x_k$ one of the sequences $\gamma_k^{\mathrm{RvM}}$, $\tau_k^{\mathrm C}$, $\tau_k^{\Cm}$, and their Ces\`aro means (centroids). Attributions are given where each item enters (the Landau coefficient and the sawtooth coefficient come from the author's Fourier notebook; the joint transform, the spectral-coherence statement, the second-difference identity, the band argument for the rearrangement term, the harmonic criterion, the reconstruction at $\omega=1$ and the collision-energy identity come from ChatGPT texts, and were re-derived here from the objects of this paper). Nothing uses RH unless stated. No symbol of the earlier sections is redefined; the new symbols are listed in the notation table (rows from \cref{def:v9-objects} on). The $220$-zero numbers of \S\ref{sec:v9-numbers} are quoted from the ChatGPT text and were not re-run here.

\subsection{Objects and exact identities}\label{sec:v9-objects}
\begin{definition}[Families, displacements, joint transform]\label{def:v9-objects}
Let $\star\in\{\mathrm{RvM},\mathrm C,\Cm\}$ and put
\begin{equation}
\mathsf x_k^{\mathrm{RvM}}=\gamma_k^{\mathrm{RvM}},\qquad
\mathsf x_k^{\mathrm C}=\tau_k^{\mathrm C}=u_k(F_h),\qquad
\mathsf x_k^{\Cm}=\tau_k^{\Cm}=u_k(F^{\Cm}),
\end{equation}
\begin{equation}
\sigma_k^\star=\gamma_k-\mathsf x_k^\star .
\end{equation}
Let $\mathfrak J_T=\{k:\ T<\gamma_k\le2T\}$; its cardinality is $M_T=N(2T)-N(T)$ of \eqref{eq:jone-AB-v3}, and $M_T=\frac{T}{2\pi}\log T\,(1+o(1))$. For real $\omega,\omega'$ set
\begin{equation}\label{eq:v9-Theta}
\Theta_T^\star(\omega,\omega')=\frac1{M_T}\sum_{k\in\mathfrak J_T}
\exp\bigl(i\omega\mathsf x_k^\star+i\omega'\sigma_k^\star\bigr),\qquad
\Xi_T^\star(\omega)=\Theta_T^\star(\omega,\omega)-\Theta_T^\star(\omega,0).
\end{equation}
Thus $\Theta_T^\star(\omega,\omega)$ is the normalized exponential sum over the true zeros, $\Theta_T^\star(\omega,0)$ the one over the approximating sequence, and $\Theta_T^\star(0,\omega')$ is the empirical characteristic function of the displacements. For $\ell\ge1$ the partial sum and the centroid (first Ces\`aro mean of the partial sums) are
\begin{equation}\label{eq:v9-Cen}
\mathrm{Par}_\ell(\omega;\mathsf y)=\sum_{k\le\ell}e^{i\omega \mathsf y_k},\qquad
\mathrm{Cen}_\ell(\omega;\mathsf y)=\frac1\ell\sum_{n\le\ell}\mathrm{Par}_n(\omega;\mathsf y)
=\frac1\ell\sum_{k=1}^\ell(\ell-k+1)\,e^{i\omega \mathsf y_k},
\end{equation}
for a real sequence $\mathsf y=(\mathsf y_k)$, in particular $\mathsf y=\gamma$ or $\mathsf y=\mathsf x^\star$. Finally put
\begin{equation}\label{eq:v9-chi}
\chi(\omega)=-\Lambda(e^{\omega})\,e^{-\omega/2}\qquad(\omega>0),
\end{equation}
which vanishes unless $e^\omega$ is a prime power; we extend it by $\chi(-\omega):=\overline{\chi(\omega)}$ ($=\chi(\omega)$, since $\chi$ is real for $\omega>0$).
\end{definition}
The function $\Theta_T^\star(\omega,\omega')$ is the two-variable transform proposed in the ChatGPT text (there called $\mathcal W_{T,x}(u,v)$); the triangular version of that text is $\mathrm{Cen}$ above, and the identities
$\Theta_T^\star(0,\omega')=$ characteristic function of the displacements and $\mathrm{Cen}_\ell(\omega;\gamma)-\mathrm{Cen}_\ell(\omega;\mathsf x^\star)=$ triangular transform at $(\omega,\omega)$ minus at $(\omega,0)$ are immediate from the definitions.

\begin{lemma}[Second differences, Stieltjes identity]\label{lem:v9-exact}
(a) With $\mathbb S_\ell(\omega;\mathsf y)=\sum_{n\le\ell}\mathrm{Par}_n(\omega;\mathsf y)=\ell\,\mathrm{Cen}_\ell(\omega;\mathsf y)$ and the second difference $(\nabla^2\mathbb S)_k=\mathbb S_k-2\mathbb S_{k-1}+\mathbb S_{k-2}$ (with $\mathbb S_0=\mathbb S_{-1}:=0$) one has $(\nabla^2\mathbb S)_k(\omega;\mathsf y)=e^{i\omega \mathsf y_k}$ and
\begin{equation}\label{eq:v9-seconddiff}
\begin{aligned}
\bigl|(\nabla^2\mathbb S)_k(\omega;\gamma)-(\nabla^2\mathbb S)_k(\omega;\mathsf x^\star)\bigr|^2
&=2\bigl(1-\cos(\omega\sigma_k^\star)\bigr),\\
\frac1{M_T}\sum_{k\in\mathfrak J_T}2\bigl(1-\cos(\omega\sigma_k^\star)\bigr)
&=2\bigl(1-\Re\,\Theta_T^\star(0,\omega)\bigr).
\end{aligned}
\end{equation}
(b) Let $N^{\mathsf x}(s)=\#\{k:\mathsf x_k\le s\}$. For $0\le s_0<s_1$ and real $\omega$,
\begin{equation}\label{eq:v9-stieltjes}
\sum_{s_0<\gamma_k\le s_1}e^{i\omega\gamma_k}-\sum_{s_0<\mathsf x_k\le s_1}e^{i\omega\mathsf x_k}
=\Bigl[e^{i\omega s}\bigl(N-N^{\mathsf x}\bigr)(s)\Bigr]_{s_0}^{s_1}
-i\omega\int_{s_0}^{s_1}e^{i\omega s}\bigl(N-N^{\mathsf x}\bigr)(s)\,ds .
\end{equation}
\end{lemma}
\begin{proof}
(a) is the definition of the second difference of a double partial sum ($\mathrm{Par}_k-\mathrm{Par}_{k-1}=e^{i\omega\mathsf y_k}$), and $|e^{is_1}-e^{is_2}|^2=2(1-\cos(s_1-s_2))$. (b) is integration by parts in the Riemann--Stieltjes integral $\int_{s_0}^{s_1}e^{i\omega s}\,d(N-N^{\mathsf x})(s)$ of two right-continuous step functions: $N$ and $N^{\mathsf x}$ are right-continuous, the integral is over $(s_0,s_1]$ so that $\{s_1\}$ carries the jump at $s_1$ and $\{s_0\}$ is excluded, and the boundary values are the right values at $s_0,s_1$. Since $e^{i\omega s}$ is $C^1$, no assumption that $s_0,s_1$ avoid the jumps is needed.
\end{proof}

\begin{lemma}[Counting function of the sublevel points]\label{lem:v9-rearr}
Let $F$ be as in \eqref{eq:sublevel}, $\mu_F(c)=\meas\{t>0:F(t)<c\}$, so $\mathsf x_k=u_k(F)=\mu_F(k-\frac12)$, and put $F^*(s)=\inf\{c:\mu_F(c)>s\}$, $E_F=\{t>0:F(t)+\frac12\in\mathbb Z\}$, $B_F=\{s:F^*(s)+\frac12\in\mathbb Z\}$.
\begin{enumerate}[label=(\alph*),leftmargin=*]
\item \emph{(Exact.)} $N^{\mathsf x}(s)=\#\{k\ge1:\mu_F(k-\frac12)\le s\}$, and with
\begin{equation}\label{eq:v9-slip}
z_F=N-\lfloor F+\tfrac12\rfloor,\qquad \varrho_F=\lfloor F+\tfrac12\rfloor-N^{\mathsf x}
\end{equation}
one has $N-N^{\mathsf x}=z_F+\varrho_F$ identically. Moreover $|z_F|=q_F$ (the integer discrepancy of \cref{thm:taucm}).
\item \emph{(Quantile form, off the boundary.)} $N^{\mathsf x}(s)=\lfloor F^*(s)+\frac12\rfloor$ for every $s\notin B_F$ (with $F^*(s)\ge\frac12$), and $N^{\mathsf x}(s)\in\{\lfloor F^*(s)+\frac12\rfloor-1,\lfloor F^*(s)+\frac12\rfloor\}$ for $s\in B_F$. Up to at most countably many points, $B_F$ is the union of the intervals of length $\meas\{F=m-\frac12\}$, so $\meas B_F=\meas E_F$.
\item \emph{(Sawtooth.)} With $D_F=N-F$ and $\mathrm{saw}(y)=y-\lfloor y+\frac12\rfloor$,
\begin{equation}\label{eq:v9-sawD}
z_F=D_F+\mathrm{saw}(F),\qquad
\mathrm{saw}(F)=-\mathrm{saw}(D_F)-\ind_{E_F},\qquad
z_F=\lfloor D_F+\tfrac12\rfloor-\ind_{E_F},
\end{equation}
so that the replacements $\mathrm{saw}(F)\to-\mathrm{saw}(D_F)$ and $z_F\to\lfloor D_F+\frac12\rfloor$ are valid exactly off $E_F$, and almost everywhere if $\meas E_F=0$ (not assumed).
\item If $F\ge-\frac12$ is continuous and nondecreasing, then $\varrho_F\equiv0$.
\end{enumerate}
\end{lemma}
\begin{proof}
(a) $N^{\mathsf x}(s)=\#\{k:\mathsf x_k\le s\}$ by definition; the rest is the definition of $\varrho_F$ and $z_F$. (b) The function $\mu_F$ is nondecreasing and left-continuous. Put $c=k-\frac12$. If $F^*(s)>c$ then $\mu_F(c)\le s$ (otherwise $F^*(s)\le c$); if $F^*(s)<c$ then $\mu_F(c)>s$. So off the boundary $\mathsf x_k\le s$ iff $k-\frac12<F^*(s)$, and $\#\{k\ge1:k<F^*(s)+\frac12\}=\lfloor F^*(s)+\frac12\rfloor$ when $F^*(s)+\frac12\notin\mathbb Z$. At the boundary the equivalence can fail only for the single $k$ with $k-\frac12=F^*(s)$, which gives the two possible values. The set $\{s:F^*(s)=c\}$ is, up to endpoints, $[\mu_F(c),\meas\{F\le c\})$, of length $\meas\{F=c\}$. (c) The first identity is $z_F=N-F+(F-\lfloor F+\frac12\rfloor)$. Since $N\in\mathbb Z$, $\mathrm{saw}(F)=-\mathrm{saw}(N-F)$ holds iff $\mathrm{saw}(-y)=-\mathrm{saw}(y)$ at $y=F$, and $-y-\lfloor-y+\frac12\rfloor=-y+\lceil y-\frac12\rceil$ equals $-y+\lfloor y+\frac12\rfloor$ iff $y+\frac12\notin\mathbb Z$. On $E_F$, writing $F+\frac12=m$: $\mathrm{saw}(F)=F-m=-\frac12$, $\mathrm{saw}(D_F)=(N-F)-(N-m+1)=-\frac12$, $z_F=N-m$ and $\lfloor D_F+\frac12\rfloor=N-m+1$, which gives the correction $-\ind_{E_F}$ in each identity. (d) For continuous nondecreasing $F$ and $c\in\mathbb R$ the set $\{F<c\}$ is an initial interval $(0,a_c)$, $\mu_F(c)=a_c$, and $F(s)\ge c$ iff $s\ge a_c$ (for $s=a_c$ by continuity). Hence $\mu_F(k-\frac12)\le s$ iff $F(s)\ge k-\frac12$, i.e.\ $N^{\mathsf x}(s)=\#\{k\ge1:k-\frac12\le F(s)\}=\lfloor F(s)+\frac12\rfloor$, so $\varrho_F=0$ with no identification of $F^*$ with $F$.
\end{proof}
The rearrangement term $\varrho_F$ is not controlled by the level-counting inequality $\sum_k|\mathsf x_k-\mathsf x_k'|\le\int|\lfloor F+\frac12\rfloor-\lfloor F'+\frac12\rfloor|$ (that inequality compares two functions, not $F$ with $F^*$). For $F=F_h$ its mean is bounded in \cref{prop:v9-rearr}; for $F=F^{\Cm}$ the band has width $O(1)$ and the same argument gives nothing, which is not needed by \cref{prop:v9-CvsCm} and \cref{lem:v9-sawcmp}.

\subsection{Sawtooth coefficient and Landau cancellation}\label{sec:v9-sawtooth}
For a phase $F$ and $\omega\neq0$ define the sawtooth coefficient
\begin{equation}\label{eq:v9-varsigma}
\varsigma_T(\omega;F)=-\frac{2\pi}{T}\int_T^{2T}e^{i\omega t}\,\mathrm{saw}\bigl(F(t)\bigr)\,dt
=\frac{2\pi}{T}\int_T^{2T}e^{i\omega t}\Bigl(\mathrm{saw}\bigl(D_F(t)\bigr)+\ind_{E_F}(t)\Bigr)\,dt,
\end{equation}
the second form by \eqref{eq:v9-sawD}, which makes the two forms identical; dropping $\ind_{E_F}$ from the second form changes it by at most $2\pi\meas(E_F\cap[T,2T])/T$, which is $0$ when $\meas E_F=0$ (this is not assumed here, and the first form is the definition). Its mean-value counterpart at $\omega=0$ is $\varsigma_T(0;F)=-\frac{2\pi}{T}\int_T^{2T}\mathrm{saw}(F)\,dt=\frac{2\pi}{T}\int_T^{2T}\bigl(\mathrm{saw}(D_F)+\ind_{E_F}\bigr)\,dt$.

Throughout this subsection and the next ones the fourth-moment input \eqref{eq:IJNTinput} of the IJNT manuscript is assumed wherever $\int_T^{2T}|\Phi^\circ|^4\ll T(\log\log T)^2$ is used (\cref{lem:v9-cancel}, \cref{cor:v9-slip}, \cref{prop:v9-CvsCm}, \cref{lem:v9-sawcmp}, \cref{prop:v9-line}, \cref{thm:v9-reduction}, \cref{cor:v9-test}, \cref{prop:v9-mean}, \cref{cor:v9-cen}(ii) and the last sentence of \cref{prop:v9-harm}). We use two external inputs, (L) for the zeros and (K) for the RvM points, for fixed $\omega>0$ (the case $\omega<0$ is the conjugate).

\medskip\noindent\textbf{(L) Landau coefficient.}
\begin{equation}\label{eq:v9-landau}
\sum_{T<\gamma_k\le2T}e^{i\omega\gamma_k}=\frac{T}{2\pi}\,\chi(\omega)+o_\omega(T).
\end{equation}
Landau's formula \cite{Landau1912} and its uniform form by Gonek \cite{Gonek1993} give $\sum_{0<\gamma\le T}\Upsilon^{\rho}=-\frac{T}{2\pi}\Lambda(\Upsilon)+O(\log T)$ for fixed $\Upsilon>1$; under RH $\Upsilon^\rho=\Upsilon^{1/2}\Upsilon^{i\gamma}$, which yields \eqref{eq:v9-landau} with error $O(\log T)$. \emph{Unconditionally}, \eqref{eq:v9-landau} with error $o(T)$ follows from the Landau-type lemma in the proof of the Ford--Zaharescu theorem \cite{FordZaharescu2005} (their relation between $\sum \Upsilon^{i\gamma}$ and $\sum \Upsilon^{\rho-1/2}$, as used in the author's notebook). Two independent readings of that paper (their eq.~(3.8)) give the sharper unconditional error $O_\omega(T/\log T)$ for $\Upsilon=e^{\omega}$ fixed, which also follows directly: the zeros $\rho$ and $1-\bar\rho$ have the same height and $\Upsilon^{\rho}+\Upsilon^{1-\bar\rho}=2\Upsilon^{1/2+i\gamma}\cosh\bigl((\beta-\tfrac12)\log\Upsilon\bigr)$, with $2(\cosh u-1)\ll u^2$ for bounded $u$; Selberg's density estimate $N(\sigma,T)\ll T^{1-(\sigma-1/2)/4}\log T$ \cite{Selberg1946} gives $\sum_{T<\gamma\le2T}(\beta-\tfrac12)^2\ll T\log T\int_0^{1/2}u\,e^{-u\log T/4}\,du\ll T/\log T$, and Landau's formula supplies the rest. The arguments below use only $o(T)$ except where $O(T/\log T)$ is stated explicitly. The unconditional error is $O(T/\log T)$, \emph{not} $O(\log T)$.

\medskip\noindent\textbf{(K) Kusmin--Landau for the RvM points.}
The map $k\mapsto\gamma_k^{\mathrm{RvM}}=N_0^{-1}(k-\frac12)$ has decreasing, positive increments $\approx2\pi/\log(\gamma_k^{\mathrm{RvM}}/2\pi)$, since $N_0$ is convex for large argument. By the Kusmin--Landau inequality \cite{GrahamKolesnik1991}, for fixed $\omega>0$ and any index interval $I\subset\{k:\gamma_k^{\mathrm{RvM}}\le U\}$ beyond a fixed index,
\begin{equation}\label{eq:v9-KL}
\Bigl|\sum_{k\in I}e^{i\omega\gamma_k^{\mathrm{RvM}}}\Bigr|\ll_\omega\log U .
\end{equation}

\begin{lemma}[Landau cancellation for the phase of the construction]\label{lem:v9-cancel}
Let $F\in\{N_0+\Phi^\circ,\ F_h,\ F^{\Cm}\}$ and $\omega\neq0$. Then
\begin{equation}\label{eq:v9-cancel}
\frac{2\pi}{T}\int_T^{2T}e^{i\omega t}\,D_F(t)\,dt\longrightarrow0\qquad(T\to\infty).
\end{equation}
\end{lemma}
\begin{proof}
\emph{Step 1 (the zero side).} Write $G=N-N_0$. Integrating by parts as in \eqref{eq:v9-stieltjes} with $N_0$ in place of $N^{\mathsf x}$ (so that $\int e^{i\omega t}dN_0$ is explicit and $O(\log T/|\omega|)$),
\begin{equation}
\frac{2\pi}{T}\int_T^{2T}e^{i\omega t}G\,dt=-\frac{2\pi}{i\omega T}\sum_{T<\gamma_k\le2T}e^{i\omega\gamma_k}+O\Bigl(\frac{\log T}{T}\Bigr)
=\frac{-\chi(\omega)}{i\omega}+o(1)
=\frac{i\,\chi(\omega)}{\omega}+o(1),
\end{equation}
using \eqref{eq:v9-landau} and $|G|\ll\log T$ for the boundary terms. For $\omega=r\log p$ one has $\chi(\omega)=-p^{-r/2}\log p$, so the limit is $-i/(rp^{r/2})$.

\emph{Step 2 (the prime side).} Each term $-\frac1\pi\frac{\sin(rt\log p)}{rp^{r/2}}$ of $\Phi_m$ in \eqref{eq:Phi-m} has $\frac{2\pi}T\int_T^{2T}e^{i\omega t}(\cdot)\,dt\to-i/(rp^{r/2})$ if $\omega=r\log p$ and $\to0$ otherwise. In \eqref{eq:interpolation} the two interpolation weights sum to $1$, and for $T$ large every fixed prime power belongs to both $\Phi_{m_j}$ and $\Phi_{m_{j+1}}$, so the coefficient of $\Phi^\circ$ at $\omega$ has the same limit as in Step 1. The terms with $\omega'=r\log p\ne\omega$ contribute $O(1/(T|\omega-\omega'|))\cdot p^{-r/2}$ each; their sum over the $\ll T^{1/2}$ available prime powers is $O(T^{-3/4+\epsilon})$.

\emph{Step 3 (clipping).} $D_{N_0+\Phi^\circ}=G-\Phi^\circ$, so \eqref{eq:v9-cancel} holds for $N_0+\Phi^\circ$. For $F_h$ and $F^{\Cm}$ the difference with the unclipped phase is bounded by $\int_T^{2T}(|\Phi^\circ|-a)_+$, $a\in\{h,H\}$: since $H\ge h$ for large $t$, $(|\Phi^\circ|-H)_+\le(|\Phi^\circ|-h)_+$, so it suffices to treat $a=h$: $\int_T^{2T}(|\Phi^\circ|-h)_+\le\int|\Phi^\circ|\ind_{\{|\Phi^\circ|>h\}}\le(\int|\Phi^\circ|^4)^{1/4}\meas\{|\Phi^\circ|>h\}^{3/4}\ll T/\log\log T$, using $\int_T^{2T}|\Phi^\circ|^4\ll T(\log\log T)^2$ (Selberg's fourth moment of $S$ \cite{Selberg1946,Fujii1974} together with \eqref{eq:IJNTinput}, since $\Phi^\circ=G-D_{N_0+\Phi^\circ}$) and Chebyshev. This is $o(T)$.
\end{proof}

\begin{corollary}[Slip coefficient]\label{cor:v9-slip}
For $F\in\{F_h,F^{\Cm}\}$ and fixed $\omega\ne0$, as $T\to\infty$,
\begin{equation}\label{eq:v9-slipcoeff}
\frac{2\pi}{T}\int_T^{2T}e^{i\omega t}z_F(t)\,dt=-\varsigma_T(\omega;F)+o(1).
\end{equation}
In particular the Fourier coefficient of the integer-valued slip $z_F$ is, up to $o(1)$, minus the sawtooth coefficient of $F$ (equivalently, $-\frac{2\pi}T\int_T^{2T}e^{i\omega t}\bigl(\mathrm{saw}(D_F)+\ind_{E_F}\bigr)\,dt$ by \eqref{eq:v9-varsigma}).
\end{corollary}
\begin{proof}
Insert $z_F=D_F+\mathrm{saw}(F)$ from \eqref{eq:v9-sawD} and use \cref{lem:v9-cancel}.
\end{proof}

\subsection{A sufficient criterion for the vanishing of the sawtooth coefficient}\label{sec:v9-harm}
\begin{proposition}[Harmonic criterion]\label{prop:v9-harm}
Let $\omega\ne0$ be fixed and $F=F_h$. Assume
\begin{itemize}[leftmargin=2em]
\item[(HC)] for every integer $\mathfrak h\ne0$: $\displaystyle\frac1T\int_T^{2T}\exp\bigl(i\omega t+2\pi i\mathfrak h F(t)\bigr)\,dt\to0$;
\item[(NC)] $\displaystyle\lim_{\mathfrak d\downarrow0}\limsup_{T\to\infty}\frac1T\meas\bigl\{t\in[T,2T]:\operatorname{dist}(F(t),\mathbb Z+\tfrac12)<\mathfrak d\bigr\}=0$.
\end{itemize}
Then $\varsigma_T(\omega;F_h)\to0$. By the measure bound in the proof of \cref{lem:v9-sawcmp} the hypotheses may equivalently be stated for $F^\circ=N_0+\Phi^\circ$ or for $F^{\Cm}$.
\end{proposition}
\begin{proof}
Let $\mathrm{Fej}_{\mathfrak n}$ be the Fej\'er mean of order $\mathfrak n$ of the $1$-periodic function $\mathrm{saw}$ (mean zero, so no constant term); $|\mathrm{Fej}_{\mathfrak n}|\le\frac12$, and $\mathrm{Fej}_{\mathfrak n}\to\mathrm{saw}$ uniformly on $\{y:\operatorname{dist}(y,\mathbb Z+\frac12)\ge\mathfrak d\}$; call $\delta(\mathfrak n,\mathfrak d)\to0$ the supremum of $|\mathrm{saw}-\mathrm{Fej}_{\mathfrak n}|$ there. Since $|\mathrm{saw}-\mathrm{Fej}_{\mathfrak n}|\le1$,
\begin{equation}
|\varsigma_T|\le\Bigl|\frac{2\pi}T\int_T^{2T}e^{i\omega t}\mathrm{Fej}_{\mathfrak n}(F)\,dt\Bigr|+2\pi\Bigl[\delta(\mathfrak n,\mathfrak d)+\frac1T\meas\{\operatorname{dist}(F,\mathbb Z+\tfrac12)<\mathfrak d\}\Bigr].
\end{equation}
The first term is a finite combination of the integrals in (HC) and tends to $0$. Take $\limsup_T$, then $\mathfrak n\to\infty$, then $\mathfrak d\downarrow0$.
\end{proof}
\begin{remark}[Arithmetic form of (HC); moving resonances]\label{rem:v9-res}
For the unclipped phase $\Phi_{\le P}$ (the sum of \eqref{eq:Phi-m} restricted to the primes $p\le P$; here $m\ge1$ is the harmonic multiplicity, not the prime count of $\Phi_m$), $e^{-2\pi i m\Phi_{\le P}}=\prod_{p\le P}\bigl((1-p^{-1/2-it})/(1-p^{-1/2+it})\bigr)^{m}$ ($m\ge1$), by taking the logarithm of each factor. Expanding the product, the harmonic $\mathfrak h=-m$ in (HC) is a sum of terms $e^{i\mathfrak P(t)}$ with $\mathfrak P(t)=\omega t-2\pi mN_0(t)+t\log(a/b)$, $a,b$ products of prime powers; $\mathfrak P'(t)=\omega-m\log(t/2\pi)+\log(a/b)$ vanishes where $a/b=e^{-\omega}(t/2\pi)^m\asymp T^me^{-\omega}$. Thus the resonant frequencies move with the height, so orthogonality for a fixed finite set of frequencies does not suffice; this identifies what (HC) requires, and is a heuristic statement about stationary points, not an estimate of the sum.
\end{remark}

\subsection{The Landau line for the sequences}\label{sec:v9-lines}
\begin{theorem}[Reduction of $\Xi_T^{\mathrm C}$, $\Xi_T^{\Cm}$ to the sawtooth coefficient]\label{thm:v9-reduction}
Let $\star\in\{\mathrm C,\Cm\}$, $F=F_h$ resp.\ $F^{\Cm}$, and fixed $\omega\ne0$. Then, as $T\to\infty$,
\begin{equation}\label{eq:v9-reduction}
\frac{2\pi}T\,M_T\,\Xi_T^\star(\omega)=i\omega\,\varsigma_T(\omega;F)-i\omega\,\frac{2\pi}{T}\int_T^{2T}e^{i\omega s}\varrho_F(s)\,ds+o(1).
\end{equation}
\end{theorem}
\begin{proof}
Apply \eqref{eq:v9-stieltjes} on $(T,2T]$, split $N-N^{\mathsf x}=z_F+\varrho_F$ by \eqref{eq:v9-slip} and insert \eqref{eq:v9-slipcoeff}. The index sets $\{k:T<\gamma_k\le2T\}$ and $\{k:T<\mathsf x_k\le2T\}$ differ in at most $|N-N^{\mathsf x}|(T)+|N-N^{\mathsf x}|(2T)$ indices, each contributing $O(1)$ to the exponential sums; this mismatch is $O(\log T)$, as is the endpoint term $\mathsf Z_T=e^{2i\omega T}(N-N^{\mathsf x})(2T)-e^{i\omega T}(N-N^{\mathsf x})(T)$ of \eqref{eq:v9-stieltjes}, as follows. By the sandwich $\{N_0+H<\lambda\}\subset\{F<\lambda\}\subset\{N_0-H<\lambda\}$ (valid for $F=F_h$ because $h\le H$ for large $t$, and for $F=F^{\Cm}$ by definition) and the eventual monotonicity of $N_0\pm H$, the rearrangement satisfies $N_0-H\le F^*\le N_0+H$ for large $s$, hence $|N^{\mathsf x}(s)-N_0(s)|\le H(s)+\frac32$; with \eqref{eq:MCOMH} this gives $|N-N^{\mathsf x}|(s)\ll\log s$.
\end{proof}

\begin{proposition}[The rearrangement term of $\tau^{\mathrm C}$]\label{prop:v9-rearr}
For $F=F_h$ and fixed $\omega$,
\begin{equation}\label{eq:v9-rearr}
\frac{2\pi}{T}\int_T^{2T}e^{i\omega s}\varrho_F(s)\,ds=O_\omega\Bigl(\frac{(\log\log T)^2}{\log T}\Bigr).
\end{equation}
\end{proposition}
\begin{proof}
Put $\lambda_k=k-\frac12$, $\mathsf U_k=\{t>0:F(t)<\lambda_k\}$ and $\mathsf V_k=[0,\tau_k^{\mathrm C})$, so $\meas\mathsf U_k=\meas\mathsf V_k$. For large $t$ the barriers $\varphi^h_\pm=N_0\pm h$ are increasing and $\{\varphi^h_+<\lambda\}\subset\{F<\lambda\}\subset\{\varphi^h_-<\lambda\}$; let $r_k^-<r_k^+$ be the solutions of $\varphi^h_+(r_k^-)=\lambda_k=\varphi^h_-(r_k^+)$. Then $\mathsf U_k$ and $\mathsf V_k$ both contain $(0,r_k^-)$ and are contained in $(0,r_k^+)$, so $\mathsf U_k\,\triangle\,\mathsf V_k\subset\mathsf I_k=[r_k^-,r_k^+)$, of length $\mathfrak w_k=\bigl(h(r_k^+)+h(r_k^-)\bigr)/N_0'(\xi_k)\ll\log\log T/\log T$ when $\tau_k^{\mathrm C}\asymp T$ (mean value theorem). Pointwise, $\varrho_F(s)=\sum_k(\ind_{\mathsf V_k}-\ind_{\mathsf U_k})(s)$ (a finite sum: $\lfloor F+\frac12\rfloor=\#\{k:\lambda_k\le F\}$ and $N^{\mathsf x}=\#\{k:\tau_k^{\mathrm C}\le s\}$). If $\mathsf I_k\subset[T,2T]$, then since $\int(\ind_{\mathsf V_k}-\ind_{\mathsf U_k})=0$,
\begin{equation}
\Bigl|\int e^{i\omega s}(\ind_{\mathsf V_k}-\ind_{\mathsf U_k})\,ds\Bigr|
=\Bigl|\int(e^{i\omega s}-e^{i\omega\tau_k^{\mathrm C}})(\ind_{\mathsf V_k}-\ind_{\mathsf U_k})\,ds\Bigr|\le|\omega|\,\mathfrak w_k\cdot\mathfrak w_k .
\end{equation}
The $O(\log\log T)$ levels with $\mathsf I_k\ni T$ or $\ni2T$ contribute at most $\mathfrak w_k$ each, and levels with $\mathsf I_k\cap[T,2T]=\varnothing$ contribute $0$. Summing over the $M_T+O(\log\log T)$ relevant levels gives $O_\omega\bigl(T(\log\log T)^2/\log T\bigr)$.
\end{proof}
For $\tau^{\Cm}$ the band has width $O(1)$ and this argument gives nothing, which is not needed by \cref{prop:v9-CvsCm} and \cref{lem:v9-sawcmp} below.

\begin{proposition}[Landau line of $\tau^{\mathrm C}$ and $\tau^{\Cm}$]\label{prop:v9-line}
For $\star\in\{\mathrm C,\Cm\}$ and fixed $\omega\ne0$, as $T\to\infty$,
\begin{equation}\label{eq:v9-line}
\frac{2\pi}{T}\sum_{T<\mathsf x_k^\star\le2T}e^{i\omega\mathsf x_k^\star}=\chi(\omega)-i\omega\,\varsigma_T(\omega;F_h)+o(1).
\end{equation}
Hence, if $\varsigma_T(\omega;F_h)\to0$ the sequence $\mathsf x^\star$ carries the same Landau line $\chi(\omega)$ as the true zeros; if $\varsigma_T$ has a nonzero limit $\varsigma_\infty$, the line is displaced by $-i\omega\varsigma_\infty$.
\end{proposition}
\begin{proof}
For $\star=\mathrm C$ subtract \eqref{eq:v9-reduction} from \eqref{eq:v9-landau} and use \cref{prop:v9-rearr}; the index block $\mathfrak J_T$ is replaced by $\{T<\mathsf x_k^\star\le2T\}$ at the cost of $O(\log T)$ indices, i.e.\ $o(1)$ after multiplication by $2\pi/T$. For $\star=\Cm$ add the bound $\bigl|\sum_{\mathfrak J_T}(e^{i\omega\tau_k^{\mathrm C}}-e^{i\omega\tau_k^{\Cm}})\bigr|\ll|\omega|T/\log\log T$ of \cref{prop:v9-CvsCm}.
\end{proof}

\begin{proposition}[The RvM sequence]\label{prop:v9-rvm}
For fixed $\omega>0$, assuming (L), $\frac{2\pi}T M_T\Xi_T^{\mathrm{RvM}}(\omega)=\chi(\omega)+o(1)$. Thus
$\Xi_T^{\mathrm{RvM}}(\omega)=\chi(\omega)/\log T+o(1/\log T)$ (more precisely $\Xi_T^{\mathrm{RvM}}=\frac{T}{2\pi M_T}\chi(\omega)+o(1/\log T)$ unconditionally, with error $O(1/\log^2T)$ if the Ford--Zaharescu error $O(T/\log T)$ of (L) is used, and $O(\log T/M_T)$ under RH), and $\Xi_T^{\mathrm{RvM}}$ does not vanish at the frequencies $\omega=r\log p$.
\end{proposition}
\begin{proof}
$\Theta_T^{\mathrm{RvM}}(\omega,0)=O(\log T/M_T)$ by \eqref{eq:v9-KL}; the other term is \eqref{eq:v9-landau}.
\end{proof}

\begin{proposition}[Same coefficient for $\tau^{\mathrm C}$ and $\tau^{\Cm}$]\label{prop:v9-CvsCm}
For fixed $\omega$,
\begin{equation}
\Bigl|\sum_{k\in\mathfrak J_T}\bigl(e^{i\omega\tau_k^{\mathrm C}}-e^{i\omega\tau_k^{\Cm}}\bigr)\Bigr|\le|\omega|\sum_{k\in\mathfrak J_T}|\tau_k^{\mathrm C}-\tau_k^{\Cm}|\ll\frac{|\omega|\,T}{\log\log T},
\end{equation}
so $\Xi_T^{\mathrm C}(\omega)-\Xi_T^{\Cm}(\omega)=O\bigl(|\omega|/(\log T\log\log T)\bigr)=o(1/\log T)$: the two corrected sequences have the same coefficient $\mathfrak b(\omega)$ whenever one of them has one. If $\tau_k^{\mathrm C}$ is replaced by the point of the unclipped phase $u_k^\circ$ of \eqref{eq:v7-untruncated} (as in the file \texttt{casanya\_ordinal\_primes\_v50}, see \cref{rem:v9-clip}), the bound improves to $O(|\omega|T/(\log T)^3)$ by \cref{thm:v7-transport}.
\end{proposition}
\begin{proof}
For the pair $(\tau^{\mathrm C},\tau^{\Cm})$: the level-counting inequality gives $\sum_k|\tau_k^{\mathrm C}-\tau_k^{\Cm}|\le\int|\lfloor F_h+\frac12\rfloor-\lfloor F^{\Cm}+\frac12\rfloor|\le\int(|F_h-F^{\Cm}|+1)\ind_{\{F_h\ne F^{\Cm}\}}\le2\int_{T/2}^{3T}|\Phi^\circ|\ind_{\{|\Phi^\circ|>h\}}\ll T/\log\log T$ as in Step 3 of \cref{lem:v9-cancel}, all sublevel sets involved lying in $(0,3T)$ and containing $(0,T/2]$ as in the proof of \cref{thm:taucm}. For the improvement via \cref{thm:v7-transport}, indices in $\mathfrak J_T$ versus $\mathcal K_T^{\mathrm{RvM}}$ differ by $O(\log T)$ terms, each of modulus at most $2$ in an exponential sum.
\end{proof}

\begin{lemma}[The three phases have the same sawtooth coefficient]\label{lem:v9-sawcmp}
For fixed $\omega\ne0$, $\varsigma_T(\omega;F_h)-\varsigma_T(\omega;F^{\Cm})=O\bigl((\log\log T)^{-2}\bigr)$ and the same holds with $F^{\circ}=N_0+\Phi^\circ$ in place of $F^{\Cm}$.
\end{lemma}
\begin{proof}
$|\mathrm{saw}|\le\frac12$, so $|\mathrm{saw}(F_h)-\mathrm{saw}(F')|\le1$, and it vanishes unless $F_h\ne F'$, which forces $|\Phi^\circ|>h$. By Chebyshev and $\int_T^{2T}|\Phi^\circ|^4\ll T(\log\log T)^2$ (Step 3 of \cref{lem:v9-cancel}) this set has measure $\ll T/(\log\log T)^2$.
\end{proof}

\begin{corollary}[A statistic for $\varsigma_T$ from the points alone]\label{cor:v9-test}
For $\omega>0$ fixed and $\star\in\{\mathrm C,\Cm\}$,
\begin{equation}\label{eq:v9-test}
\varsigma_T(\omega;F_h)=\frac{i}{\omega}\Bigl[\frac{2\pi}{T}\sum_{T<\mathsf x_k^\star\le2T}e^{i\omega\mathsf x_k^\star}+\Lambda(e^{\omega})\,e^{-\omega/2}\Bigr]+o(1).
\end{equation}
In particular, for $\omega=1$ ($\Lambda(e)=0$): $\varsigma_T(1;F_h)\to0$ if and only if $\frac{2\pi}T\sum_{T<\tau_k^{\mathrm C}\le2T}e^{i\tau_k^{\mathrm C}}\to0$; and for $\omega=\log p$: if and only if $\frac{2\pi}T\sum e^{i\tau_k^{\mathrm C}\log p}+\frac{\log p}{\sqrt p}\to0$.
\end{corollary}
\begin{proof}
Solve \eqref{eq:v9-line} for $\varsigma_T$, with $\chi=-\Lambda(e^\omega)e^{-\omega/2}$.
\end{proof}
This is the formulation, in terms of the points $\tau_k^{\mathrm C}$ and the RvM block only, of the test in the ChatGPT text; the index block $\{T<\mathsf x_k\le2T\}$ may be replaced by $\{k:T<N_0^{-1}(k-\tfrac12)\le2T\}$ at the cost of $O(\log T)$ indices (MCOM).

\begin{proposition}[Mean displacement]\label{prop:v9-mean}
$\displaystyle\frac1{M_T}\sum_{k\in\mathfrak J_T}\bigl(\gamma_k-\tau_k^{\mathrm C}\bigr)=\frac{\varsigma_T(0;F_h)}{\log T}+o\Bigl(\frac1{\log T}\Bigr)$, with $\varsigma_T(0;F_h)=-\frac{2\pi}T\int_T^{2T}\mathrm{saw}(F_h)$.
\end{proposition}
\begin{proof}
$\sum_k(\gamma_k-\tau_k^{\mathrm C})=\int(N^{\mathsf x}-N)\,ds$ over the block up to $O(\log T)\cdot O(1)$ from straddling indices, and $N-N^{\mathsf x}=D_F+\mathrm{saw}(F)+\varrho_F$. Here $\int_T^{2T}\varrho_F=o(T)$ by \cref{prop:v9-rearr}, $\int_T^{2T}(N-N_0)=\int S+O(1)=O(\log T)$ (classical, \cite{Titchmarsh}), and $\int_T^{2T}\Phi^\circ=O(T^{1/4+\epsilon})$ (each Euler term integrates to $O(1/\log p)$), while the clip changes $\int_T^{2T}\Phi^\circ$ by at most $\int_T^{2T}(|\Phi^\circ|-h)_+\ll T/\log\log T$ (Step 3 of \cref{lem:v9-cancel}). Hence $\int_T^{2T}D_F=o(T)$ and the sum equals $-\int_T^{2T}\mathrm{saw}(F_h)+o(T)=\frac{T}{2\pi}\varsigma_T(0;F_h)+o(T)$; divide by $M_T=\frac T{2\pi}\log T(1+o(1))$.
\end{proof}

\begin{remark}[$\tau^{\mathrm C}$ clipped or unclipped]\label{rem:v9-clip}
In this note $\tau_k^{\mathrm C}=u_k(F_h)$ with the loglog clip, \eqref{eq:old}, which is also the definition of the current IJNT manuscript (\texttt{casanya\_\allowbreak ordinal\_\allowbreak primes\_\allowbreak ijnt\_\allowbreak v5\_\allowbreak final}, truncated phase at height $\log\log\max(t,e^e)$, sublevel construction of that manuscript). The file \texttt{casanya\_ordinal\_primes\_v50} defines $\tau^{\mathrm C}$ from the unclipped phase $N_0+\Phi^\circ$; that definition survives here only through the auxiliary points $u_k^\circ$ of \eqref{eq:v7-untruncated}. A comparison between $\tau^{\mathrm C}$ and $\tau^{\Cm}$ that uses only the set $\{|\Phi^\circ|\ge H\}$ is correct for the unclipped points only; for $F_h$ the set $\{|\Phi^\circ|>h\}$ must be used, which is why \cref{prop:v9-CvsCm} has $T/\log\log T$ and not $T/(\log T)^3$.
\end{remark}

\subsection{Gaussian scale for RvM and coherence of the corrected sequences}\label{sec:v9-coherence}
Put $\varpi_T=\sqrt{2\log\log T}/\log T$ and $\omega_T(\psi)=\psi/\varpi_T=\psi\log T/\sqrt{2\log\log T}$ for fixed $\psi\in\mathbb R$.

\begin{theorem}[Coherence of the corrected sequences]\label{thm:v9-coherence}
Under the inputs of \cref{thm:taucm}, for $\star\in\{\mathrm C,\Cm\}$ and every $\omega=\omega_T$ with $|\omega_T|=o(\log T)$ (in particular $\omega=\omega_T(\psi)$),
\begin{equation}
\Theta_T^\star(0,\omega_T)\longrightarrow1,\qquad
\frac1{M_T}\sum_{k\in\mathfrak J_T}\bigl|(\nabla^2\mathbb S)_k(\omega_T;\gamma)-(\nabla^2\mathbb S)_k(\omega_T;\mathsf x^\star)\bigr|^2\longrightarrow0 .
\end{equation}
The first statement holds with rate $O(|\omega_T|/\log T)$.
\end{theorem}
\begin{proof}
By \eqref{eq:cm-mean} and Cauchy--Schwarz, $\sum_{\mathfrak J_T}|\sigma_k^{\Cm}|\le(M_T\sum\sigma_k^2)^{1/2}\ll(T\log T\cdot T/\log T)^{1/2}=T$, i.e.\ the mean of $|\sigma_k^{\Cm}|$ over $\mathfrak J_T$ is $O(1/\log T)$. For $\mathrm C$ add $\sum|\tau^{\mathrm C}-\tau^{\Cm}|\ll T/\log\log T$ (\cref{prop:v9-CvsCm}). Now $|e^{i\omega\sigma}-1|\le\min(2,|\omega||\sigma|)$ gives $|1-\Theta_T^\star(0,\omega)|\le|\omega|\cdot O(1/\log T)$, and by \eqref{eq:v9-seconddiff} the second quantity is $2(1-\Re\Theta)\le2|1-\Theta|$.
\end{proof}

\begin{proposition}[Gaussian scale for the RvM points]\label{prop:v9-gauss}
Assume the Selberg central limit theorem holds for $\mathcal S_k:=k-N_0(\gamma_k)$ (which equals $S(\gamma_k)+R(\gamma_k)$ when $\gamma_k<\gamma_{k+1}$, i.e.\ when $N(\gamma_k)=k$) when $k$ ranges over $\mathfrak J_T$ with the variance $\frac1{2\pi^2}\log\log T$ (the passage from the uniform $t\in[T,2T]$ to the points $\gamma_k$ is \emph{not proved} here; the result \cite{Selberg1946} is stated for $t$). Then $\Theta_T^{\mathrm{RvM}}(0,\omega_T(\psi))\to e^{-\psi^2/2}$. Consequently
\begin{equation}
\frac1{M_T}\sum_{k\in\mathfrak J_T}\bigl|(\nabla^2\mathbb S)_k(\omega_T(\psi);\gamma)-(\nabla^2\mathbb S)_k(\omega_T(\psi);\gamma^{\mathrm{RvM}})\bigr|^2\to2\bigl(1-e^{-\psi^2/2}\bigr).
\end{equation}
\end{proposition}
\begin{proof}
By definition $N_0(\gamma_k)-(k-\frac12)=\frac12-\mathcal S_k$, so $\sigma_k^{\mathrm{RvM}}=\frac{2\pi}{\log(\gamma_k/2\pi)}\bigl(\frac12-\mathcal S_k\bigr)(1+o(1))$. The standard deviation of this is $\varpi_T(1+O(1/\log T))$, and the shift $\frac12$ is multiplied by $\omega_T\cdot\pi/\log T=O(1/\sqrt{\log\log T})\to0$.
\end{proof}
Thus, under the hypothesis of \cref{prop:v9-gauss}, on the scale $\omega\asymp\varpi_T^{-1}$ the RvM sequence has a non-trivial Gaussian dispersion of its displacements, whereas, under the inputs of \cref{thm:taucm}, $\tau^{\mathrm C},\tau^{\Cm}$ keep coherent phases. The scale $\omega\asymp\log T$ is not reached by the second-moment input, which only gives $O(1)$ for $\omega\sigma$ there (open, \S\ref{sec:v9-open-centroid}). Kargin's Gaussian law for $\gamma_k-\gamma_k^{\mathrm{RvM}}$ with $k$ uniform in a window of size $T^{\theta}$, $\frac12<\theta\le1$ \cite{Kargin2013}, is the natural source for \cref{prop:v9-gauss}; whether $\theta=1$ (the full block) and the centering and scale used there coincide with $\gamma_k^{\mathrm{RvM}}$ and $\varpi_T$ has not been checked here.

\subsection{Reconstruction at frequency \texorpdfstring{$1$}{1} and collision energy}\label{sec:v9-recon}
Write $\Delta^{\mathsf x}_k=\mathsf x_{k+1}-\mathsf x_k$ (for $\mathsf x=\tau^{\mathrm{Cms}}(X;W)$ this is the $\Delta_k(X;W)$ of the earlier sections), $h_k^{\mathrm{RvM}}=\gamma_{k+1}^{\mathrm{RvM}}-\gamma_k^{\mathrm{RvM}}$, and for a cutoff $\omega_\sharp>0$
\begin{equation}\label{eq:v9-energy}
\mathfrak E_k(\omega_\sharp;\mathsf x)=\frac1{2\omega_\sharp}\int_{-\omega_\sharp}^{\omega_\sharp}\bigl|(\nabla^2\mathbb S)_{k+1}(\omega;\mathsf x)-(\nabla^2\mathbb S)_k(\omega;\mathsf x)\bigr|^2d\omega .
\end{equation}
\begin{proposition}[Exact reconstruction and collision energy]\label{prop:v9-recon}
Let $\mathsf x\in\{\gamma,\gamma^{\mathrm{RvM}},\tau^{\mathrm C},\tau^{\Cm}\}$ and put $\mathfrak c_{\mathsf x}=0$ for RvM, $\mathfrak c_{\mathsf x}=0.08382$ for $\tau^{\mathrm C},\tau^{\Cm}$, and for $\mathsf x=\gamma$ any $\mathfrak c_{\mathsf x}>C_{\mathrm H}$ (\eqref{eq:asymptMCOM}). Then for every $\epsilon>0$ and $k\ge k_0(\epsilon)$:
\begin{enumerate}[label=(\roman*),leftmargin=*]
\item $|\mathsf x_k-\gamma_k^{\mathrm{RvM}}|\le2\pi\mathfrak c_{\mathsf x}+\epsilon$; numerically $<0.5267+\epsilon$ for $\tau^{\mathrm C},\tau^{\Cm}$ and $<0.4704$ for $\gamma$ (with $\mathfrak c_\gamma=C_{\mathrm H}+$ small); in all cases $<\pi$.
\item $\displaystyle\mathsf x_k=\gamma_k^{\mathrm{RvM}}+\operatorname{Arg}\bigl[e^{-i\gamma_k^{\mathrm{RvM}}}(\nabla^2\mathbb S)_k(1;\mathsf x)\bigr]$.
\item $|\Delta_k^{\mathsf x}|\le h_k^{\mathrm{RvM}}+4\pi\mathfrak c_{\mathsf x}+\epsilon<\pi$ (at most $1.0534$ for $\tau^{\mathrm C},\tau^{\Cm}$, $0.9408$ for $\gamma$), hence $\Delta_k^{\mathsf x}=\operatorname{Arg}\bigl[(\nabla^2\mathbb S)_{k+1}(1;\mathsf x)\,\overline{(\nabla^2\mathbb S)_k(1;\mathsf x)}\bigr]$.
\item $\mathfrak E_k(\omega_\sharp;\mathsf x)=2\bigl[1-\sin(\omega_\sharp\Delta_k^{\mathsf x})/(\omega_\sharp\Delta_k^{\mathsf x})\bigr]$ (value $0$ at $\Delta_k^{\mathsf x}=0$). Hence, for fixed $\omega_\sharp$ and fixed $k$, and for every family $\mathsf x^{(X)}$ of nondecreasing real sequences (the formula for $\mathfrak E_k$ holds for every real sequence), $\liminf_{X\to\infty}\mathfrak E_k(\omega_\sharp;\mathsf x^{(X)})>0$ if and only if $\liminf_{X\to\infty}\Delta_k^{\mathsf x^{(X)}}>0$ (the function $2(1-\sin y/y)$ vanishes only at $y=0$ and is bounded below away from it).
\end{enumerate}
\end{proposition}
\begin{proof}
(i) For $\tau^{\Cm}$, the sandwich by $N_0\pm H$ (as in the proof of \cref{thm:v9-reduction}) gives $|\tau_k^{\Cm}-\gamma_k^{\mathrm{RvM}}|\le(H/N_0')(\gamma_k^{\mathrm{RvM}})(1+o(1))=2\pi\bigl(0.08382+5.5/\log t\bigr)(1+o(1))\to2\pi\cdot0.08382$; $\tau^{\mathrm C}$ has $h\le H$. For the zeros, $N(\gamma_k^-)\le k-1<k\le N(\gamma_k)$ and \eqref{eq:RVM} (applied at $\gamma_k^-$ and at $\gamma_k$) give $S(\gamma_k^-)+\frac12\le k-\frac12-N_0(\gamma_k)-R(\gamma_k)\le S(\gamma_k)-\frac12$, hence $|\gamma_k-\gamma_k^{\mathrm{RvM}}|\le\bigl(\max|S(\gamma_k^{\pm})|+\frac12+|R(\gamma_k)|\bigr)/N_0'(\xi_k)$ for some $\xi_k$ between $\gamma_k^{\mathrm{RvM}}$ and $\gamma_k$, and $\limsup|S|/\log t<\mathfrak c_\gamma$ applies to the one-sided values $S(\gamma_k^\pm)$ as well. (ii) is \cref{lem:v9-exact}(a) at $\omega=1$ together with (i): the real number $\mathsf x_k-\gamma_k^{\mathrm{RvM}}$ lies in $(-\pi,\pi)$, so it is the principal argument of its exponential. (iii) follows from (i) and $h_k^{\mathrm{RvM}}\to0$. (iv): $|e^{i\omega\mathsf x_{k+1}}-e^{i\omega\mathsf x_k}|^2=2(1-\cos\omega\Delta_k^{\mathsf x})$, averaged over $[-\omega_\sharp,\omega_\sharp]$.
\end{proof}
This is an exact reconstruction of the points (and spacings) from the sequence of centroids $\mathrm{Cen}_\ell(1;\mathsf x)$ and the RvM points; it does not say that the centroid built from primes has the same limit as the one built from the zeros. The natural frequency in (iv) is $\omega_\sharp=1/h_k^{\mathrm{RvM}}\sim\log\gamma_k^{\mathrm{RvM}}/2\pi$, the microscopic scale; the arithmetic input still missing for non-collision as $X\to\infty$ is the one already stated in the earlier sections.

\begin{remark}[Cauchy--Schwarz transport]\label{rem:v9-cs}
By $|e^{i\omega\gamma}-e^{i\omega\mathsf x}|\le\min\{2,|\omega|\,|\gamma-\mathsf x|\}$ and Cauchy--Schwarz, $|\Xi_T^\star(\omega)|\le|\omega|\bigl(M_T^{-1}\sum_{\mathfrak J_T}\sigma_k^2\bigr)^{1/2}$. The IJNT mean-square asymptotics for the RvM points (\cite{CasanaIJNT}; in the notation of this paper $\frac1{M_T}\sum(\gamma_k^{\mathrm{RvM}}-\gamma_k)^2=\frac{2\log\log T}{\log^2T}(1+o(1))$) give $|\Xi_T^{\mathrm{RvM}}|\le|\omega|\varpi_T(1+o(1))$, and \eqref{eq:cm-mean} gives $|\Xi_T^{\Cm}|\ll|\omega|/\log T$ (same for $\mathrm C$ with the IJNT bound). Hence $\sup_{|\omega|\le U(T)}|\Xi_T^\star(\omega)|\to0$ for $U(T)=o(\log T)$ for the corrected sequences and $U(T)=o(\log T/\sqrt{\log\log T})$ for RvM. These are upper bounds, much weaker than \cref{prop:v9-rvm} at fixed $\omega$.
\end{remark}

\subsection{Centroids}\label{sec:v9-cen}
\begin{theorem}[Centroid difference from a block coefficient]\label{thm:v9-cen}
Let $\star\in\{\mathrm{RvM},\mathrm C,\Cm\}$, $\omega\ne0$ and $\mathfrak b\in\mathbb C$. Assume that for \emph{every} large real $T$
\begin{equation}
\Xi_T^\star(\omega)=\frac{\mathfrak b}{\log T}+o\Bigl(\frac1{\log T}\Bigr)
\end{equation}
(a single coefficient $\mathfrak b=\mathfrak b^\star(\omega)$ for all $T$, not only dyadic ones). Then, as $\ell\to\infty$,
\begin{equation}\label{eq:v9-cen}
\mathrm{Cen}_\ell(\omega;\gamma)-\mathrm{Cen}_\ell(\omega;\mathsf x^\star)=\frac{\mathfrak b\,\ell}{2\log\ell}+o\Bigl(\frac\ell{\log\ell}\Bigr),
\qquad
\sum_{k\le\ell}\bigl(e^{i\omega\gamma_k}-e^{i\omega\mathsf x_k^\star}\bigr)=\frac{\mathfrak b\,\ell}{\log\ell}+o\Bigl(\frac\ell{\log\ell}\Bigr).
\end{equation}
\end{theorem}
\begin{proof}
Put $\mathsf f_k=e^{i\omega\gamma_k}-e^{i\omega\mathsf x_k^\star}$ and $\mathbb Y(U)=\sum_{\gamma_k\le U}\mathsf f_k$. Since $M_T=\frac{T\log T}{2\pi}(1+o(1))$, the hypothesis reads $\mathbb Y(2T)-\mathbb Y(T)=\frac{T}{2\pi}(\mathfrak b+o(1))$. Telescoping over $T,T/2,T/4,\dots$ (the contribution of $T\le T_0$ is $O(1)$) gives $\mathbb Y(U)=\frac U{2\pi}(\mathfrak b+o(1))$. As $\gamma_n=\frac{2\pi n}{\log n}(1+o(1))$, $\sum_{k\le n}\mathsf f_k=\frac{\mathfrak b\,n}{\log n}(1+o(1))$, which is the second formula in \eqref{eq:v9-cen}. Summation over $n\le\ell$ and division by $\ell$ give the first, since $\sum_{n\le\ell}n/\log n=\frac{\ell^2}{2\log\ell}(1+o(1))$.
\end{proof}

\begin{corollary}[What is determined]\label{cor:v9-cen}
(i) For RvM and $\omega>0$: $\mathfrak b^{\mathrm{RvM}}(\omega)=\chi(\omega)$ by \cref{prop:v9-rvm}; hence $\mathrm{Cen}_\ell(\omega;\gamma)-\mathrm{Cen}_\ell(\omega;\gamma^{\mathrm{RvM}})=\chi(\omega)\ell/(2\log\ell)+o(\ell/\log\ell)$, which is non-zero for $\omega=r\log p$: the centroids of the zeros and of the RvM points differ at the order $\ell/\log\ell$ at the frequencies $\omega=r\log p$. For $\ell=N(T)$, $\ell/(2\log\ell)=\frac T{4\pi}(1+o(1))$, so the difference is $-\frac{T}{4\pi}\Lambda(e^\omega)e^{-\omega/2}+o(T)$, e.g.\ $-\frac{T\log p}{4\pi p^{r/2}}+o(T)$ at $\omega=r\log p$; the same formula appears in the ChatGPT text.

(ii) For $\tau^{\mathrm C},\tau^{\Cm}$: $\mathfrak b^{\mathrm C}(\omega)=\mathfrak b^{\Cm}(\omega)=i\omega\lim_T\varsigma_T(\omega;F_h)$ whenever this limit exists (no further hypothesis, by \cref{prop:v9-rearr} and \cref{prop:v9-CvsCm}). In particular $\mathfrak b=0$ if and only if the sawtooth coefficient tends to $0$, and then $\mathrm{Cen}_\ell(\omega;\gamma)-\mathrm{Cen}_\ell(\omega;\tau^{\mathrm C})=o(\ell/\log\ell)$.

(iii) For $\omega=1$: $e$ is not a prime power, $\chi(1)=0$. By (K), $\mathrm{Par}_n(1;\gamma^{\mathrm{RvM}})=O(\log n)$, hence $\mathrm{Cen}_\ell(1;\gamma^{\mathrm{RvM}})=O(\log\ell)$, and by \cref{thm:v9-cen} and (L) (unconditional error $o(T)$)
\begin{equation}
\mathrm{Cen}_\ell(1;\gamma)=o\bigl(\ell/\log\ell\bigr).
\end{equation}
Under RH the Landau error is $O(\log T)$ and one gets $\mathrm{Par}_n(1;\gamma)=O(\log n)$, $\mathrm{Cen}_\ell(1;\gamma)=O(\log\ell)$. The unconditional bound is therefore $o(\ell/\log\ell)$; the bound $O(\log\ell)$ is obtained here only from the $O(\log T)$ Landau error just quoted. With the Ford--Zaharescu error $O(T/\log T)$ of (L), $\mathbb Y(U)=O(U/\log U)$ and the same summation gives the sharper $\mathrm{Cen}_\ell(1;\gamma)=O(\ell/\log^2\ell)$.
\end{corollary}
\begin{proof}
(i) is \cref{thm:v9-cen} with \cref{prop:v9-rvm}. (ii) combines \cref{prop:v9-line}, \cref{prop:v9-rearr}, \cref{prop:v9-CvsCm} and \cref{thm:v9-cen}. (iii): $\mathrm{Cen}_\ell(1;\gamma)=\mathrm{Cen}_\ell(1;\gamma^{\mathrm{RvM}})+[\mathrm{Cen}_\ell(1;\gamma)-\mathrm{Cen}_\ell(1;\gamma^{\mathrm{RvM}})]$ and the bracket is $0\cdot\ell/(2\log\ell)+o(\ell/\log\ell)$; the RH variant uses $\sum_{0<\gamma\le T}e^{i\gamma}=O(\log T)$, which follows from Landau--Gonek when $\Lambda(e)=0$ and all $\beta=\frac12$.
\end{proof}

\subsection{The four centroids: formulas, differences, and convergence with and without RH}\label{sec:v11-four}
This subsection collects, in the notation of \cref{def:v9-objects}, what is known about the centroids of $\gamma$, $\gamma^{\mathrm{RvM}}$, $\tau^{\mathrm C}$ and $\tau^{\Cm}$ and about their differences. Where a statement comes from the author's earlier work it is marked at that place: the Riemann-hypothesis limit, the Tauberian criterion and the window equivalence are from the author's notes on the zeta--RvM paper \cite{CasanaV26} (quoted, proofs not reproduced here); the values $K^0$ and the non-convergence record are from the README of the repository directory \texttt{KC\_\allowbreak tauC\_\allowbreak 2026-11-02}; the convergence of the RvM centroid is stated in the author's notes and is proved here. No computation is made in this subsection.

\paragraph{Formulas.} For real $\omega$ and $\ell\ge1$ the four centroids are $\mathrm{Cen}_\ell(\omega;\mathsf y)$ of \eqref{eq:v9-Cen} for $\mathsf y=\gamma,\ \gamma^{\mathrm{RvM}},\ \tau^{\mathrm C},\ \tau^{\Cm}$; at the frequency of the author's earlier work, $\omega=1$,
\begin{equation}
\mathrm{Cen}_\ell(1;\mathsf y)=\frac1\ell\sum_{k=1}^{\ell}(\ell-k+1)\,e^{i\mathsf y_k}.
\end{equation}
For $\star\in\{\mathrm{RvM},\mathrm C,\Cm\}$ put
\begin{equation}\label{eq:v11-diff}
\mathsf D_\ell^\star(\omega)=\mathrm{Cen}_\ell(\omega;\gamma)-\mathrm{Cen}_\ell(\omega;\mathsf x^\star)
=\frac1\ell\sum_{k\le\ell}(\ell-k+1)\,e^{i\omega\mathsf x_k^\star}\bigl(e^{i\omega\sigma_k^\star}-1\bigr)
=\frac1\ell\sum_{n\le\ell}\ \sum_{k\le n}\mathsf f_k ,
\end{equation}
with $\sigma_k^\star=\gamma_k-\mathsf x_k^\star$ and $\mathsf f_k=e^{i\omega\gamma_k}-e^{i\omega\mathsf x_k^\star}$ as in \cref{thm:v9-cen}. Splitting $e^{i\omega\sigma}-1=i\sin\omega\sigma+(\cos\omega\sigma-1)$ separates the angular part (weights $i\sin\omega\sigma_k^\star$) from the radial part (weights $\cos\omega\sigma_k^\star-1$) of the centroid defect, the decomposition used in the author's drift analysis. The centroids of two approximating sequences differ by
\begin{equation}\label{eq:v11-between}
\mathrm{Cen}_\ell(\omega;\mathsf x^{\star})-\mathrm{Cen}_\ell(\omega;\mathsf x^{\star'})=\mathsf D_\ell^{\star'}(\omega)-\mathsf D_\ell^{\star}(\omega),
\end{equation}
and the second differences of $\ell\,\mathrm{Cen}_\ell$ are $e^{i\omega\mathsf y_k}$ itself (\cref{lem:v9-exact}). Everything below follows from \eqref{eq:v11-diff}, \eqref{eq:v11-between} and the block coefficients of \S\ref{sec:v9-lines}.

\begin{proposition}[Centroid of the RvM points converges, unconditionally]\label{prop:v11-rvm}
For fixed $\omega>0$ the limit
\begin{equation}\label{eq:v11-mid}
\mathfrak m(\omega)=\lim_{n\to\infty}\Bigl(\mathrm{Par}_n(\omega;\gamma^{\mathrm{RvM}})+\frac i\omega\,N_0'(\gamma_n^{\mathrm{RvM}})\,e^{i\omega\gamma_n^{\mathrm{RvM}}}-\frac12\,e^{i\omega\gamma_n^{\mathrm{RvM}}}\Bigr),\qquad N_0'(t)=\frac1{2\pi}\log\frac t{2\pi},
\end{equation}
exists, and $\mathrm{Cen}_\ell(\omega;\gamma^{\mathrm{RvM}})=\mathfrak m(\omega)+O_\omega(1/\log\ell)$. For $\omega=1$ the author's value of the limit is $\mathfrak m(1)=0.04157\ldots+1.75754\ldots\,i$ (quoted from the notes; not recomputed).
\end{proposition}
\begin{proof}
Put $t(s)=N_0^{-1}(s)$ and fix an integer $m_0$ so large that $N_0$ is increasing on $[t(m_0),\infty)$ and $\omega\,t'(s)\le\pi$ for $s\ge m_0$, where $t'(s)=1/N_0'(t(s))=2\pi/\log(t(s)/2\pi)$ is positive and decreasing. Put $g^{\mathrm{RvM}}_\omega(s)=e^{i\omega t(s)}$ and $t_n=\gamma_n^{\mathrm{RvM}}=t(n-\frac12)$, so that $\mathrm{Par}_n=\mathrm{Par}_{m_0}+\sum_{m_0<k\le n}g^{\mathrm{RvM}}_\omega(k-\frac12)$ for $n>m_0$, and $(g^{\mathrm{RvM}}_\omega)'(s)=i\omega\,t'(s)\,g^{\mathrm{RvM}}_\omega(s)$.

\emph{Exact midpoint identity.} For $f\in C^1[k-1,k]$, integration by parts on $[k-1,k-\frac12]$ and on $[k-\frac12,k]$ gives $\int_{k-1}^kf\,ds=f(k-\frac12)-\int_{k-1}^k\mathrm{saw}(s)f'(s)\,ds$, with $\mathrm{saw}(s)=s-\lfloor s+\frac12\rfloor$ as in \eqref{eq:v9-sawD}. Summing over $k$,
\begin{equation}\label{eq:v11-mid-exact}
\sum_{m_0<k\le n}g^{\mathrm{RvM}}_\omega(k-\tfrac12)=\int_{m_0}^ng^{\mathrm{RvM}}_\omega\,ds+\int_{m_0}^n\mathrm{saw}(s)\,(g^{\mathrm{RvM}}_\omega)'(s)\,ds .
\end{equation}

\emph{The sawtooth integral.} One has $\mathrm{saw}(s)=-\sum_{\nu\ge1}(-1)^\nu\sin(2\pi\nu s)/(\pi\nu)$ for $s\notin\frac12+\mathbb Z$, the series converging in $L^2$ on every bounded interval; since $(g^{\mathrm{RvM}}_\omega)'$ is bounded on bounded intervals, it may be integrated termwise against $(g^{\mathrm{RvM}}_\omega)'$ there. For $\nu\ge1$ put $\Psi^{\mathrm{RvM}}_{\nu,\pm}(s)=\omega t(s)\pm2\pi\nu s$; then $\sin(2\pi\nu s)\,\omega t'(s)e^{i\omega t(s)}=\frac1{2i}\bigl(\omega t'e^{i\Psi^{\mathrm{RvM}}_{\nu,+}}-\omega t'e^{i\Psi^{\mathrm{RvM}}_{\nu,-}}\bigr)$, $|(\Psi^{\mathrm{RvM}}_{\nu,\pm})'|=|\omega t'\pm2\pi\nu|\ge2\pi\nu-\pi\ge\pi\nu$ on $[m_0,\infty)$, and $\omega t'/(\Psi^{\mathrm{RvM}}_{\nu,\pm})'$ is monotone there, of modulus at most $\omega t'(m)/(\pi\nu)$ on $[m,\infty)$, because $\omega t'$ is decreasing. Writing $\omega t'e^{i\Psi}=\frac{\omega t'}{i\Psi'}\frac d{ds}e^{i\Psi}$ with $\Psi=\Psi^{\mathrm{RvM}}_{\nu,\pm}$ and integrating by parts, for $m_0\le m<M$,
\begin{equation}
\Bigl|\int_m^M\omega t'(s)e^{i\Psi^{\mathrm{RvM}}_{\nu,\pm}(s)}\,ds\Bigr|\le\frac{3\,\omega t'(m)}{\pi\nu},
\end{equation}
hence
\begin{equation}
\Bigl|\int_m^M\mathrm{saw}(s)\,(g^{\mathrm{RvM}}_\omega)'(s)\,ds\Bigr|\le\sum_{\nu\ge1}\frac1{\pi\nu}\cdot\frac{3\,\omega t'(m)}{\pi\nu}=\frac{\omega\,t'(m)}2=\frac{\pi\omega}{\log(t(m)/2\pi)} .
\end{equation}
Consequently $\int_{m_0}^n\mathrm{saw}\,(g^{\mathrm{RvM}}_\omega)'\,ds$ has a limit as $n\to\infty$, and differs from it by at most $\pi\omega/\log(t(n)/2\pi)=O_\omega(1/\log n)$.

\emph{The main integral.} Changing variable $s=N_0(t)$,
$\int_{m_0}^ng^{\mathrm{RvM}}_\omega\,ds=\int_{t(m_0)}^{t(n)}e^{i\omega t}N_0'(t)\,dt=\frac1{i\omega}\bigl[e^{i\omega t}N_0'(t)\bigr]_{t(m_0)}^{t(n)}-\frac1{i\omega}\int_{t(m_0)}^{t(n)}\frac{e^{i\omega t}}{2\pi t}\,dt$, and the last integral converges as $n\to\infty$, with tail $O(1/(\omega\,t(n)))$ by one more integration by parts. By \eqref{eq:v11-mid-exact} there is therefore a constant $\mathfrak m(\omega)$ with
\begin{equation}
\mathrm{Par}_n=\mathfrak m(\omega)+\frac1{i\omega}e^{i\omega t(n)}N_0'(t(n))+O_\omega\Bigl(\frac1{\log n}\Bigr).
\end{equation}
The upper limit of the integral is $s=n$, that is $t(n)$, and not $t_n=t(n-\frac12)$. Since $N_0(t(n))-N_0(t_n)=\frac12$ and $N_0''(t)=1/(2\pi t)$, $\varDelta_n:=t(n)-t_n$ satisfies $\varDelta_nN_0'(t_n)=\frac12+O(\varDelta_n^2/t_n)$ and $\varDelta_n\asymp1/\log n$, so
\begin{equation}
\begin{aligned}
e^{i\omega t(n)}N_0'(t(n))&=e^{i\omega t_n}\bigl(1+i\omega\varDelta_n+O(\omega^2\varDelta_n^2)\bigr)\bigl(N_0'(t_n)+O(\varDelta_n/t_n)\bigr)\\
&=e^{i\omega t_n}N_0'(t_n)+\tfrac{i\omega}2e^{i\omega t_n}+O_\omega\Bigl(\frac1{\log n}\Bigr).
\end{aligned}
\end{equation}
Hence, since $\frac1{i\omega}=-\frac i\omega$,
$\mathrm{Par}_n=\mathfrak m(\omega)-\frac i\omega N_0'(t_n)e^{i\omega t_n}+\frac12e^{i\omega t_n}+O_\omega(1/\log n)$, which is \eqref{eq:v11-mid}. Without the term $-\frac12e^{i\omega t_n}$ the bracket in \eqref{eq:v11-mid} would keep an oscillation of amplitude $\frac12$ and have no limit.

\emph{The centroid.} Averaging over $n\le\ell$, the $O_\omega(1/\log n)$ errors give $O_\omega(1/\log\ell)$, and the two oscillating terms contribute $\frac1\ell\sum_{n\le\ell}\bigl(\frac12-\frac i\omega N_0'(t_n)\bigr)e^{i\omega t_n}=O_\omega(\log^2\ell/\ell)$ by Abel summation from \eqref{eq:v9-KL} (the weights $N_0'(t_n)$ increase by $O(1/n)$ per step and are $O(\log n)$). So the Ces\`aro mean converges to $\mathfrak m(\omega)$ with the rate stated.
\end{proof}

\begin{proposition}[The centroid of the zeros $\gamma$]\label{prop:v11-gamma}
Let $\omega>0$.
\begin{enumerate}[label=(\alph*),leftmargin=*]
\item \emph{(Unconditional, from \cref{cor:v9-cen}(i) and \cref{prop:v11-rvm}.)} $\mathrm{Cen}_\ell(\omega;\gamma)=\mathfrak m(\omega)+\mathsf D_\ell^{\mathrm{RvM}}(\omega)+O(1/\log\ell)$, and $\mathsf D_\ell^{\mathrm{RvM}}(\omega)=\chi(\omega)\,\ell/(2\log\ell)+o(\ell/\log\ell)$. Consequently $\mathrm{Cen}_\ell(\omega;\gamma)$ is unbounded when $\omega=r\log p$ (it grows like $-\Lambda(e^{\omega})e^{-\omega/2}\,\ell/(2\log\ell)$), and is $o(\ell/\log\ell)$ when $\Lambda(e^\omega)=0$, in particular for $\omega=1$. At $\omega=1$ the centroid of $\gamma$ converges if and only if $\mathsf D_\ell^{\mathrm{RvM}}(1)$ converges.
\item \emph{(Under RH; the author's notes, Level~1: Landau's formula at $\Upsilon=e$, Cram\'er's continuation \cite{Cramer1919}, local Tauberian theorem \cite{DebruyneVindas2019}; in the author's manuscript \cite{CasanaV26} this is Theorem~5.15, resting on Theorem~5.5 (window theorem) and Lemma~5.14 (local Tauberian lemma); quoted, proof not reproduced.)} $\mathrm{Par}_n(1;\gamma)=O(\log n)$ and
\begin{equation}
\lim_{\ell\to\infty}\mathrm{Cen}_\ell(1;\gamma)=e^{-1/2}\,\mathfrak V(1),\qquad \mathfrak V(z)=\sum_{\operatorname{Im}\rho>0}e^{\rho z},
\end{equation}
where $\mathfrak V$ denotes Cram\'er's meromorphic continuation (simple poles at $\pm\log p^m$); the limit $e^{-1/2}\mathfrak V(1)$ has the value $0.77689426079325\ldots+0.76785566886455\ldots\,i$ and real part $\frac1{2\sqrt e}\bigl(e-\frac1{e^2-1}\bigr)$. Therefore, under RH, $\lim_\ell\mathsf D_\ell^{\mathrm{RvM}}(1)=K^0$ exists and
\begin{equation}
K^0=e^{-1/2}\mathfrak V(1)-\mathfrak m(1)\qquad(\text{RH}).
\end{equation}
The two quoted constants give $0.77689-0.04157=0.73532$ and $0.76786-1.75754=-0.98968$, in agreement to four decimals with the value $K^0=0.735337-0.989697\,i$ of the repository README (stable to six digits from $N\approx10^5$); this is a comparison of quoted numbers, not a new computation.
\item \emph{(Without RH; the author's notes.)} Neither the convergence of $\mathrm{Cen}_\ell(1;\gamma)$ nor its value is proved. The exact Tauberian criterion of the notes is: the centroid converges to $\mathsf c$ if and only if $\int_s^\infty\mathbb G_1(\kappa)\kappa^{-2}\,d\kappa-\mathsf c/s$ has a local pseudofunction boundary value at $s=0$, where $\mathbb G_1(s)=\sum_ke^{i\gamma_k}e^{-s\gamma_k}$ (in the notes this function is called $G$ and the constant $b$; the letters are changed here because $G=N-N_0$ and $\mathfrak b$ are used in this paper). The criterion is proved in \cref{thm:v14-taub} (under (Rep), automatic for simple zeros). In the variable $u$ (centroid of $e^{iu\gamma_k}$) the notes show that RH is equivalent to $\lim_\ell\mathrm{Cen}_\ell(u;\gamma)=e^{-u/2}\mathfrak V(u)$ in $\mathcal D'(\log2,\log3)$; $u=1$ lies in this window, but a single point is not sufficient for that equivalence (this equivalence is \cite[Thm.~5.5]{CasanaV26}, and is proved here in the form of \cref{prop:v13-def}). The transfer from the Laplace transform to the centroid of $e^{i\gamma_k}$, which \cite{CasanaV26} proves under RH, is proved here without RH (for simple zeros, or under (Rep)) in \cref{thm:v14-transfer}; the criterion thus characterises convergence to a prescribed value $\mathsf c$, and the identification of that value is a separate question (\cref{rem:v13-which}).
\end{enumerate}
\end{proposition}

\begin{proposition}[The centroids of $\tau^{\mathrm C}$ and $\tau^{\Cm}$]\label{prop:v11-tau}
Let $\omega>0$. (Statements on $\tau^{\mathrm C},\tau^{\Cm}$ use \eqref{eq:IJNTinput} through \cref{prop:v9-line} and \cref{prop:v9-CvsCm}.)
\begin{enumerate}[label=(\alph*),leftmargin=*]
\item $\mathrm{Cen}_\ell(\omega;\tau^{\mathrm C})-\mathrm{Cen}_\ell(\omega;\tau^{\Cm})=o(\ell/\log\ell)$, without RH and under the input \eqref{eq:IJNTinput}.
\item If $\varsigma_T(\omega;F_h)\to\varsigma_\infty$ as $T\to\infty$ through real $T$, then for $\star\in\{\mathrm C,\Cm\}$
\begin{equation}
\mathsf D_\ell^\star(\omega)=\frac{i\omega\,\varsigma_\infty\,\ell}{2\log\ell}+o\Bigl(\frac\ell{\log\ell}\Bigr),\qquad
\mathrm{Cen}_\ell(\omega;\tau^\star)=\frac{\bigl(\chi(\omega)-i\omega\varsigma_\infty\bigr)\ell}{2\log\ell}+o\Bigl(\frac\ell{\log\ell}\Bigr).
\end{equation}
At $\omega=1$ ($\chi(1)=0$) the centroids of $\tau^{\mathrm C},\tau^{\Cm}$ are $o(\ell/\log\ell)$ if and only if $\varsigma_\infty=0$; at $\omega=r\log p$ they have the same leading term as $\gamma$ (\cref{prop:v11-gamma}(a)) if and only if $\varsigma_\infty=0$, and are $o(\ell/\log\ell)$ if and only if $\varsigma_\infty=\chi(\omega)/(i\omega)$.
\item The convergence of $\mathrm{Cen}_\ell(1;\tau^{\mathrm C})$ and $\mathrm{Cen}_\ell(1;\tau^{\Cm})$ is not proved, with or without RH. The exact relation is $\mathrm{Cen}_\ell(1;\tau^\star)=\mathrm{Cen}_\ell(1;\gamma)-\mathsf D_\ell^\star(1)$: under RH, $\mathrm{Cen}_\ell(1;\tau^\star)=e^{-1/2}\mathfrak V(1)-\mathsf D_\ell^\star(1)+o(1)$, so the centroid of $\tau^\star$ converges if and only if $\mathsf D_\ell^\star(1)$ does. The repository record (quoted) shows $|\mathsf D_N^{\mathrm C}(1)|=1.39,\ 1.97,\ 10.8,\ 14.7$ at $N=10^4,10^5,10^6,2\cdot10^6$ and no convergence for the other prime-selection rules tried there; it is compatible with (b) for $\varsigma_\infty=0$ and does not decide convergence.
\end{enumerate}
\end{proposition}
\begin{proof}
(a) By \cref{prop:v9-CvsCm}, $\sum_{k\in\mathfrak J_T}(e^{i\omega\tau_k^{\mathrm C}}-e^{i\omega\tau_k^{\Cm}})=O(T/\log\log T)=o(T)$, i.e.\ the block coefficient of the difference is $0$; the telescoping argument of \cref{thm:v9-cen} applies verbatim to $\mathsf f_k=e^{i\omega\tau_k^{\mathrm C}}-e^{i\omega\tau_k^{\Cm}}$. (b) \Cref{prop:v9-line} together with (L) gives $\frac{2\pi}{T}M_T\Xi_T^\star(\omega)=i\omega\varsigma_T+o(1)$, i.e.\ $\Xi_T^\star(\omega)=i\omega\varsigma_\infty/\log T+o(1/\log T)$, which is the hypothesis of \cref{thm:v9-cen} with $\mathfrak b=i\omega\varsigma_\infty$; then use \cref{prop:v11-gamma}(a) and \eqref{eq:v11-diff}. (c) is \eqref{eq:v11-diff} together with \cref{prop:v11-gamma}(b).
\end{proof}

\begin{table}[ht]
\centering\small
\renewcommand{\arraystretch}{1.2}
\begin{tabular}{@{}p{0.11\textwidth}p{0.28\textwidth}p{0.24\textwidth}p{0.26\textwidth}@{}}
\toprule
 & $\mathrm{Cen}_\ell(1;\cdot)$, unconditional & $\mathrm{Cen}_\ell(1;\cdot)$, under RH & $\mathrm{Cen}_\ell(r\log p;\cdot)$, unconditional\\
\midrule
$\gamma$ & $o(\ell/\log\ell)$ ($O(\ell/\log^2\ell)$ with the Ford--Zaharescu error); convergence open; $=\mathfrak m(1)+\mathsf D^{\mathrm{RvM}}_\ell(1)+o(1)$ & $\to e^{-1/2}\mathfrak V(1)=0.7769+0.7679\,i$ (notes) & $\sim\chi\,\ell/(2\log\ell)$: unbounded\\
$\gamma^{\mathrm{RvM}}$ & $\to\mathfrak m(1)=0.0416+1.7575\,i$ (\cref{prop:v11-rvm}) & same & $\to\mathfrak m(\omega)$\\
$\tau^{\mathrm C},\tau^{\Cm}$ & $\sim-i\varsigma_\infty\ell/(2\log\ell)$ if $\varsigma_T\to\varsigma_\infty$; $o(\ell/\log\ell)$ iff $\varsigma_\infty=0$; convergence open & $e^{-1/2}\mathfrak V(1)-\mathsf D^\star_\ell(1)+o(1)$ & $\sim(\chi-i\omega\varsigma_\infty)\ell/(2\log\ell)$\\
\midrule
$\mathsf D^{\mathrm{RvM}}_\ell$ & $o(\ell/\log\ell)$; converges iff $\mathrm{Cen}_\ell(1;\gamma)$ does & $\to K^0=e^{-1/2}\mathfrak V(1)-\mathfrak m(1)$ & $\sim\chi\,\ell/(2\log\ell)$\\
$\mathsf D^{\mathrm C}_\ell,\mathsf D^{\Cm}_\ell$ & $\sim i\varsigma_\infty\ell/(2\log\ell)$ (same for both up to $o(\ell/\log\ell)$) & as unconditional & $\sim i\omega\varsigma_\infty\ell/(2\log\ell)$\\
\bottomrule
\end{tabular}
\caption{Status of the four centroids and of their differences with $\gamma$. $\varsigma_\infty$ is the limit of the sawtooth coefficient $\varsigma_T(\cdot;F_h)$ when it exists; whether it vanishes is open (\S\ref{sec:v9-open-centroid}). Rows for $\tau^{\mathrm C},\tau^{\Cm},\mathsf D^{\mathrm C},\mathsf D^{\Cm}$ use \eqref{eq:IJNTinput}; ``unconditional'' means without RH.}
\label{tab:v11-centroids}
\end{table}

\subsection{What is provable without RH: the real part of Cram\'er's function, and distributional statements}\label{sec:v13-uncond}
Throughout, $J=(\log2,\log3)$, $\mathcal D'(J)$ is the space of distributions, and for a sequence $\mathsf x=(\mathsf x_k)$ of positive numbers $C^{(\mathsf x)}_\ell(u)=\sum_{k\le \ell}(1-\frac{k-1}{\ell})e^{iu\mathsf x_k}=\mathrm{Cen}_\ell(u;\mathsf x)$. For $f\in C_c^\infty(J)$ put $\hat f(t)=\int_Jf(u)e^{iut}\,du$; repeated integration by parts gives $|\hat f(t)|\le C_m t^{-m}$ for every $m$. The hypothesis
\phantomsection
\begin{equation}
\mathsf x_k\ge c\,k/\log(k+2)\quad(k\ge1,\ \text{for some }c>0)\tag{Sp}\label{eq:v13-G}
\end{equation}
holds for $\gamma$ and $\gamma^{\mathrm{RvM}}$ by the Riemann--von Mangoldt formula, and for $\tau^{\mathrm C},\tau^{\Cm}$ because $|\tau_k-\gamma_k^{\mathrm{RvM}}|\le2\pi(0.08382+5.5/\log t)(1+o(1))$ is bounded (\cref{prop:v9-recon}(i)). The letter $u$ stands for the frequency $\omega$; the letter $\ell$ is the length of the centroid (the counting function is always written $N(t)$). Numbered references to \cite{CasanaV26} are to Section~5 of that manuscript, with its numbering.

\begin{proposition}[Unconditional distributional convergence, with the $1/\ell$ term]\label{prop:v13-dist}
Let $\mathsf x$ satisfy \eqref{eq:v13-G} and $f\in C_c^\infty(J)$. Then $\mathscr N_\infty^{(\mathsf x)}(f)=\sum_{k\ge1}\hat f(\mathsf x_k)$ and $\mathscr O_{\mathsf x}(f)=\sum_{k\ge1}(k-1)\hat f(\mathsf x_k)$ converge absolutely, and
\begin{equation}
\int_JfC_\ell^{(\mathsf x)}\,du=\mathscr N_\infty^{(\mathsf x)}(f)-\frac{\mathscr O_{\mathsf x}(f)}\ell+o_f\!\left(\frac1\ell\right).
\end{equation}
In particular the centroids of $\gamma$, $\gamma^{\mathrm{RvM}}$, $\tau^{\mathrm C}$ and $\tau^{\Cm}$ all converge in $\mathcal D'(J)$, with no hypothesis on the zeros.
\end{proposition}
\begin{proof}
$\int_JfC_\ell^{(\mathsf x)}=\sum_{k\le \ell}\hat f(\mathsf x_k)-\frac1\ell\sum_{k\le \ell}(k-1)\hat f(\mathsf x_k)$. By \eqref{eq:v13-G} and $|\hat f(t)|\le C_3t^{-3}$, $|\hat f(\mathsf x_k)|\ll(\log(k+2)/k)^3$, so both series converge absolutely, and the tails are $\sum_{k>\ell}|\hat f(\mathsf x_k)|\ll \ell^{-2}\log^3\ell$ and $\frac1\ell\sum_{k>\ell}k|\hat f(\mathsf x_k)|\ll \ell^{-2}\log^3\ell$, both $o(1/\ell)$.
\end{proof}

\begin{proposition}[Real part of Cram\'er's function, without RH]\label{prop:v13-real}
Let $\mathfrak V(u)=\sum_{\operatorname{Im}\rho>0}e^{\rho u}$ (zeros with multiplicity), a distribution on $(0,\infty)$: for $\varphi\in C_c^\infty(0,\infty)$, $\int\varphi e^{\rho u}du=(-\rho)^{-2}\int\varphi''e^{\rho u}du=O(|\rho|^{-2})$ and $\sum|\rho|^{-2}<\infty$. Then, unconditionally, in $\mathcal D'(0,\infty)$
\begin{equation}
2\operatorname{Re}\mathfrak V(u)=e^{u}-\frac1{e^{2u}-1}-\sum_{n\ge2}\Lambda(n)\,\delta_{\log n}(u),
\end{equation}
so on $J$ (which contains no $\log n$) $\;2\operatorname{Re}\mathfrak V(u)=e^u-(e^{2u}-1)^{-1}$, a smooth function that involves neither a non-trivial zero nor a prime. In particular
\begin{equation}
\operatorname{Re}\bigl(e^{-1/2}\mathfrak V(1)\bigr)=\frac1{2\sqrt e}\Bigl(e-\frac1{e^2-1}\Bigr)=0.77689\ldots,\qquad 1\in J,
\end{equation}
for the value of the (meromorphically continued) Cram\'er function.
\end{proposition}
\begin{proof}
The zeros with $\operatorname{Im}\rho<0$ are the conjugates of those with $\operatorname{Im}\rho>0$, and $e^{\bar\rho u}=\overline{e^{\rho u}}$ for real $u$, so $2\operatorname{Re}\mathfrak V(u)=\sum_\rho e^{\rho u}$ over all non-trivial zeros, summed in $\mathcal D'$ (absolutely, as above). Von Mangoldt's formula $\psi_0(x)=x-\sum_\rho x^\rho/\rho-\log2\pi-\frac12\log(1-x^{-2})$ holds for every $x>1$ with symmetric partial sums that are bounded in $L^1_{\mathrm{loc}}$ \cite[Ch.~17]{Davenport2000}, hence in $\mathcal D'(1,\infty)$; differentiating, and using $\psi_0'=\sum_n\Lambda(n)\delta_n$ and $\frac{d}{dx}\bigl(-\frac12\log(1-x^{-2})\bigr)=-\frac1{x(x^2-1)}$,
\begin{equation}
\sum_\rho x^{\rho-1}=1-\sum_n\Lambda(n)\delta_n(x)-\frac1{x(x^2-1)}\quad\text{in }\mathcal D'(1,\infty).
\end{equation}
Pairing with $\varphi(\log x)$, $x=e^u$, $dx=e^udu$, gives the displayed identity in $u$. Nothing in this argument uses the real parts of the zeros. The value at $u=1$ is the one quoted in \cref{prop:v11-gamma}(b) (real-analyticity of $\mathscr C=e^{-w/2}\mathfrak V$ on $J$ is \cref{thm:v17-cont}).
\end{proof}

\begin{remark}[Which of the two parts is unconditional, and what it does not say]\label{rem:v13-which}
(i) The \emph{real} part of the Cram\'er function is determined by the primes alone (\cref{prop:v13-real}); the \emph{imaginary} part is not local: since $\mathfrak V$ is the boundary value of a holomorphic function in $\operatorname{Im}w>0$, $\operatorname{Im}\mathfrak V$ is tied to $\operatorname{Re}\mathfrak V$ by a Hilbert-type transform, after the usual regularization (\cite[Rem.~5.9]{CasanaV26}; $\operatorname{Re}\mathfrak V(u)\sim e^u/2$ grows, so the integral over the whole line cannot be taken without it); this transform is not local and involves the values of $\operatorname{Re}\mathfrak V$ away from $u$, hence all the primes. This is the sense in which one of the two parts is proved without RH.\\
(ii) This is a statement about $e^{-u/2}\mathfrak V(u)$, the value to which the centroid of $e^{iu\gamma_k}$ converges \emph{under RH} (\cref{prop:v11-gamma}(b)). It does not assert that $\operatorname{Re}\mathrm{Cen}_\ell(1;\gamma)$ converges to this value without RH: by \cite[Cor.~5.8]{CasanaV26} the equality of the real parts on $J$ is \emph{equivalent} to the identity of the ``cell $(2,3)$'', and by \cite[Prop.~5.10]{CasanaV26} that identity fails if only finitely many zeros are off the critical line. The reason is the exact pair relation of \cref{prop:v13-def} below: off the line the two zeros $\rho,1-\bar\rho$ share the imaginary part but contribute $e^{iu\gamma}\cdot2\cosh(ud)$ to $e^{-u/2}\mathfrak V$ and $2e^{iu\gamma}$ to the centroid of $e^{iu\gamma_k}$, and $\cosh(ud)\cos(u\gamma)$ and $\cos(u\gamma)$ cannot be separated.\\
(iii) \emph{Unverified:} a statement that the real or the imaginary part of the centroid of $e^{i\gamma_k}$ is proved without RH is recalled from earlier work of the author; it was not located in the sources available for this note, and the unconditional statements re-derived here on this point are \cref{prop:v13-real} and \cref{thm:v17-cont} (for $\mathscr C$), and \cref{prop:v13-dist}, \cref{prop:v13-def}, \cref{prop:v13-rig} (distributional). If the intended statement is a different one it is not claimed here.
\end{remark}

\begin{proposition}[Exact defect between the $\gamma$-centroid and Cram\'er's function]\label{prop:v13-def}
Put $\mathsf o_k=|\beta_k-\frac12|$ for the zero $\rho_k=\beta_k+i\gamma_k$ (all zeros with $\gamma_k>0$, with multiplicity). For $f\in C_c^\infty(J)$, unconditionally,
\begin{equation}
\lim_{\ell\to\infty}\int_JfC_\ell^{(\gamma)}-\int_Jf(u)\,e^{-u/2}\mathfrak V(u)\,du=-\sum_k\int_Jf(u)\bigl(\cosh(u\mathsf o_k)-1\bigr)e^{iu\gamma_k}\,du ,
\end{equation}
and the right side vanishes for every $f$ if and only if RH holds (\cite[Thm.~5.5]{CasanaV26}; the window theorem).
\end{proposition}
\begin{proof}
The zeros $\rho$ and $1-\bar\rho$ have the same imaginary part, so off the line the zeros with $\gamma>0$ come in pairs with $\beta-\frac12=\pm\mathsf o$, and $\sum_{\gamma>0}e^{(\rho-1/2)u}=\sum_ke^{iu\gamma_k}e^{u(\beta_k-1/2)}$ with $e^{u\mathsf o}+e^{-u\mathsf o}-2=2(\cosh u\mathsf o-1)$ per pair, which is $(\cosh u\mathsf o_k-1)$ for each of the two members. This gives the identity, all series converging absolutely against $f$. If RH holds all $\mathsf o_k=0$. Conversely, let $\mathfrak X(w)=\sum_k(\cosh(w\mathsf o_k)-1)e^{iw\gamma_k}$, holomorphic in $\operatorname{Im}w>0$ and of polynomial growth at the real axis; if its boundary value vanishes on $J$, extending $\mathfrak X$ by $0$ to $\operatorname{Im}w<0$ gives a function with zero jump across $J$, holomorphic across $J$ by the edge-of-the-wedge argument (\cite[Lemma~5.4]{CasanaV26}: repeated primitives and Schwarz reflection), so $\mathfrak X\equiv0$. At $w=iy$ with $0<y\le4\pi$, $\cosh(iy\mathsf o_k)-1=\cos(y\mathsf o_k)-1\le0$, with strict inequality when $\mathsf o_k>0$ because $y\mathsf o_k<2\pi$ ($\mathsf o_k<\frac12$); so $\mathfrak X(iy)=0$ forces $\mathsf o_k=0$ for all $k$.
\end{proof}

\begin{theorem}[Spectral rigidity on the window]\label{prop:v13-rig}
Let $\mathsf x$ satisfy \eqref{eq:v13-G}, $\mathsf x_k$ nondecreasing. Then
\begin{equation}
\begin{gathered}
\int_Jf\,\bigl(C_\ell^{(\gamma)}-C_\ell^{(\mathsf x)}\bigr)\to0\ \ (\ell\to\infty)\ \text{for all }f\in C_c^\infty(J)\\
\iff \{\gamma_k\}=\{\mathsf x_k\}\text{ as multisets}\iff\gamma_k=\mathsf x_k\ \forall k .
\end{gathered}
\end{equation}
\end{theorem}
\begin{proof}
($\Leftarrow$) is trivial. ($\Rightarrow$) Put $\mathscr U_{\mathsf x}(w)=\sum_k(e^{iw\gamma_k}-e^{iw\mathsf x_k})$, holomorphic in $\operatorname{Im}w>0$ since $\gamma_k,\mathsf x_k\gg k/\log k$, with $|\mathscr U_{\mathsf x}(u+iy)|\ll y^{-2}$ for $0<y\le1$. By \cref{prop:v13-dist} and dominated convergence, $\lim_{y\downarrow0}\int f\,\mathscr U_{\mathsf x}(\cdot+iy)=\sum_k(\hat f(\gamma_k)-\hat f(\mathsf x_k))=\lim_\ell\int f(C^{(\gamma)}_\ell-C^{(\mathsf x)}_\ell)=0$. As in the proof of \cref{prop:v13-def} (\cite[Lemma~5.4]{CasanaV26}), $\mathscr U_{\mathsf x}\equiv0$: a function holomorphic in $\operatorname{Im}w>0$ with zero distributional boundary value on $J$ extends by $0$ across $J$ (edge of the wedge, \cref{lem:v17-refl} with $\mathfrak g=0$), hence vanishes identically. Then $\int_0^\infty e^{-yt}\,d(\mathsf n_\gamma-\mathsf n_{\mathsf x})(t)=\mathscr U_{\mathsf x}(iy)=0$ for all $y>0$ for the counting measures $\mathsf n_\gamma=\sum\delta_{\gamma_k}$, $\mathsf n_{\mathsf x}=\sum\delta_{\mathsf x_k}$, which have finite Laplace transforms; uniqueness of the Laplace transform of positive measures gives $\mathsf n_\gamma=\mathsf n_{\mathsf x}$. Both sequences being nondecreasing, $\gamma_k=\mathsf x_k$ for all $k$.
\end{proof}

\begin{corollary}[No exact transfer from $\gamma$ to the corrected points on $J$]\label{cor:v13-notransfer}
(i) The difference of the centroids of $\gamma$ and of $\gamma^{\mathrm{RvM}}$ tends to $0$ in $\mathcal D'(J)$ only if $\gamma_k=\gamma_k^{\mathrm{RvM}}$ for all $k$, which is false: $N_0(\gamma_1)\approx0.45\ne\frac12=N_0(\gamma_1^{\mathrm{RvM}})$. (ii) The same difference for $\tau^{\mathrm C}$ (or $\tau^{\Cm}$) tends to $0$ in $\mathcal D'(J)$ iff $\tau_k^{\mathrm C}=\gamma_k$ for every $k$, which is not the case ($\mathrm{RMS}(\gamma_k-\tau_k^{\mathrm C})=0.0689$ on the first $220$ zeros, quoted in \S\ref{sec:v9-numbers}). Hence, \emph{under RH}, the distributional limit of the centroid of $\tau^{\mathrm C}$ on $J$ is \emph{not} $e^{-u/2}\mathfrak V(u)$ (the limit of the centroid of $\gamma$), and equals it only if the corrected points are exactly the zeros; the distributional limit of the $\tau^{\mathrm C}$-centroid is $\mathscr N_\infty^{(\mathrm{RvM})}+\mathscr V_{\mathrm C}(\cdot+i0)$ with $\mathscr V_{\mathrm C}(w)=\sum_k\bigl(e^{iw\tau_k^{\mathrm C}}-e^{iw\gamma_k^{\mathrm{RvM}}}\bigr)$ and $\mathscr N_\infty^{(\mathrm{RvM})}$ the distributional limit of the RvM centroid (it is $\mathfrak m$ wherever the pointwise limit of \cref{prop:v11-rvm} is dominated), and what remains open is the actual value of this boundary distribution and its pointwise behaviour at $u=1$ (\cref{prop:v11-tau}).
\end{corollary}
\begin{proof}
\cref{prop:v13-rig} with $\mathsf x=\gamma^{\mathrm{RvM}},\tau^{\mathrm C}$. The computation $N_0(\gamma_1)=\frac{\gamma_1}{2\pi}\log\frac{\gamma_1}{2\pi e}+\frac78\approx0.4493$ uses $\gamma_1\approx14.1347$ only.
\end{proof}

\subsection{The Tauberian criterion for the centroid without RH}\label{sec:v14-taub}
This subsection develops part~(c) of \cref{prop:v11-gamma} (from the author's notes, whose steps are re-derived here; where a step could not be confirmed as stated it is flagged). Notation: $u>0$ is the frequency $\omega$ of the section (the sublevel points $u_k(F)$ always carry an index and an argument); $N(t)$ is the zero-counting function, $G=N-N_0$ (\S\ref{sec:v9-objects}); the centroid length is $\ell$ and the time is $T$; $\mathsf L(T)=\log\frac T{2\pi}$; $\mathrm{mult}(\gamma)$ is the multiplicity of the zero with imaginary part $\gamma$, and
\begin{gather}
\mathbb B_u(t)=\sum_{0<\gamma_k\le t}e^{iu\gamma_k},\quad
\mathbb A_u(T)=\int_0^T\mathbb B_u,\quad
\mathbb M_u(T)=\frac{\mathbb A_u(T)}T,\\
\overline{\mathbb M}_u(T)=\frac1T\int_0^T\mathbb M_u,\quad
\mathsf W_{\mathrm{rep}}(T)=\sum_{\gamma_k\le T}\bigl(\mathrm{mult}(\gamma_k)-1\bigr).
\end{gather}
Two hypotheses are used. \emph{(Rep)} $\mathsf W_{\mathrm{rep}}(T)=o(T\log T)$ (automatic if the zeros are simple; no unconditional proof is known to the author that multiple zeros are rare in this sense). \emph{(Pa)} for the sequence $\mathsf x$ and the frequency $u$: $\mathrm{Par}_\ell(u;\mathsf x)=o(\ell)$. By (L) with error $o(T)$ (\eqref{eq:v9-landau}) and $\gamma_\ell\sim2\pi\ell/\log\ell$, (Pa) holds for $\gamma$ for every $u>0$; it holds for $\gamma^{\mathrm{RvM}}$ by (K); and for $\tau^{\mathrm C},\tau^{\Cm}$ (under the input \eqref{eq:IJNTinput}) by \cref{prop:v9-line}: since $|\varsigma_T|\le\pi$, the sums over the dyadic blocks $(T,2T]$ are $O(T)$, hence $\mathrm{Par}_\ell(u;\tau)=O(\ell/\log\ell)$. The finer relation $\mathrm{Par}_\ell(u;\tau)=\mathrm{Par}_\ell(u;\gamma)-\mathfrak b\ell/\log\ell+o(\ell/\log\ell)$, with $\mathfrak b=\mathfrak b^{\mathrm C}(u)=iu\lim_T\varsigma_T(u;F_h)$, holds whenever $\varsigma_T(u;F_h)$ has a limit, by \cref{thm:v9-cen} and \cref{cor:v9-cen}(ii).

\begin{lemma}[Stieltjes identity for the centroid]\label{lem:v14-stielt}
For every $T>0$, with $\ell=N(T)$ and $u>0$,
\begin{multline}
\sum_{n\le\ell}\mathrm{Par}_n(u;\gamma)=\int_0^T\mathbb B_u(t)N_0'(t)\,dt+\mathbb B_u(T)G(T)-\int_0^TG(t)e^{iut}N_0'(t)\,dt\\
-\tfrac12\int_{(0,T]}e^{iut}\,d(G^2)+\tfrac12\mathbb B_u(T)+O\bigl(\mathsf W_{\mathrm{rep}}(T)+1\bigr).
\end{multline}
\end{lemma}
\begin{proof}
$\mathbb B_u$ is right-continuous with $d\mathbb B_u=e^{iut}dN$, and $\sum_{n\le\ell}\mathbb B_u(\gamma_n)=\int_{(0,T]}\mathbb B_u\,dN$. In a block of $m$ equal heights $\gamma$ occupying the indices $j+1,\dots,j+m$, $\mathbb B_u(\gamma_n)-\mathrm{Par}_n(u;\gamma)=(j+m-n)e^{iu\gamma}$, so the block contributes at most $\frac12m(m-1)$ to $\sum_n(\mathbb B_u(\gamma_n)-\mathrm{Par}_n)$; the sum over blocks is at most $\frac12\mathsf W_{\mathrm{rep}}(T)$. Write $dN=N_0'dt+dG$. Integrating by parts,
\begin{equation}
\int_{(0,T]}\mathbb B_u\,dG=\mathbb B_u(T)G(T)-\int_{(0,T]}G(t^-)\,d\mathbb B_u,\qquad d\mathbb B_u=e^{iut}(N_0'\,dt+dG),
\end{equation}
and, since $G^2$ jumps by $2G(t^-)\mathrm{mult}+\mathrm{mult}^2$ at a zero height and has derivative $-2GN_0'$ elsewhere, $\int e^{iut}G(t^-)dG=\frac12\int e^{iut}d(G^2)-\frac12\sum_{\gamma\le T}\mathrm{mult}(\gamma)^2e^{iu\gamma}$ (sum over distinct heights). Finally $\sum_{\rm distinct}\mathrm{mult}^2e^{iu\gamma}=\mathbb B_u(T)+\sum_{\rm distinct}\mathrm{mult}(\mathrm{mult}-1)e^{iu\gamma}=\mathbb B_u(T)+O(\mathsf W_{\mathrm{rep}}(T))$.
\end{proof}

\begin{theorem}[Time centroid and index centroid agree, without RH]\label{thm:v14-transfer}
Let $u>0$ with $\Lambda(e^u)=0$ (that is, $u\ne\log p^m$) and assume (Rep). For $\mathsf c\in\mathbb C$,
\begin{equation}
\lim_{\ell\to\infty}\mathrm{Cen}_\ell(u;\gamma)=\mathsf c\iff\lim_{T\to\infty}\mathbb M_u(T)=\mathsf c .
\end{equation}
\end{theorem}
\begin{proof}
\emph{Step 0.} By (L) (error $o(T)$), $\mathbb B_u(T)=o(T)$. Put $\mathbb H_u(T)=\int_0^TG(t)e^{iut}dt$. Integrating $\int_{(0,T]}e^{iut}dN$ by parts and using $\int_0^Te^{iut}N_0'dt=O_u(\log T)$ gives $iu\mathbb H_u(T)=O_u(\log T)+e^{iuT}G(T)-G(0^+)-\mathbb B_u(T)$, hence $\mathbb H_u(T)=o(T)$ because $G(T)=O(\log T)$ (\eqref{eq:MCOMH}).

\emph{Step 1.} In \cref{lem:v14-stielt}: $\mathbb B_uG=o(T\log T)$; $\int_0^TGe^{iut}N_0'=N_0'(T)\mathbb H_u(T)-\int_0^T\mathbb H_uN_0''=o(T\log T)$; $\int_{(0,T]}e^{iut}d(G^2)=e^{iuT}G(T)^2-G(0^+)^2-iu\int_0^TG^2e^{iut}=O(\log^2T+T\log\log T)$ by Selberg's mean-square bound $\int_0^TG^2\ll T\log\log T$ \cite{Selberg1946} ($G=S+O(1/t)$). By (Rep), $\mathsf W_{\mathrm{rep}}(T)=o(T\log T)$, and $\mathbb B_u(T)=o(T)$, so the remaining terms are $o(T\log T)$. Thus
\begin{equation}
\sum_{n\le N(T)}\mathrm{Par}_n(u;\gamma)=\int_0^T\mathbb B_u(t)N_0'(t)\,dt+o(T\log T).
\end{equation}
Integrating by parts once more, $\int_0^T\mathbb B_uN_0'=\frac T{2\pi}\bigl[\mathsf L(T)\mathbb M_u(T)-\overline{\mathbb M}_u(T)\bigr]$ (use $\mathbb A_u'=\mathbb B_u$). With $N(T)=\frac T{2\pi}(\mathsf L(T)-1)+\frac78+O(\log T)$ and $\mathbb M_u=o(T)$,
\begin{equation}\label{eq:v14-keyid}
\mathrm{Cen}_{N(T)}(u;\gamma)=\frac{\mathsf L(T)\mathbb M_u(T)-\overline{\mathbb M}_u(T)}{\mathsf L(T)-1}+o(1)\qquad(T\to\infty).
\end{equation}

\emph{Step 2 ($\Leftarrow$).} If $\mathbb M_u\to\mathsf c$ then $\overline{\mathbb M}_u\to\mathsf c$, and \eqref{eq:v14-keyid} gives $\mathrm{Cen}_{N(T)}\to\mathsf c$. Every $\ell$ lies in a block $N(T^-)<\ell\le N(T)$ of length $\le C\log T$, and $|\mathrm{Cen}_{\ell+1}-\mathrm{Cen}_\ell|=|\mathrm{Par}_{\ell+1}-\mathrm{Cen}_\ell|/(\ell+1)=o(1/\log\ell)$ because $\mathrm{Par}_n,\mathrm{Cen}_n=o(n/\log n)$ (this is (L): $\mathbb B_u(T)=o(T)$ with $T\asymp n/\log n$); so $\mathrm{Cen}_\ell\to\mathsf c$ for all $\ell$.

\emph{Step 3 ($\Rightarrow$).} Assume $\mathrm{Cen}_\ell\to\mathsf c$. Put $\mathbb I_u(T)=\int_0^T\mathbb M_u$, so $\mathbb I_u'=\mathbb M_u$ and $\overline{\mathbb M}_u=\mathbb I_u/T$. By \eqref{eq:v14-keyid} (valid for every $T$), $\mathsf L\,\mathbb I_u'-\mathbb I_u/T=\mathsf c(\mathsf L-1)+o(\mathsf L)$, i.e. $\mathbb I_u'-\mathbb I_u/(T\mathsf L)=\mathsf c(1-1/\mathsf L)+o(1)$. With $\mathbb K_u(T)=\mathbb I_u(T)-\mathsf cT$: $\mathbb K_u'-\mathbb K_u/(T\mathsf L)=o(1)$, and since $\mathsf L'=1/T$, $(\mathbb K_u/\mathsf L)'=o(1)/\mathsf L$. As $\int^T\!dt/\mathsf L(t)\sim T/\mathsf L(T)$, $\mathbb K_u(T)=o(T)$, and $\mathbb M_u-\mathsf c=\mathbb K_u'=\mathbb K_u/(T\mathsf L)+o(1)\to0$.
\end{proof}

\begin{remark}
(i) The proof of \cref{thm:v14-transfer} was derived here independently of the text received; the identity of \cref{lem:v14-stielt} agrees with the one in that text for simple zeros (there the term $\frac12\mathbb B_u(T)$ appears as $(G(T)+\frac12)\mathbb B_u(T)$), and the treatment of multiplicities in that text (``symmetric counting'') could not be confirmed, which is why (Rep) is stated. (ii) At a resonant frequency $u=\log p^m$, $\mathbb B_u(T)\sim-\frac T{2\pi}\Lambda(e^u)e^{-u/2}$, so $\mathbb M_u\to\infty$ and the centroid is unbounded (\cref{prop:v11-gamma}(a)); the theorem is not stated there.
\end{remark}

\begin{definition}[Local pseudofunction at $s=0$]
Let $\Phi$ be holomorphic in $\operatorname{Re}s>0$ with $\Phi(\eta+i\xi)\to\Phi(i\xi)$ in $\mathcal D'(-\delta,\delta)$ as $\eta\downarrow0$. $\Phi$ has a \emph{local pseudofunction boundary value at $s=0$} if for every $\phi\in C_c^\infty(-\delta,\delta)$ the function $T\mapsto\frac1{2\pi}\langle\phi\,\Phi(i\,\cdot),e^{i\,\cdot\,T}\rangle$ tends to $0$ as $T\to\infty$. Cf.\ \cite{DebruyneVindas2019}. The condition concerns the distributional boundary values for $\xi$ in an open interval around $0$, not the behaviour along the positive real axis.
\end{definition}

\begin{theorem}[Tauberian criterion for part (c)]\label{thm:v14-taub}
Under the hypotheses of \cref{thm:v14-transfer} let $\mathbb G_u(s)=\sum_{k\ge1}e^{iu\gamma_k}e^{-s\gamma_k}$ ($\operatorname{Re}s>0$) and, for $\mathsf c\in\mathbb C$,
\begin{equation}
\mathscr P^{\mathrm t}_{u,\mathsf c}(s)=\int_s^\infty\frac{\mathbb G_u(\kappa)}{\kappa^2}\,d\kappa-\frac{\mathsf c}s .
\end{equation}
Then $\mathscr P^{\mathrm t}_{u,\mathsf c}$ is the Laplace transform of $(\mathbb M_u-\mathsf c)\mathbf 1_{(0,\infty)}$, and
\begin{equation}
\lim_{\ell\to\infty}\mathrm{Cen}_\ell(u;\gamma)=\mathsf c\iff\mathscr P^{\mathrm t}_{u,\mathsf c}\ \text{has a local pseudofunction boundary value at }s=0 .
\end{equation}
In particular, for $u=1$ the criterion of \cref{prop:v11-gamma}(c) is justified without RH, for simple zeros (or under (Rep)): $\mathsf c$ is the value of the limit; the existence of the limit is the existence of some $\mathsf c$ for which the condition holds, and the identification $\mathsf c=e^{-1/2}\mathfrak V(1)$ is a separate matter (\cref{prop:v13-def}, \cref{rem:v13-which}).
\end{theorem}
\begin{proof}
$\mathbb G_u(\kappa)=\kappa\int_0^\infty\mathbb B_ue^{-\kappa t}dt$, so $\mathcal L\mathbb A_u=\mathbb G_u/\kappa^2$ and, dividing by $T$, $\mathcal L\mathbb M_u(s)=\int_s^\infty\mathcal L\mathbb A_u(\kappa)d\kappa=\int_s^\infty\mathbb G_u\kappa^{-2}d\kappa$; $\mathcal L[\mathsf c\mathbf 1_{(0,\infty)}]=\mathsf c/s$. By \cref{thm:v14-transfer} it suffices to show $\mathbb M_u\to\mathsf c\iff$ pseudofunction. ($\Rightarrow$) $f=\mathbb M_u-\mathsf c$ is bounded and $\to0$; for $\phi\in C_c^\infty$, $\frac1{2\pi}\langle\phi\hat f,e^{i\cdot T}\rangle=(f*\check\phi)(T)\to0$. ($\Leftarrow$) $f=\mathbb M_u-\mathsf c$ has $f=o(T)$ and $f'(T)=(\mathbb B_u(T)-\mathbb M_u(T))/T=o(1)$ a.e. (Step~0 and $\mathbb M_u=o(T)$). Choose $K$ Schwartz with $\int K=1$ and $\hat K\in C_c^\infty(-\delta,\delta)$. The hypothesis gives $(f*K)(T)\to0$. Also $|f(T)-(f*K)(T)|\le\int|K(y)||f(T)-f(T-y)|dy\le o(1)\int|K(y)||y|dy+\int_{|y|>T/2}|K(y)|(|f(T)|+|f(T-y)|)dy\to0$, using $f'=o(1)$ on $|y|\le T/2$ and $|f(\cdot)|\ll T$ with rapid decay of $K$ for $|y|>T/2$. Hence $f\to0$.
\end{proof}

\begin{proposition}[Exact index-based identity and criterion for the four centroids]\label{prop:v14-index}
Let $\mathsf x$ be any sequence of real numbers and $\mathbb P_{\mathsf x}(s;u)=\sum_{k\ge1}e^{-sk+iu\mathsf x_k}$, $\mathbb Q_{\mathsf x}(s;u)=\sum_{\ell\ge1}\mathrm{Cen}_\ell(u;\mathsf x)e^{-s\ell}$ ($\operatorname{Re}s>0$). Then, without any hypothesis,
\begin{equation}\label{eq:v14-Q}
\mathbb Q_{\mathsf x}(s;u)=\int_s^\infty\frac{\mathbb P_{\mathsf x}(\kappa;u)}{(1-e^{-\kappa})^2}\,d\kappa,\qquad\frac1{(1-e^{-s})^2}=\frac1{s^2}+\frac1s+\frac5{12}+O(s).
\end{equation}
If moreover $\mathrm{Par}_\ell(u;\mathsf x)=o(\ell)$ (Pa), then for $\mathsf c\in\mathbb C$
\begin{equation}
\begin{gathered}
\lim_{\ell\to\infty}\mathrm{Cen}_\ell(u;\mathsf x)=\mathsf c\iff\mathscr P^{\mathrm o}_{\mathsf x,\mathsf c}(s;u)\ \text{has a local pseudofunction b.v.\ at }s=0,\\
\mathscr P^{\mathrm o}_{\mathsf x,\mathsf c}(s;u)=\int_s^\infty\frac{\mathbb P_{\mathsf x}(\kappa;u)}{(1-e^{-\kappa})^2}d\kappa-\frac{\mathsf c}{e^s-1}.
\end{gathered}
\end{equation}
For two such sequences $\mathsf x,\mathsf x'$ and $\mathsf q\in\mathbb C$: $\mathrm{Cen}_\ell(u;\mathsf x)-\mathrm{Cen}_\ell(u;\mathsf x')\to\mathsf q$ iff $\mathscr X(s;u)-\mathsf q/(e^s-1)$ has such a boundary value, where $\mathscr X(s;u)=\sum_\ell[\mathrm{Cen}_\ell(u;\mathsf x)-\mathrm{Cen}_\ell(u;\mathsf x')]e^{-s\ell}=\int_s^\infty\frac{\mathbb P_{\mathsf x}-\mathbb P_{\mathsf x'}}{(1-e^{-\kappa})^2}d\kappa$; in particular for $\mathsf x=\gamma$ and $\mathsf x'=\mathsf x^\star$ one may write $\mathbb P_\gamma-\mathbb P_{\mathsf x^\star}=\sum_ke^{-sk+iu\mathsf x^\star_k}\bigl(e^{iu\sigma_k^\star}-1\bigr)$ with $\sigma_k^\star=\gamma_k-\mathsf x^\star_k$.
\end{proposition}
\begin{proof}
$\sum_\ell\ell\,\mathrm{Cen}_\ell z^\ell=\sum_ke^{iu\mathsf x_k}z^k\sum_{m\ge0}(m+1)z^m=\mathbb P_{\mathsf x}(z)/(1-z)^2$ with $z=e^{-s}$, and $\int_s^\infty\ell e^{-\ell\kappa}d\kappa=e^{-\ell s}$; the Laurent expansion is $(s^{-1}+\frac12+\frac s{12}+\dots)^2$. Under (Pa), $\mathrm{Cen}_{\ell+1}-\mathrm{Cen}_\ell=(\mathrm{Par}_{\ell+1}-\mathrm{Cen}_\ell)/(\ell+1)=o(1)$ and $\mathrm{Cen}_\ell=o(\ell)$; the step function $\ell\mapsto\mathrm{Cen}_{\lceil t\rceil}$ has Laplace transform $e^{s}\frac{1-e^{-s}}s\mathbb Q_{\mathsf x}(s;u)$, the factor being holomorphic and non-zero at $s=0$, so pseudofunction boundary values are preserved, and $\sum_\ell\mathsf c\,e^{-s\ell}=\mathsf c/(e^s-1)$. The argument of the proof of \cref{thm:v14-taub} applies verbatim.
\end{proof}

\begin{corollary}[The window criterion in index variables]\label{cor:v14-window}
Let $J=(\log2,\log3)$. Then RH holds if and only if, for every $f\in C_c^\infty(J)$,
\begin{equation}
\lim_{s\downarrow0}\ s\int_Jf(u)\,\mathbb Q_\gamma(s;u)\,du=\int_Jf(u)\,e^{-u/2}\mathfrak V(u)\,du .
\end{equation}
\end{corollary}
\begin{proof}
By \cref{prop:v13-dist}, $\int_Jf\,\mathrm{Cen}_\ell(\cdot;\gamma)\to\mathscr N^{(\gamma)}_\infty(f)$, so its Abel mean $s\sum_\ell e^{-s\ell}\int_Jf\,\mathrm{Cen}_\ell$ has the same limit as $s\downarrow0$, and $\sum_\ell e^{-s\ell}\int_Jf\,\mathrm{Cen}_\ell=\int_Jf\,\mathbb Q_\gamma$. The equivalence with RH is \cref{prop:v13-def}.
\end{proof}

\begin{remark}[Heuristic, not used above]
If the block coefficients $\varsigma_T(1;F_h)$ decay only like $T^{-\alpha}$, $0<\alpha<1$ (the quantity called $\rho_1(T)$ in the notes received, which equals $\varsigma_T(1;F_h)$ for $\langle\cdot\rangle=\mathrm{saw}$), coherent accumulation over the earlier blocks could produce contributions of size $T^{1-\alpha}$ to the centroid difference. This would be consistent with a good RMS of $\gamma_k-\tau^{\mathrm C}_k$ coexisting with growing differences of centroids, but nothing in this paper proves it: what \cref{prop:v14-index} localises is the behaviour of $\mathscr X$ at $s=0$, and it asks for a decomposition $\mathscr X(s;1)=\mathsf q/(e^s-1)+(\text{local pseudofunction})$ for $\mathsf x=\gamma$, $\mathsf x'=\tau^{\mathrm C}$, which is open.
\end{remark}

\begin{table}[ht]
\centering\small
\begin{tabular}{@{}p{0.58\linewidth}p{0.36\linewidth}@{}}
\toprule
Statement & Status\\
\midrule
$\mathrm{Cen}_\ell(u;\gamma)\to\mathsf c$ iff $\mathbb M_u(T)\to\mathsf c$ ($u\ne\log p^m$) & \cref{thm:v14-transfer}; no RH; (Rep)\\
Pseudofunction criterion for $u=1$ (part (c)) & \cref{thm:v14-taub}; no RH; (Rep)\\
Identity \eqref{eq:v14-Q} & exact, any real sequence\\
Index-based criterion and criterion for differences of centroids & \cref{prop:v14-index}; needs (Pa)\\
(Pa) for $\gamma$; for $\gamma^{\mathrm{RvM}}$; for $\tau^{\mathrm C},\tau^{\Cm}$ & (L); (K); \cref{prop:v9-line} with $|\varsigma_T|\le\pi$ (under \eqref{eq:IJNTinput})\\
Convergence of $\mathrm{Cen}_\ell(1;\gamma)$ without RH; identification of the limit with $e^{-1/2}\mathfrak V(1)$ & open\\
Window equivalence with RH & \cref{prop:v13-def}, \cref{cor:v14-window}\\
\bottomrule
\end{tabular}
\caption{Status of the Tauberian formulation.}
\label{tab:v14-status}
\end{table}

\subsection{One test function and a family of modulations: reduction of the window criterion}\label{sec:v15-family}
This subsection records a reduction of the window criterion of \cref{prop:v13-def} to one fixed test function (from the author's notes; the steps are re-derived here). Notation: $J=(\log2,\log3)$ as in \S\ref{sec:v13-uncond}; $J_0=(0.94,1.06)\subset J$ (any open subinterval would do); $f\in C_c^\infty(J)$ is real, $f\ge0$, $f>0$ on $J_0$; $\mathsf a$ is a real or complex modulation parameter and $\hat f(t)=\int_Jf(u)e^{iut}du$ (entire in $t$). Put
\begin{gather}
\mathscr I_f(\mathsf a)=\int_Jf(u)\,e^{i\mathsf au}\,e^{-u/2}\mathfrak V(u)\,du,\qquad
\mathscr Y_f(\mathsf a)=\sum_{k\ge1}\hat f(\gamma_k+\mathsf a)-\mathscr I_f(\mathsf a),\\
\mathscr W_f=\sum_{m\in\mathbb Z}2^{-|m|}\,|\mathscr Y_f(m)|^2 .
\end{gather}

\begin{lemma}[Local form of the window criterion]\label{lem:v15-local}
If $\sum_ke^{iu\gamma_k}=e^{-u/2}\mathfrak V(u)$ in $\mathcal D'(J_0)$ for some open interval $J_0\subset J$, then RH holds.
\end{lemma}
\begin{proof}
In the proof of \cref{prop:v13-def} the function $\mathfrak X(w)=\sum_k(\cosh(w\mathsf o_k)-1)e^{iw\gamma_k}$ is holomorphic in $\operatorname{Im}w>0$ with polynomial growth at the real axis, and its boundary value on $J_0$ is $e^{-u/2}\mathfrak V(u)-\sum_ke^{iu\gamma_k}$ (the identity of that proof). If this vanishes on $J_0$, extension by $0$ to $\operatorname{Im}w<0$ is holomorphic across $J_0$ (edge of the wedge, \cite[Lemma~5.4]{CasanaV26}, which is local), so $\mathfrak X\equiv0$, and the argument at $w=iy$ gives $\mathsf o_k=0$ for all $k$.
\end{proof}

\begin{theorem}[Reduction to one test function]\label{thm:v15-reduction}
The following are equivalent: (i) RH; (ii) $\mathscr Y_f(\mathsf a)=0$ for every real $\mathsf a$; (iii) $\mathscr Y_f(\mathsf a_j)=0$ for a sequence of pairwise distinct complex $\mathsf a_j$ with a finite accumulation point (for example $\mathsf a_j=1/j$); (iv) $\mathscr W_f=0$.
\end{theorem}
\begin{proof}
Let $\mathscr T=\sum_ke^{iu\gamma_k}-e^{-u/2}\mathfrak V(u)$ in $\mathcal D'(J)$ (the first sum exists by \cref{prop:v13-dist}). Then $f\mathscr T$ has compact support in $J$ and $\mathscr Y_f(\mathsf a)=\langle f\mathscr T,e^{i\mathsf au}\rangle$ is its Fourier--Laplace transform: an entire function of exponential type with polynomial growth on $\mathbb R$ (the series for $\mathscr Y_f$ converges uniformly on compacts of $\mathbb C$). (i)$\Rightarrow$(ii) is \cref{prop:v13-def}. (ii)$\Rightarrow$(iii) is trivial for real $\mathsf a_j$ and in general follows from $\mathscr Y_f\equiv0$ (identity theorem applied to (ii)), (iii)$\Rightarrow$(ii) is the identity theorem for entire functions, and (iv)$\iff$(ii)-at-integers because every summand of $\mathscr W_f$ is $\ge0$. (ii) $\Rightarrow f\mathscr T=0$ by injectivity of the Fourier transform on compactly supported distributions; (iv) gives the same, because $\mathscr Y_f(m)$ are the Fourier coefficients of the $2\pi$-periodisation of $f\mathscr T$ and $|J|=\log\frac32<2\pi$, so no overlap occurs and a distribution on the circle with vanishing coefficients is zero. Hence $\mathscr T=0$ on the open set $\{f>0\}\supseteq J_0$, and \cref{lem:v15-local} gives (i).
\end{proof}
\begin{remark}
\cref{thm:v15-reduction} is an exact reformulation, not a reduction of difficulty: (iv) is the vanishing of a norm of $\mathscr Y_f$, with weights $2^{-|m|}$ chosen arbitrarily, and is equivalent to $\mathscr Y_f\equiv0$. The tail estimate below shows that no small value of $\mathscr W_f$ obtained from finitely many zeros decides RH: for $|m|\le T/2$ the sum $\sum_k\hat f(\gamma_k+m)$ in $\mathscr Y_f(m)$ is determined, up to $O(T^{1-M}\log T)$, by the zeros with $\gamma_k\le T$ (the term $\mathscr I_f(m)$ does not depend on $T$, and the terms with $|m|>T/2$ are absorbed by the weights $2^{-|m|}$).
\end{remark}

\begin{lemma}[Tails of the modulated sums]\label{lem:v15-tail}
Let $\mathsf x$ satisfy \eqref{eq:v13-G}. For every $M>2$, every $T\ge2$ and every real $\mathsf a$ with $|\mathsf a|\le T/2$ (or complex $\mathsf a$ with $|\operatorname{Re}\mathsf a|\le T/2$ and $|\operatorname{Im}\mathsf a|\le A$, the implied constant then depending on $A$), $\sum_{\mathsf x_k>T}|\hat f(\mathsf x_k+\mathsf a)|\ll_{f,M}T^{1-M}\log T$, uniformly in $\mathsf a$. Hence for $\mathsf x=\gamma^{\mathrm{RvM}},\tau^{\mathrm C},\tau^{\Cm}$ (for $\mathsf x=\gamma$ one has $\mathscr Z_f^\star=0$),
\begin{equation}
\mathscr Z_f^{\star}(\mathsf a):=\sum_k\bigl[\hat f(\gamma_k+\mathsf a)-\hat f(\mathsf x_k+\mathsf a)\bigr]=\sum_{\gamma_k\le T}\bigl[\hat f(\gamma_k+\mathsf a)-\hat f(\mathsf x_k+\mathsf a)\bigr]+O_{f,M}(T^{1-M}\log T),
\end{equation}
and $\mathscr Y_f=\mathscr Y_f^{\star}+\mathscr Z_f^{\star}$ with $\mathscr Y_f^{\star}(\mathsf a)=\sum_k\hat f(\mathsf x_k+\mathsf a)-\mathscr I_f(\mathsf a)$, an identity by definition.
\end{lemma}
\begin{proof}
For $\mathsf x_k>T$ and real $|\mathsf a|\le T/2$, $\mathsf x_k+\mathsf a\ge\mathsf x_k/2$ (for complex $\mathsf a$ as in the statement, $\operatorname{Re}(\mathsf x_k+\mathsf a)\ge\mathsf x_k/2$ and $|\hat f(t)|\ll_{M,A}|t|^{-M}$ uniformly for $|\operatorname{Im}t|\le A$, $|t|\ge1$), so $|\hat f(\mathsf x_k+\mathsf a)|\ll\mathsf x_k^{-M}$; with $\mathsf x_k\gg k/\log k$ the sum over $\mathsf x_k>T$ is $\ll\int_T^\infty t^{-M}\log t\,dt$. Moreover $\gamma_k-\mathsf x_k=O(1)$ for $\mathsf x=\gamma^{\mathrm{RvM}}$ (since $N-N_0=O(\log t)$ and $N_0'(t)\asymp\log t$) and for $\mathsf x=\tau^{\mathrm C},\tau^{\Cm}$ (by the bound on $\tau_k-\gamma_k^{\mathrm{RvM}}$ used for \eqref{eq:v13-G}); hence the terms of $\sum_{\gamma_k\le T}$ and of $\sum_{\mathsf x_k\le T}$ differ only by terms with $\min(\gamma_k,\mathsf x_k)>T-C$, which are $O(T^{1-M}\log T)$ by the first estimate (applied at $T-C$). So the sum over $\gamma_k\le T$ may replace the full sum with the stated error.
\end{proof}
\begin{remark}
Only the decay of $\hat f$, \eqref{eq:v13-G} and $\gamma_k-\mathsf x_k=O(1)$ are used. The quadratic estimate $\sum_{T<\gamma_k\le2T}|\gamma_k-\tau_k|^2\ll T/\log T$ attributed to the author's earlier results in the notes received is \emph{not} needed here; it is established for $\tau^{\Cm}$ in \eqref{eq:cm-mean} (under the inputs of \cref{thm:taucm}), but not for $\tau^{\mathrm C}$, for which only \cref{prop:v9-CvsCm} is available; the splitting of $\mathscr Y_f$ in \cref{lem:v15-tail} is a definition, and $\mathscr Y_f^{\mathrm C}(\mathsf a)$ is computed from the points $\tau_k^{\mathrm C}$ and the Cram\'er representation of $\mathscr I_f$, not from a quadratic estimate. What is not controlled is $\mathscr Z_f^{\mathrm C}$, which depends on the true zeros.
\end{remark}

\begin{remark}[Rate in \cref{prop:v13-dist}]
Taking $m$ large in the proof of \cref{prop:v13-dist} (and replacing $f$ by $fe^{i\mathsf au}$) gives $\int_Jf\,e^{i\mathsf au}\,\mathrm{Cen}_\ell(u;\mathsf x)\,du=\mathscr N_\infty^{(\mathsf x)}-\mathscr O_{\mathsf x}/\ell+O_{f,M,\mathsf a}(\ell^{-M})$ for every $M$, with $\mathscr N_\infty^{(\mathsf x)}(\mathsf a)=\sum_k\hat f(\mathsf x_k+\mathsf a)$ and $\mathscr O_{\mathsf x}(\mathsf a)=\sum_k(k-1)\hat f(\mathsf x_k+\mathsf a)$ entire in $\mathsf a$, uniformly on compacts.
\end{remark}

\begin{proposition}[Smooth averages of the sawtooth coefficient vanish automatically]\label{prop:v15-auto}
For any real measurable $F$, $f\in C_c^\infty(J)$ and $M\ge1$,
\begin{equation}
\Bigl|\int_Jf(u)\,\varsigma_T(u;F)\,du\Bigr|=\frac{2\pi}T\Bigl|\int_T^{2T}\mathrm{saw}(F(t))\,\hat f(t)\,dt\Bigr|\le\pi\,C_{f,M}\,T^{-M} .
\end{equation}
\end{proposition}
\begin{proof}
Fubini and $|\mathrm{saw}|\le\frac12$, $|\hat f(t)|\le C_{f,M}(1+t)^{-M}$.
\end{proof}
\begin{remark}
Consequently $\varsigma_T(\cdot;F)\to0$ in $\mathcal D'(J)$ faster than any power of $1/T$ for \emph{every} $F$, so this convergence neither requires nor implies the pointwise statement $\varsigma_T(1;F)\to0$ (open, \S\ref{sec:v9-open-centroid}(a)), nor RH. The coefficient called $\rho_u(T)$ in the notes received coincides with $\varsigma_T(u;F)$ of \eqref{eq:v9-varsigma} when $\langle\cdot\rangle=\mathrm{saw}$. This also shows that the heuristic remark of \S\ref{sec:v14-taub} concerns pointwise accumulation only.
\end{remark}

\subsection{Phase transport: exact moment expansion of the centroid difference, and the band-occupation identity}\label{sec:v16-transport}
This subsection proves the algebraic and measure-theoretic statements behind a computation of the phase correction $\mathsf D_\ell^{\mathrm C}$ (numbers in \S\ref{sec:v9-numbers}). Notation is that of \eqref{eq:v9-Cen} and \eqref{eq:v11-diff}: $\sigma_k=\sigma_k^{\mathrm C}=\gamma_k-\tau_k^{\mathrm C}$ and $\mathsf D_\ell=\mathsf D_\ell^{\mathrm C}$. New symbols (none occurs earlier in this paper): $\mathsf M_\ell^{(j)}(\omega)$, $\mathsf s_n(\omega)$, $\mathsf m$ (order of a Taylor expansion), $\mathsf R_{\ell,\mathsf m}(\omega)$, $\operatorname{Rem}_{\mathsf m}$, and, for \cref{prop:v16-counter}, $\mathsf E$, $\mathsf w$, $\mathsf v$, $\mathfrak o$, $\mathsf K_{\mathfrak o}$, $t_\natural$, $t_\lambda$.

\begin{definition}[Moments of the displacement]\label{def:v16-moments}
For $\omega\in\mathbb R$, $\ell\ge1$ and an integer $j\ge1$ put
\begin{equation}
\mathsf M_\ell^{(j)}(\omega)=\frac1\ell\sum_{k=1}^{\ell}(\ell-k+1)\,\sigma_k^{\,j}\,e^{i\omega\tau_k^{\mathrm C}},\qquad
\mathsf s_n(\omega)=\sum_{k=1}^{n}\sigma_k\,e^{i\omega\tau_k^{\mathrm C}} .
\end{equation}
\end{definition}

\begin{proposition}[Exact expansion of $\mathsf D_\ell^{\mathrm C}$ in moments of the displacement]\label{prop:v16-moment}
For every $\ell\ge1$, every real $\omega$ and every integer $\mathsf m\ge0$,
\begin{equation}\label{eq:v16-moment}
\mathsf D_\ell^{\mathrm C}(\omega)=\sum_{j=1}^{\mathsf m}\frac{(i\omega)^j}{j!}\,\mathsf M_\ell^{(j)}(\omega)+\mathsf R_{\ell,\mathsf m}(\omega),\qquad
|\mathsf R_{\ell,\mathsf m}(\omega)|\le\frac{|\omega|^{\mathsf m+1}}{(\mathsf m+1)!}\cdot\frac1\ell\sum_{k=1}^{\ell}(\ell-k+1)\,|\sigma_k|^{\mathsf m+1}.
\end{equation}
Moreover
\begin{equation}\label{eq:v16-abel}
\mathsf M_\ell^{(1)}(\omega)=\frac1\ell\sum_{n=1}^{\ell}\mathsf s_n(\omega),
\end{equation}
so the first-order term of \eqref{eq:v16-moment} is $i\omega$ times the Ces\`aro mean of the partial sums $\mathsf s_n$.
\end{proposition}
\begin{proof}
By \eqref{eq:v11-diff}, $\mathsf D_\ell^{\mathrm C}(\omega)=\frac1\ell\sum_k(\ell-k+1)e^{i\omega\tau_k^{\mathrm C}}(e^{i\omega\sigma_k}-1)$. For real $y$ put $\operatorname{Rem}_{\mathsf m}(y)=e^{iy}-\sum_{j=0}^{\mathsf m}(iy)^j/j!$. Taylor's formula with integral remainder gives $\operatorname{Rem}_{\mathsf m}(y)=i^{\mathsf m+1}\int_0^y\frac{(y-v)^{\mathsf m}}{\mathsf m!}e^{iv}\,dv$, hence $|\operatorname{Rem}_{\mathsf m}(y)|\le|y|^{\mathsf m+1}/(\mathsf m+1)!$ because $|e^{iv}|=1$. Apply this with $y=\omega\sigma_k$, multiply by $\frac{\ell-k+1}\ell e^{i\omega\tau_k^{\mathrm C}}$ (modulus $\frac{\ell-k+1}\ell\le1$) and sum over $k$; the $j=0$ term cancels the $-1$, and $\mathsf R_{\ell,\mathsf m}$ is the sum of the remainders. For \eqref{eq:v16-abel}: $\sum_{n=1}^\ell\mathsf s_n=\sum_{n=1}^\ell\sum_{k\le n}a_k=\sum_{k=1}^\ell(\ell-k+1)a_k$ with $a_k=\sigma_ke^{i\omega\tau_k^{\mathrm C}}$.
\end{proof}

\begin{remark}[What \cref{prop:v16-moment} does and does not give]\label{rem:v16-diag}
(i) The moments $\mathsf M_\ell^{(j)}$ contain $\sigma_k=\gamma_k-\tau_k^{\mathrm C}$, hence the zeros: \eqref{eq:v16-moment} is an exact rewriting of $\mathsf D_\ell^{\mathrm C}$, \emph{not} a reconstruction of it from RvM and primes. A reconstruction would require expressing $\mathsf M_\ell^{(j)}$ through $\tau^{\mathrm C}$ alone, which is not done here. (ii) For $\omega=1$ and $\mathsf m=3$ the remainder is at most $\frac1{24}\cdot\frac1\ell\sum_k(\ell-k+1)\sigma_k^4$. (iii) The block coefficient $\mathfrak b^{\mathrm C}(\omega)$ of \cref{thm:v9-cen} is the coefficient of the \emph{whole} sum $\sum_k(e^{i\omega\gamma_k}-e^{i\omega\tau_k^{\mathrm C}})$, which is the sum of all the orders in \eqref{eq:v16-moment}; no single $\mathsf M_\ell^{(j)}$ is identified with it here, and which order carries the size $\ell/\log\ell$ (if any) is not decided. No limit statement for $\mathsf D_\ell^{\mathrm C}(1)$ follows.
\end{remark}

\begin{proposition}[Band-occupation identity for the sublevel points]\label{prop:v16-band}
Let $F:(0,\infty)\to\mathbb R$ be measurable with $F(t)\to+\infty$ as $t\to\infty$. For each level $\lambda$ the set $\{t>0:F(t)<\lambda\}$ is bounded, and put $u_k(F)=\meas\{t>0:F(t)<k-\tfrac12\}$. Then for every $k\ge1$
\begin{equation}\label{eq:v16-band}
u_{k+1}(F)-u_k(F)=\meas\Bigl\{t>0:\ k-\tfrac12\le F(t)<k+\tfrac12\Bigr\}.
\end{equation}
Consequently, for a family $(F_X)_X$ of such functions, $\inf_X\bigl(u_{k+1}(F_X)-u_k(F_X)\bigr)>0$ if and only if $\inf_X\meas\{k-\frac12\le F_X<k+\frac12\}>0$. In particular, for $F_X=F_X^W$ and $\tau_k^{\Cms}(X;W)=u_k(F_X^W)$, the gap $\tau_{k+1}^{\Cms}-\tau_k^{\Cms}$ is the measure of the band of width $1$ at level $k$.
\end{proposition}
\begin{proof}
$F\to\infty$ gives a $t_0$ with $F\ge\lambda$ on $(t_0,\infty)$, so $\{F<\lambda\}\subset(0,t_0]$ has finite measure. $\{F<k-\frac12\}\subset\{F<k+\frac12\}$ and the difference is $\{k-\frac12\le F<k+\frac12\}$; since both have finite measure, the measures subtract. The statement about families is \eqref{eq:v16-band} read for each $X$.
\end{proof}

\begin{proposition}[Bounded perturbation of the smooth count and regular spacing do not force separation]\label{prop:v16-counter}
Let $k\ge3$ and let $t_\natural>0$ be the point with $N_0(t_\natural)=k$ (it exists and is unique, and $t_\natural>2\pi$, see the proof). For $\mathsf E>0$ put $F_{\mathsf E}(t)=N_0(t)+\tanh\bigl(\mathsf E(t-t_\natural)\bigr)$. Then $|F_{\mathsf E}-N_0|\le1$, $F_{\mathsf E}$ is strictly increasing on $(2\pi,\infty)$, $u_k(F_{\mathsf E})$ and $u_{k+1}(F_{\mathsf E})$ are the roots in $(2\pi,\infty)$ of $F_{\mathsf E}=k-\tfrac12$ and $F_{\mathsf E}=k+\tfrac12$ respectively, and
\begin{equation}
u_{k+1}(F_{\mathsf E})-u_k(F_{\mathsf E})=\frac{\log3}{\mathsf E}+O_k\!\bigl(\mathsf E^{-2}\bigr)\qquad(\mathsf E\to\infty).
\end{equation}
In particular the gap tends to $0$ although $|F_{\mathsf E}-N_0|\le1$ and the points $\gamma_j^{\mathrm{RvM}}$, defined by $N_0(\gamma_j^{\mathrm{RvM}})=j-\frac12$, are regularly spaced.
\end{proposition}
\begin{proof}
\emph{Monotonicity and location of the roots.} $|\tanh|\le1$ gives $|F_{\mathsf E}-N_0|\le1$. $N_0'(t)=\frac1{2\pi}\log\frac t{2\pi}$ is $<0$ on $(0,2\pi)$ and $>0$ on $(2\pi,\infty)$; $N_0(0^+)=\frac78$, $N_0(2\pi)=-\frac18$, $N_0(t)\to\infty$. Hence $N_0\le\frac78<k$ on $(0,2\pi]$ and $N_0$ increases from $-\frac18$ to $\infty$ on $(2\pi,\infty)$, so $N_0(t_\natural)=k$ has exactly one solution, and $t_\natural>2\pi$. On $(2\pi,\infty)$, $F_{\mathsf E}'=N_0'+\mathsf E\operatorname{sech}^2(\mathsf E(t-t_\natural))>0$. On $(0,2\pi]$, $F_{\mathsf E}<\frac78+1<\frac52\le k-\frac12$. So for $\lambda\in\{k-\frac12,k+\frac12\}$, $\{t>0:F_{\mathsf E}<\lambda\}=(0,t_\lambda)$, where $t_\lambda$ is the unique root of $F_{\mathsf E}=\lambda$ in $(2\pi,\infty)$ (it exists because $F_{\mathsf E}\to\infty$ and $F_{\mathsf E}(2\pi)<\lambda$), and $u_k=t_{k-1/2}$, $u_{k+1}=t_{k+1/2}$.

\emph{Asymptotics.} Put $\mathsf w=1/\mathsf E$, $t=t_\natural+\mathsf w\mathsf v$ and, for $\mathfrak o\in\{-\frac12,\frac12\}$,
\begin{equation}
\mathsf K_{\mathfrak o}(\mathsf v,\mathsf w)=\tanh\mathsf v+\bigl(N_0(t_\natural+\mathsf w\mathsf v)-k\bigr)-\mathfrak o .
\end{equation}
Then $F_{\mathsf E}(t)=k+\mathfrak o$ iff $\mathsf K_{\mathfrak o}(\mathsf v,\mathsf w)=0$. $\mathsf K_{\mathfrak o}$ is real-analytic in $(\mathsf v,\mathsf w)$ near $(\mathsf v_{\mathfrak o},0)$, where $\mathsf v_{\mathfrak o}=\operatorname{artanh}\mathfrak o=\pm\frac12\log3$ and $\mathsf K_{\mathfrak o}(\mathsf v,0)=\tanh\mathsf v-\mathfrak o$. Since $\partial_{\mathsf v}\mathsf K_{\mathfrak o}(\mathsf v_{\mathfrak o},0)=1-\mathfrak o^2=\frac34\ne0$, the implicit function theorem gives a real-analytic $\mathsf v_{\mathfrak o}(\mathsf w)$ with $\mathsf v_{\mathfrak o}(0)=\mathsf v_{\mathfrak o}$ and $\mathsf K_{\mathfrak o}(\mathsf v_{\mathfrak o}(\mathsf w),\mathsf w)=0$. For small $\mathsf w$ the point $t_\natural+\mathsf w\mathsf v_{\mathfrak o}(\mathsf w)$ lies in $(2\pi,\infty)$ and is a root of $F_{\mathsf E}=k+\mathfrak o$, hence is $t_{k+\mathfrak o}$ by uniqueness. Therefore
\begin{equation}
u_{k+1}-u_k=\mathsf w\bigl(\mathsf v_{1/2}(\mathsf w)-\mathsf v_{-1/2}(\mathsf w)\bigr)=\mathsf w\bigl(\log3+O_k(\mathsf w)\bigr)=\frac{\log 3}{\mathsf E}+O_k\!\bigl(\mathsf E^{-2}\bigr),
\end{equation}
the $O_k$ being the size of $\mathsf v_{\mathfrak o}'$ near $\mathsf w=0$, which depends on $k$ through $N_0'(t_\natural)$.
\end{proof}

\begin{remark}[Scope of \cref{prop:v16-counter}]\label{rem:v16-scope}
\Cref{prop:v16-counter} says nothing about the prime-built $F_X^W$ or $F_h$: it shows that no argument may use only $|F-N_0|\le H$ together with the monotonicity of $N_0$. Whatever separates $\tau_{k+1}$ from $\tau_k$ must come from the arithmetic structure of the phase, through \eqref{eq:v16-band}. Whether $\liminf_{X}\meas\{k-\frac12\le F_X^W<k+\frac12\}>0$ holds is open here.
\end{remark}

\subsection{The target function: Cram\'er's function extends across the window, with explicit principal parts}\label{sec:v17-target}
Notation (new symbols, except $\mathfrak V$, which is the function of \cref{prop:v13-real}): $\mathfrak V(w)=\sum_{\operatorname{Im}\rho>0}e^{\rho w}$ for $\operatorname{Im}w>0$ (as in \cref{prop:v13-real}; the series converges absolutely there because $|e^{\rho w}|\le e^{\operatorname{Re}\rho\,\operatorname{Re}w-\gamma\operatorname{Im}w}$ and $N(T)\ll T\log T$), and
\begin{equation}
\mathscr C(w)=e^{-w/2}\mathfrak V(w),\qquad
\mathfrak p(w)=\tfrac12\Bigl(e^{w/2}-\frac{e^{-w/2}}{e^{2w}-1}\Bigr),\qquad
\mathfrak R=\{w:\operatorname{Re}w>0,\ |\operatorname{Im}w|<\pi\}.
\end{equation}
$\mathfrak p$ is holomorphic on $\mathfrak R$ (the zeros of $e^{2w}-1$ are $w=i\pi n$, $n\in\mathbb Z$, all with $\operatorname{Re}w=0$, so none lies in $\mathfrak R$) and real on $\mathfrak R\cap\mathbb R$, so $\mathfrak p^\natural=\mathfrak p$, where for a function $\mathsf H$ we write $\mathsf H^\natural(w)=\overline{\mathsf H(\bar w)}$ (defined where $\mathsf H$ is defined at $\bar w$). For a distribution $\mathfrak u$ on an interval, $\operatorname{Re}\mathfrak u=\frac12(\mathfrak u+\bar{\mathfrak u})$ and $\operatorname{Im}\mathfrak u=\frac1{2i}(\mathfrak u-\bar{\mathfrak u})$.

\begin{lemma}[Reflection across an interval where the real part is analytic]\label{lem:v17-refl}
Let $U\subset\mathbb C$ be open and symmetric under $w\mapsto\bar w$, with $I=U\cap\mathbb R$ an interval and $U\cap\{\operatorname{Im}w>0\}$ connected. Let $\mathsf H$ be holomorphic on $U\cap\{\operatorname{Im}w>0\}$ with $\mathsf H(\cdot+iy)\to\mathfrak u$ in $\mathcal D'(I)$ as $y\downarrow0$, and let $\mathfrak g$ be holomorphic on $U$ with $\mathfrak g^\natural=\mathfrak g$. If $\operatorname{Re}\mathfrak u=\mathfrak g|_I$ in $\mathcal D'(I)$, then $\mathsf H$ extends holomorphically to $U$, and the extension satisfies $\mathsf H+\mathsf H^\natural=2\mathfrak g$.
\end{lemma}
\begin{proof}
Put $\mathsf H_0=\mathsf H-\mathfrak g$ on $U\cap\{\operatorname{Im}w>0\}$; its boundary value is $\mathfrak u-\mathfrak g=\mathfrak u-\operatorname{Re}\mathfrak u=i\operatorname{Im}\mathfrak u$. Define $\widehat{\mathsf H}_0(w)=-\overline{\mathsf H_0(\bar w)}$ on $U\cap\{\operatorname{Im}w<0\}$; it is holomorphic, and its boundary value at $I$ from below is $-\overline{i\operatorname{Im}\mathfrak u}=i\operatorname{Im}\mathfrak u$ (as $\operatorname{Im}\mathfrak u$ is a real distribution), the same distribution. A function holomorphic on both sides of an interval with the same distributional boundary value on it is the restriction of a function holomorphic across it (one-variable edge-of-the-wedge theorem, Painlevé; \cite[pp.~63--66, 343--344]{Hormander1990}, pages taken from a secondary listing, books not opened; see also \cite{Rudin1971,BOT1958}). So $\mathsf H_0$ extends holomorphically to $U$ (shrinking $U$ about $I$ if needed) with $\mathsf H_0^\natural=-\mathsf H_0$ off $\mathbb R$, hence everywhere by continuity. Then $\mathsf H=\mathfrak g+\mathsf H_0$ extends, and $\mathsf H+\mathsf H^\natural=2\mathfrak g+\mathsf H_0+\mathsf H_0^\natural=2\mathfrak g$.
\end{proof}

\begin{theorem}[Meromorphic continuation of Cram\'er's function with principal parts, without RH]\label{thm:v17-cont}
$\mathscr C$ extends from $\{\operatorname{Im}w>0\}$ to a meromorphic function on $\{\operatorname{Im}w>0\}\cup\mathfrak R$ with
\begin{equation}
\mathscr C(w)+\mathscr C^\natural(w)=2\,\mathfrak p(w)\qquad(w\in\mathfrak R),
\end{equation}
whose poles are exactly the points $w=\log n$ with $\Lambda(n)\ne0$, all simple, with principal part
\begin{equation}\label{eq:v17-pp}
\mathfrak s_n(w)=-\frac{i}{2\pi}\,\frac{\Lambda(n)}{\sqrt n}\,\frac1{w-\log n}.
\end{equation}
Equivalently $\mathfrak V(w)=e^{w/2}\mathscr C(w)$ has the principal part $-\frac{i}{2\pi}\Lambda(n)/(w-\log n)$ at $w=\log n$.
\end{theorem}
\begin{proof}
\emph{Boundary value.} For $\varphi\in C_c^\infty(0,\infty)$ and a zero $\rho=\beta+i\gamma$ with $\gamma>0$, $\int\varphi(u)e^{\rho(u+iy)}du=e^{i\rho y}\int\varphi e^{\rho u}du$ with $|e^{i\rho y}|=e^{-\gamma y}\le1$ and $\int\varphi e^{\rho u}du=O_\varphi(|\rho|^{-2})$ (two integrations by parts), and $\sum|\rho|^{-2}<\infty$. By dominated convergence $\mathfrak V(\cdot+iy)\to\mathfrak V$ in $\mathcal D'(0,\infty)$, the distribution of \cref{prop:v13-real}. Multiplication by the smooth function $e^{-u/2}$ gives the boundary value $\mathfrak u$ of $\mathscr C$ in $\mathcal D'(0,\infty)$, and $\operatorname{Re}\mathfrak u=e^{-u/2}\operatorname{Re}\mathfrak V$, because the factor is real.

\emph{Real part of the boundary value.} By \cref{prop:v13-real}, with $e^{-u/2}\delta_{\log n}=n^{-1/2}\delta_{\log n}$, in $\mathcal D'(0,\infty)$
\begin{equation}
\operatorname{Re}\mathfrak u=\mathfrak p(u)-\tfrac12\sum_{n\ge2}\frac{\Lambda(n)}{\sqrt n}\,\delta_{\log n}(u).
\end{equation}

\emph{Definition in the lower half.} For $w\in\mathfrak R$ with $\operatorname{Im}w<0$ put $\mathscr C(w):=2\mathfrak p(w)-\overline{\mathscr C(\bar w)}$; this is holomorphic there (because $\bar w$ lies in the upper half plane where $\mathscr C$ is defined, and $\mathfrak p$ is holomorphic on $\mathfrak R$) and has no poles.

\emph{Across the real axis away from the atoms.} On an interval $I\subset(0,\infty)\setminus\{\log n:n\ge2\}$, $\operatorname{Re}\mathfrak u=\mathfrak p|_I$. Take $U\subset\mathfrak R$ a small symmetric neighbourhood of $I$ with $U\cap\{\operatorname{Im}w>0\}$ connected. \cref{lem:v17-refl} with $\mathfrak g=\mathfrak p$ extends $\mathscr C$ holomorphically to $U$ with $\mathscr C+\mathscr C^\natural=2\mathfrak p$; the extension below $I$ coincides with the definition above (both are the reflection of the same function). 

\emph{At an atom $a=\log n$, $\Lambda(n)\ne0$.} The boundary value of $\mathfrak s_n$ from above is $-\frac{i}{2\pi}\frac{\Lambda(n)}{\sqrt n}\bigl(\mathrm{PV}\frac1{u-a}-i\pi\delta_a\bigr)$, with real part $-\frac12\frac{\Lambda(n)}{\sqrt n}\delta_a$, and $\mathfrak s_n+\mathfrak s_n^\natural=0$ on $\mathbb C\setminus\{a\}$. Hence $\mathsf H=\mathscr C-\mathfrak s_n$ has, near $a$ (where no other atom lies), boundary value with real part $\mathfrak p$; \cref{lem:v17-refl} extends $\mathsf H$ holomorphically to a neighbourhood of $a$ with $\mathsf H+\mathsf H^\natural=2\mathfrak p$. So $\mathscr C=\mathsf H+\mathfrak s_n$ has a simple pole at $a$ with principal part \eqref{eq:v17-pp}, and $\mathscr C+\mathscr C^\natural=2\mathfrak p$ near $a$ (off the pole). If $\Lambda(n)=0$ there is no atom and the previous case applies.

\emph{Conclusion.} The function so defined is holomorphic on $\{\operatorname{Im}w>0\}$ and on $\{-\pi<\operatorname{Im}w<0,\operatorname{Re}w>0\}$, extends holomorphically across every real interval between atoms and meromorphically at every atom, hence is meromorphic on $\{\operatorname{Im}w>0\}\cup\mathfrak R$ with the stated poles, and $\mathscr C+\mathscr C^\natural=2\mathfrak p$ holds on $\mathfrak R$ by construction.
\end{proof}

\begin{corollary}[Zero-free local form of the target at $u=1$]\label{cor:v17-local}
Put $\mathscr C^\sharp(w)=\mathscr C(w)-\mathfrak s_2(w)-\mathfrak s_3(w)$. Then $\mathscr C^\sharp$ is holomorphic on the disc $|w-1|<\log4-1=0.3863\ldots$, $\operatorname{Re}\mathscr C^\sharp(u)=\mathfrak p(u)$ for real $u$ in it, and $\mathfrak q(u)=\operatorname{Im}\mathscr C^\sharp(u)$ is real-analytic there. For real $u\in(\log2,\log3)$,
\begin{equation}\label{eq:v17-local}
\mathscr C(u)=\mathfrak p(u)+i\Bigl[\frac1{2\pi}\Bigl(\frac{\log3}{\sqrt3\,(\log3-u)}-\frac{\log2}{\sqrt2\,(u-\log2)}\Bigr)+\mathfrak q(u)\Bigr].
\end{equation}
At $u=1$ the explicit part of the bracket equals $0.7694859\ldots$ and $\mathfrak p(1)=0.7768943\ldots$.
\end{corollary}
\begin{proof}
The disc is contained in $\mathfrak R$ (its points have $\operatorname{Re}w>0.6137$ and $|\operatorname{Im}w|<0.3863<\pi$), its real points form $(0.6137,1.3863)$, and the only $\log n$ with $n\ge2$ in the open disc are $\log2$ and $\log3$ ($\log4$ lies on the boundary circle). By \cref{thm:v17-cont}, $\mathscr C$ is meromorphic on the disc with simple poles at $\log2$ and $\log3$ only; subtracting their principal parts leaves $\mathscr C^\sharp$ holomorphic on the whole disc. For real $u\ne\log2,\log3$ in the disc, $\mathfrak s_2(u),\mathfrak s_3(u)$ are purely imaginary, so $\operatorname{Re}\mathscr C^\sharp(u)=\operatorname{Re}\mathscr C(u)=\mathfrak p(u)$ (the latter by $\mathscr C+\mathscr C^\natural=2\mathfrak p$), and by continuity at the two points. Formula \eqref{eq:v17-local} is $\mathscr C=\mathscr C^\sharp+\mathfrak s_2+\mathfrak s_3$ with $\mathscr C^\sharp=\mathfrak p+i\mathfrak q$. Values: $\frac1{2\pi}\frac{\log3}{\sqrt3(\log3-1)}=1.0237005\ldots$ and $\frac1{2\pi}\frac{\log2}{\sqrt2(1-\log2)}=0.2542146\ldots$, difference $0.7694859\ldots$.
\end{proof}

\begin{remark}[What \cref{thm:v17-cont} contains and what it leaves open]\label{rem:v17-scope}
(i) \Cref{thm:v17-cont} is Cram\'er's continuation \cite{Cramer1919} in the region $\mathfrak R$ (attribution recalled, not verified), here with the residues derived from \cref{prop:v13-real} and no assumption on the zeros. It shows why the phase of $\mathscr C$ is smooth on $J$: $\operatorname{Re}\mathscr C=\mathfrak p>0$ on $J$ (as $e^{u/2}>1>e^{-u/2}/(e^{2u}-1)$ there), so $\mathscr C$ has no zero on $J$ and $\arg\mathscr C$ is real-analytic on $J$; numerically $\mathscr C(1)=0.7769+0.7679\,i$ (the value quoted in \cref{prop:v11-gamma}(b), cross-checked in \S\ref{sec:v9-numbers}). (ii) The function $\mathfrak q$ has no explicit formula here: the series of all principal parts does not converge absolutely, and the value $\mathfrak q(1)=0.767855669-0.769485915=-0.0016302\ldots$ is obtained from the quoted $\operatorname{Im}\bigl(e^{-1/2}\mathfrak V(1)\bigr)=0.76785566886\ldots$ (\cref{prop:v11-gamma}(b)). (iii) Under RH, $\mathscr C(1)$ is the limit of the centroid of $\gamma$ at $u=1$ (\cref{prop:v11-gamma}(b), quoted); on $J$ the relation holds in $\mathcal D'(J)$ up to the defect of \cref{prop:v13-def}; the pointwise convergence at $u=1$ is quoted without proof, and for other $u$ it is not proved here.
\end{remark}

\begin{corollary}[Zero-free target and the zeros]\label{cor:v17-rig}
Let $\mathsf x$ satisfy \eqref{eq:v13-G}, $\mathsf x_k$ nondecreasing, and suppose that the centroid of $\mathsf x$ tends to $\mathscr C$ in $\mathcal D'(J)$, i.e.\ $\mathscr N_\infty^{(\mathsf x)}(f)=\int_Jf(u)\,\mathscr C(u)\,du$ for all $f\in C_c^\infty(J)$ (on $J$, $\mathscr C$ is the real-analytic function of \cref{thm:v17-cont}, which is the distribution $e^{-u/2}\mathfrak V$ of \cref{prop:v13-def}, since a holomorphic function's boundary value is the function itself where it extends analytically). Then
\begin{equation}
\mathsf x_k=\gamma_k\ \ \forall k\iff\text{RH}.
\end{equation}
\end{corollary}
\begin{proof}
By \cref{prop:v13-dist} applied to $\mathsf x$ and to $\gamma$, and by the hypothesis, $\lim_\ell\int_Jf\,(C_\ell^{(\gamma)}-C_\ell^{(\mathsf x)})=\mathscr N_\infty^{(\gamma)}(f)-\int_Jf\,\mathscr C=-\sum_k\int_Jf(\cosh(u\mathsf o_k)-1)e^{iu\gamma_k}du$ by \cref{prop:v13-def}. By \cref{prop:v13-rig}, the left side tends to $0$ for every $f$ iff $\mathsf x_k=\gamma_k$ for all $k$; by \cref{prop:v13-def}, the right side vanishes for every $f$ iff RH.
\end{proof}

\begin{remark}[Use of \cref{cor:v17-rig}]\label{rem:v17-use}
The statement has content on the side of the sequence: the target $\mathscr C$ on $J$ is fixed by the primes up to the function $\mathfrak q$ (\cref{thm:v17-cont}, \cref{cor:v17-local}), and any $\mathsf x$ satisfying (Sp) with that centroid limit is the zero sequence exactly when RH holds. No such prime-built $\mathsf x$ is available: under RH the points $\tau^{\mathrm C}$ do not have this limit (\cref{cor:v13-notransfer}(ii), since $\tau^{\mathrm C}\ne\gamma$); without RH the question is open. Matching the phase of $\mathscr C$ at finitely many frequencies, or one real function of $u$ such as $\arg\mathscr C$, does not identify the $\mathsf x_k$; \cref{prop:v13-rig} needs the complex limit on all of $J$.
\end{remark}

\subsection{Exploratory numbers (not proofs)}\label{sec:v9-numbers}
\emph{Model computation (this note, Python, not the IJNT phase).} With a model of $F$ built from the primes $p^k\le\sqrt{1.5T}$ with a log-linear taper (full weight up to $X^{1/2}$, $X$ being the cutoff parameter of the model, not reproduced here), the quantity $\mathsf r_\omega=-2\pi\,\mathrm{mean}\bigl(e^{i\omega t}\,\mathrm{saw}(F)\bigr)$ (the model's $\varsigma_T$) has modulus decreasing with $T$ for $\omega=\log2,\log3,\log5,\log7$ (e.g.\ $0.0142,0.0088,0.0054,0.0040$ for $\omega=\log2$ at $T=10^3,\dots,6.4\cdot10^4$) and is of order $10^{-3}$ for $\omega=0$. The sorted-sublevel sequence at $T=10^4$ reproduces the Landau lines $\Xi=\frac{T}{2\pi M_T}\chi(\omega)$ (\cref{prop:v9-rvm}; $M_T=N_0(2T)-N_0(T)$, so $\Xi$ differs by $0.3\%$ from $\chi(\omega)/\log(1.5T/2\pi)$) to within $\approx7\%$ ($\omega=\log2$: $-0.0632$ predicted vs $-0.0637$ observed; $\log3$: $-0.0818$ vs $-0.0836$; $\log 5$: $-0.0928$ vs $-0.0970$; $\log 7$: $-0.0948$ vs $-0.1010$) and is $\approx10^{-3}$ at $\omega=\log6,\log10$ where $\Lambda=0$. This supports $\varsigma_T\to0$ in the model but proves nothing for the actual construction.

\emph{First $220$ zeros, $\omega=1$ (ChatGPT text, recomputed there with \texttt{mpmath}, step $\approx0.004$; interval arithmetic not used).} RMS of $\gamma_k-\mathsf x_k$: $0.4226$ (RvM), $0.0689$ ($\tau^{\mathrm C}$ and $\tau^{\Cm}$ coincide at this height); $\max|\gamma_k-\mathsf x_k|$: $0.9937$, $0.5853$; $|\mathrm{Cen}^{(\gamma)}-\mathrm{Cen}^{(x)}|$: $1.2032$, $0.6060$; $\Theta^{\mathrm{RvM}}_T(0,1)\approx0.9137-0.0006i$, $\Theta^{\mathrm C}(0,1)\approx0.9976+0.0038i$; $\mathrm{Cen}^{(\gamma)}_{220}(1)=0.7739+0.8022i$, $\mathrm{Cen}^{(\mathrm C)}_{220}(1)=1.3167+0.5330i$. These agree in sign and size with \cref{thm:v9-coherence} and \cref{cor:v9-cen}(iii) (the values $|\mathrm{Cen}|\approx1$ are far below $\ell/\log\ell\approx40$), but $220$ zeros reach only $T\approx430$ and settle no asymptotic question. Not re-run here.

\emph{Repository record up to $2\cdot10^6$ zeros (README of the repository directory \texttt{KC\_\allowbreak tauC\_\allowbreak 2026-11-02}, quoted, not re-run here).} $\mathrm{Cen}_N(1;\gamma)-\mathrm{Cen}_N(1;\gamma^{\mathrm{RvM}})$ is stable at $K^0=0.735337-0.989697\,i$ to six digits from $N\approx10^5$; $|\mathrm{Cen}_N(1;\gamma)-\mathrm{Cen}_N(1;\tau^{\mathrm C})|=1.39,\ 1.97,\ 10.8,\ 14.7$ at $N=10^4,10^5,10^6,2\cdot10^6$. The first is compatible with $\mathrm{Cen}_\ell(1;\gamma^{\mathrm{RvM}})\to\mathfrak m(1)$ (\cref{prop:v11-rvm}) and $\mathrm{Cen}_\ell(1;\gamma)=o(\ell/\log\ell)$ (\cref{cor:v9-cen}); the second grows with $N$ but stays far below $\ell/\log\ell\approx1.4\cdot10^5$ at $\ell=2\cdot10^6$, so it is compatible with \cref{thm:v9-cen} for $\mathfrak b=0$, and neither shows nor excludes convergence. Even $\varsigma_T\to0$ would give only $o(\ell/\log\ell)$ for the difference of centroids.

\emph{Recomputation from the zero file and the RvM file (this note; C and Python, double precision; files \texttt{zeross.bin}, \texttt{zLamb.bin}, $2\,001\,052$ zeros each).} The corrected points $\tau^{\mathrm C}$ were recomputed from the definition (sublevel measures of $F_h$, cubic Hermite cells with root bisection; source \texttt{tauC-c.c} in the repository) and compared with the supplied file \texttt{zTaucs2e6.bin}: $\max|\tau^{\mathrm C}_{\rm new}-\tau^{\mathrm C}_{\rm file}|=3.0\cdot10^{-6}$ (at $k=1\,818\,451$), rms $2.0\cdot10^{-7}$, none above $10^{-3}$; the RvM points recomputed by Newton's method from $N_0(t)=k-\frac12$ agree with \texttt{zLamb.bin} to $2.3\cdot10^{-10}$. The file \texttt{zTaucs2e6c.bin} contains the recomputed points. The recomputed points are strictly increasing.

\begin{center}
\small
\begin{tabular}{rrrrrr}
$N$ & $\mathrm{RMS}(\sigma_k^{\mathrm{RvM}})$ & $\mathrm{RMS}(\sigma_k^{\mathrm C})$ & $|\mathrm{Par}_N(\gamma)|$ & $|\mathrm{Par}_N(\tau^{\mathrm C})|$ & $|\mathrm{Par}_N(\gamma^{\mathrm{RvM}})|$\\\hline
$10^3$ & -- & $0.0554$ & $0.30$ & $0.78$ & $0.86$\\
$10^4$ & $0.2838$ & $0.0446$ & $1.71$ & $2.51$ & $1.40$\\
$10^5$ & $0.2327$ & $0.0378$ & $2.65$ & $6.34$ & $0.92$\\
$10^6$ & -- & $0.0328$ & $2.29$ & $13.15$ & $0.82$\\
$2\,001\,052$ & $0.1881$ & $0.0315$ & $1.42$ & $26.31$ & $3.67$
\end{tabular}
\end{center}
(RMS over the first $N$ indices, $\omega=1$; dashes: not computed.) Over $10^3\le N<2\cdot10^6$, $|\mathrm{Par}_N(\gamma)|\le4.74$ throughout, while $|\mathrm{Par}_N(\tau^{\mathrm C})|/\sqrt N$ has mean $0.027,\ 0.017,\ 0.019,\ 0.016$ on $[10^3,10^4)$, $[10^4,10^5)$, $[10^5,10^6)$, $[10^6,2\cdot10^6)$. Hence $|\mathrm{Par}_N(\gamma)-\mathrm{Par}_N(\tau^{\mathrm C})|/(\mathrm{RMS}(\sigma^{\mathrm C})\sqrt N)=0.61,\,0.49,\,0.37,\,0.43,\,0.58$ at the five values of $N$: the partial sums of the $\tau^{\mathrm C}$-phases differ from those of the zeros by an amount of the order of $\sqrt N$ times the rms displacement, and not by $o(\sqrt N)$. The displacements $\sigma_k^{\mathrm C}$ have lag-$1$ autocorrelation $-0.400$ and autocorrelations $-0.029,\ -0.006,\ -0.001$ at lags $2,\,5,\,50$.

\emph{Phase transport at $N=2\,001\,052$.} $\mathrm{Cen}_N(1;\gamma)=0.776907+0.767845\,i$, $\mathrm{Cen}_N(1;\gamma^{\mathrm{RvM}})=0.041570+1.757542\,i$, $\mathrm{Cen}_N(1;\tau^{\mathrm C})=15.209902+3.533042\,i$, so $\mathsf D_N^{\mathrm C}(1)=-14.432995-2.765198\,i$. In \eqref{eq:v16-moment} at $\omega=1$: $\mathsf M_N^{(1)}=-2.3708+14.4155\,i$, $\mathsf M_N^{(2)}=0.1035+0.7562\,i$, $\mathsf M_N^{(3)}=0.1255+0.1930\,i$ (moduli $14.609,\ 0.763,\ 0.230$); the orders $m=1,2,3$ leave errors $0.3948,\ 0.0380,\ 0.00508$; the error for $m=3$ lies below the bound of \eqref{eq:v16-moment}, $\frac1{24N}\sum_k(N-k+1)\sigma_k^4=0.1273$. The partial sums satisfy $\max_{n\le N}|\mathsf s_n(1)|=27.9$, $|\mathsf s_N(1)|=25.1$, and their mean is $\mathsf M_N^{(1)}$ by \eqref{eq:v16-abel}. On the grid $u_m=0.75+0.0003\,m$, $m=0,\dots,1000$ ($1001$ frequencies), the rms over the grid is $12.703$ for $|\mathsf D_N^{\mathrm C}|$ and $0.6321,\ 0.02676,\ 0.003610$ for the errors of the expansions of order $1,2,3$ (that is $5.0\%$, $0.21\%$, $0.028\%$ of the rms). The value $u=1$ is not a grid point of this progression.

\emph{Block coefficient and separation.} With $\mathfrak b^{\rm est}(T)=\frac{2\pi}T\sum_{T\le\tau_k^{\mathrm C}<2T}\bigl(e^{i\gamma_k}-e^{i\tau_k^{\mathrm C}}\bigr)$ (blocks cut by $\tau^{\mathrm C}$, $\omega=1$), $|\mathfrak b^{\rm est}(T)|=0.0082,\ 0.0106,\ 0.0076,\ 0.0036,\ 0.0016,\ 0.0003,\ 0.0012,\ 0.0010,\ 0.0009$ for $T=2^8,\dots,2^{16}$, $<5\cdot10^{-5}$ for $T=2^{17}$, and $0.0001$ for $T=2^{18},2^{19}$. For the spacing of $\tau^{\mathrm C}$: minimum gap $0.1260$; ratio of consecutive gaps of $\tau^{\mathrm C}$ to those of $\gamma^{\mathrm{RvM}}$ has minimum $0.2260$ (at $k=716\,289$) and maximum $3.2228$, with minima $0.2576,\ 0.2574,\ 0.2260,\ 0.2352$ on $k<10^3$, $[10^3,10^5)$, $[10^5,10^6)$, $[10^6,N)$. These are the quantities \eqref{eq:v16-band} measures in the range reached; they show no collapse of consecutive corrected points up to $N=2\,001\,052$ and decide nothing beyond it.

\emph{Test of \cref{thm:v17-cont} and \cref{cor:v17-local} on the $2\,001\,052$ zeros (this note, Python).} On the grid $u_m=0.75+0.0003\,m$, $m=0,\dots,1000$, the centroid $\mathrm{Cen}_N(u_m;\gamma)$ has $\max_m|\operatorname{Re}\mathrm{Cen}_N-\mathfrak p|=5.0\cdot10^{-4}$ and rms $1.9\cdot10^{-5}$. For the imaginary part, the rms residual of a least-squares polynomial of degree $d$ in $u$ (on $[0.75,1.05]$) is, for $d=1,\dots,6$: $1.5\cdot10^{-1},\ 1.3\cdot10^{-1},\ 3.1\cdot10^{-2},\ 2.6\cdot10^{-2},\ 6.6\cdot10^{-3},\ 5.4\cdot10^{-3}$ when fitting $\operatorname{Im}\mathrm{Cen}_N$ alone, and $6.6\cdot10^{-3},\ 1.0\cdot10^{-3},\ 1.2\cdot10^{-4},\ 3.3\cdot10^{-5},\ 2.9\cdot10^{-5},\ 2.9\cdot10^{-5}$ after subtracting the two principal parts $\mathfrak s_2,\mathfrak s_3$ of \eqref{eq:v17-pp} (no fitted coefficient: the coefficients are $\Lambda(n)/(2\pi\sqrt n)$). Adding the principal part $\mathfrak s_4$ gives $3.0\cdot10^{-3},\ 4.4\cdot10^{-4},\ 4.1\cdot10^{-5},\ 2.9\cdot10^{-5}$ for $d=1,\dots,4$. In every model with the two principal parts and $d\ge5$, or with the three principal parts and $d\ge4$, the residual stays at $2.9\cdot10^{-5}$; the origin of this floor was not analysed. The fitted value of the smooth remainder at $u=1$ is $\mathfrak q(1)\approx-0.0016$, against $-0.0016302$ obtained from \eqref{eq:v17-local} and the quoted value of $\mathfrak V(1)$.

\emph{Direct evaluation of $\mathscr C$ at $u=1$ from the zeros.} Let $\mathsf S(y)=\sum_{k}e^{i\gamma_k(1+iy)}$ over the $2\,001\,052$ zeros, $y\in\{0.0004,0.0008,\dots,0.0024\}$ (the tail beyond $\gamma_{\max}=1\,132\,490.66$ is $<10^{-190}$ in size). This equals $\mathscr C(1+iy)$ when all $\beta_k=\frac12$, which is known to hold for these zeros by the numerical verifications of RH in the literature (recalled, not re-verified here). Extrapolating $y\to0$ with polynomials of degree $2,3,4$ (legitimate since $\mathscr C$ is holomorphic near $1$ by \cref{thm:v17-cont}) gives $0.7768960+0.7678557\,i$, $0.7768943+0.7678556\,i$, $0.7768943+0.7678556\,i$, against $0.7768943+0.7678557\,i$ of \cref{prop:v11-gamma}(b): degrees $3$ and $4$ agree with it to about $10^{-7}$, degree $2$ to $2\cdot10^{-6}$.

\subsection{What remains open in this section}\label{sec:v9-open-centroid}
\begin{enumerate}[label=(\alph*),leftmargin=*]
\item $\varsigma_T(\omega;F_h)\to0$ for $\omega=r\log p$ and $\omega=1$ (decides whether $\mathfrak b^{\mathrm C}=\mathfrak b^{\Cm}=0$, \cref{cor:v9-cen}(ii)); \cref{prop:v9-harm} reduces it to (HC) and (NC), i.e.\ to the cancellation of the harmonics $e^{i\omega t+2\pi i\mathfrak hF}$ with moving prime resonances (\cref{rem:v9-res}) and to non-concentration of $F$ modulo $1$ at the half-integers; \cref{cor:v9-test} is the corresponding statistic on the points $\tau_k^{\mathrm C}$.
\item The Gaussian law at the points $\gamma_k$ (\cref{prop:v9-gauss}) and the scale $\omega\asymp\log T$ (second-moment input insufficient).
\item Convergence of $\mathrm{Cen}_\ell(1;\gamma)$ without RH, and of $\mathrm{Cen}_\ell(1;\tau^{\mathrm C})$, $\mathrm{Cen}_\ell(1;\tau^{\Cm})$ with or without RH (\cref{prop:v11-gamma}, \cref{prop:v11-tau}, \cref{tab:v11-centroids}); for $\gamma$ the Tauberian criterion of \cref{prop:v11-gamma}(c) (\cref{thm:v14-taub}) is the exact reformulation.
\item Pointwise convergence at $u=1$ of $\operatorname{Re}\mathrm{Cen}_\ell(u;\gamma)$ or of $\operatorname{Im}\mathrm{Cen}_\ell(u;\gamma)$ without RH (\cref{rem:v13-which}(iii), unverified recollection), and the value of the boundary distribution $\mathscr V_{\mathrm C}(\cdot+i0)$ of \cref{cor:v13-notransfer}.
\end{enumerate}

\section{The taucw points}\label{sec:taucw}
This section concerns a single object: the sublevel points of $N_0$ plus the Selberg-weighted Euler phase with prime support up to $t^{2\theta}$, called taucw. It is self-contained except for the analytic inputs listed here. It uses neither the dyadic interpolation and prime budget of the original taucs \eqref{eq:old} nor the fourth-moment inputs that \cref{sec:taucm} attributes to the IJNT manuscript. The analytic inputs are the Riemann--von Mangoldt formula \eqref{eq:RVM}, the Selberg--Fujii increment estimate \eqref{eq:FujiiB} (in the form stated in \S\ref{sec:definitions}), Selberg's fourth moment of $S$ \cite{Selberg1946,Titchmarsh}, and the explicit bound \eqref{eq:MCOMH}, which enters only through the truncation \eqref{eq:taucwH}. The subsections \ref{sec:taucw-exist}--\ref{sec:taucw-S} use in addition only the mean value theorem of Montgomery and Vaughan \cite{MontgomeryVaughan}, Mertens' theorem \cite{Mertens1874} and the prime number theorem in the form $\pi(x)\ll x/\log x$ with its partial-summation consequences \cite{Davenport2000}; they use neither \eqref{eq:FujiiB} nor \eqref{eq:MCOMH}.

\begin{remark}[Nature of the construction]\label{rem:taucw-nature}
The function $F^{\mathrm w}_\theta$ of \cref{def:taucw} is a finite explicit expression in the primes and prime powers up to $t^{2\theta}$, the weight $W$ and the exponent $\theta$; the truncated version adds only the explicit function $H$ of \eqref{eq:MCOMH}. No zero of $\zeta$, no zero-free region and no information on the location of the zeros enters its definition or its computation. The zeros are the target with which the sublevel points are compared, and $\zeta$ enters the proofs below through the explicit formula \eqref{eq:taucw-expl}, the Riemann--von Mangoldt formula \eqref{eq:RVM}, the cited moment estimates for $S$ (\eqref{eq:FujiiB} and Selberg's fourth moment) and the explicit bound \eqref{eq:MCOMH} for $S$ (through the truncation \eqref{eq:taucwH}), and nowhere else. The taucw points are therefore not a reformulation of RH: RH is a statement about $\zeta$, and $\zeta$ is absent from the object. Every parameter that controls the object ($W$, $\theta$, $H$) is explicit and finite. The use of Selberg smoothing and of finite Euler products is not claimed as new \cite{Gonek2012}.
\end{remark}

\subsection{Definition}\label{sec:taucw-def}
\begin{definition}[taucw]\label{def:taucw}
Let $W(u)=\min\{1,2-u\}_+$ be the Selberg weight, the case $q\equiv1$ of \eqref{eq:weights}. For $\theta>0$ and $t\ge2$ put $X=t^\theta$ and
\begin{align}
\Phi^{\mathrm w}_\theta(t)&=-\frac1\pi\sum_{2\le n\le t^{2\theta}}
\frac{\Lambda(n)\,W\!\left(\frac{\log n}{\theta\log t}\right)}{\sqrt n\,\log n}\,\sin(t\log n),\label{eq:taucw-phase}\\
F^{\mathrm w}_\theta(t)&=N_0(t)+\Phi^{\mathrm w}_\theta(t),\qquad
\tau_k^{\mathrm w}=2\pi+\meas\{t\ge2\pi:\ F^{\mathrm w}_\theta(t)<k-\tfrac12\}.\label{eq:taucw-F}
\end{align}
Here $k\ge1$. For $0<\theta\le1$ the measure in \eqref{eq:taucw-F} is finite because $|\Phi^{\mathrm w}_\theta(t)|\ll t^\theta/\log t=o(N_0(t))$; for $\theta>1$ finiteness is part of the statement wherever $\tau_k^{\mathrm w}$ is used. The points $\tau_k^{\mathrm w}$ are the taucw points; the case $\theta=1$ is the one computed in \cref{sec:taucw-data}. The sublevel set is measured from $2\pi$, the minimum of $N_0$, as in \eqref{eq:v8-flat}, because $F^{\mathrm w}_\theta$ is defined for $t\ge2$ only. The coefficient of $\sin(t\log n)$ contains the weight $1$ for $n\le t^\theta$ and $\log(t^{2\theta}/n)/(\theta\log t)$ for $t^\theta<n\le t^{2\theta}$. With $H$ of \eqref{eq:MCOMH} the truncated version is
\begin{equation}\label{eq:taucwH}
F^{\mathrm w,H}_\theta(t)=N_0(t)+\clip_{H(t)}\bigl(\Phi^{\mathrm w}_\theta(t)\bigr),\qquad
\tau_k^{\mathrm w,H}=2\pi+\meas\{t\ge2\pi:\ F^{\mathrm w,H}_\theta(t)<k-\tfrac12\}.
\end{equation}
\end{definition}
For fixed $t$ with $t^\theta\ge2$, $\Phi^{\mathrm w}_\theta(t)$ is the value at $X=t^\theta$ of the phase $\Phi_X^W(t)$ of \eqref{eq:PhiW}, because $\Im\,n^{-1/2-it}=-n^{-1/2}\sin(t\log n)$.

\begin{lemma}[The compensator is negligible in the regime $X=t^\theta$, $\theta\le1$]\label{lem:taucw-pole}
For the Selberg weight, uniformly in $X\ge2$ and $t\ge1$,
\begin{equation}\label{eq:taucw-pole}
\Pi_X^W(t)\ll\frac1t+\frac{X}{t^2\log^2X}.
\end{equation}
In particular, without any hypothesis on the zeros,
\begin{equation}\label{eq:taucw-pole-t}
\Pi^W_{t^\theta}(t)\ll\frac1t\quad(0<\theta\le1,\ t^\theta\ge2),\qquad\text{hence}\qquad
\Phi^{\mathrm w}_\theta(t)=P^W_{t^\theta}(t)+O\Bigl(\frac1t\Bigr),
\end{equation}
and for the cutoff $X=t$ of the original taucw phase, $\Pi_t^W(t)\ll1/t\ll(\log t)^{-2}$.
\end{lemma}
\begin{proof}
Put $L=\log X\ge\log2$ and $a(v)=W(v/L)e^{v/2}/v$ for $\log2\le v\le2L$, $a=0$ for $v>2L$. Substituting $u=e^v$ in \eqref{eq:PiW} gives $\pi\Pi_X^W(t)=\Im\int_{\log2}^{2L}a(v)e^{-itv}\,dv$. The function $a$ is continuous, $a(\log2)=\sqrt2/\log2$ because $W=1$ on $[0,1]$, $a(2L)=0$, and $a$ is piecewise $C^2$ with breaks at $v=L$ and $v=2L$ only. There $a'$ jumps by $\sqrt X/L^2$ and $X/(2L^2)$ respectively, and $a'(\log2)=O(1)$. Away from the breaks $a''=W(v/L)g''+(2/L)W'(v/L)g'$ with $g(v)=e^{v/2}/v$ and $W'\in\{0,-1\}$, so that
\begin{equation}\label{eq:taucw-a2}
\int|a''|\,dv\ll\int_{\log2}^{L}|g''|\,dv+\frac1L\int_L^{2L}|g'|\,dv+\int_L^{2L}W\Bigl(\frac vL\Bigr)|g''|\,dv\ll1+\frac{\sqrt X}L+\frac{X}{L^2}+\frac X{L^2},
\end{equation}
the last integral being of order $X/L^2$ because $W(v/L)=(2L-v)/L$ vanishes linearly at $v=2L$ while $g''(v)\asymp e^{v/2}/v$ there. Two integrations by parts against $e^{-itv}$ give
$\int a\,e^{-itv}dv\ll t^{-1}a(\log2)+t^{-2}\bigl(|a'(\log2)|+\sqrt X/L^2+X/L^2+\int|a''|\bigr)$. Since $e^L\ge L^2$ for all $L>0$, one has $1\le X/L^2$ and $\sqrt X/L\le X/L^2$, and this is \eqref{eq:taucw-pole}. For $X=t^\theta$ with $\theta\le1$, $X/(t^2\log^2X)\le1/(t\log^2X)$, which gives \eqref{eq:taucw-pole-t}.
\end{proof}

\begin{remark}[Relation with \cref{lem:pole}]\label{rem:taucw-pole}
The bound $(1+X)/t$ of \cref{lem:pole} is of size $O(1)$ at $X=t$ and does not exclude a compensator of constant size in the regime of taucw. \Cref{lem:taucw-pole} shows that there is none: one integration by parts that keeps the exact variation of the amplitude on $[L,2L]$ already gives $\Pi_t^W(t)\ll(\log t)^{-2}$, and the second one, which uses that $W$ vanishes linearly at $2L$, gives $\ll1/t$. For $t$ fixed and $X\to\infty$ the compensator does not tend to zero (the bound \eqref{eq:taucw-pole} grows with $X$), so \eqref{eq:taucw-pole-t} is a statement on the regime $X\le t$, not on the limit $X\to\infty$ at fixed $t$. A computation of $\Pi_X^W(t)$ in closed form, by means of the exponential integral $E_1$, gives $t\,|\Pi^W_t(t)|\le0.66$ at $t=50,100,300,10^3,3\cdot10^3,10^4,10^5,10^6,10^7$ (output file of \texttt{check\_pi.py} in the archive; the values are $0.65,0.64,0.54,0.26,0.62,0.28,0.12,0.20,0.65$).
\end{remark}

\subsection{An explicit formula, with remainder \texorpdfstring{$O(1/t)$}{O(1/t)}, for the error of the phase}\label{sec:taucw-exact}
For $L>0$ put
\begin{equation}\label{eq:kernel-G}
\mathfrak G_L(w)=\frac{e^{Lw}-e^{2Lw}}{Lw^2},\qquad
\mathfrak K_L(\delta)=\frac{\cos L\delta-\cos2L\delta}{\pi L\delta^2}\quad(\delta\neq0),
\end{equation}
and for $\Im w\ne0$ let $\mathfrak I_L(w)=\int_{-\infty}^{w}\mathfrak G_L(u)\,du$, the integral along the horizontal ray ending at $w$; it converges because $|\mathfrak G_L(u)|\le2e^{L\Re u}/(L|u|^2)$ for $\Re u\le0$. Let $\mathfrak H_L(v)=\int_{-\infty}^v\mathfrak K_L(\delta)\,d\delta$ and $C_{\mathfrak K}=\int_{\mathbb R}|\mathfrak K_1|<\infty$.

\begin{lemma}[The kernels]\label{lem:taucw-kernel}
For every $L>0$:
\begin{enumerate}[label=(\alph*),leftmargin=*]
\item $\mathfrak K_L$ is even, $|\mathfrak K_L(\delta)|\le L\min\{1,(L\delta)^{-2}\}$, $\int_{\mathbb R}\mathfrak K_L=1$ and $\int_{\mathbb R}|\mathfrak K_L|=C_{\mathfrak K}$.
\item For $v\neq0$, $\ \bigl|\mathfrak H_L(v)-\mathbf 1_{v>0}\bigr|\le\min\{C_{\mathfrak K},\,3/(\pi L^2v^2)\}$.
\item For $v\ne0$, $\ \frac1\pi\Im\,\mathfrak I_L(iv)=\mathbf 1_{v>0}-\mathfrak H_L(v)$.
\item For $\Re w\le0$, $|\mathfrak I_L(w)|\le2e^{L\Re w}/(L^2(\Im w)^2)$.
\end{enumerate}
\end{lemma}
\begin{proof}
(a) $\mathfrak K_L(\delta)=L\mathfrak K_1(L\delta)$, and $\mathfrak K_1(u)=\frac{2\sin(3u/2)\sin(u/2)}{\pi u^2}$ gives $|\mathfrak K_1|\le3/(2\pi)$ and $|\mathfrak K_1(u)|\le2/(\pi u^2)$. From $\int(1-\cos au)u^{-2}du=\pi|a|$ one gets $\int\mathfrak K_1=\frac1\pi(2\pi-\pi)=1$. (b) For $v>0$, $\mathfrak H_L(v)-1=-\int_v^\infty\mathfrak K_L$ and one integration by parts gives $|\int_v^\infty\cos(a\delta)\delta^{-2}d\delta|\le2/(av^2)$; with $a=L$ and $a=2L$ this yields $\bigl|\int_v^\infty\mathfrak K_L\bigr|\le3/(\pi L^2v^2)$; the case $v<0$ is symmetric, and the bound $C_{\mathfrak K}$ is trivial. (c) The function $f(v)=\frac1\pi\Im\,\mathfrak I_L(iv)$ satisfies $f'(v)=\frac1\pi\Im(i\,\mathfrak G_L(iv))=\frac1\pi\Re\,\mathfrak G_L(iv)=-\mathfrak K_L(v)$ on each half-line, and $f(v)\to0$ as $|v|\to\infty$ by (d). Hence $f(v)=\int_v^\infty\mathfrak K_L=1-\mathfrak H_L(v)$ for $v>0$ and $f(v)=-\mathfrak H_L(v)$ for $v<0$. (d) For $\Re u\le0$, $|\mathfrak G_L(u)|\le2e^{L\Re u}/(L|u|^2)\le2e^{L\Re u}/(L(\Im w)^2)$ on the ray, and $\int_{-\infty}^{\Re w}e^{Lx}dx=e^{L\Re w}/L$.
\end{proof}

\begin{lemma}[Smoothed explicit formula]\label{lem:taucw-expl}
Let $X=e^L\ge2$ and $W$ the Selberg weight.
\begin{enumerate}[label=(\alph*),leftmargin=*]
\item For $s$ distinct from $1$, the zeros and the trivial zeros,
\begin{equation}\label{eq:taucw-zl}
\frac{\zeta'}{\zeta}(s)=-\sum_{n\ge2}\Lambda(n)W\!\Bigl(\frac{\log n}{L}\Bigr)n^{-s}-\mathfrak G_L(1-s)+\sum_\rho\mathfrak G_L(\rho-s)+\sum_{m\ge1}\mathfrak G_L(-2m-s),
\end{equation}
the sum over $\rho$ running over all nontrivial zeros with multiplicity.
\item For $t\ge t_0$ not the imaginary part of a zero, uniformly in $X\ge2$ (the identities \eqref{eq:taucw-piS} and $\Im\,\mathfrak I_L(\frac12-it)=-\pi\Pi^W_X(t)+\mathscr A(t)$ of the proof hold for every $t>0$ that is not the imaginary part of a zero),
\begin{equation}\label{eq:taucw-expl}
N_0(t)+P_X^W(t)-N(t)=\frac1\pi\Im\sum_\rho\mathfrak I_L\bigl(\rho-\tfrac12-it\bigr)+O\Bigl(\frac1t\Bigr).
\end{equation}
\end{enumerate}
\end{lemma}
\begin{proof}
(a) For $c>0$ and $y\ge1$, $\frac1{2\pi i}\int_{(c)}\mathfrak G_L(w)y^{-w}dw=\frac1L\bigl[(L-\log y)_+-(2L-\log y)_+\bigr]=-W(\log y/L)$, since $\frac1{2\pi i}\int_{(c)}e^{w\xi}w^{-2}dw=\xi_+$. Hence for $\Re s>1$ the Dirichlet series is $\sum_n\Lambda(n)W(\tfrac{\log n}L)n^{-s}=\frac1{2\pi i}\int_{(c)}\mathfrak G_L(w)\frac{\zeta'}{\zeta}(s+w)\,dw$. Move the line of integration to the left. The poles are $w=0$ (simple, with residue $-1$ for $\mathfrak G_L$, since $\mathfrak G_L(w)=-1/w+O(1)$), $w=1-s$ (residue $-\mathfrak G_L(1-s)$), $w=\rho-s$ and $w=-2m-s$ (residues $\mathfrak G_L(\rho-s)$, $\mathfrak G_L(-2m-s)$). The horizontal integrals tend to zero on heights avoiding the zeros, because $\mathfrak G_L(w)\ll_L|w|^{-2}$ and $\zeta'/\zeta\ll\log^2$ there, and the integral on $\Re w=-U$ tends to zero because $\mathfrak G_L$ decays like $e^{L\Re w}$. The sums over zeros converge absolutely since $\sum_\rho|\rho|^{-2}<\infty$. This is \eqref{eq:taucw-zl}, which extends by analytic continuation.

(b) For $t>0$ not the imaginary part of a zero, $\pi S(t)=\Im\log\zeta(\tfrac12+it)$ with the branch tending to $0$ as $\sigma\to\infty$ on $\Im s=t$, and $\log\zeta(\tfrac12+it)=-\int_{1/2}^\infty\frac{\zeta'}{\zeta}(\sigma+it)\,d\sigma$. Integrate \eqref{eq:taucw-zl} termwise (the interchange of $\sum_\rho$ and $\int d\sigma$ is justified as follows: with $u=\rho-\sigma-it$ one has $|u|\ge|\gamma-t|$, $|\mathfrak G_L(u)|\le2e^{2L(\Re u)_+}/(L|u|^2)\le2X/(L(\gamma-t)^2)$ for $\frac12\le\sigma\le1$, since $\Re u=\beta-\sigma<\frac12$, and $|\mathfrak G_L(u)|\le2e^{L\Re u}/(L|u|^2)$ for $\sigma\ge1$, since then $\Re u\le0$; hence $\sum_\rho\int_{1/2}^\infty|\mathfrak G_L(\rho-\sigma-it)|\,d\sigma\le(X/L+2/L^2)\sum_\rho(\gamma-t)^{-2}<\infty$ for such $t$, while for the trivial zeros $\sum_{m\ge1}\int_{1/2}^\infty|\mathfrak G_L(-2m-\sigma-it)|\,d\sigma\le\sum_{m\ge1}2e^{-L(2m+1/2)}/(L^2(2m+\frac12)^2)<\infty$): $\int_{1/2}^\infty n^{-\sigma-it}d\sigma=n^{-1/2-it}/\log n$, so the Dirichlet series gives $\pi\Phi_X^W(t)$ after taking imaginary parts; and $\int_{1/2}^\infty\mathfrak G_L(\rho-\sigma-it)\,d\sigma=\mathfrak I_L(\rho-\tfrac12-it)$ and $\int_{1/2}^\infty\mathfrak G_L(1-\sigma-it)\,d\sigma=\mathfrak I_L(\tfrac12-it)$. Thus
\begin{equation}\label{eq:taucw-piS}
\pi S(t)=\pi\Phi_X^W(t)+\Im\,\mathfrak I_L(\tfrac12-it)-\Im\sum_\rho\mathfrak I_L(\rho-\tfrac12-it)-\Im\sum_{m\ge1}\mathfrak I_L(-2m-\tfrac12-it).
\end{equation}
By \cref{lem:taucw-kernel}(d) the trivial-zero sum is $O(1/t^2)$. For the pole term write $\mathfrak G_L(1-s)=\frac1{s-1}-\int_1^{X^2}W(\tfrac{\log u}L)u^{-s}du$, which holds for $\Re s>1$ and, both sides being meromorphic, everywhere. Integrating over $\sigma\in[\tfrac12,B]$, using $W=1$ on $[1,2]$ and letting $B\to\infty$, one finds $\Im\,\mathfrak I_L(\tfrac12-it)=-\pi\Pi_X^W(t)+\mathscr A(t)$, where $\mathscr A$ does not depend on $X$: with $z=\tfrac12-it$ and $B^\star=B-\tfrac12$ the contribution of $u\in[1,2]$ is $\int_0^{\log2}(e^{zv}-e^{(z-B^\star)v})v^{-1}dv=-\mathrm{Ein}(-z\log2)+\log\bigl((B^\star-z)\log2\bigr)+\gamma_E+o(1)$ with $\mathrm{Ein}(x)=\int_0^x(1-e^{-s})s^{-1}ds=E_1(x)+\log x+\gamma_E$, and it combines with $\log\frac{B-1+it}{-1/2+it}$ to $\mathscr A(t)=\Im\,E_1\bigl((-\tfrac12+it)\log2\bigr)$. From $E_1(w)=e^{-w}\int_0^\infty e^{-s}(w+s)^{-1}ds$ and $|w+s|\ge\Im w$ one gets $|\mathscr A(t)|\le\sqrt2/(t\log2)$. Finally $N=N_0+S+R$ with $R=O(1/t)$, and $P_X^W=\Phi_X^W-\Pi_X^W$, which gives \eqref{eq:taucw-expl}.
\end{proof}

Write $y_\rho=\beta-\frac12$; the real part of $\rho-\frac12-it$ is $y_\rho$ and the imaginary part is $\gamma-t$. Split the sum in \eqref{eq:taucw-expl} into the contribution of the same zeros placed on the critical line and the remainder:
\begin{equation}\label{eq:taucw-on-off}
\mathcal O_L(t)=\sum_\rho\bigl(\mathfrak H_L(t-\gamma)-\mathbf 1_{\gamma\le t}\bigr),\qquad
\mathcal Y_L(t)=\frac1\pi\Im\sum_{\rho}\bigl(\mathfrak I_L(y_\rho+i(\gamma-t))-\mathfrak I_L(i(\gamma-t))\bigr).
\end{equation}
By \cref{lem:taucw-kernel}(b),(c) the series $\mathcal O_L(t)$ converges absolutely and equals $\frac1\pi\Im\sum_\rho\mathfrak I_L(i(\gamma-t))$, using $\mathfrak H_L(-v)=1-\mathfrak H_L(v)$. It involves only the imaginary parts of the zeros, whereas $\mathcal Y_L$ vanishes term by term for zeros on the critical line.

\begin{proposition}[Decomposition]\label{prop:taucw-decomp}
Let $X=e^L\ge2$ and let $t\ge t_0$ not be the imaginary part of a zero. Then, uniformly in $X$,
\begin{align}
N_0(t)+P_X^W(t)-N(t)&=\mathcal O_L(t)+\mathcal Y_L(t)+O\Bigl(\frac1t\Bigr),\label{eq:taucw-decomp}\\
\mathcal O_L(t)&=\int_{|\delta|\le t/2}\bigl(N(t-\delta)-N(t)\bigr)\mathfrak K_L(\delta)\,d\delta+O\Bigl(\frac{\log t}{L}\Bigr).\label{eq:taucw-on}
\end{align}
\end{proposition}
\begin{proof}
\eqref{eq:taucw-decomp} is \eqref{eq:taucw-expl} with the splitting \eqref{eq:taucw-on-off}. For \eqref{eq:taucw-on}, the zeros with $|\gamma-t|>t/2$ contribute at most $\sum3/(\pi L^2(t-\gamma)^2)\ll\log t/(L^2t)$ by \cref{lem:taucw-kernel}(b) and $N(s)\ll s\log s$. For the zeros with $\gamma\in(t/2,3t/2]$ write $\mathfrak H_L(t-\gamma)-\mathbf 1_{\gamma\le t}=\int_{\mathbb R}\mathfrak K_L(\delta)\bigl(\mathbf 1_{\gamma\le t-\delta}-\mathbf 1_{\gamma\le t}\bigr)d\delta$ and sum over $\gamma$: for $|\delta|\le t/2$ the sum of the bracket is $N(t-\delta)-N(t)$, and the ranges $\delta<-t/2$ and $\delta>t/2$ contribute at most $\bigl(N(\tfrac{3t}2)-N(\tfrac t2)+N(t)\bigr)\int_{|\delta|>t/2}|\mathfrak K_L|\ll t\log t\cdot(Lt)^{-1}$.
\end{proof}

\subsection{The fourth moment}\label{sec:taucw-moment}
\begin{theorem}[Fourth moment of the on-line part]\label{thm:taucw-online}
Let $\theta_0>0$. There are $C=C(\theta_0)$ and $T_0=T_0(\theta_0)$ such that for $T\ge T_0$ and every $\theta\ge\theta_0$,
\begin{equation}\label{eq:taucw-L4}
\int_T^{2T}\bigl|\mathcal O_{L(t)}(t)\bigr|^4\dd\le C\,T,\qquad L(t)=\theta\log t.
\end{equation}
The same bound holds for a fixed $L\ge\theta_0\log T$ in place of $L(t)$. No hypothesis on the zeros is used; the proof uses \eqref{eq:FujiiB} and the fourth moment of $S$ \cite{Selberg1946,Titchmarsh}.
\end{theorem}
\begin{proof}
Let $T\le t\le2T$. By \eqref{eq:taucw-on} with $L=L(t)\ge\theta_0\log T$, the error $O(\log t/L)$ is $O(1/\theta_0)$. On $|\delta|\le t/2$ write $N(t-\delta)-N(t)=[N_0(t-\delta)-N_0(t)]+[S(t-\delta)-S(t)]+O(1/t)$ by \eqref{eq:RVM}.

\emph{The term $N_0$.} Taylor's formula gives $N_0(t-\delta)-N_0(t)=-N_0'(t)\delta+\tfrac12N_0''(\xi)\delta^2$ with $|N_0''|\le1/(\pi t)$ on $[t/2,3t/2]$. The linear term integrates to zero against the even kernel on the symmetric range $|\delta|\le t/2$. Since $\delta^2|\mathfrak K_L(\delta)|\le2/(\pi L)$, the quadratic term is at most $\frac12\cdot\frac1{\pi t}\cdot\frac2{\pi L}\cdot t=\frac1{\pi^2L}$.

\emph{The term $S$.} For $\delta$ with $|\delta|\le t/2\le T$, Minkowski's inequality gives
\begin{equation}\label{eq:taucw-mink}
\Bigl\|\int_{|\delta|\le t/2}\bigl(S(t-\delta)-S(t)\bigr)\mathfrak K_{L(t)}(\delta)\,d\delta\Bigr\|_{L^4[T,2T]}\le\int_{|\delta|\le T}\sup_t|\mathfrak K_{L(t)}(\delta)|\;\bigl\|S(\cdot-\delta)-S\bigr\|_{L^4(E_\delta)}d\delta,
\end{equation}
where $E_\delta=\{t\in[T,2T]:|\delta|\le t/2\}$. Because $L(t)\in[\theta\log T,\theta\log 2T]$ and $\theta\log2T\le2\theta\log T$ for large $T$, \cref{lem:taucw-kernel}(a) gives $\sup_t|\mathfrak K_{L(t)}(\delta)|\le2L\min\{1,(L\delta)^{-2}\}$ with $L=\theta\log T$. For $|\delta|\le1$, \eqref{eq:FujiiB} with $\mathsf k=2$, $\alpha_{\rm F}=\tfrac12$ and Fujii's length parameter equal to $T$, applied on the bounded number of intervals covering $[T/2,3T]$, gives $\|S(\cdot-\delta)-S\|_{L^4}\ll T^{1/4}(\log(2+|\delta|\log T))^{1/2}$. Substituting $u=L\delta$ and using $\log T/L=1/\theta\le1/\theta_0$, the part $|\delta|\le1$ of the integral is at most
\begin{equation}\label{eq:taucw-fujii}
C_1\,T^{1/4}\int_0^\infty\min\{1,u^{-2}\}\Bigl(\log\bigl(2+u/\theta_0\bigr)\Bigr)^{1/2}du,
\end{equation}
which is finite and depends on $\theta_0$ only. For $1<|\delta|\le T$ one uses $\|S(\cdot-\delta)-S\|_{L^4}\le\|S(\cdot-\delta)\|_{L^4}+\|S\|_{L^4}\ll T^{1/4}(\log\log T)^{1/2}$ (Selberg's fourth moment on $O(1)$ dyadic intervals, $\int_{T/2}^{3T}|S|^4\,dt\ll T(\log\log T)^2$), and $\int_{|\delta|>1}2L\,(L\delta)^{-2}d\delta\le4/L$; this part is $O\bigl(T^{1/4}(\log\log T)^{1/2}/\log T\bigr)$. The remainders $O(1/t)$ and $O(1/\theta_0)$ are bounded. Combining the $L^4$ norms proves \eqref{eq:taucw-L4}; the case of a fixed $L$ is the same with $L(t)\equiv L$.
\end{proof}


\begin{remark}[Where the off-line zeros enter]\label{rem:taucw-off}
Everything in $N_0+P^W_X-N$ that \cref{thm:taucw-online} does not control is contained in $\mathcal Y_L$. A zero $\rho=\frac12+y+i\gamma$ with $y>0$ contributes to $\mathcal Y_L(t)$ at most $\frac1\pi\int_0^y|\mathfrak G_L(x+i(\gamma-t))|dx\le(X^{2y}-1)/(\pi L^2(\gamma-t)^2)$ for $\gamma\ne t$, since $|\mathfrak G_L(x+iv)|\le2e^{2Lx}/(Lv^2)$ for $x\ge0$, and the reflected zero $\frac12-y+i\gamma$ at most $2/(\pi L^2(\gamma-t)^2)$. The location of an off-line zero thus enters only through the explicit factor $X^{2y}$. For $X$ in the range $T^\varepsilon\le X\le T^{1/4}$, uniformly in $X$, the fourth moment of the truncated $F_X^W-N$ is bounded by \cref{thm:moderate} (subject to the Fujii input \eqref{eq:FujiiA}); outside that range the control of $N_0+P^W_X-N$ is the control of $\mathcal Y_L$. The kernels are rescalings of one function: $\mathfrak I_L(w)=\mathfrak I_1(Lw)$, so that $N_0+P^W_X-N=\frac1\pi\Im\sum_\rho\mathfrak I_1\bigl(L(\rho-\frac12-it)\bigr)+O(1/t)$ with $\mathfrak I_1(z)=\int_{-\infty}^z(e^v-e^{2v})v^{-2}dv$.
\end{remark}

\begin{remark}[Comparison with $\tau^{\mathrm C}$]\label{rem:taucw-tauc}
A bound for taucw deduced from a bound for $\tau^{\mathrm C}$ and a bound for the difference of the two phases, by the triangle inequality, keeps the order of the larger of the two terms and discards the cancellation between the error of $\tau^{\mathrm C}$ and the additional arithmetic correction of taucw. It is a safety bound. The accuracy of taucw is read from $F^{\mathrm w}_\theta$ directly, as in \cref{thm:taucw-online}.
\end{remark}

\subsection{Existence, order and coarse localization of the taucw points}\label{sec:taucw-exist}
Throughout this and the following subsections $0<\theta\le1$, $\lambda_k=k-\frac12$, and $\tau_k^{\mathrm w}(\theta)$ denotes the point of \cref{def:taucw}, so that $\tau_k^{\mathrm w}=\tau_k^{\mathrm w}(1)$ for the cutoff $X=t$. For brevity put $\mathbf w_{\theta,t}(n)=W\bigl(\frac{\log n}{\theta\log t}\bigr)$, the weight in \eqref{eq:taucw-phase}, and extend it to real $n$ by the same formula. The counting functions are
\begin{equation}\label{eq:taucw-count}
N^{\mathrm w}_\theta(t)=\#\{k\ge1:\ \tau_k^{\mathrm w}(\theta)\le t\},\qquad N^{\mathrm w}=N^{\mathrm w}_1,
\end{equation}
and $S^{\mathrm w}(t)=N^{\mathrm w}(t)-N_0(t)-R(t)$ with $R$ of \eqref{eq:RVM}. The synthetic points $\frac12+i\tau_k^{\mathrm w}$ are defined by this formula only; $p^{1/2+i\tau_k^{\mathrm w}}=\sqrt p\,e^{i\tau_k^{\mathrm w}\log p}$ is an algebraic identity.

\begin{lemma}[Amplitude of the weighted phase]\label{lem:taucw-amp}
For $t\ge3$ and fixed $0<\theta\le1$,
\begin{equation}\label{eq:taucw-amp}
\bigl|\Phi^{\mathrm w}_\theta(t)\bigr|\ll_\theta1+\frac{t^\theta}{\log^2t}+\frac{t^{\theta/2}}{\log t}.
\end{equation}
In particular $|\Phi^{\mathrm w}_1(t)|\ll t/\log^2t$, and $F^{\mathrm w}_\theta(t)\to+\infty$.
\end{lemma}
\begin{proof}
Put $X=t^\theta$. The primes $p\le X$ carry weight $1$, and by partial summation with $\pi(x)\ll x/\log x$ their contribution is at most $\frac1\pi\sum_{p\le X}p^{-1/2}\ll1+X^{1/2}/\log X$. For $X<p\le X^2$ the weight is $\mathbf w_{\theta,t}(p)=(2\log X-\log p)/\log X$, decreasing from $1$ to $0$, and partial summation gives
\begin{equation}
\sum_{X<p\le X^2}\frac{\mathbf w_{\theta,t}(p)}{\sqrt p}\le\int_X^{X^2}\pi(x)\Bigl(\frac{\mathbf w_{\theta,t}(x)}{2x^{3/2}}+\frac{1}{x^{3/2}\log X}\Bigr)\mathrm dx
\ll\int_X^{X^2}\Bigl(\mathbf w_{\theta,t}(x)+\frac1{\log X}\Bigr)\frac{\mathrm dx}{\sqrt x\,\log x}.
\end{equation}
The part $x\le X^{3/2}$ is $\ll X^{3/4}/\log X$. On $x>X^{3/2}$ write $x=X^2e^{-s}$, so that $\mathbf w_{\theta,t}=s/\log X$, $\mathrm dx/\sqrt x=Xe^{-s/2}\mathrm ds$ and $\log x\ge\frac32\log X$; that part is $\ll X(\log X)^{-2}\int_0^\infty(1+s)e^{-s/2}\,\mathrm ds\ll X/\log^2X$. Since $X^{3/4}/\log X\ll X/\log^2X$ for large $X$, the primes in $(X,X^2]$ contribute $\ll X/\log^2X$. The prime squares $p\le X$ contribute at most $\frac1{2\pi}\sum_{p\le X}p^{-1}\ll\log\log(3t)$, and the powers with exponent at least $3$ at most $\frac1\pi\sum_p\sum_{r\ge3}(rp^{r/2})^{-1}<\infty$. The term $\log\log(3t)$ is absorbed in $1+t^{\theta/2}/\log t$. Finally $N_0(t)\asymp t\log t$.
\end{proof}

Let $\mathfrak A_\theta(t)=C_\theta\bigl(1+t^\theta/\log^2t+t^{\theta/2}/\log t\bigr)$ with $C_\theta$ so large that $|\Phi^{\mathrm w}_\theta|\le\mathfrak A_\theta$ for $t\ge3$. For $t\ge t_0(\theta)$ the function $\mathfrak A_\theta$ is increasing and $\mathfrak A_\theta'(t)=o(N_0'(t))$, so that $N_0\pm\mathfrak A_\theta$ are increasing on $[t_0,\infty)$.

\begin{proposition}[Ordered points]\label{prop:taucw-ord}
For every $0<\theta\le1$ and $k\ge1$ the measure in \eqref{eq:taucw-F} is finite and $k\mapsto\tau^{\mathrm w}_k(\theta)$ is nondecreasing. It is strictly increasing on the levels $\lambda_k>F^{\mathrm w}_\theta(2\pi)$. For $k\ge k_0(\theta)$, with $\mathsf T_k^-<\mathsf T_k^+$ defined by $(N_0+\mathfrak A_\theta)(\mathsf T_k^-)=\lambda_k=(N_0-\mathfrak A_\theta)(\mathsf T_k^+)$,
\begin{equation}\label{eq:taucw-ord}
\mathsf T_k^-\le\tau_k^{\mathrm w}(\theta)\le\mathsf T_k^+,\qquad
N^{\mathrm w}_\theta(t)=N_0(t)+O_\theta\Bigl(1+\frac{t^\theta}{\log^2t}+\frac{t^{\theta/2}}{\log t}\Bigr).
\end{equation}
In particular $N^{\mathrm w}(t)-N_0(t)\ll t/\log^2t$.
\end{proposition}
\begin{proof}
Each summand of \eqref{eq:taucw-phase} vanishes where it enters or leaves its moving support and only finitely many summands are active on a compact interval, so $F^{\mathrm w}_\theta$ is continuous on $[2\pi,\infty)$. By \cref{lem:taucw-amp}, $F^{\mathrm w}_\theta\ge N_0-\mathfrak A_\theta\to\infty$; the sublevel sets are therefore bounded and nested in $k$, which gives finiteness and monotonicity. If $F^{\mathrm w}_\theta(2\pi)<\lambda_k<\lambda_{k+1}$, the intermediate value theorem gives a nonempty open set on which $F^{\mathrm w}_\theta\in(\lambda_k,\lambda_{k+1})$; it has positive measure, so $\tau_k^{\mathrm w}<\tau_{k+1}^{\mathrm w}$.
Choose $k_0$ with $\lambda_{k_0}>\max\{(N_0+\mathfrak A_\theta)(t_0),\max_{[2\pi,t_0]}F^{\mathrm w}_\theta\}$. For $k\ge k_0$ the sublevel set contains $[2\pi,t_0]$, and on $[t_0,\infty)$ it contains $\{N_0+\mathfrak A_\theta<\lambda_k\}=[t_0,\mathsf T_k^-)$ and is contained in $\{N_0-\mathfrak A_\theta<\lambda_k\}=[t_0,\mathsf T_k^+)$. Taking measures gives $\mathsf T_k^-\le\tau^{\mathrm w}_k(\theta)\le\mathsf T_k^+$. Hence for $t\ge t_0$, by monotonicity of $N_0\pm\mathfrak A_\theta$, $\mathsf T_k^+\le t$ iff $k\le(N_0-\mathfrak A_\theta)(t)+\frac12$, and likewise for $\mathsf T_k^-$, so that $\lfloor(N_0-\mathfrak A_\theta)(t)+\frac12\rfloor-k_0\le N^{\mathrm w}_\theta(t)-(k_0-1)\le\lfloor(N_0+\mathfrak A_\theta)(t)+\frac12\rfloor-k_0+1$; this is \eqref{eq:taucw-ord}.
\end{proof}

\begin{proposition}[Coarse displacement from the Riemann--von Mangoldt points]\label{prop:taucw-coarse}
Let $\gamma_k^{\mathrm{RvM}}$ be the solution of $N_0(\gamma_k^{\mathrm{RvM}})=k-\frac12$ with $\gamma_k^{\mathrm{RvM}}>2\pi$. As $k\to\infty$,
\begin{equation}\label{eq:taucw-coarse}
\bigl|\tau_k^{\mathrm w}-\gamma_k^{\mathrm{RvM}}\bigr|\ll\frac{\gamma_k^{\mathrm{RvM}}}{\log^3\gamma_k^{\mathrm{RvM}}}.
\end{equation}
\end{proposition}
\begin{proof}
Let $\theta=1$. Since $(N_0+\mathfrak A_1)(\gamma_k^{\mathrm{RvM}})>\lambda_k>(N_0-\mathfrak A_1)(\gamma_k^{\mathrm{RvM}})$ and both functions increase, $\mathsf T_k^-<\gamma_k^{\mathrm{RvM}}<\mathsf T_k^+$. Then $N_0(\mathsf T_k^+)-N_0(\mathsf T_k^-)=\mathfrak A_1(\mathsf T_k^-)+\mathfrak A_1(\mathsf T_k^+)\ll\gamma_k^{\mathrm{RvM}}/\log^2\gamma_k^{\mathrm{RvM}}$ because $\mathfrak A_1(t)\ll t/\log^2t$ and $\mathsf T_k^\pm\asymp\gamma_k^{\mathrm{RvM}}$, while $N_0'\asymp\log\gamma_k^{\mathrm{RvM}}$ on $[\mathsf T_k^-,\mathsf T_k^+]$. Hence $\mathsf T_k^+-\mathsf T_k^-\ll\gamma_k^{\mathrm{RvM}}/\log^3\gamma_k^{\mathrm{RvM}}$, and \eqref{eq:taucw-ord} gives \eqref{eq:taucw-coarse}.
\end{proof}

\subsection{Sublevel transport and contraction}\label{sec:taucw-transport}
For a measurable function $f$ on a finite interval $I$ put
\begin{equation}\label{eq:taucw-vol}
\vol_{f,I}(v)=\meas\{t\in I:\ f(t)<v\},\qquad v\in\mathbb R,
\end{equation}
the sublevel distribution of $f$ on $I$.

\begin{lemma}[Level crossing count]\label{lem:taucw-cross}
Let $y\in\mathbb R$ and let $n\ge0$ be an integer. Then
\begin{equation}\label{eq:taucw-cross}
\sum_{k\ge1}\bigl|\ind_{\{y<\lambda_k\}}-\ind_{\{n<\lambda_k\}}\bigr|=\bigl|\max\{0,\lfloor y+\tfrac12\rfloor\}-n\bigr|\le\bigl|\lfloor y+\tfrac12\rfloor-n\bigr|.
\end{equation}
\end{lemma}
\begin{proof}
For every real $x$, $\ind_{\{x<\lambda_k\}}=1-\ind_{\{\lambda_k\le x\}}$, and $\{k\ge1:\lambda_k\le x\}=\{1,\dots,\max\{0,\lfloor x+\frac12\rfloor\}\}$ because $k-\frac12\le x$ iff $k\le x+\frac12$. For $x=n$ this initial segment has $n$ elements. The left side of \eqref{eq:taucw-cross} is the number of elements of the symmetric difference of two initial segments, that is, the difference of their lengths. The inequality holds because $|\max\{0,m\}-n|\le|m-n|$ for $n\ge0$.
\end{proof}

\begin{theorem}[Absolute transport]\label{thm:taucw-L1}
Let $0<\theta\le1$, $F=F^{\mathrm w}_\theta$, and let $\mathcal U\subset\mathbb N$ be finite. Let $I\subset[2\pi,\infty)$ be a finite interval such that, for every $k\in\mathcal U$, the two sets $\{t\ge2\pi:F(t)<\lambda_k\}$ and $\{t\ge2\pi:N(t)<\lambda_k\}$ coincide outside $I$ up to a null set. Then
\begin{equation}\label{eq:taucw-L1}
\sum_{k\in\mathcal U}\bigl|\tau_k^{\mathrm w}(\theta)-\gamma_k\bigr|\le\int_Iq_F(t)\dd,\qquad q_F=\bigl|\lfloor F+\tfrac12\rfloor-N\bigr|.
\end{equation}
The same holds for the sublevel points, measured from $2\pi$, of any measurable function with finite sublevel measures.
\end{theorem}
\begin{proof}
Since $N=0$ on $[0,\gamma_1)$ and $\gamma_1>2\pi$, $\{t\ge2\pi:N(t)<\lambda_k\}=[2\pi,\gamma_k)$ up to a null set, so that $\gamma_k=2\pi+\meas\{t\ge2\pi:N(t)<\lambda_k\}$. Hence
\begin{equation}
\tau_k^{\mathrm w}(\theta)-\gamma_k=\int_I\bigl(\ind_{\{F(t)<\lambda_k\}}-\ind_{\{N(t)<\lambda_k\}}\bigr)\dd.
\end{equation}
Take absolute values inside, sum over $k\in\mathcal U$, enlarge the sum to all $k\ge1$ and apply \cref{lem:taucw-cross} with $y=F(t)$ and $n=N(t)$. Neither a derivative of $F$ nor the simplicity of any crossing is used.
\end{proof}

\begin{theorem}[$L^1$ contraction of sublevel distributions]\label{thm:taucw-contr}
Let $I$ be a finite interval and $f_1,f_2\in L^1(I)$ real valued. Then
\begin{align}
\int_{\mathbb R}\bigl|\vol_{f_2,I}(v)-\vol_{f_1,I}(v)\bigr|\,\mathrm dv&\le\int_I|f_2(t)-f_1(t)|\dd,\label{eq:taucw-contr}\\
\int_{\mathbb R}\bigl(\vol_{f_2,I}(v)-\vol_{f_1,I}(v)\bigr)\,\mathrm dv&=\int_I\bigl(f_1(t)-f_2(t)\bigr)\dd.\label{eq:taucw-signed}
\end{align}
The left side of \eqref{eq:taucw-signed} converges absolutely.
\end{theorem}
\begin{proof}
For real $x,y$ the function $v\mapsto\ind_{\{x<v\}}-\ind_{\{y<v\}}$ vanishes outside the interval with endpoints $x,y$, and
\begin{equation}
\int_{\mathbb R}\bigl|\ind_{\{x<v\}}-\ind_{\{y<v\}}\bigr|\,\mathrm dv=|x-y|,\qquad
\int_{\mathbb R}\bigl(\ind_{\{x<v\}}-\ind_{\{y<v\}}\bigr)\,\mathrm dv=y-x.
\end{equation}
Put $x=f_2(t)$, $y=f_1(t)$, integrate over $t\in I$, and use Tonelli's theorem together with the triangle inequality for \eqref{eq:taucw-contr}, and Fubini's theorem for \eqref{eq:taucw-signed}; the latter interchange is justified by $f_1,f_2\in L^1(I)$.
\end{proof}

\begin{corollary}[Mean-square form]\label{cor:taucw-L2}
For $f_1,f_2\in L^2(I)$, $\ \int_{\mathbb R}|\vol_{f_2,I}-\vol_{f_1,I}|\,\mathrm dv\le|I|^{1/2}\bigl(\int_I|f_2-f_1|^2\dd\bigr)^{1/2}$.
\end{corollary}
\begin{proof}
\Cref{thm:taucw-contr} and the Cauchy--Schwarz inequality.
\end{proof}
Both bounds compare sublevel distributions at all real levels simultaneously. For two phases $F^{\mathrm w}_{\theta_-}$ and $F^{\mathrm w}_{\theta_+}$ the common term $N_0$ cancels in $f_2-f_1$, which is how \cref{thm:taucw-contr} is used in the next subsection.

\subsection{The cutoff parameter: differentiation, variance and increments}\label{sec:taucw-theta}
\begin{lemma}[Derivative in the cutoff exponent]\label{lem:taucw-theta}
For fixed $t\ge2\pi$ and $0<\theta\le1$, away from the finitely many values of $\theta$ at which a prime power equals $t^\theta$ or $t^{2\theta}$,
\begin{equation}\label{eq:taucw-theta}
\frac{\partial F^{\mathrm w}_\theta(t)}{\partial\theta}=-\frac1{\pi\theta^2\log t}\sum_{t^\theta<n<t^{2\theta}}\frac{\Lambda(n)}{\sqrt n}\,\sin(t\log n).
\end{equation}
The map $\theta\mapsto F^{\mathrm w}_\theta(t)$ is locally Lipschitz on $(0,1]$ and \eqref{eq:taucw-theta} holds almost everywhere.
\end{lemma}
\begin{proof}
On $t^\theta<n\le t^{2\theta}$ the weight is $2-\log n/(\theta\log t)$, with $\theta$-derivative $\log n/(\theta^2\log t)$, and on $n\le t^\theta$ it is constant. The weight is continuous at both break points, so no point mass appears, and for $\theta$ in a compact subinterval of $(0,1]$ only finitely many terms are active.
\end{proof}

For $T\ge3$ and $0<\theta<\frac12$ define the phase with weights frozen at height $T$,
\begin{equation}\label{eq:taucw-frozen}
\Phi^{\mathrm{fr}}_{\theta,T}(t)=-\frac1\pi\sum_{2\le n\le(2T)^{2\theta}}\frac{\Lambda(n)\,W\!\left(\frac{\log n}{\theta\log T}\right)}{\sqrt n\,\log n}\sin(t\log n).
\end{equation}

\begin{lemma}[Frozen weights on a dyadic block]\label{lem:taucw-freeze}
For fixed $0<\theta<\frac12$,
\begin{equation}\label{eq:taucw-freeze}
\frac1T\int_T^{2T}\bigl|\Phi^{\mathrm w}_\theta(t)-\Phi^{\mathrm{fr}}_{\theta,T}(t)\bigr|^2\dd\ll_\theta\frac{1+\log\log T}{\log^2T}.
\end{equation}
\end{lemma}
\begin{proof}
Put $\mathfrak k_n=\Lambda(n)/(\pi\sqrt n\log n)$ and $\mathfrak L=(2T)^{2\theta}$. For $T\le s\le2T$ and $n\le\mathfrak L$ the function $s\mapsto W(\log n/(\theta\log s))$ is Lipschitz with derivative $\mathbf d_n(s)$ of modulus at most $\log n/(\theta s\log^2s)\ll_\theta1/(T\log T)$, since $\log n\le2\theta\log(2T)$. Writing the weight difference as the integral of this derivative over $[T,t]$, Minkowski's integral inequality gives
\begin{equation}
\bigl\|\Phi^{\mathrm w}_\theta-\Phi^{\mathrm{fr}}_{\theta,T}\bigr\|_{L^2(T,2T)}\le\int_T^{2T}\Bigl\|\mathbf 1_{(s,2T]}(t)\sum_{n\le\mathfrak L}\mathfrak k_n\mathbf d_n(s)\sin(t\log n)\Bigr\|_{L^2(\mathrm dt)}\mathrm ds,\qquad|\mathbf d_n(s)|\ll_\theta\frac1{T\log T}.
\end{equation}
By the mean value theorem of Montgomery and Vaughan \cite{MontgomeryVaughan}, for coefficients $\iota_n$ and an interval of length $H\le T$, $\int|\sum_{n\le\mathfrak L}\iota_nn^{it}|^2\mathrm dt=\sum_{n\le\mathfrak L}|\iota_n|^2(H+O(n))$, and $\mathfrak L\ll T$ here. Hence the inner norm is $\ll_\theta T^{1/2}(T\log T)^{-1}(\sum_{n\le\mathfrak L}\mathfrak k_n^2)^{1/2}$, and
$\sum_{n\le\mathfrak L}\mathfrak k_n^2\ll1+\sum_{p\le\mathfrak L}p^{-1}\ll1+\log\log T$ by Mertens' theorem and $\sum_p\sum_{r\ge2}(r^2p^r)^{-1}<\infty$. Integrating over $s$ and squaring gives \eqref{eq:taucw-freeze}.
\end{proof}

\begin{theorem}[Variance of a fixed moderate cutoff]\label{thm:taucw-var}
For every fixed $0<\theta<\frac12$,
\begin{equation}\label{eq:taucw-var}
\frac1T\int_T^{2T}\bigl|\Phi^{\mathrm w}_\theta(t)\bigr|^2\dd=\frac1{2\pi^2}\log\log T+O_\theta(1).
\end{equation}
\end{theorem}
\begin{proof}
Let $\iota_n=\mathfrak k_nW(\log n/(\theta\log T))$ for $n\le\mathfrak L=(2T)^{2\theta}$ and $\mathsf X(t)=\sum_n\iota_nn^{it}$, so that $\Phi^{\mathrm{fr}}_{\theta,T}=-\Im\mathsf X$ and $|\Im\mathsf X|^2=\frac12|\mathsf X|^2-\frac12\Re\mathsf X^2$. By the mean value theorem recalled above, $\int_T^{2T}|\mathsf X|^2\dd=T\sum_n\iota_n^2+O(\sum_nn\iota_n^2)$ and $\sum_nn\iota_n^2\ll\mathfrak L=o(T)$ because $\Lambda(n)\le\log n$. In $\int_T^{2T}\mathsf X^2\dd=\sum_{n,m}\iota_n\iota_m\int_T^{2T}(nm)^{it}\dd$ every integral is $O(1/\log(nm))$, so that this term is $\ll(\sum_n\iota_n)^2\ll\mathfrak L/\log^2\mathfrak L=o(T)$; here $\theta<\frac12$ is used. Hence $\int_T^{2T}|\Phi^{\mathrm{fr}}_{\theta,T}|^2\dd=\frac T2\sum_n\iota_n^2+o(T)$. The primes carry $\sum_pp^{-1}W(\log p/(\theta\log T))^2=\log\log T+O_\theta(1)$ by Mertens' theorem, since $W=1$ for $p\le T^\theta$ and $\sum_{T^\theta<p\le T^{2\theta}}p^{-1}=\log2+o(1)$, and the prime powers with exponent at least $2$ give $O(1)$. By \cref{lem:taucw-freeze} and the Cauchy--Schwarz inequality, replacing $\Phi^{\mathrm{fr}}_{\theta,T}$ by $\Phi^{\mathrm w}_\theta$ changes the norm in $L^2(T,2T)$ by $O(\sqrt{T}\sqrt{1+\log\log T}/\log T)$, which changes the square by $o(T)$.
\end{proof}

\begin{theorem}[Variance of an increment in the cutoff]\label{thm:taucw-inc}
For fixed $0<\theta_-<\theta_+<\frac12$ let
\begin{equation}\label{eq:taucw-S}
\mathfrak S(\theta_-,\theta_+)=\int_0^\infty\Bigl[W\!\Bigl(\frac v{\theta_+}\Bigr)-W\!\Bigl(\frac v{\theta_-}\Bigr)\Bigr]^2\frac{\mathrm dv}v ,
\end{equation}
a finite positive number. Then
\begin{equation}\label{eq:taucw-inc}
\lim_{T\to\infty}\frac1T\int_T^{2T}\bigl|\Phi^{\mathrm w}_{\theta_+}(t)-\Phi^{\mathrm w}_{\theta_-}(t)\bigr|^2\dd=\frac{\mathfrak S(\theta_-,\theta_+)}{2\pi^2}.
\end{equation}
For $\theta_+=\theta(1+\varepsilon)<\frac12$ with $\theta>0$ fixed and $\varepsilon\to0$,
\begin{equation}\label{eq:taucw-inc-small}
\frac{\mathfrak S(\theta,\theta(1+\varepsilon))}{2\pi^2}=\frac3{4\pi^2}\varepsilon^2+o(\varepsilon^2).
\end{equation}
\end{theorem}
\begin{proof}
The integrand of \eqref{eq:taucw-S} vanishes for $v\le\theta_-$ and $v\ge2\theta_+$, so $\mathfrak S<\infty$, and it is positive on $(\theta_-,\theta_+)$, where the first weight is $1$ and the second is smaller than $1$. By \cref{lem:taucw-freeze} applied to both cutoffs, the integral in \eqref{eq:taucw-inc} equals that of the difference $\Phi^{\mathrm{fr}}_{\theta_+,T}-\Phi^{\mathrm{fr}}_{\theta_-,T}$ up to $o(T)$, the latter being $O(T)$. As in the proof of \cref{thm:taucw-var}, with $\mathfrak L=(2T)^{2\theta_+}=o(T)$, this integral equals $\frac T2\sum_n\iota_n^2+o(T)$ for the coefficients $\iota_n=\mathfrak k_n[W(\frac{\log n}{\theta_+\log T})-W(\frac{\log n}{\theta_-\log T})]$. The prime powers with exponent $r\ge2$ contribute $o(1)$ by dominated convergence: for fixed $p^r$ both weights equal $1$ once $\log T\ge\log(p^r)/\theta_-$, and the terms are dominated by $(r^2p^r)^{-1}$. For primes, Mertens' theorem and partial summation turn $\sum_pp^{-1}[\cdots]^2$ into the Riemann integral \eqref{eq:taucw-S}, with $v=\log p/\log T$; this proves \eqref{eq:taucw-inc}.
For \eqref{eq:taucw-inc-small}, $W$ is $1$-Lipschitz, so $|W(\frac v{\theta(1+\varepsilon)})-W(\frac v\theta)|\le\varepsilon v/\theta$, supported in $\theta\le v\le4\theta$, and for $v\notin\{\theta,2\theta\}$ the quotient by $\varepsilon$ tends to $\frac v\theta\mathbf 1_{(\theta,2\theta)}(v)$. Dominated convergence in $L^2(\mathrm dv/v)$ gives $\varepsilon^{-2}\mathfrak S(\theta,\theta(1+\varepsilon))\to\int_\theta^{2\theta}(v/\theta)^2\,\mathrm dv/v=\frac32$.
\end{proof}

\begin{corollary}[Absolute transport between moderate cutoffs]\label{cor:taucw-mod}
For fixed $0<\theta_-<\theta_+<\frac12$,
\begin{equation}\label{eq:taucw-mod}
\int_{\mathbb R}\bigl|\vol_{F^{\mathrm w}_{\theta_+},[T,2T]}(v)-\vol_{F^{\mathrm w}_{\theta_-},[T,2T]}(v)\bigr|\,\mathrm dv\le T\Bigl(\frac{\mathfrak S(\theta_-,\theta_+)}{2\pi^2}+o(1)\Bigr)^{1/2}.
\end{equation}
\end{corollary}
\begin{proof}
Apply \cref{cor:taucw-L2} to $f_i=F^{\mathrm w}_{\theta_i}$, for which $f_2-f_1=\Phi^{\mathrm w}_{\theta_+}-\Phi^{\mathrm w}_{\theta_-}$, and insert \eqref{eq:taucw-inc}.
\end{proof}

\begin{theorem}[Signed transport up to the full cutoff]\label{thm:taucw-signed}
For fixed $0<\theta_-<\theta_+\le1$ and $I_T=[T,2T]$,
\begin{equation}\label{eq:taucw-signed-main}
\int_{\mathbb R}\bigl(\vol_{F^{\mathrm w}_{\theta_+},I_T}(v)-\vol_{F^{\mathrm w}_{\theta_-},I_T}(v)\bigr)\,\mathrm dv=O_{\theta_+}\Bigl(\frac{T^{\theta_+}}{\log^2T}\Bigr).
\end{equation}
\end{theorem}
\begin{proof}
By \eqref{eq:taucw-signed} the left side equals $\int_T^{2T}(\Phi^{\mathrm w}_{\theta_-}-\Phi^{\mathrm w}_{\theta_+})\dd$. For fixed $n$ the weight $t\mapsto W(\log n/(\theta\log t))$ is nondecreasing with values in $[0,1]$, so $\mathscr K_n(t)=W(\frac{\log n}{\theta_+\log t})-W(\frac{\log n}{\theta_-\log t})$ satisfies $|\mathscr K_n|\le1$ and has total variation at most $2$ on $[T,2T]$, uniformly in $n$. One integration by parts against $\sin(t\log n)$ gives $|\int_T^{2T}\mathscr K_n(t)\sin(t\log n)\dd|\le4/\log n$. Hence the left side of \eqref{eq:taucw-signed-main} is at most $\frac4\pi\sum_{2\le n\le(2T)^{2\theta_+}}\Lambda(n)/(\sqrt n\log^2n)$. The primes give $\ll\sqrt{\mathfrak L}/\log^2\mathfrak L$ with $\mathfrak L=(2T)^{2\theta_+}$, by partial summation and the prime number theorem, and the prime powers with exponent at least $2$ give $\sum_p\sum_{r\ge2}(r^2p^{r/2}\log p)^{-1}<\infty$.
\end{proof}

\begin{proposition}[Half-integer sampling and the sawtooth]\label{prop:taucw-saw}
For $0<\theta_-,\theta_+\le1$ and $I=[T,2T]$, with $\mathrm{saw}(y)=y-\lfloor y+\frac12\rfloor$ as in \cref{lem:v9-rearr},
\begin{equation}\label{eq:taucw-saw}
\sum_{k\in\mathbb Z}\Bigl[\vol_{F^{\mathrm w}_{\theta_+},I}(k-\tfrac12)-\vol_{F^{\mathrm w}_{\theta_-},I}(k-\tfrac12)\Bigr]
=\int_I\bigl(F^{\mathrm w}_{\theta_-}-F^{\mathrm w}_{\theta_+}\bigr)\dd+\int_I\Bigl[\mathrm{saw}\bigl(F^{\mathrm w}_{\theta_+}\bigr)-\mathrm{saw}\bigl(F^{\mathrm w}_{\theta_-}\bigr)\Bigr]\dd.
\end{equation}
Only finitely many summands are nonzero, and the first integral on the right is $O(T^{\theta_*}/\log^2T)$ with $\theta_*=\max\{\theta_-,\theta_+\}$ and the implied constant depending on $\theta_*$ only, by \cref{thm:taucw-signed}.
\end{proposition}
\begin{proof}
For real $x,y$, the number of half-integers $\lambda$ with $x<\lambda\le y$ is $\lfloor y+\frac12\rfloor-\lfloor x+\frac12\rfloor$ (negative when $y<x$), and it equals $\sum_{k\in\mathbb Z}(\ind_{\{x<k-1/2\}}-\ind_{\{y<k-1/2\}})$. Put $x=F^{\mathrm w}_{\theta_+}(t)$ and $y=F^{\mathrm w}_{\theta_-}(t)$ and integrate over $I$. Both phases are bounded on $I$, so only a common finite set of levels is active and summation and integration commute. Finally $\lfloor y+\frac12\rfloor=y-\mathrm{saw}(y)$.
\end{proof}
The two sides of \eqref{eq:taucw-saw} are exact; the estimate of the first integral involves no information on any single point $\tau_k^{\mathrm w}(\theta)$, since the positive and negative contributions of the individual levels cancel in the sum.

\subsection{Prime powers and the long-prime tail}\label{sec:taucw-powers}
Let $T\ge3$ and let $\mathfrak a_T$ be the $\lfloor\sqrt T\rfloor$-th prime, so that $\mathfrak a_T\sim\frac12\sqrt T\log T<T$ for large $T$. For $t\in[T,2T]$ the weight of every prime $p\le\mathfrak a_T$ in \eqref{eq:taucw-phase} with $\theta=1$ equals $1$, and
\begin{equation}\label{eq:taucw-split}
\Phi^{\mathrm w}_1=\Phi^{\mathrm{sm}}_T+\Phi^{\mathrm{pw}}_T+\Phi^{\mathrm{lg}}_T,
\end{equation}
where, with $\mathbf w_{1,t}(n)=W(\log n/\log t)$,
\begin{align}
\Phi^{\mathrm{sm}}_T(t)&=-\frac1\pi\sum_{p\le\mathfrak a_T}\frac{\sin(t\log p)}{\sqrt p},\qquad
\Phi^{\mathrm{pw}}_T(t)=-\frac1\pi\sum_{\substack{r\ge2,\ p\ \mathrm{prime}\\p^r\le t^2}}\frac{\mathbf w_{1,t}(p^r)}{r\,p^{r/2}}\sin(rt\log p),\label{eq:taucw-sm}\\
\Phi^{\mathrm{lg}}_T(t)&=-\frac1\pi\sum_{\mathfrak a_T<p\le t^2}\frac{\mathbf w_{1,t}(p)}{\sqrt p}\sin(t\log p).\label{eq:taucw-lg}
\end{align}
Each prime power occurs in exactly one of the three sums, because $\Lambda(p^r)/(\sqrt{p^r}\log p^r)=1/(rp^{r/2})$.

\begin{theorem}[Fourth moment of all prime powers]\label{thm:taucw-pow}
With an absolute implied constant,
\begin{equation}\label{eq:taucw-pow}
\int_T^{2T}\bigl|\Phi^{\mathrm{pw}}_T(t)\bigr|^4\dd\ll T.
\end{equation}
\end{theorem}
\begin{proof}
The terms with $r\ge3$ are bounded by $\frac1\pi\sum_p\sum_{r\ge3}(rp^{r/2})^{-1}<\infty$. For $r=2$ let $\mathfrak Y=(T/2)^{1/4}$. If $p\le\mathfrak Y$ then $p^2\le(T/2)^{1/2}<t$, so $\mathbf w_{1,t}(p^2)=1$, and these terms form $\mathsf Y(t)=-\frac1{2\pi}\Im\mathsf F(t)$ with $\mathsf F(t)=\sum_{p\le\mathfrak Y}p^{-1}p^{2it}$. The function $\mathsf F^2=\sum_m\upsilon_mm^{it}$ is a Dirichlet polynomial in $m=(pq)^2\le\mathfrak Y^4=T/2$ with $|\upsilon_m|\le2/(pq)$, hence $\sum_m|\upsilon_m|^2\le4(\sum_pp^{-2})^2$, and the mean value theorem \cite{MontgomeryVaughan} gives $\int_T^{2T}|\mathsf Y|^4\dd\ll\int_T^{2T}|\mathsf F^2|^2\dd\ll T$. The primes $\mathfrak Y<p\le2T$ with $r=2$ contribute at most $\frac1{2\pi}\sum_{\mathfrak Y<p\le2T}p^{-1}=O(1)$ pointwise by Mertens' theorem \cite{Mertens1874}. Combining the three parts gives \eqref{eq:taucw-pow}.
\end{proof}

\begin{lemma}[Mean of the long-prime tail]\label{lem:taucw-mean}
\begin{equation}\label{eq:taucw-mean}
\Bigl|\int_T^{2T}\Phi^{\mathrm{lg}}_T(t)\dd\Bigr|\ll\frac T{\log^2T}.
\end{equation}
\end{lemma}
\begin{proof}
For fixed $p$ the function $t\mapsto\mathbf w_{1,t}(p)$ is nondecreasing with values in $[0,1]$, so one integration by parts gives $|\int_T^{2T}\mathbf w_{1,t}(p)\sin(t\log p)\dd|\le3/\log p$. Since only $p\le4T^2$ occur, $|\int_T^{2T}\Phi^{\mathrm{lg}}_T|\le\frac3\pi\sum_{\mathfrak a_T<p\le4T^2}1/(\sqrt p\log p)\ll T/\log^2T$ by partial summation and the prime number theorem.
\end{proof}

\begin{proposition}[The continuous density is negligible]\label{prop:taucw-cont}
For $T\le t\le2T$ put
\begin{equation}\label{eq:taucw-cont}
\mathscr J_T(t)=\int_{\mathfrak a_T}^{t^2}\frac{W\bigl(\frac{\log x}{\log t}\bigr)}{\sqrt x\,\log x}\,e^{it\log x}\,\mathrm dx .
\end{equation}
Then
\begin{equation}\label{eq:taucw-cont-bound}
\sup_{T\le t\le2T}\bigl|\mathscr J_T(t)\bigr|\ll T^{-3/4}(\log T)^{-1/2}.
\end{equation}
\end{proposition}
\begin{proof}
Put $\mathfrak f_t(x)=\sqrt x\,W(\log x/\log t)/\log x$ and $\mathscr J_T(t)=\int\mathfrak f_t(x)e^{it\log x}\,\mathrm dx/x$. Since $\mathfrak f_t(t^2)=0$, one integration by parts gives
\begin{equation}
\mathscr J_T(t)=-\frac{\mathfrak f_t(\mathfrak a_T)e^{it\log\mathfrak a_T}}{it}-\frac1{it}\int_{\mathfrak a_T}^{t^2}\mathfrak j_t(x)e^{it\log x}\frac{\mathrm dx}x,\qquad \mathfrak j_t(x)=x\,\mathfrak f_t'(x).
\end{equation}
Direct differentiation gives the following. On $[\mathfrak a_T,t]$, $\mathfrak j_t(x)=\frac{\sqrt x}{2\log x}-\frac{\sqrt x}{\log^2x}$, so that $\int|\mathfrak j_t'|\ll\int\mathrm dx/(\sqrt x\log x)\ll\sqrt t/\log t$. On $(t,t^2]$, $\mathfrak j_t(x)=\sqrt x\bigl(\frac1{\log x}-\frac1{2\log t}\bigr)-2\sqrt x/\log^2x$, and $|\mathfrak j_t'(x)|\le\bigl(\frac1{\log x}-\frac1{2\log t}\bigr)/(2\sqrt x)+3/(\sqrt x\log^2x)$ for $t$ large; at $x=t^2e^{-s}$ one has $0\le\frac1{\log x}-\frac1{2\log t}=\frac s{2\log t(2\log t-s)}\le s/(2\log^2t)$ and $\mathrm dx/\sqrt x=te^{-s/2}\mathrm ds$, so that $\int_t^{t^2}|\mathfrak j_t'|\ll t\log^{-2}t$. At $x=t$ there is a single jump of size $\sqrt t/\log^2t$, and $|\mathfrak j_t(\mathfrak a_T)|\ll\sqrt{\mathfrak a_T}/\log\mathfrak a_T$, $|\mathfrak j_t(t^2-)|\ll t/\log^2t$. Hence $|\mathfrak j_t(\mathfrak a_T)|+|\mathfrak j_t(t^2-)|+\operatorname{Var}_{[\mathfrak a_T,t^2]}\mathfrak j_t\ll t/\log^2t$. A second integration by parts against $e^{it\log x}\,\mathrm dx/x$ bounds the last integral by $O(\log^{-2}t)$. Since $\mathfrak f_t(\mathfrak a_T)=\sqrt{\mathfrak a_T}/\log\mathfrak a_T$ and $\mathfrak a_T\asymp\sqrt T\log T$, $|\mathscr J_T(t)|\ll\sqrt{\mathfrak a_T}/(t\log\mathfrak a_T)+1/(t\log^2t)\ll T^{-3/4}(\log T)^{-1/2}$.
\end{proof}

\begin{corollary}[Chebyshev representation of the long-prime tail]\label{cor:taucw-cheb}
Let $\mathfrak T(x)=\sum_{p\le x}\log p$, $\mathfrak e(x)=\mathfrak T(x)-x$, and for $T\le t\le2T$
\begin{equation}\label{eq:taucw-G}
\mathscr G_T(t)=\int_{(\mathfrak a_T,t^2]}\frac{W\bigl(\frac{\log x}{\log t}\bigr)}{\sqrt x\,\log x}\,e^{it\log x}\,\mathrm d\mathfrak e(x).
\end{equation}
Then, uniformly in $t\in[T,2T]$,
\begin{equation}\label{eq:taucw-cheb}
\Phi^{\mathrm{lg}}_T(t)=-\frac1\pi\Im\mathscr G_T(t)+O\bigl(T^{-3/4}(\log T)^{-1/2}\bigr).
\end{equation}
\end{corollary}
\begin{proof}
$\Phi^{\mathrm{lg}}_T(t)=-\frac1\pi\Im\int_{(\mathfrak a_T,t^2]}W(\frac{\log x}{\log t})(\sqrt x\log x)^{-1}e^{it\log x}\,\mathrm d\mathfrak T(x)$, and $\mathrm d\mathfrak T=\mathrm dx+\mathrm d\mathfrak e$ splits the integral into $\mathscr J_T(t)$ and $\mathscr G_T(t)$. Apply \cref{prop:taucw-cont}.
\end{proof}

\subsection{Fixed-frequency coefficients and the spectrum of the points}\label{sec:taucw-spec}
\begin{theorem}[Fixed-frequency coefficient of the phase]\label{thm:taucw-coef}
Fix $\omega>0$ and $0<\theta\le1$ and put $\Upsilon=e^\omega$. Then
\begin{equation}\label{eq:taucw-coef}
\frac{2\pi}T\int_T^{2T}e^{i\omega t}\,\Phi^{\mathrm w}_\theta(t)\dd=-\,i\,\frac{\Lambda(\Upsilon)}{\sqrt\Upsilon\,\log\Upsilon}+O_{\omega,\theta}\bigl(\log^{-2}T\bigr),
\end{equation}
where $\Lambda(\Upsilon)=0$ unless $\Upsilon$ is a prime power.
\end{theorem}
\begin{proof}
Put $\mathfrak k_n=\Lambda(n)/(\pi\sqrt n\log n)$, so that $\Phi^{\mathrm w}_\theta=-\sum_n\mathfrak k_n\mathbf w_{\theta,t}(n)\sin(t\log n)$. For fixed $n$ the weight $t\mapsto\mathbf w_{\theta,t}(n)$ is nondecreasing with values in $[0,1]$, so integration by parts gives, for $\log n\ne\omega$,
\begin{equation}
\Bigl|\int_T^{2T}e^{i\omega t}\mathbf w_{\theta,t}(n)\sin(t\log n)\dd\Bigr|\le\frac3{2|\omega-\log n|}+\frac3{2(\omega+\log n)}.
\end{equation}
The finitely many $n\ne\Upsilon$ with $\log n<2\omega$ contribute $O_{\omega,\theta}(1)$. For $\log n\ge2\omega$ both denominators are $\gg\log n$, and the sum over $n\le(2T)^{2\theta}$ of $\mathfrak k_n/\log n$ is $\ll T^\theta/\log^2T\le T/\log^2T$ by the prime number theorem and the convergence of the series over prime powers with exponent at least $2$. If $\Upsilon$ is a prime power then $\mathbf w_{\theta,t}(\Upsilon)=1$ for all $t\in[T,2T]$ once $T^\theta\ge\Upsilon$, and $\int_T^{2T}e^{i\omega t}\sin(\omega t)\dd=\frac{iT}2+O_{\omega,\theta}(1)$. The resonant term is therefore $-\mathfrak k_\Upsilon(\frac{iT}2+O_{\omega,\theta}(1))$. Multiply the sum of these contributions by $2\pi/T$ and use $\pi\mathfrak k_\Upsilon=\Lambda(\Upsilon)/(\sqrt\Upsilon\log\Upsilon)$.
\end{proof}

With $\chi(\omega)=-\Lambda(e^\omega)e^{-\omega/2}$ of \eqref{eq:v9-chi}, \cref{thm:taucw-coef} reads $\frac{2\pi}T\int_T^{2T}e^{i\omega t}\Phi^{\mathrm w}_\theta\dd=i\chi(\omega)/\omega+O_{\omega,\theta}(\log^{-2}T)$, the same limit as for the zeros in Step~1 of the proof of \cref{lem:v9-cancel}.

\begin{proposition}[Spectral formula for the taucw points]\label{prop:taucw-spec}
For $\omega>0$ and $T>2\pi$,
\begin{equation}\label{eq:taucw-spec}
\frac{2\pi}T\sum_{T<\tau_k^{\mathrm w}\le2T}e^{i\omega\tau_k^{\mathrm w}}=\chi(\omega)-i\omega\,\frac{2\pi}T\int_T^{2T}e^{i\omega t}\bigl(N^{\mathrm w}(t)-F^{\mathrm w}_1(t)\bigr)\dd+O_\omega\bigl(\log^{-2}T\bigr).
\end{equation}
In particular, for each fixed prime $p$,
\begin{equation}\label{eq:taucw-spec-p}
\frac{2\pi}T\sum_{T<\tau_k^{\mathrm w}\le2T}p^{\frac12+i\tau_k^{\mathrm w}}=-\log p-i\sqrt p\,\log p\,\frac{2\pi}T\int_T^{2T}p^{it}\bigl(N^{\mathrm w}(t)-F^{\mathrm w}_1(t)\bigr)\dd+O_p\bigl(\log^{-2}T\bigr).
\end{equation}
With $\varrho_F=\lfloor F+\frac12\rfloor-N^{\mathrm w}$ as in \eqref{eq:v9-slip} and $F=F^{\mathrm w}_1$, one has $N^{\mathrm w}-F=-\varrho_F-\mathrm{saw}(F)$, so that the integral in \eqref{eq:taucw-spec} equals $-\int_T^{2T}e^{i\omega t}\varrho_F\dd+\frac T{2\pi}\varsigma_T(\omega;F)$ with $\varsigma_T$ of \eqref{eq:v9-varsigma}.
\end{proposition}
\begin{proof}
Integration by parts in the Stieltjes integral of $e^{i\omega t}$ against the step function $N^{\mathrm w}$ gives $\sum_{T<\tau_k^{\mathrm w}\le2T}e^{i\omega\tau_k^{\mathrm w}}=[e^{i\omega t}N^{\mathrm w}(t)]_T^{2T}-i\omega\int_T^{2T}e^{i\omega t}N^{\mathrm w}(t)\dd$. Insert $N^{\mathrm w}=N_0+\Phi^{\mathrm w}_1+(N^{\mathrm w}-F^{\mathrm w}_1)$. One integration by parts gives $-i\omega\int_T^{2T}e^{i\omega t}N_0\dd=-[e^{i\omega t}N_0]_T^{2T}+\int_T^{2T}e^{i\omega t}N_0'\dd$, and the last integral is $O_\omega(\log T)$. Hence the boundary terms combine to $[e^{i\omega t}(N^{\mathrm w}-N_0)]_T^{2T}=O(T/\log^2T)$ by \eqref{eq:taucw-ord}. By \cref{thm:taucw-coef}, $-i\omega\int_T^{2T}e^{i\omega t}\Phi^{\mathrm w}_1\dd=\frac T{2\pi}\chi(\omega)+O_\omega(T/\log^2T)$. Multiplying by $2\pi/T$ gives \eqref{eq:taucw-spec}, and \eqref{eq:taucw-spec-p} follows with $\omega=\log p$, since $\sqrt p\,\chi(\log p)=-\log p$. The last assertion is $N^{\mathrm w}-F=(N^{\mathrm w}-\lfloor F+\frac12\rfloor)+(\lfloor F+\frac12\rfloor-F)$.
\end{proof}

For a bounded measurable function $f$ on a finite interval $I=[a,b]$ and a half-integer $\lambda$ put
\begin{equation}\label{eq:taucw-local}
\begin{aligned}
\mathsf O_\lambda&=\{t\in I:\ f(t)<\lambda\},\qquad\mathfrak M_\lambda=a+|\mathsf O_\lambda|,\\
\varrho_{f;I}(t)&=\sum_{\lambda\in\mathbb Z+\frac12}\bigl(\ind_{\{t<\mathfrak M_\lambda\}}-\ind_{\{f(t)<\lambda\}}\bigr),\qquad a<t<b.
\end{aligned}
\end{equation}
The point $\mathfrak M_\lambda$ is the position at which the sublevel set $\mathsf O_\lambda$ would end if it were moved to the left end of $I$ without change of measure. For $\lambda\le\operatorname{ess\,inf}f$ and for $\lambda>\sup f$ the summand vanishes on $I$ (up to a null set), so the sum is finite. The function $\varrho_{f;I}$ is determined by the sublevel sets of $f$ on $I$ alone; it is not asserted to coincide with $\varrho_F$ of \eqref{eq:v9-slip}, which is defined by the sublevel sets on $(0,\infty)$.

\begin{theorem}[Quadratic estimate for the local reordering defect]\label{thm:taucw-quad}
Let $I=[a,b]$ be a finite interval, $g\in C^1(I)$ with $g'\ge\mathfrak l>0$, $\phi$ a bounded measurable function on $I$ and $f=g+\phi$. For every Lipschitz function $\mathsf d:I\to\mathbb C$,
\begin{equation}\label{eq:taucw-quad}
\Bigl|\int_I\mathsf d(t)\,\varrho_{f;I}(t)\dd\Bigr|\le\frac{2\operatorname{Lip}(\mathsf d)}{\mathfrak l}\int_I\bigl(|\phi(t)|^2+|\phi(t)|\bigr)\dd.
\end{equation}
No differentiability of $f$, no simplicity of its level crossings and no lower bound for $|f'|$ is used.
\end{theorem}
\begin{proof}
Fix a half-integer $\lambda$ and put $\mathsf B_\lambda=\{t\in I:g(t)<\lambda\}$, $\mathfrak N_\lambda=a+|\mathsf B_\lambda|$, $\mathsf J_\lambda=[a,\mathfrak M_\lambda)$. As $g$ is strictly increasing and continuous, $\mathsf B_\lambda=[a,\mathfrak N_\lambda)$ (possibly empty or all of $I$), and $|\mathsf J_\lambda|=|\mathsf O_\lambda|$. The summand of \eqref{eq:taucw-local} is $\ind_{\mathsf J_\lambda}-\ind_{\mathsf O_\lambda}$, and $\int_I(\ind_{\mathsf J_\lambda}-\ind_{\mathsf O_\lambda})\dd=0$, so
\begin{equation}
\int_I\mathsf d\,(\ind_{\mathsf J_\lambda}-\ind_{\mathsf O_\lambda})\dd=\int_I\bigl(\mathsf d(t)-\mathsf d(\mathfrak N_\lambda)\bigr)(\ind_{\mathsf J_\lambda}-\ind_{\mathsf O_\lambda})\dd .
\end{equation}
Let $\mathsf Q_\lambda=\int_I|t-\mathfrak N_\lambda|\,|\ind_{\mathsf O_\lambda}-\ind_{\mathsf B_\lambda}|\dd$, and let $\mathsf e_+=|\mathsf O_\lambda\setminus\mathsf B_\lambda|$, $\mathsf e_-=|\mathsf B_\lambda\setminus\mathsf O_\lambda|$, so that $\mathfrak M_\lambda-\mathfrak N_\lambda=\mathsf e_+-\mathsf e_-$. The set $\mathsf O_\lambda\setminus\mathsf B_\lambda$ lies to the right of $\mathfrak N_\lambda$ and $\mathsf B_\lambda\setminus\mathsf O_\lambda$ to the left, so their first moments about $\mathfrak N_\lambda$ are at least $\mathsf e_+^2/2$ and $\mathsf e_-^2/2$. Hence $\mathsf Q_\lambda\ge(\mathsf e_+^2+\mathsf e_-^2)/2\ge(\mathfrak M_\lambda-\mathfrak N_\lambda)^2/2=\int_I|t-\mathfrak N_\lambda|\,|\ind_{\mathsf J_\lambda}-\ind_{\mathsf B_\lambda}|\dd$, and the triangle inequality gives
\begin{equation}
\Bigl|\int_I\mathsf d\,(\ind_{\mathsf J_\lambda}-\ind_{\mathsf O_\lambda})\dd\Bigr|\le2\operatorname{Lip}(\mathsf d)\,\mathsf Q_\lambda .
\end{equation}
If $\ind_{\mathsf O_\lambda}(t)\ne\ind_{\mathsf B_\lambda}(t)$ then $\lambda$ lies between $f(t)$ and $g(t)$, so $|g(t)-\lambda|\le|\phi(t)|$; since $g'\ge\mathfrak l$, this gives $|t-\mathfrak N_\lambda|\le|\phi(t)|/\mathfrak l$, also when $\mathsf B_\lambda$ is empty or equal to $I$ (then $\mathfrak N_\lambda$ is an endpoint and $g$ lies on one side of $\lambda$ throughout $I$). For each $t$ at most $|\phi(t)|+1$ half-integers lie between $f(t)$ and $g(t)$, so $\sum_\lambda \mathsf Q_\lambda\le\mathfrak l^{-1}\int_I|\phi|(|\phi|+1)\dd$. Summing over $\lambda$ proves \eqref{eq:taucw-quad}.
\end{proof}

\begin{corollary}[Low frequencies for a moderate cutoff]\label{cor:taucw-lowfreq}
Fix $0<\theta<\frac12$ and a smooth $\mathsf b$ supported in $(1,2)$, let $I_T=[T,2T]$ and $\mathsf d_u(t)=\mathsf b(t/T)e^{iut}$, $u\in\mathbb R$. Then, for $f=F^{\mathrm w}_\theta$,
\begin{equation}\label{eq:taucw-lowfreq}
\Bigl|\frac1T\int_{I_T}\mathsf d_u(t)\,\varrho_{f;I_T}(t)\dd\Bigr|\ll_{\theta,\mathsf b}\frac{(|u|+T^{-1})(1+\log\log T)}{\log T}.
\end{equation}
In particular the normalized Fourier coefficient tends to $0$ uniformly in $|u|=o(\log T/\log\log T)$.
\end{corollary}
\begin{proof}
Take $g=N_0$ and $\phi=\Phi^{\mathrm w}_\theta$, bounded on $I_T$ by \cref{lem:taucw-amp}. Then $N_0'\ge\mathfrak l=\frac1{2\pi}\log\frac T{2\pi}\gg\log T$ on $I_T$ and $\operatorname{Lip}(\mathsf d_u)\ll_{\mathsf b}|u|+T^{-1}$. By \cref{thm:taucw-var}, $T^{-1}\int_{I_T}|\Phi^{\mathrm w}_\theta|^2\ll_\theta1+\log\log T$, and the Cauchy--Schwarz inequality bounds $T^{-1}\int_{I_T}|\Phi^{\mathrm w}_\theta|$ by the square root of this. Insert both into \eqref{eq:taucw-quad}.
\end{proof}
For $|u|\asymp\log T$ the factor $\operatorname{Lip}(\mathsf d_u)$ cancels the gain $1/\mathfrak l$ and the right side of \eqref{eq:taucw-lowfreq} is of order $\log\log T$, so the estimate carries no information at that frequency.

\subsection{Fourier identities for counting discrepancies}\label{sec:taucw-plancherel}
Let $\mathfrak x_1\le\dots\le\mathfrak x_\ell$ and $\mathfrak y_1\le\dots\le\mathfrak y_\ell$ be real numbers, repetitions allowed, and put
\begin{equation}\label{eq:taucw-M}
\mathscr M_\ell(t)=\sum_{j=1}^\ell\bigl(\ind_{\{\mathfrak x_j\le t\}}-\ind_{\{\mathfrak y_j\le t\}}\bigr),\qquad
\mathsf C_\ell(u)=\sum_{j=1}^\ell\bigl(e^{iu\mathfrak x_j}-e^{iu\mathfrak y_j}\bigr).
\end{equation}
With $\mathfrak x_j=\gamma_j$ and $\mathfrak y_j=\tau_j^{\mathrm w}$ these are the difference of the two counting functions of the first $\ell$ points and the difference of the two exponential sums.

\begin{theorem}[Plancherel identity for two ordered lists]\label{thm:taucw-plan}
\begin{align}
\int_{\mathbb R}|\mathscr M_\ell(t)|^2\dd&=\frac1{2\pi}\int_{\mathbb R}\frac{|\mathsf C_\ell(u)|^2}{u^2}\,\mathrm du,\label{eq:taucw-plan}\\
\int_{\mathbb R}|\mathscr M_\ell(t)|\dd&=\sum_{j=1}^\ell|\mathfrak x_j-\mathfrak y_j|.\label{eq:taucw-L1sorted}
\end{align}
The integrand in \eqref{eq:taucw-plan} is continuous at $u=0$, where $\mathsf C_\ell(0)=0$. For every real $u$, $|\mathsf C_\ell(u)|\le|u|\sum_j|\mathfrak x_j-\mathfrak y_j|$.
\end{theorem}
\begin{proof}
$\mathscr M_\ell$ is bounded, piecewise constant and compactly supported, and its derivative in the sense of distributions is the sum of the unit point masses at the $\mathfrak x_j$ minus the sum of the unit point masses at the $\mathfrak y_j$. With $\widehat f(u)=\int f(t)e^{-iut}\dd$ this gives $iu\widehat{\mathscr M_\ell}(u)=\sum_j(e^{-iu\mathfrak x_j}-e^{-iu\mathfrak y_j})=\mathsf C_\ell(-u)=\overline{\mathsf C_\ell(u)}$, and \eqref{eq:taucw-plan} is Plancherel's identity. For \eqref{eq:taucw-L1sorted}, because both lists are sorted, $\{j:\mathfrak x_j\le t\}=\{1,\dots,\mathsf j_1(t)\}$ and $\{j:\mathfrak y_j\le t\}=\{1,\dots,\mathsf j_2(t)\}$ are initial segments, so $|\mathscr M_\ell(t)|=|\mathsf j_1(t)-\mathsf j_2(t)|=\sum_j|\ind_{\{\mathfrak x_j\le t\}}-\ind_{\{\mathfrak y_j\le t\}}|$, and $\int_{\mathbb R}|\ind_{\{\mathfrak x_j\le t\}}-\ind_{\{\mathfrak y_j\le t\}}|\dd=|\mathfrak x_j-\mathfrak y_j|$. The last assertion follows from $|e^{iux}-e^{iuy}|\le|u||x-y|$.
\end{proof}

For a bandwidth $\mathfrak B>0$ let $\mathfrak F_{\mathfrak B}(t)=\frac{\mathfrak B}{2\pi}\bigl(\frac{\sin(\mathfrak Bt/2)}{\mathfrak Bt/2}\bigr)^2$, the Fej\'er kernel, which is nonnegative with $\int\mathfrak F_{\mathfrak B}=1$ and $\widehat{\mathfrak F_{\mathfrak B}}(u)=(1-|u|/\mathfrak B)_+$.

\begin{theorem}[Band-limited form]\label{thm:taucw-fejer}
For $\mathscr M_\ell,\mathsf C_\ell$ as in \eqref{eq:taucw-M},
\begin{equation}\label{eq:taucw-fejer}
\int_{\mathbb R}\bigl|\mathfrak F_{\mathfrak B}*\mathscr M_\ell(t)\bigr|^2\dd=\frac1{2\pi}\int_{-\mathfrak B}^{\mathfrak B}\Bigl(1-\frac{|u|}{\mathfrak B}\Bigr)^2\frac{|\mathsf C_\ell(u)|^2}{u^2}\,\mathrm du ,
\end{equation}
and for every real $t$
\begin{equation}\label{eq:taucw-pointwise}
|\mathscr M_\ell(t)|\le\Bigl[\frac{\mathfrak B}{2\pi^2}\int_{-\mathfrak B}^{\mathfrak B}\Bigl(1-\frac{|u|}{\mathfrak B}\Bigr)^2\frac{|\mathsf C_\ell(u)|^2}{u^2}\mathrm du\Bigr]^{1/2}+\sum_{j=1}^\ell\min\Bigl\{1,\frac2{\pi\mathfrak B|t-\mathfrak x_j|}\Bigr\}+\sum_{j=1}^\ell\min\Bigl\{1,\frac2{\pi\mathfrak B|t-\mathfrak y_j|}\Bigr\}.
\end{equation}
\end{theorem}
\begin{proof}
The transform of $\mathfrak F_{\mathfrak B}*\mathscr M_\ell$ is $(1-|u|/\mathfrak B)_+\widehat{\mathscr M_\ell}(u)$, supported in $[-\mathfrak B,\mathfrak B]$, and \eqref{eq:taucw-fejer} is Plancherel's identity combined with $|\widehat{\mathscr M_\ell}(u)|^2=|\mathsf C_\ell(u)|^2/u^2$. By Fourier inversion and the Cauchy--Schwarz inequality, $\|\mathfrak F_{\mathfrak B}*\mathscr M_\ell\|_\infty\le\sqrt{\mathfrak B/\pi}\,\|\mathfrak F_{\mathfrak B}*\mathscr M_\ell\|_2$. For a Heaviside function $\ind_{\{c\le t\}}$, the difference from its convolution with $\mathfrak F_{\mathfrak B}$ equals $\pm\int_{|t-c|}^\infty\mathfrak F_{\mathfrak B}$, which is at most $\min\{1,2/(\pi\mathfrak B|t-c|)\}$ because $\mathfrak F_{\mathfrak B}(y)\le2/(\pi\mathfrak By^2)$. Summing over the $2\ell$ Heaviside functions of $\mathscr M_\ell$ gives \eqref{eq:taucw-pointwise}.
\end{proof}

\subsection{The counting surrogate \texorpdfstring{$S^{\mathrm w}$}{Sw} and the quadratic comparison of phases}\label{sec:taucw-S}
The function $S^{\mathrm w}=N^{\mathrm w}-N_0-R$ of \S\ref{sec:taucw-exist} is determined by the taucw points and the explicit function $R$ of \eqref{eq:RVM}; neither the imaginary parts nor the real parts of the zeros enter its definition.

\begin{proposition}[Exact counting discrepancy]\label{prop:taucw-S}
For every $t>2\pi$,
\begin{equation}\label{eq:taucw-S-exact}
S(t)-S^{\mathrm w}(t)=N(t)-N^{\mathrm w}(t)\in\mathbb Z .
\end{equation}
Let $N_{[\ell]}(t)=\#\{j\le\ell:\gamma_j\le t\}$ and $N^{\mathrm w}_{[\ell]}(t)=\#\{j\le\ell:\tau_j^{\mathrm w}\le t\}$. Then
\begin{equation}\label{eq:taucw-exc}
\meas\{t:\ N_{[\ell]}(t)\ne N^{\mathrm w}_{[\ell]}(t)\}\le\int_{\mathbb R}\bigl|N_{[\ell]}-N^{\mathrm w}_{[\ell]}\bigr|\dd=\sum_{j=1}^\ell\bigl|\gamma_j-\tau_j^{\mathrm w}\bigr| .
\end{equation}
\end{proposition}
\begin{proof}
Subtract the definition of $S^{\mathrm w}$ from \eqref{eq:RVM}. The difference of two integer-valued functions is $0$ or has modulus at least $1$, which gives the first inequality in \eqref{eq:taucw-exc}, and the equality is \eqref{eq:taucw-L1sorted} with $\mathfrak x_j=\gamma_j$, $\mathfrak y_j=\tau_j^{\mathrm w}$; both lists are sorted by \cref{prop:taucw-ord}.
\end{proof}

\begin{proposition}[Transfer of coefficients needs counts]\label{prop:taucw-transfer}
Let
\begin{equation}\label{eq:taucw-gh}
\mathsf g^{\mathrm w}=\limsup_{t\to\infty}\frac{|N^{\mathrm w}(t)-N_0(t)|}{\log t},\qquad
\mathsf h^{\mathrm w}=\limsup_{t\to\infty}\frac{|N(t)-N^{\mathrm w}(t)|}{\log t},
\end{equation}
with the value $+\infty$ allowed. Then
\begin{equation}\label{eq:taucw-transfer}
\limsup_{t\to\infty}\frac{|S(t)|}{\log t}\le\mathsf g^{\mathrm w}+\mathsf h^{\mathrm w}.
\end{equation}
\end{proposition}
\begin{proof}
By \eqref{eq:taucw-S-exact}, $S=S^{\mathrm w}+(N-N^{\mathrm w})$, and $|S^{\mathrm w}|\le|N^{\mathrm w}-N_0|+|R|$ with $R=O(1/t)$.
\end{proof}
Thus a bound on $\mathsf h^{\mathrm w}$ does not require a uniform bound on $|\gamma_k-\tau_k^{\mathrm w}|$; it requires only that the integer-valued function $N-N^{\mathrm w}$ be small compared with $\log t$.

\begin{proposition}[Exact quadratic comparison of two phases]\label{prop:taucw-polar}
Let $I$ be a finite interval, $F$ and $F'$ real functions in $L^2(I)$, and $D_{F'}=N-F'$. Then
\begin{equation}\label{eq:taucw-polar}
\int_I|F-N|^2\dd-\int_I|F'-N|^2\dd=-2\int_ID_{F'}\,(F-F')\dd+\int_I|F-F'|^2\dd .
\end{equation}
In particular $F$ has a smaller quadratic discrepancy than $F'$ on $I$ if and only if
\begin{equation}\label{eq:taucw-polar-crit}
2\int_ID_{F'}\,(F-F')\dd>\int_I|F-F'|^2\dd .
\end{equation}
\end{proposition}
\begin{proof}
Write $F-N=(F'-N)+(F-F')$ and expand the square; $F'-N=-D_{F'}$.
\end{proof}
The identity retains the sign of the mixed term $\int_ID_{F'}(F-F')$, which a comparison through the triangle inequality discards. For $F=F^{\mathrm w}_\theta$ and $F'$ any of the phases of the preceding sections, $F-F'$ is a difference of two explicit prime sums.

\subsection{Computed taucw points: closeness to the zeros, centroids and exponential sums}\label{sec:taucw-data}
The points $\tau_k^{\mathrm w}$ for $\theta=1$ and $1\le k\le3000$ were computed from $F^{\mathrm w}_1$ alone: the prime powers up to $t^2$ with the weights of \eqref{eq:taucw-phase}, the level $k-\frac12$ of $F^{\mathrm w}_1$ located by a search on a grid of step $0.1$ and refined (file \texttt{taucw\_independent\_3000.csv}, one crossing detected for each level). The first $3000$ zeros $\gamma_k$ were computed separately (mpmath, $25$ digits) and enter only the comparisons below. Write $\mathsf y=(\mathsf y_k)$ for either sequence and $\mathrm{Cen}_\ell(\omega;\mathsf y)$ as in \eqref{eq:v9-Cen}.

\begin{enumerate}[label=(\roman*),leftmargin=*]
\item \emph{Points.} $\tau_1^{\mathrm w}=14.13526$ against $\gamma_1=14.13473$, and $\tau_{3000}^{\mathrm w}=3533.32873$ against $\gamma_{3000}=3533.32824$. The root-mean-square of $\tau^{\mathrm w}_k-\gamma_k$ is $0.00411$, that is $0.00362$ in units of the local mean spacing $2\pi/\log(\gamma_k/2\pi)$; the maximum of $|\tau^{\mathrm w}_k-\gamma_k|$ is $0.0279$ (at $k=1496$) and $\max_n|\sum_{k\le n}(\tau^{\mathrm w}_k-\gamma_k)|=0.0281$.
\item \emph{Truncation.} Since $F^{\mathrm w}_1(\tau_k^{\mathrm w})=k-\frac12$, the phase at the points is $\Phi^{\mathrm w}_1(\tau_k^{\mathrm w})=k-\frac12-N_0(\tau_k^{\mathrm w})$, and $\max_{k\le3000}|\Phi^{\mathrm w}_1(\tau_k^{\mathrm w})|=0.828$, whereas $\min_k H(\tau_k^{\mathrm w})=5.72$. The truncation of \eqref{eq:taucwH} is therefore inactive at all $3000$ points; this alone does not give $\tau^{\mathrm w,H}_k=\tau^{\mathrm w}_k$, which needs $|\Phi^{\mathrm w}_1(t)|<H(t)$ on all of $[2\pi,\tau^{\mathrm w}_{3000}]$.
\item \emph{Centroids.} \Cref{tab:taucw-cen} gives $\mathrm{Cen}_\ell(1;\mathsf y)$ for $\mathsf y=\gamma$ and $\mathsf y=\tau^{\mathrm w}$; \cref{fig:taucw-cen3d,fig:taucw-cenplanar} show the two centroid curves, which coincide to the resolution of the plots. One has $\max_{10\le\ell\le3000}|\mathrm{Cen}_\ell(1;\tau^{\mathrm w})-\mathrm{Cen}_\ell(1;\gamma)|=0.0176$ (at $\ell=1717$) and the mean of this difference over $1\le\ell\le3000$ is $0.0122$.
\item \emph{Exponential sums in the frequency $u$.} Put $\mathsf p_{\mathsf y}(u)=\mathrm{Par}_{3000}(u;\mathsf y)=\sum_{k\le3000}e^{iu\mathsf y_k}$ (\eqref{eq:v9-Cen}; for $\mathsf y=\tau^{\mathrm w}$ we write $\mathsf p_{\mathrm{taucw}}$). On the grid $0.25\le u\le5$ of step $4\cdot10^{-5}$, $\max_u|\mathsf p_\gamma(u)-\mathsf p_{\mathrm{taucw}}(u)|=2.49$ (at $u=4.7613$), the median of the difference is $0.32$ and $\max_u|\mathsf p_\gamma(u)|=412.8$ (the grid does not contain $u=\log7$, where $\Re\mathsf p_\gamma=-412.86$, \cref{tab:taucw-cen}). The peaks of both sums sit at the prime powers $u=\log p^m$ and have the height of the Landau term $-\frac T{2\pi}\Lambda(n)n^{-1/2}$ of \eqref{eq:v9-landau}, with $T=\tau^{\mathrm w}_{3000}$ (\cref{tab:taucw-cen}, lower part; \cref{fig:taucw-exp}). At $u=\log6,\log10,\log12$, where $\Lambda=0$, the real parts are $2.03,\ 1.87,\ 2.38$ for $\mathsf p_\gamma$ and $2.14,\ 1.52,\ 2.54$ for $\mathsf p_{\mathrm{taucw}}$.
\end{enumerate}

\begin{table}[ht]
\centering\small
\begin{tabular}{@{}rrrr@{}}
\toprule
$\ell$ & $\mathrm{Cen}_\ell(1;\gamma)$ & $\mathrm{Cen}_\ell(1;\tau^{\mathrm w})$ & difference\\
\midrule
$10$ & $0.998352+1.186415\,i$ & $1.002120+1.186981\,i$ & $0.0038$\\
$30$ & $0.714294+0.939104\,i$ & $0.714425+0.937074\,i$ & $0.0020$\\
$100$ & $0.769068+0.842240\,i$ & $0.763560+0.844983\,i$ & $0.0062$\\
$300$ & $0.753707+0.781773\,i$ & $0.749720+0.786638\,i$ & $0.0063$\\
$1000$ & $0.772074+0.762675\,i$ & $0.762925+0.758071\,i$ & $0.0102$\\
$3000$ & $0.773945+0.767306\,i$ & $0.761620+0.778420\,i$ & $0.0166$\\
\bottomrule
\end{tabular}

\medskip
\begin{tabular}{@{}rrrrr@{}}
\toprule
$n$ & $-\frac T{2\pi}\Lambda(n)n^{-1/2}$ & $\Re\mathsf p_\gamma(\log n)$ & $\Re\mathsf p_{\mathrm{taucw}}(\log n)$ & $|\mathsf p_\gamma-\mathsf p_{\mathrm{taucw}}|$\\
\midrule
$2$ & $-275.62$ & $-276.26$ & $-276.29$ & $0.03$\\
$3$ & $-356.69$ & $-356.46$ & $-356.52$ & $0.07$\\
$4$ & $-194.89$ & $-193.80$ & $-193.91$ & $0.12$\\
$5$ & $-404.76$ & $-403.95$ & $-403.55$ & $0.42$\\
$7$ & $-413.60$ & $-412.86$ & $-412.20$ & $0.67$\\
$8$ & $-137.81$ & $-137.58$ & $-137.52$ & $0.11$\\
$9$ & $-205.93$ & $-204.72$ & $-204.84$ & $0.27$\\
$11$ & $-406.57$ & $-404.77$ & $-404.04$ & $0.73$\\
$49$ & $-156.33$ & $-156.37$ & $-156.43$ & $0.34$\\
$125$ & $-80.95$ & $-74.65$ & $-75.27$ & $1.28$\\
\bottomrule
\end{tabular}
\caption{Upper part: centroids $\mathrm{Cen}_\ell(1;\mathsf y)$ for $\ell\in\{10,30,100,300,1000,3000\}$, computed from the first $3000$ zeros and from the first $3000$ taucw points ($\theta=1$). Lower part: exponential sums $\mathsf p_{\mathsf y}(u)=\mathrm{Par}_{3000}(u;\mathsf y)$ at the frequencies $u=\log n$ and the Landau term, $T=3533.3287$.}\label{tab:taucw-cen}
\end{table}

\begin{figure}[ht]
\centering
\includegraphics[width=0.92\linewidth]{fig_taucw_centroid_3d_v27.png}
\caption{Centroids $\mathrm{Cen}_\ell(1;\mathsf y)$, $\ell\le3000$, with axes $(\Re,\Im,\log_{10}\ell)$: zeros $\gamma_k$ (black, solid) and taucw points (green, dashed; computed without zeros).}\label{fig:taucw-cen3d}
\end{figure}

\begin{figure}[ht]
\centering
\includegraphics[width=\linewidth]{fig_taucw_centroid_planar_v27.png}
\caption{Real part, imaginary part, modulus of the difference (log--log) and trajectory in the plane of $\mathrm{Cen}_\ell(1;\mathsf y)$ for the zeros and for the taucw points.}\label{fig:taucw-cenplanar}
\end{figure}

\begin{figure}[ht]
\centering
\includegraphics[width=\linewidth]{fig_taucw_exp_sums_u_v27.png}
\caption{Top: $\Re\sum_{k\le3000}e^{iu\gamma_k}$ and $\Re\sum_{k\le3000}e^{iu\,\mathrm{taucw}_k}$ for $0.25\le u\le5$; dotted lines at $u=\log p^m$. Bottom: $|\mathsf p_\gamma(u)|$ and $|\mathsf p_\gamma(u)-\mathsf p_{\mathrm{taucw}}(u)|$ (logarithmic scale).}\label{fig:taucw-exp}
\end{figure}

\subsection{Detection of zeros off the critical line by the phase}\label{sec:taucw-detect}
Every statement of this subsection is unconditional. It isolates what the explicit formula of \cref{lem:taucw-expl} says about a zero $\rho=\frac12+y+i\gamma$ with $y\ne0$: the real part of a zero enters $N_0+P^W_X-N$ only through $\mathcal Y_L$ of \eqref{eq:taucw-on-off}; the contribution of the pair $\{\rho,1-\bar\rho\}$ is computed in closed form, with its spectral density and its $L^2$ mass (\cref{prop:taucw-pair}, \cref{cor:taucw-pair}); up to $O(\log t)$, uniformly in the cutoff, the prime sum with its compensator \emph{is} $\mathcal Y_L$ (\cref{thm:taucw-P}); the absence of zeros off the line is equivalent to the boundedness of the weighted $L^2$ norm of this prime sum, and of $\mathcal Y_L$, along any unbounded sequence of cutoffs (\cref{thm:taucw-equiv}); and a zero off the line, like a multiple zero, forces the strictly increasing taucw points to leave the imaginary parts, at an explicit cost in $L^1$ (\cref{prop:taucw-shared}). The zeros with $\beta\ne\frac12$ occur in pairs $\{\rho,1-\bar\rho\}$ with the same imaginary part, the same multiplicity and real parts $\frac12\pm y$, $y>0$, so that each pair has exactly one member with $\beta>\frac12$. As in \eqref{eq:weights}, $W(u)=\min\{1,2-u\}_+$ for all real $u$.

\begin{lemma}[Laplace representations of the kernel]\label{lem:taucw-laplace}
For every $L>0$,
\begin{align}
\mathfrak G_L(w)&=-\int_{-\infty}^{2L}W\!\Bigl(\frac sL\Bigr)e^{sw}\,ds\qquad(\Re w>0),\label{eq:taucw-lap+}\\
\mathfrak G_L(w)&=\int_{L}^{\infty}\Bigl(1-W\!\Bigl(\frac sL\Bigr)\Bigr)e^{sw}\,ds\qquad(\Re w<0).\label{eq:taucw-lap-}
\end{align}
\end{lemma}
\begin{proof}
For $\Re w>0$, $\int_{-\infty}^Le^{sw}ds=e^{Lw}/w$ and, integrating by parts, $\int_L^{2L}(2-\frac sL)e^{sw}ds=-\frac{e^{Lw}}w+\frac{e^{2Lw}-e^{Lw}}{Lw^2}$; the sum is $-\mathfrak G_L(w)$. For $\Re w<0$, $\int_{2L}^\infty e^{sw}ds=-e^{2Lw}/w$ and $\int_L^{2L}(\frac sL-1)e^{sw}ds=\frac{e^{2Lw}}w-\frac{e^{2Lw}-e^{Lw}}{Lw^2}$; the sum is $\mathfrak G_L(w)$.
\end{proof}

For $L>0$ and $y>0$ put
\begin{equation}\label{eq:taucw-J}
\mathcal J_L(y;v)=-\int_0^{2L}\mathscr R_{L,y}(u)\sin(uv)\,du,\qquad
\mathscr R_{L,y}(u)=\frac2\pi\,W\!\Bigl(\frac uL\Bigr)\frac{\cosh(yu)-1}{u},
\end{equation}
and $\mathcal J_L(0;v)=0$.

\begin{proposition}[Exact profile of an off-line pair]\label{prop:taucw-pair}
Let $L>0$, $y>0$, $\gamma\in\mathbb R$ and $t\ne\gamma$. For the zero $\rho=\frac12+y+i\gamma$ and its reflection $1-\bar\rho=\frac12-y+i\gamma$, the terms of $\mathcal Y_L$ in \eqref{eq:taucw-on-off} add up to
\begin{equation}\label{eq:taucw-pair}
\mathcal Y_L[\rho]+\mathcal Y_L[1-\bar\rho]=\mathcal J_L(y;\gamma-t).
\end{equation}
Here $\mathcal Y_L[\rho]=\frac1\pi\Im\bigl(\mathfrak I_L(y_\rho+i(\gamma-t))-\mathfrak I_L(i(\gamma-t))\bigr)$ is the term of the zero $\rho$. The function $v\mapsto\mathcal J_L(y;v)$ is entire, odd and of exponential type $2L$.
\end{proposition}
\begin{proof}
Fix $v=\gamma-t\neq0$. The segment $\{x+iv\}$ avoids the pole $w=0$ of $\mathfrak G_L$, and $\mathfrak I_L(\pm y+iv)-\mathfrak I_L(iv)=\int_0^{\pm y}\mathfrak G_L(x+iv)\,dx$. Put $\Psi(y)=\int_0^{y}\mathfrak G_L(x+iv)dx+\int_0^{-y}\mathfrak G_L(x+iv)dx$, so that $\Psi(0)=0$ and $\Psi'(y)=\mathfrak G_L(y+iv)-\mathfrak G_L(-y+iv)$. For $y>0$ apply \eqref{eq:taucw-lap+} at $y+iv$ and \eqref{eq:taucw-lap-} at $-y+iv$ and take imaginary parts:
\begin{equation}
\Im\,\Psi'(y)=-\int_{-\infty}^{2L}W\!\Bigl(\frac sL\Bigr)e^{sy}\sin(sv)\,ds-\int_L^\infty\Bigl(1-W\!\Bigl(\frac sL\Bigr)\Bigr)e^{-sy}\sin(sv)\,ds .
\end{equation}
Since $W=1$ on $(-\infty,L]$, the part $s<0$ of the first integral equals $+\int_0^\infty e^{-sy}\sin(sv)\,ds$ after $s\mapsto-s$; because $1-W(s/L)=0$ for $s\le L$, this combines with the second integral into $\int_0^\infty W(\frac sL)e^{-sy}\sin(sv)\,ds=\int_0^{2L}W(\frac sL)e^{-sy}\sin(sv)\,ds$. Hence
\begin{equation}
\Im\,\Psi'(y)=-\int_0^{2L}W\!\Bigl(\frac sL\Bigr)\bigl(e^{sy}-e^{-sy}\bigr)\sin(sv)\,ds=-2\int_0^{2L}W\!\Bigl(\frac sL\Bigr)\sinh(sy)\sin(sv)\,ds .
\end{equation}
Integrating over $[0,y]$ and using $\int_0^y\sinh(sy')\,dy'=(\cosh(sy)-1)/s$ gives $\frac1\pi\Im\Psi(y)=\mathcal J_L(y;v)$. The integrand of $\mathcal J_L$ is an entire function of $v$ on the compact interval $[0,2L]$, which gives the last assertions; at $v=0$ the integral defines the continuous extension.
\end{proof}

For a zero with $\beta>\frac12$ and $y=\beta-\frac12$, \eqref{eq:taucw-on-off} and \cref{prop:taucw-pair} give, after pairing,
\begin{equation}\label{eq:taucw-Ypairs}
\mathcal Y_L(t)=\sum_{\rho:\ \beta>1/2}\mathcal J_L\bigl(\beta-\tfrac12;\gamma-t\bigr),
\end{equation}
with multiplicity, the sum running over all zeros with $\beta>\frac12$ (of either sign of $\gamma$). The series converges absolutely and locally uniformly in $t$ by \cref{cor:taucw-pair}(c), because $\sum_\rho(\gamma-t)^{-2}<\infty$ and $e^{2Ly}\le X$.

\begin{corollary}[Spectrum, mass and size of a pair]\label{cor:taucw-pair}
Let $L>0$ and $y>0$.
\begin{enumerate}[label=(\alph*),leftmargin=*]
\item \emph{Spectrum.} $\mathcal J_L(y;v)=-\Im\int_0^{2L}\mathscr R_{L,y}(u)e^{iuv}du$: the pair has the frequencies $0<u<2L$ only, with density $\mathscr R_{L,y}(u)>0$, which is strictly increasing in $y$ for every fixed $u$ and tends to $0$ as $y\to0$. In particular $\mathscr R_{L,y}(u)\le\frac{y^2}{\pi}\,u\,W(\tfrac uL)\cosh(yu)$.
\item \emph{Exact mass.}
\begin{equation}\label{eq:taucw-mass}
\int_{\mathbb R}\bigl|\mathcal J_L(y;v)\bigr|^2dv=\pi\int_0^{2L}\mathscr R_{L,y}(u)^2du=\frac4\pi\int_0^{2L}W\!\Bigl(\frac uL\Bigr)^2\frac{(\cosh yu-1)^2}{u^2}\,du,
\end{equation}
which is strictly increasing in $y$, and
\begin{align}
\int_{\mathbb R}|\mathcal J_L(y;v)|^2dv&\ge\frac{(\cosh yL-1)^2}{3\pi L}\qquad\text{for all }y>0,\label{eq:taucw-mass-lb}\\
\int_{\mathbb R}|\mathcal J_L(y;v)|^2dv&\ge\frac{(\cosh(2Ly-1)-1)^2}{3\pi L^4y^3}\qquad\text{if }Ly\ge1 .\label{eq:taucw-mass-lb2}
\end{align}
\item \emph{Size.} For $v\ne0$,
\begin{equation}\label{eq:taucw-pair-size}
|\mathcal J_L(y;v)|\le\min\Bigl\{\frac2\pi\int_0^{2L}\frac{\cosh yu-1}{u}\,du,\ \frac{e^{2Ly}+1}{\pi L^2v^2}\Bigr\}.
\end{equation}
\item \emph{Localization.} If $Ly\ge2$, then $\displaystyle\int_{|v|\le4y}|\mathcal J_L(y;v)|^2dv\ge\frac{(\cosh(2Ly-1)-1)^2}{6\pi L^4y^3}$.
\end{enumerate}
\end{corollary}
\begin{proof}
(a) $\cosh z-1=\frac{z^2}2\cdot\frac{2(\cosh z-1)}{z^2}$ and $2(\cosh z-1)\le z^2\cosh z$; monotonicity in $y$ is that of $\cosh$ on $[0,\infty)$. (b) For $f(v)=\int_0^{2L}a(u)\sin(uv)\,du$ with $a\in L^2(0,2L)$, extend $a$ to an odd function $b$ on $\mathbb R$; then $2if=\int_{\mathbb R}b(u)e^{iuv}du$ and Plancherel gives $\int|f|^2=\frac14\cdot2\pi\cdot2\int_0^{2L}a^2=\pi\int_0^{2L}a^2$. For \eqref{eq:taucw-mass-lb} restrict the integral to $u\in[L,\frac43L]$, where $W(u/L)\ge\frac23$, $u^{-2}\ge\frac9{16L^2}$ and $\cosh yu-1\ge\cosh yL-1$: this gives $\int W^2(\cosh yu-1)^2u^{-2}du\ge\frac49\cdot\frac9{16L^2}\cdot\frac L3(\cosh yL-1)^2=\frac{(\cosh yL-1)^2}{12L}$, and $\frac4\pi\cdot\frac1{12L}=\frac1{3\pi L}$. For \eqref{eq:taucw-mass-lb2} put $\eta_0=1/y\le L$ and restrict to $u\in[2L-\eta_0,2L]\subset[L,2L]$, where $W(u/L)=(2L-u)/L$, $u^{-2}\ge(2L)^{-2}$ and $\cosh yu-1\ge\cosh(2Ly-1)-1$: the integral is at least $(\cosh(2Ly-1)-1)^2\frac1{4L^2}\cdot\frac{\eta_0^3}{3L^2}=\frac{(\cosh(2Ly-1)-1)^2}{12L^4y^3}$. (c) The first bound is $|\sin|\le1$ and $W\le1$. For the second, $\mathcal J_L(y;v)=\frac1\pi\Im\Psi(y)$ and, for $v\ne0$, $|\mathfrak G_L(x+iv)|\le2e^{2Lx}/(Lv^2)$ if $x\ge0$ and $\le2e^{Lx}/(Lv^2)$ if $x\le0$; hence $|\Psi(y)|\le\frac{e^{2Ly}-1}{L^2v^2}+\frac2{L^2v^2}$. (d) By \eqref{eq:taucw-pair-size}, $\int_{|v|>V}|\mathcal J_L|^2\le\frac{2(e^{2Ly}+1)^2}{3\pi^2L^4V^3}$. With $z=2Ly\ge4$ one has $(e^z+1)/(\cosh(z-1)-1)\le6.2$, so the ratio of this tail, for $V=4y$, to the lower bound in \eqref{eq:taucw-mass-lb2} is at most $\frac{2\cdot6.2^2y^3}{\pi\,(4y)^3}<\frac12$; subtract it from the total mass.
\end{proof}

\begin{proposition}[Energy of finitely many pairs]\label{prop:taucw-energy}
Let $\mathcal N$ be a finite set of pairs, indexed by $\mathbf p$ in a fixed order, $(y_{\mathbf p},\gamma_{\mathbf p})$ with $y_{\mathbf p}>0$, and $\mathcal J_{\mathcal N}(t)=\sum_{\mathbf p\in\mathcal N}\mathcal J_L(y_{\mathbf p};\gamma_{\mathbf p}-t)$. Then
\begin{align}
\int_{\mathbb R}|\mathcal J_{\mathcal N}(t)|^2\dd&=\pi\int_0^{2L}\Bigl|\sum_{\mathbf p\in\mathcal N}\mathscr R_{L,y_{\mathbf p}}(u)e^{iu\gamma_{\mathbf p}}\Bigr|^2du\label{eq:taucw-energy}\\
&=\sum_{\mathbf p}\int_{\mathbb R}|\mathcal J_L(y_{\mathbf p};v)|^2dv+2\pi\sum_{\mathbf p<\mathbf q}\int_0^{2L}\mathscr R_{L,y_{\mathbf p}}\mathscr R_{L,y_{\mathbf q}}\cos\bigl(u(\gamma_{\mathbf p}-\gamma_{\mathbf q})\bigr)du .\label{eq:taucw-energy-x}
\end{align}
\end{proposition}
\begin{proof}
$\mathcal J_{\mathcal N}(t)=-\Im\mathcal F_{\mathcal N}(t)$ with $\mathcal F_{\mathcal N}(t)=\int_0^{2L}\mathcal H(u)e^{-iut}du$, $\mathcal H(u)=\sum_{\mathbf p}\mathscr R_{L,y_{\mathbf p}}(u)e^{iu\gamma_{\mathbf p}}$. Since $\widehat{\mathcal F_{\mathcal N}}$ and $\widehat{\overline{\mathcal F_{\mathcal N}}}$ are supported on $(0,2L)$ and $(-2L,0)$, $\int\mathcal F_{\mathcal N}^2=0$, while $\int|\mathcal F_{\mathcal N}|^2=2\pi\int|\mathcal H|^2$. From $|\Im z|^2=\frac12(|z|^2-\Re z^2)$ one gets the first equality; the second is the expansion of $|\mathcal H|^2$.
\end{proof}
Thus the total $L^2$ mass of finitely many pairs is the energy, on $0<u<2L$, of the weighted exponential sum $\sum_{\mathbf p}\mathscr R_{L,y_{\mathbf p}}(u)e^{iu\gamma_{\mathbf p}}$ over the imaginary parts; the cross terms in \eqref{eq:taucw-energy-x} are cosine transforms of the product of two densities at the difference of the two imaginary parts.

\begin{theorem}[The prime sum equals $\mathcal Y_L$ up to $O(\log t)$]\label{thm:taucw-P}
There are absolute constants $C_0,C$ such that, for every $X\ge2$, $L=\log X$, and every $t>0$ that is not the imaginary part of a zero,
\begin{equation}\label{eq:taucw-P-exact}
P^W_X(t)=S(t)+\mathcal O_L(t)+\mathcal Y_L(t)+\mathcal E_L(t),\qquad|\mathcal E_L(t)|\le C_0,
\end{equation}
and $|S(t)|+|\mathcal O_L(t)|\le C\log(2+t)$. Consequently
\begin{equation}\label{eq:taucw-P-Y}
\bigl|P^W_X(t)-\mathcal Y_L(t)\bigr|\le C\log(2+t)\qquad\text{uniformly in }X\ge2 .
\end{equation}
No hypothesis on the zeros is used.
\end{theorem}
\begin{proof}
The proof of \cref{lem:taucw-expl}(b) gives, for every $t>0$ that is not the imaginary part of a zero, \eqref{eq:taucw-piS} and $\Im\,\mathfrak I_L(\frac12-it)=-\pi\Pi^W_X(t)+\mathscr A(t)$. Eliminating $\Pi_X^W$ in $P^W_X=\Phi^W_X-\Pi^W_X$,
\begin{equation}
P^W_X(t)=S(t)+\frac1\pi\Im\sum_\rho\mathfrak I_L\bigl(\rho-\tfrac12-it\bigr)+\frac1\pi\Im\sum_{m\ge1}\mathfrak I_L\bigl(-2m-\tfrac12-it\bigr)-\frac{\mathscr A(t)}\pi .
\end{equation}
By \cref{lem:taucw-kernel}(c) and \eqref{eq:taucw-on-off} the sum over $\rho$ splits into $\mathcal O_L+\mathcal Y_L$. For the trivial zeros, $|\mathfrak G_L(u)|\le2e^{L\Re u}/(L|u|^2)$ with $|u|\ge|\Re u|\ge\frac52$ on the ray ending at $-2m-\frac12-it$ gives $|\mathfrak I_L(-2m-\frac12-it)|\le8e^{-L(2m+1/2)}/L^2$, summable in $m$ and bounded for $L\ge\log2$. The function $\mathscr A(t)=\Im E_1((-\frac12+it)\log2)$ is bounded on $(0,\infty)$: by $\sqrt2/(t\log2)$ for $t\ge1$ (proof of \cref{lem:taucw-expl}), and by continuity of $E_1$ up to the cut on $(0,1]$. This is \eqref{eq:taucw-P-exact}. The bound $|S(t)|\ll\log(2+t)$ is classical \cite{Titchmarsh}. For $\mathcal O_L$, \cref{lem:taucw-kernel}(b) gives $|\mathfrak H_L(v)-\mathbf 1_{v>0}|\le\min\{C_{\mathfrak K},3/(\pi L^2v^2)\}$, hence
\begin{equation}
|\mathcal O_L(t)|\le C_{\mathfrak K}\,\#\{\rho:|\gamma-t|\le1\}+\frac3{\pi\log^22}\sum_{|\gamma-t|>1}\frac1{(\gamma-t)^2}\ll\log(2+t),
\end{equation}
because every interval of length $1$ contains $O(\log(2+|\gamma|))$ imaginary parts of zeros; this bound is uniform in $L\ge\log2$.
\end{proof}

\begin{theorem}[Zeros off the line are exactly what the weighted norm sees]\label{thm:taucw-equiv}
Let $X_j\to\infty$ be any unbounded sequence and $L_j=\log X_j$. The following are equivalent.
\begin{enumerate}[label=(\roman*),leftmargin=*]
\item Every nontrivial zero of $\zeta$ has real part $\frac12$.
\item $\displaystyle\sup_j\int_0^\infty\frac{|P^W_{X_j}(t)|^2}{(1+t)^3}\dd<\infty$.
\item $\displaystyle\sup_j\int_0^\infty\frac{|\mathcal Y_{L_j}(t)|^2}{(1+t)^3}\dd<\infty$.
\end{enumerate}
\end{theorem}
\begin{proof}
(i)$\Rightarrow$(iii): if all zeros are on the line, every term of $\mathcal Y_L$ in \eqref{eq:taucw-on-off} vanishes, so $\mathcal Y_{L}\equiv0$. (iii)$\Rightarrow$(ii): by \eqref{eq:taucw-P-Y}, $|P^W_{X_j}|^2\le2|\mathcal Y_{L_j}|^2+2C^2\log^2(2+t)$ outside the imaginary parts of the zeros and $\int_0^\infty\log^2(2+t)(1+t)^{-3}dt<\infty$. (ii)$\Rightarrow$(i): \cref{thm:v7-normality}.
\end{proof}
The hypothesis (ii) is that of \cref{thm:v7-normality} for the Selberg weight; \cref{thm:taucw-equiv} shows that it is the same as boundedness of the weighted norm of $\mathcal Y_{L_j}$, which by \eqref{eq:taucw-Ypairs} is the sum of the profiles of \cref{prop:taucw-pair}. The statement (iii) is not proved in this note.

\begin{remark}[Reading for the taucw phase]\label{rem:taucw-detect}
For $0<\theta\le1$, $X=t^\theta\ge2$ and $t\ge t_0$ not the imaginary part of a zero, \cref{lem:taucw-pole} and \cref{prop:taucw-decomp} give
\begin{equation}
F^{\mathrm w}_\theta(t)-N(t)=\mathcal O_{L(t)}(t)+\mathcal Y_{L(t)}(t)+O\Bigl(\frac1t\Bigr),\qquad L(t)=\theta\log t .
\end{equation}
The phase that defines the taucw points is built from the primes alone; the real parts of the zeros appear in its error exactly through \eqref{eq:taucw-Ypairs}, each off-line pair with the profile \eqref{eq:taucw-pair}, whose frequencies fill $(0,2L)$ with the positive density $\mathscr R_{L,y}$, whose mass is \eqref{eq:taucw-mass}, and whose size at distance $|\gamma-t|$ from the imaginary part is bounded by \eqref{eq:taucw-pair-size}. The first-order term in $y$ cancels in the pair (the density is $O(y^2)$), the factor $e^{2Ly}=X^{2y}$ is the amplification, and $\mathcal O_L$, which depends on the imaginary parts only, is not affected by the real parts (\cref{rem:taucw-off}). \Cref{thm:taucw-P} shows that this holds for the prime sum $P^W_X$ at every height and for every cutoff $X\ge2$, not only at $X=t^\theta$.
\end{remark}

\begin{proposition}[Shared imaginary parts cost displacement]\label{prop:taucw-shared}
Let $0<\theta\le1$, let $\gamma_1\le\gamma_2\le\cdots$ be the imaginary parts of the zeros with $\gamma>0$, repeated with multiplicity, and put $\tau_j=\tau_j^{\mathrm w}(\theta)$.
\begin{enumerate}[label=(\alph*),leftmargin=*]
\item If $\rho=\beta+i\gamma$ with $\gamma>0$ has $\beta\ne\frac12$ or has multiplicity at least $2$, then $\gamma=\gamma_j=\gamma_{j+1}$ for some $j$. Conversely, if $\gamma_{j-1}<\gamma_j<\gamma_{j+1}$, the zero with imaginary part $\gamma_j$ is simple and lies on the critical line.
\item If $\gamma_j=\gamma_{j+1}$ and $j-\frac12>F^{\mathrm w}_\theta(2\pi)$, then
\begin{equation}\label{eq:taucw-shared}
0<\tau_{j+1}-\tau_j\le|\tau_j-\gamma_j|+|\tau_{j+1}-\gamma_{j+1}| .
\end{equation}
\item For $\ell\ge1$ let $\mathcal Q_\ell$ be the set of $j<\ell$ with $\gamma_j=\gamma_{j+1}$ and $j-\frac12>F^{\mathrm w}_\theta(2\pi)$. Then
\begin{equation}\label{eq:taucw-shared-sum}
\sum_{j\in\mathcal Q_\ell}(\tau_{j+1}-\tau_j)\le2\sum_{j\le\ell}|\tau_j-\gamma_j|=2\int_{\mathbb R}\bigl|N_{[\ell]}-N^{\mathrm w}_{[\ell]}\bigr|\dd,\qquad N^{\mathrm w}_{[\ell]}(t)=\#\{j\le\ell:\tau_j\le t\},
\end{equation}
and, for every $\mathsf s_\circ>0$ and every $c>0$,
\begin{align}
\#\{j\in\mathcal Q_\ell:\ \tau_{j+1}-\tau_j\ge\mathsf s_\circ\}&\le\frac2{\mathsf s_\circ}\sum_{j\le\ell}|\tau_j-\gamma_j|,\label{eq:taucw-shared-count}\\
\#\Bigl\{j\in\mathcal Q_\ell:\ \frac{\tau_{j+1}-\tau_j}{2\pi}\log\frac{\tau_j}{2\pi}\ge c\Bigr\}&\le\frac{\log(\tau_\ell/2\pi)}{\pi c}\sum_{j\le\ell}|\tau_j-\gamma_j| .\label{eq:taucw-shared-count2}
\end{align}
\end{enumerate}
\end{proposition}
\begin{proof}
(a) If $\beta\ne\frac12$, then $1-\bar\rho\ne\rho$ is a zero with the same imaginary part, so the imaginary part occurs at least twice in the list; a zero of multiplicity $m\ge2$ occurs $m$ times. If the imaginary part $\gamma_j$ occurs once, the zero is simple, and $1-\bar\rho$, having the same imaginary part, equals $\rho$, so $\beta=\frac12$. (b) $\tau_j<\tau_{j+1}$ by \cref{prop:taucw-ord}, and $\tau_{j+1}-\tau_j=(\tau_{j+1}-\gamma_{j+1})+(\gamma_j-\tau_j)$ because $\gamma_j=\gamma_{j+1}$. (c) Sum \eqref{eq:taucw-shared} over $\mathcal Q_\ell$; each index occurs in at most two of the pairs $(j,j+1)$. The equality is the identity $\int|N_{[\ell]}-N^{\mathrm w}_{[\ell]}|\dd=\sum_{j\le\ell}|\gamma_j-\tau_j|$, which holds for any two nondecreasing lists of length $\ell$, as in \eqref{eq:taucw-exc}. \eqref{eq:taucw-shared-count} is Chebyshev's inequality for \eqref{eq:taucw-shared-sum}. For \eqref{eq:taucw-shared-count2}, $\tau_j\ge2\pi$ and $j\le\ell$ give $\log(\tau_j/2\pi)\le\log(\tau_\ell/2\pi)$, so every $j$ in the set satisfies $\tau_{j+1}-\tau_j\ge2\pi c/\log(\tau_\ell/2\pi)$; apply \eqref{eq:taucw-shared-count} with this $\mathsf s_\circ$.
\end{proof}
The taucw points are strictly increasing and defined without the zeros, while a zero off the line, like a multiple zero, makes an imaginary part occur at least twice (\cref{prop:taucw-shared}(a)); by \eqref{eq:taucw-shared-sum} the separation of the points at such an index is paid for by displacement, and the total displacement is the $L^1$ mass of the integer-valued discrepancy $N_{[\ell]}-N^{\mathrm w}_{[\ell]}$ of \cref{prop:taucw-S} (with $\tau_j^{\mathrm w}(\theta)$ in place of $\tau_j^{\mathrm w}$). Where the points are within a small fraction of the mean spacing of the corresponding imaginary parts, \eqref{eq:taucw-shared-count2} bounds the number of shared imaginary parts whose taucw gap exceeds that fraction.


\subsection{Moments beyond the support \texorpdfstring{$n\le t$}{n <= t}, the increment to the full cutoff, and the harmonics}\label{sec:taucw-beyond}
Every statement of this subsection is unconditional. Here $W$ is the Selberg weight, except in \cref{lem:taucw-tail} and \cref{lem:taucw-gram}, where $W$ is any weight of \eqref{eq:weights}. It extends the second moment of \cref{thm:taucw-var} past the support $n\le t$ by logarithmic factors, using that the Selberg weight vanishes linearly at the end of its support (\cref{lem:taucw-mvsel}, \cref{prop:taucw-beyond}); computes the diagonal part of the increment of the phase from $\theta=\frac12$ to $\theta=1$ (\cref{prop:taucw-inc1}); writes the Selberg weight as a difference of two Fej\'er-type autocorrelations and identifies the kernel $\mathfrak K_L$ of \eqref{eq:kernel-G} as $(2\pi)^{-1}$ times its Fourier transform (\cref{lem:taucw-fejer}); bounds the energy of $\mathcal Y_L$ at moderate cutoffs from the primes (\cref{prop:taucw-Yenergy}); shows that zeros of bounded height leave the taucw error on dyadic intervals untouched up to $O_\theta(\#\mathscr Z^\diamond\,T^{\theta-2}\log^{-2}T)$ (\cref{prop:taucw-bh}); reduces the weighted norm of \cref{thm:taucw-equiv} to a finite range of $t$ and writes it as a Gram form (\cref{lem:taucw-tail}, \cref{cor:taucw-trunc}, \cref{lem:taucw-gram}); gives the exact arithmetic form of the harmonics $e^{2\pi i\mathfrak hF^{\mathrm w}_\theta}$ (\cref{lem:taucw-harmar}); and records what the rank--trace inequality of Alp\"oge and Furman \cite{AlpogeFurman2026} would give at twice their bandwidth (\cref{prop:taucw-AF2}). The coefficients $\mathfrak k_n=\Lambda(n)/(\pi\sqrt n\log n)$ are those of \cref{lem:taucw-freeze}, and $\mathrm B_1$ denotes the Meissel--Mertens constant, $\sum_{p\le x}p^{-1}=\log\log x+\mathrm B_1+o(1)$ \cite{Mertens1874}.

\begin{lemma}[Mean value with the Selberg endpoint]\label{lem:taucw-mvsel}
For $X\ge2$ put $A^{\mathrm w}_X(t)=\pi\sum_{n\le X^2}\mathfrak k_nW\bigl(\frac{\log n}{\log X}\bigr)n^{it}$. Then, for $T\ge1$,
\begin{equation}\label{eq:taucw-mvsel}
\int_T^{2T}|A^{\mathrm w}_X(t)|^2\dd=T\sum_{n\le X^2}\pi^2\mathfrak k_n^2W\Bigl(\frac{\log n}{\log X}\Bigr)^2+O\Bigl(\frac{X^2}{\log^3X}\Bigr),\qquad
\Bigl|\int_T^{2T}A^{\mathrm w}_X(t)^2\dd\Bigr|\ll\frac{X^2}{\log^4X},
\end{equation}
with absolute implied constants, and, as $X\to\infty$,
\begin{equation}\label{eq:taucw-mvsel-diag}
\sum_{n\le X^2}\pi^2\mathfrak k_n^2W\Bigl(\frac{\log n}{\log X}\Bigr)^2=\log\log X+\mathrm B_1+4\log2-\frac52+\sum_p\sum_{r\ge2}\frac1{r^2p^r}+o(1).
\end{equation}
\end{lemma}
\begin{proof}
Write $c^{\mathrm w}_n=\pi\mathfrak k_nW(\log n/\log X)$, so that $0\le c^{\mathrm w}_n\le\Lambda(n)/(\sqrt n\log n)$. The mean value theorem of Montgomery and Vaughan \cite{MontgomeryVaughan} gives $\int_T^{2T}|A^{\mathrm w}_X|^2\dd=\sum_n(c^{\mathrm w}_n)^2(T+O(n))$. Here $n(c^{\mathrm w}_n)^2=W(\log p/\log X)^2$ for $n=p$, and $n(c^{\mathrm w}_n)^2\le r^{-2}\le1$ for $n=p^r$, $r\ge2$, of which there are $\ll X/\log X$ with $p^r\le X^2$. The primes $p\le X$ contribute $\pi(X)\ll X/\log X$ to $\sum_nn(c^{\mathrm w}_n)^2$, and the primes in $(X,X^2]$ contribute $\ll X^2/\log^3X$ by the elementary estimate \eqref{eq:wprime2} (Chebyshev's bound and partial summation); since $X/\log X\ll X^2/\log^3X$, this is the first estimate of \eqref{eq:taucw-mvsel}. For the second, $\int_T^{2T}(nm)^{it}\dd\ll1/\log(nm)\le1/\log4$, so $|\int_T^{2T}(A^{\mathrm w}_X)^2\dd|\ll(\sum_nc^{\mathrm w}_n)^2$, and $\sum_nc^{\mathrm w}_n\ll X/\log^2X$ by the computation in the proof of \cref{lem:taucw-amp} (with $X$ in place of $t^\theta$). For \eqref{eq:taucw-mvsel-diag}, the primes $p\le X$ give $\log\log X+\mathrm B_1+o(1)$, and Mertens' theorem with partial summation turns $\sum_{X<p\le X^2}(2-\log p/\log X)^2p^{-1}$ into $\int_1^2(2-v)^2\,dv/v+o(1)=4\log2-\frac52+o(1)$. A prime power $p^r$, $r\ge2$, has $(c^{\mathrm w}_{p^r})^2=(r^2p^r)^{-1}W(r\log p/\log X)^2$, which tends to $(r^2p^r)^{-1}$ for fixed $p^r$, and the series of $(r^2p^r)^{-1}$ converges; dominated convergence gives the last term.
\end{proof}

\begin{proposition}[Second moment past the support $n\le t$]\label{prop:taucw-beyond}
Fix $0\le\mathsf z<\frac32$, put $X(t)=t^{1/2}(\log t)^{\mathsf z}$ for $t\ge t_0(\mathsf z)$, where $t_0(\mathsf z)$ is so large that $X(t)\ge2$ and $X$ is increasing on $[t_0(\mathsf z),\infty)$, and
\begin{equation}\label{eq:taucw-beyond-def}
\Phi^{\mathrm w}_{\langle\mathsf z\rangle}(t)=-\sum_{n\le X(t)^2}\mathfrak k_n\,W\Bigl(\frac{\log n}{\log X(t)}\Bigr)\sin(t\log n),
\end{equation}
whose prime-power support is $n\le X(t)^2=t(\log t)^{2\mathsf z}$. Then
\begin{equation}\label{eq:taucw-beyond}
\frac1T\int_T^{2T}\bigl|\Phi^{\mathrm w}_{\langle\mathsf z\rangle}(t)\bigr|^2\dd=\frac1{2\pi^2}\Bigl[\log\log T+\mathrm B_1+3\log2-\frac52+\sum_p\sum_{r\ge2}\frac1{r^2p^r}\Bigr]+o(1).
\end{equation}
For $\mathsf z=0$ this is the phase $\Phi^{\mathrm w}_{1/2}$ of \cref{def:taucw}.
\end{proposition}
\begin{proof}
Put $L_s=\log X(s)=\frac12\log s+\mathsf z\log\log s$, and let $T_0(\mathsf z)\ge t_0(\mathsf z)$ be so large that $L_s\asymp\log T$ and $0<L_s'\le1/s$ for $s\in[T,2T]$ and $T\ge T_0(\mathsf z)$. Let $\Phi^{\mathrm{fr}}_{\langle\mathsf z\rangle,T}(t)$ be \eqref{eq:taucw-beyond-def} with $L_T$ in place of $\log X(t)$ in the weight and with the sum over $n\le X(T)^2$. Since $\sin(t\log n)=\Im n^{it}$, $\Phi^{\mathrm{fr}}_{\langle\mathsf z\rangle,T}=-\pi^{-1}\Im A^{\mathrm w}_{X(T)}$, and $|\Im z|^2=\frac12|z|^2-\frac12\Re z^2$. By \cref{lem:taucw-mvsel} with $X=X(T)$, $X(T)^2/\log^3X(T)\ll_{\mathsf z}T(\log T)^{2\mathsf z-3}=o(T)$ and $X(T)^2/\log^4X(T)=o(T)$, so
\begin{equation}
\int_T^{2T}|\Phi^{\mathrm{fr}}_{\langle\mathsf z\rangle,T}|^2\dd=\frac T{2\pi^2}\Bigl[\log\log X(T)+\mathrm B_1+4\log2-\frac52+\sum_p\sum_{r\ge2}\frac1{r^2p^r}\Bigr]+o(T),
\end{equation}
and $\log\log X(T)=\log\log T-\log2+o(1)$; this is the right side of \eqref{eq:taucw-beyond} times $T$.

It remains to compare $\Phi^{\mathrm w}_{\langle\mathsf z\rangle}$ with $\Phi^{\mathrm{fr}}_{\langle\mathsf z\rangle,T}$ on $[T,2T]$. For fixed $n$ the function $s\mapsto W(\log n/L_s)$ is Lipschitz, with derivative $(\log n)L_s'/L_s^2$ when $L_s<\log n<2L_s$ and $0$ when $\log n<L_s$ or $\log n>2L_s$ (at the two break points it is continuous, and terms entering the support enter with weight $0$). Hence, for $T\le t\le2T$,
\begin{equation}
\Phi^{\mathrm w}_{\langle\mathsf z\rangle}(t)-\Phi^{\mathrm{fr}}_{\langle\mathsf z\rangle,T}(t)=-\int_T^{2T}\mathbf 1_{(s,2T]}(t)\sum_{n}b^{\langle\mathsf z\rangle}_n(s)\sin(t\log n)\,\mathrm ds,
\end{equation}
where
\begin{equation}
b^{\langle\mathsf z\rangle}_n(s)=\frac{\Lambda(n)L_s'}{\pi\sqrt n\,L_s^2}\mathbf 1_{\{L_s<\log n<2L_s\}} .
\end{equation}
By Minkowski's integral inequality and the mean value theorem of Montgomery and Vaughan on $(s,2T]$, the $L^2(T,2T)$ norm of the difference is at most $\int_T^{2T}\bigl(\sum_nb^{\langle\mathsf z\rangle}_n(s)^2(T+O(n))\bigr)^{1/2}\mathrm ds$. Here $\sum_nb^{\langle\mathsf z\rangle}_n(s)^2\le(\pi L_s^2)^{-2}L_s'^2\sum_{n\le X(s)^2}\Lambda(n)^2/n\ll1/(T^2\log^2T)$ and $\sum_nnb^{\langle\mathsf z\rangle}_n(s)^2\le(\pi L_s^2)^{-2}L_s'^2\sum_{n\le X(2T)^2}\Lambda(n)^2\ll X(2T)^2/(T^2\log^3T)\ll_{\mathsf z}(\log T)^{2\mathsf z-3}/T$. Therefore
\begin{equation}
\bigl\|\Phi^{\mathrm w}_{\langle\mathsf z\rangle}-\Phi^{\mathrm{fr}}_{\langle\mathsf z\rangle,T}\bigr\|_{L^2(T,2T)}\ll_{\mathsf z}\sqrt T\Bigl(\frac1{\log T}+(\log T)^{\mathsf z-3/2}\Bigr),
\end{equation}
and, since $\|\Phi^{\mathrm{fr}}_{\langle\mathsf z\rangle,T}\|_{L^2(T,2T)}\ll\sqrt{T\log\log T}$, the squares of the two norms differ by $O_{\mathsf z}\bigl(T\sqrt{\log\log T}\,((\log T)^{-1}+(\log T)^{\mathsf z-3/2})\bigr)=o(T)$, because $\mathsf z<\frac32$.
\end{proof}
\begin{remark}[What the extension is]\label{rem:taucw-beyond}
\Cref{thm:taucw-var} covers fixed $\theta<\frac12$, that is, support $n\le t^{2\theta}$ with $2\theta<1$. \Cref{prop:taucw-beyond} reaches $\theta=\frac12$ and support $n\le t(\log t)^{3-\varepsilon}$ for every $\varepsilon>0$, with the same diagonal asymptotics: the Montgomery--Vaughan remainder $\sum_nn(c^{\mathrm w}_n)^2$ is $\ll X^2/\log^3X$, instead of $\asymp X^2/\log X$ for the sharp cutoff, because the Selberg weight vanishes linearly at $\log n=2\log X$, the mechanism of \cref{thm:moderate-log}. A support $n\le t^{1+\delta}$ with fixed $\delta>0$ is not reached: there the bound of the mean value theorem for the remainder is $\gg T$, and the off-diagonal sum is not controlled by this argument.
\end{remark}

For the increment of the phase from $\theta=\frac12$ to $\theta=1$ put, as in the proof of \cref{thm:taucw-signed} with $(\theta_-,\theta_+)=(\frac12,1)$,
\begin{equation}\label{eq:taucw-inc1-def}
\begin{gathered}
\mathscr K_n(t)=W\Bigl(\frac{\log n}{\log t}\Bigr)-W\Bigl(\frac{2\log n}{\log t}\Bigr),\\
\Phi^{\mathrm w}_{[1/2,1]}(t)=\Phi^{\mathrm w}_1(t)-\Phi^{\mathrm w}_{1/2}(t)=-\sum_{n\le t^2}\mathfrak k_n\mathscr K_n(t)\sin(t\log n),
\end{gathered}
\end{equation}
and define the off-diagonal part
\begin{equation}\label{eq:taucw-inc1-off}
\mathscr E^{\mathrm{off}}_T=\frac1T\sum_{\substack{n,m\le4T^2\\n\ne m}}\mathfrak k_n\mathfrak k_m\int_T^{2T}\mathscr K_n(t)\mathscr K_m(t)\sin(t\log n)\sin(t\log m)\dd ,
\end{equation}
a finite sum (the coefficients $\mathscr K_n(t)$ vanish for $n>t^2$).

\begin{proposition}[Diagonal of the increment to the full cutoff]\label{prop:taucw-inc1}
As $T\to\infty$,
\begin{equation}\label{eq:taucw-inc1}
\begin{gathered}
\frac1T\int_T^{2T}\bigl|\Phi^{\mathrm w}_{[1/2,1]}(t)\bigr|^2\dd=\frac{\mathfrak S(\frac12,1)}{2\pi^2}+\mathscr E^{\mathrm{off}}_T+o(1),\\
\mathfrak S\bigl(\tfrac12,1\bigr)=\int_0^\infty\bigl[W(v)-W(2v)\bigr]^2\frac{dv}v=5\log2-3,
\end{gathered}
\end{equation}
so that the diagonal constant is $\mathfrak S(\frac12,1)/(2\pi^2)=0.0235944564682\ldots$. The primes $T<p\le T^2$ contribute $\frac1{2\pi^2}\int_1^2(2-v)^2\frac{dv}v=\frac{4\log2-5/2}{2\pi^2}=0.0138095059924\ldots$ to it.
\end{proposition}
\begin{proof}
Expanding the square, $\frac1T\int_T^{2T}|\Phi^{\mathrm w}_{[1/2,1]}|^2\dd$ is the sum of $\mathscr E^{\mathrm{off}}_T$ and the diagonal $\frac1T\sum_n\mathfrak k_n^2\int_T^{2T}\mathscr K_n(t)^2\sin^2(t\log n)\dd$. Write $\sin^2=\frac12-\frac12\cos(2t\log n)$. For fixed $n$, $t\mapsto\mathscr K_n(t)$ is the difference of two nondecreasing functions with values in $[0,1]$, so $|\mathscr K_n|\le1$ and $\mathscr K_n^2$ has total variation at most $4$ on $[T,2T]$; one integration by parts gives $|\int_T^{2T}\mathscr K_n^2\cos(2t\log n)\dd|\le(2+4)/(2\log n)=3/\log n$, and $\sum_n\mathfrak k_n^2/\log n<\infty$, so the cosine terms contribute $O(1/T)$. For $t\in[T,2T]$ and $n\le4T^2$, the Lipschitz property of $W$ gives $|\mathscr K_n(t)-\mathscr K_n(T)|\le3\log n\,\bigl(\frac1{\log T}-\frac1{\log t}\bigr)\ll1/\log T$; moreover $\mathscr K_n(t)=0$ unless $t^{1/2}<n\le t^2$, and $\sum_{T^{1/2}<n\le4T^2}\mathfrak k_n^2=O(1)$ by Mertens' theorem. The diagonal is therefore $\frac12\sum_n\mathfrak k_n^2\mathscr K_n(T)^2+O(1/\log T)$. The prime powers $p^r$ with $r\ge2$ contribute at most $\frac1{2\pi^2}\sum_{p^r>T^{1/2}}(r^2p^r)^{-1}=o(1)$. For the primes, $\mathscr K_p(T)=W_{[1/2,1]}(\log p/\log T)$ with $W_{[1/2,1]}(v)=W(v)-W(2v)$, which is continuous, piecewise linear and supported in $[\frac12,2]$: $W_{[1/2,1]}(v)=0$ for $v\le\frac12$, $2v-1$ for $\frac12<v\le1$, $2-v$ for $1<v\le2$, and $0$ for $v>2$. Mertens' theorem and Stieltjes integration by parts, as in the proof of \cref{thm:v20-anticorr}, give $\frac1{2\pi^2}\sum_pW_{[1/2,1]}(\log p/\log T)^2/p=\frac1{2\pi^2}\int_0^\infty W_{[1/2,1]}(v)^2dv/v+o(1)$, and
\begin{equation}
\int_{1/2}^1\frac{(2v-1)^2}v\,dv+\int_1^2\frac{(2-v)^2}v\,dv=\Bigl(\log2-\frac12\Bigr)+\Bigl(4\log2-\frac52\Bigr)=5\log2-3 .
\end{equation}
The primes $T<p\le T^2$ are those with $1<v\le2$, which gives the second constant.
\end{proof}
\begin{remark}[Scope of \cref{prop:taucw-inc1}]\label{rem:taucw-inc1}
The constant $\mathfrak S(\frac12,1)$ is the value at $(\theta_-,\theta_+)=(\frac12,1)$ of the integral \eqref{eq:taucw-S}. For $\theta_+<\frac12$, \cref{thm:taucw-inc} shows that the corresponding off-diagonal part tends to $0$; for $\theta_+=1$ the frequencies $\log n$ with $n$ up to $4T^2$ are not separated on $[T,2T]$, the bound of the mean value theorem for the remainder is not small, and no bound for $\mathscr E^{\mathrm{off}}_T$ is proved here. The formula \eqref{eq:taucw-inc1} splits the left side, up to $o(1)$, into its diagonal constant and $\mathscr E^{\mathrm{off}}_T$.
\end{remark}

\begin{lemma}[The Selberg weight as a difference of Fej\'er-type autocorrelations]\label{lem:taucw-fejer}
Put $W^{\mathrm{ev}}(v)=W(|v|)$ and, for $L>0$, $\mathbb T_L(y)=(L-|y|)_+$. Then
\begin{gather}
W^{\mathrm{ev}}(v)=(2-|v|)_+-(1-|v|)_+,\qquad L\,W^{\mathrm{ev}}\Bigl(\frac yL\Bigr)=\mathbb T_{2L}(y)-\mathbb T_L(y),\label{eq:taucw-fejer-W}\\
\mathbb T_L(y)=\int_{\mathbb R}\ind_{[-L/2,L/2]}(u)\,\ind_{[-L/2,L/2]}(u+y)\,du,\label{eq:taucw-fejer-auto}\\
\int_{\mathbb R}W^{\mathrm{ev}}(v)e^{-iuv}\,dv=\frac{2(\cos u-\cos2u)}{u^2},\qquad
\mathfrak K_L(\delta)=\frac1{2\pi}\int_{\mathbb R}W^{\mathrm{ev}}\Bigl(\frac yL\Bigr)e^{-i\delta y}\,dy .\label{eq:taucw-fejer-ft}
\end{gather}
In particular $\mathfrak K_L(\pi/L)=-2L/\pi^3<0$, so $W^{\mathrm{ev}}$ is not the autocorrelation $\int f(u)\overline{f(u+v)}\,du$ of any $f\in L^2(\mathbb R)$.
\end{lemma}
\begin{proof}
For $|v|\le1$ the right side of the first identity is $(2-|v|)-(1-|v|)=1$, for $1<|v|\le2$ it is $2-|v|$, and for $|v|>2$ it is $0$; these are the values of $W(|v|)$. The second identity is the first at $v=y/L$, multiplied by $L$. The autocorrelation of $\ind_{[-L/2,L/2]}$ is $(L-|y|)_+$, which is \eqref{eq:taucw-fejer-auto}. The Fourier transforms of $(2-|v|)_+$ and $(1-|v|)_+$ are $2(1-\cos2u)/u^2$ and $2(1-\cos u)/u^2$; their difference is the first formula of \eqref{eq:taucw-fejer-ft}. Substituting $v=y/L$ gives $\frac1{2\pi}\int W^{\mathrm{ev}}(y/L)e^{-i\delta y}dy=\frac{L}{2\pi}\cdot\frac{2(\cos L\delta-\cos2L\delta)}{L^2\delta^2}$, which is $\mathfrak K_L(\delta)$ of \eqref{eq:kernel-G}. At $\delta=\pi/L$, $\cos\pi-\cos2\pi=-2$. An autocorrelation of an $L^2$ function has nonnegative Fourier transform $|\hat f|^2$.
\end{proof}
The triangles $\mathbb T_L$ are the autocorrelations of indicators that bound the kernels $\phi^2\star\phi^2$ and $\phi\star\phi$ of \cite[\S5.1, (5.1)--(5.2)]{AlpogeFurman2026}, which lie between $(L-2w-|y|)_+$ and $(L-|y|)_+$, with $w=1$ as in \cite[\S5.1]{AlpogeFurman2026}. By \cref{lem:taucw-fejer}, $\mathfrak K_L$ is $(2\pi L)^{-1}$ times the Fourier transform of $\mathbb T_{2L}-\mathbb T_L$, a difference of two such triangles of widths $2L$ and $L$, and this difference is not positive definite; the negative values of $\mathfrak K_L$ are the reason why $C_{\mathfrak K}=\int|\mathfrak K_1|>1=\int\mathfrak K_1$ in \cref{lem:taucw-kernel}(a).

\begin{proposition}[Energy of the off-line part at moderate cutoffs]\label{prop:taucw-Yenergy}
Fix $0<\theta\le1$ and put $L(t)=\theta\log t$. For $T\ge T_0(\theta)$,
\begin{equation}\label{eq:taucw-Yenergy}
\Bigl|\,\bigl\|\Phi^{\mathrm w}_\theta\bigr\|_{L^2(T,2T)}-\bigl\|\mathcal Y_{L(t)}(t)\bigr\|_{L^2(T,2T)}\Bigr|\ll_\theta\sqrt{T\log\log T},
\end{equation}
where $\mathcal Y_{L(t)}(t)$ is defined outside the imaginary parts of the zeros. Consequently, for $0<\theta<\frac12$,
\begin{equation}\label{eq:taucw-Yenergy-mod}
\int_T^{2T}\bigl|\mathcal Y_{\theta\log t}(t)\bigr|^2\dd\ll_\theta T\log\log T .
\end{equation}
\end{proposition}
\begin{proof}
Outside the imaginary parts of the zeros, \cref{lem:taucw-pole}, \cref{prop:taucw-decomp} with $X=t^\theta$ and \eqref{eq:RVM} give $\Phi^{\mathrm w}_\theta=S+R+\mathcal O_{L(t)}+\mathcal Y_{L(t)}+O(1/t)$. By Selberg's second moment, $\int_{T/2}^{3T}S^2\dd\ll T\log\log T$ \cite{Selberg1946,Titchmarsh}, and $R=O(1/t)$. For $\mathcal O_{L(t)}$ use \eqref{eq:taucw-on}: its error term is $O(\log t/L(t))=O(1/\theta)$, the contribution of $N_0$ is at most $1/(\pi^2L(t))$ (as in the proof of \cref{thm:taucw-online}), and for the contribution of $S$ Minkowski's inequality and the bound $\sup_{t\in[T,2T]}|\mathfrak K_{L(t)}(\delta)|\le2L_\theta\min\{1,(L_\theta\delta)^{-2}\}$, $L_\theta=\theta\log T$, used there give
\begin{equation}
\begin{aligned}
\Bigl\|\int_{|\delta|\le t/2}\bigl(S(t-\delta)-S(t)\bigr)\mathfrak K_{L(t)}(\delta)\,d\delta\Bigr\|_{L^2(T,2T)}&\le\int_{\mathbb R}2L_\theta\min\{1,(L_\theta\delta)^{-2}\}\,d\delta\cdot2\|S\|_{L^2(T/2,3T)}\\
&=16\,\|S\|_{L^2(T/2,3T)} ,
\end{aligned}
\end{equation}
since $t-\delta\in[T/2,3T]$ for $T\le t\le2T$ and $|\delta|\le t/2$. Hence $\|S+R+\mathcal O_{L(t)}\|_{L^2(T,2T)}\ll_\theta\sqrt{T\log\log T}$, and the triangle inequality gives \eqref{eq:taucw-Yenergy}. For $\theta<\frac12$, \cref{thm:taucw-var} gives $\|\Phi^{\mathrm w}_\theta\|_{L^2(T,2T)}^2=\frac T{2\pi^2}\log\log T+O_\theta(T)$, and \eqref{eq:taucw-Yenergy-mod} follows.
\end{proof}
\begin{remark}[Reading of \cref{prop:taucw-Yenergy}]\label{rem:taucw-Yenergy}
\eqref{eq:taucw-Yenergy-mod} is an upper bound for the energy on $[T,2T]$ of all the zeros off the line at bandwidth $2\theta\log t<\log t$, proved from \cref{thm:taucw-var}, Selberg's second moment of $S$ and the explicit formula. For a fixed $L$ and finitely many pairs, \cref{prop:taucw-energy} gives the energy on all of $\mathbb R$, which dominates the energy on $[T,2T]$ at that $L$. For $\theta\ge\frac12$, \eqref{eq:taucw-Yenergy} says that a bound $\int_T^{2T}|\Phi^{\mathrm w}_\theta|^2\ll T\log\log T$ is equivalent to the same bound for $\mathcal Y_{\theta\log t}$; neither is proved here.
\end{remark}

\begin{proposition}[Zeros of bounded height]\label{prop:taucw-bh}
Let $\gamma^\diamond>0$ and let $\mathscr Z^\diamond$ be the finite set of zeros $\rho=\beta+i\gamma$ with $\beta\ne\frac12$ and $|\gamma|\le\gamma^\diamond$, counted with multiplicity. Let $0<\theta\le1$, $t\ge2\gamma^\diamond$ with $t^\theta\ge2$, and $L=\theta\log t$. The sum of the terms of $\mathcal Y_L(t)$ in \eqref{eq:taucw-on-off} over $\rho\in\mathscr Z^\diamond$ has modulus at most
\begin{equation}\label{eq:taucw-bh}
\#\mathscr Z^\diamond\cdot\frac{8\,t^{\theta-2}}{\pi\theta^2\log^2t}\le\#\mathscr Z^\diamond\cdot\frac{8}{\pi\theta^2\,t\log^2t}.
\end{equation}
\end{proposition}
\begin{proof}
For $\rho\in\mathscr Z^\diamond$, $|\gamma-t|\ge t/2$. By \cref{rem:taucw-off}, a zero $\frac12+y+i\gamma$ with $0<y<\frac12$ contributes at most $(X^{2y}-1)/(\pi L^2(\gamma-t)^2)\le4t^{2\theta y}/(\pi\theta^2t^2\log^2t)\le4t^{\theta-2}/(\pi\theta^2\log^2t)$, and a zero $\frac12-y+i\gamma$ at most $2/(\pi L^2(\gamma-t)^2)\le8t^{-2}/(\pi\theta^2\log^2t)\le8t^{\theta-2}/(\pi\theta^2\log^2t)$. The larger of the two bounds, summed over $\mathscr Z^\diamond$, is at most \eqref{eq:taucw-bh}; the second inequality uses $\theta\le1$.
\end{proof}
\begin{remark}[Two regimes]\label{rem:taucw-bh}
By \cref{prop:taucw-bh} and \cref{rem:taucw-detect}, on $[T,2T]$ with $X=t^\theta$, $\theta\le1$, the zeros off the line of bounded height change $F^{\mathrm w}_\theta-N$ only by $O_{\theta}(\#\mathscr Z^\diamond\,T^{-1}\log^{-2}T)$; any moment of $F^{\mathrm w}_\theta-N$ on dyadic intervals is insensitive to them. In the regime of \cref{thm:taucw-equiv}, $t$ ranges over $(0,\infty)$ while $X_j\to\infty$, and by \cref{cor:taucw-pair}(d) the profile of a single pair of zeros off the line has weighted mass that diverges with $L_j$, and for $\mathcal Y_{L_j}$ itself the divergence when a zero lies off the line follows from \cref{thm:taucw-equiv}. The detection of zeros off the line by the phase therefore takes place at cutoffs $X$ that are large compared with the height $t$.
\end{remark}

Let $\mathsf t_X=\bigl(X^2(1+\log\log3X)\bigr)^{1/3}$.

\begin{lemma}[Tail of the weighted norm]\label{lem:taucw-tail}
There is an absolute constant $C_{\mathrm{tail}}$ such that, for every $W\in\mathcal W$ and every $X\ge2$,
\begin{equation}\label{eq:taucw-tail}
\int_{\mathsf t_X}^\infty\frac{|P^W_X(t)|^2}{(1+t)^3}\dd\le C_{\mathrm{tail}} .
\end{equation}
\end{lemma}
\begin{proof}
Write $P^W_X=\Phi^W_X-\Pi^W_X$. The coefficients of the Dirichlet polynomial $\pi\Phi^W_X$ satisfy $|a_n|\le\Lambda(n)/(\sqrt n\log n)$ for $n\le X^2$ and vanish for $n>X^2$, so $\sum_n|a_n|^2\ll1+\log\log3X$ by Mertens' theorem and $\sum_nn|a_n|^2\le\sum_{n\le X^2}\Lambda(n)^2/\log^2n\ll X^2$. By the mean value theorem of Montgomery and Vaughan, $\int_U^{2U}|\Phi^W_X|^2\dd\ll U(1+\log\log3X)+X^2$ for every $U\ge1$. On $[2^m\mathsf t_X,2^{m+1}\mathsf t_X]$ one has $(1+t)^{-3}\le(2^m\mathsf t_X)^{-3}$, so
\begin{equation}
\begin{aligned}
\int_{\mathsf t_X}^\infty\frac{|\Phi^W_X|^2}{(1+t)^3}\dd&\ll\sum_{m\ge0}\frac{2^m\mathsf t_X(1+\log\log3X)+X^2}{(2^m\mathsf t_X)^3}\ll\frac{1+\log\log3X}{\mathsf t_X^2}+\frac{X^2}{\mathsf t_X^3}\\
&=\frac{(1+\log\log3X)^{1/3}}{X^{4/3}}+\frac1{1+\log\log3X}\ll1 .
\end{aligned}
\end{equation}
By \cref{lem:pole}, $|\Pi^W_X(t)|\ll(1+X)/t$ for $t\ge1$, uniformly in $W$, and $\mathsf t_X\ge1$, so $\int_{\mathsf t_X}^\infty|\Pi^W_X|^2(1+t)^{-3}\dd\ll X^2\mathsf t_X^{-4}\le X^{-2/3}\ll1$. Finally $|P^W_X|^2\le2|\Phi^W_X|^2+2|\Pi^W_X|^2$.
\end{proof}

\begin{corollary}[A finite range of heights]\label{cor:taucw-trunc}
For the Selberg weight and every unbounded sequence $X_j$, each of the conditions of \cref{thm:taucw-equiv} is equivalent to
\begin{equation}\label{eq:taucw-trunc}
\sup_j\int_0^{\mathsf t_{X_j}}\frac{|P^W_{X_j}(t)|^2}{(1+t)^3}\dd<\infty,
\end{equation}
and to the same statement with $\mathcal Y_{L_j}$ in place of $P^W_{X_j}$.
\end{corollary}
\begin{proof}
\Cref{lem:taucw-tail} and \cref{thm:taucw-equiv}(ii); for $\mathcal Y_{L_j}$ use in addition \eqref{eq:taucw-P-Y} and $\int_0^\infty\log^2(2+t)(1+t)^{-3}\dd<\infty$.
\end{proof}

\begin{lemma}[The weighted norm as a Gram form]\label{lem:taucw-gram}
For $W\in\mathcal W$ and $X\ge2$ let $\nu_X$ be the signed measure on $[2,X^2]$
\begin{equation}\label{eq:taucw-nu}
\int f\,\mathrm d\nu_X=\sum_{2\le n\le X^2}\frac{\Lambda(n)\,W(\frac{\log n}{\log X})}{\sqrt n\,\log n}f(n)-\int_2^{X^2}\frac{W(\frac{\log u}{\log X})}{\sqrt u\,\log u}f(u)\,\mathrm du\qquad(f\in C[2,X^2]),
\end{equation}
and, for real $z$ and $u,v\ge2$,
\begin{equation}\label{eq:taucw-gram-kernel}
\mathfrak C(z)=\frac12\int_0^\infty\frac{\cos(tz)}{(1+t)^3}\dd=\frac14\int_0^\infty\frac{r^3e^{-r}}{r^2+z^2}\,dr,\qquad
\mathbb J(u,v)=\mathfrak C\Bigl(\log\frac uv\Bigr)-\mathfrak C(\log uv).
\end{equation}
Then $P^W_X(t)=-\frac1\pi\int\sin(t\log u)\,\mathrm d\nu_X(u)$ for every $t$, and
\begin{equation}\label{eq:taucw-gram}
\int_0^\infty\frac{|P^W_X(t)|^2}{(1+t)^3}\dd=\frac1{\pi^2}\iint\mathbb J(u,v)\,\mathrm d\nu_X(u)\,\mathrm d\nu_X(v).
\end{equation}
The kernel $\mathbb J$ is positive on $[2,\infty)^2$ and positive semidefinite.
\end{lemma}
\begin{proof}
The representation of $P^W_X$ is \eqref{eq:PhiW}--\eqref{eq:PW}, because $\Im n^{-1/2-it}=-n^{-1/2}\sin(t\log n)$ and $\Im u^{-1/2-it}=-u^{-1/2}\sin(t\log u)$. The measure $\nu_X$ has finite total variation and $|\sin|\le1$, so Fubini's theorem gives \eqref{eq:taucw-gram} with $\mathbb J(u,v)=\int_0^\infty\sin(t\log u)\sin(t\log v)(1+t)^{-3}\dd$, and $\sin x\sin y=\frac12[\cos(x-y)-\cos(x+y)]$ gives the form \eqref{eq:taucw-gram-kernel}. The second expression for $\mathfrak C$ follows from $(1+t)^{-3}=\frac12\int_0^\infty r^2e^{-r(1+t)}dr$ and $\int_0^\infty\cos(tz)e^{-rt}\dd=r/(r^2+z^2)$. The function $\mathfrak C$ is even and strictly decreasing in $|z|$, and $|\log(u/v)|<\log uv$ for $u,v\ge2$, so $\mathbb J>0$ there. For finitely many points $u_i$ and real $c_i$, $\sum_{i,j}c_ic_j\mathbb J(u_i,u_j)=\int_0^\infty|\sum_ic_i\sin(t\log u_i)|^2(1+t)^{-3}\dd\ge0$.
\end{proof}

\begin{lemma}[Arithmetic form of the harmonics of the taucw phase]\label{lem:taucw-harmar}
For $0<\theta\le1$, $t\ge2$ with $t^\theta\ge2$, and complex $s$ put
\begin{equation}\label{eq:taucw-Euler}
\mathcal P_\theta(s;t)=\exp\Bigl(\sum_{n\le t^{2\theta}}\frac{\Lambda(n)\,\mathbf w_{\theta,t}(n)}{n^s\log n}\Bigr),
\end{equation}
an entire function of $s$ without zeros. Then, for every integer $\mathfrak h$,
\begin{equation}\label{eq:taucw-harmar}
e^{2\pi i\mathfrak hF^{\mathrm w}_\theta(t)}=e^{2i\mathfrak h\vartheta(t)}\,e^{-2\pi i\mathfrak hR(t)}\Bigl(\frac{\mathcal P_\theta(\frac12+it;t)}{\mathcal P_\theta(\frac12-it;t)}\Bigr)^{\mathfrak h},
\end{equation}
with $\vartheta$ and $R$ as in \eqref{eq:RVM}, and $|e^{-2\pi i\mathfrak hR(t)}-1|\le1.06\,|\mathfrak h|/t$ for $t\ge1$. Consequently, for real $\omega\ne0$, condition (HC) of \cref{prop:v9-harm} with $F^{\mathrm w}_\theta$ in place of $F_h$ holds if and only if, for every integer $\mathfrak h\ne0$,
\begin{equation}\label{eq:taucw-harmar-HC}
\frac1T\int_T^{2T}e^{i\omega t+2i\mathfrak h\vartheta(t)}\Bigl(\frac{\mathcal P_\theta(\frac12+it;t)}{\mathcal P_\theta(\frac12-it;t)}\Bigr)^{\mathfrak h}\dd\longrightarrow0 .
\end{equation}
\end{lemma}
\begin{proof}
Let $\ell^{\mathrm w}_\theta(t)$ be the finite sum in \eqref{eq:taucw-Euler} at $s=\frac12+it$. Its coefficients are real, so the sum at $s=\frac12-it$ is $\overline{\ell^{\mathrm w}_\theta(t)}$, and $\Im\ell^{\mathrm w}_\theta(t)=-\sum_n\Lambda(n)\mathbf w_{\theta,t}(n)n^{-1/2}(\log n)^{-1}\sin(t\log n)=\pi\Phi^{\mathrm w}_\theta(t)$ by \eqref{eq:taucw-phase}. Hence $e^{2\pi i\mathfrak h\Phi^{\mathrm w}_\theta}=e^{\mathfrak h(\ell^{\mathrm w}_\theta-\overline{\ell^{\mathrm w}_\theta})}=(\mathcal P_\theta(\frac12+it;t)/\mathcal P_\theta(\frac12-it;t))^{\mathfrak h}$. By \eqref{eq:RVM}, $N_0=\vartheta/\pi+1-R$, so $e^{2\pi i\mathfrak hN_0}=e^{2i\mathfrak h\vartheta}e^{-2\pi i\mathfrak hR}$. Multiplying gives \eqref{eq:taucw-harmar}. The bound for the factor with $R$ follows from $|e^{ix}-1|\le|x|$ and $|R(t)|\le0.53/(\pi t)$. The modulus of the remaining factors is $1$, so replacing $e^{-2\pi i\mathfrak hR}$ by $1$ changes the mean in (HC) by at most $1.06\,|\mathfrak h|/T$.
\end{proof}
\begin{remark}[The two factors]\label{rem:taucw-harmar}
In \eqref{eq:taucw-harmar} the factor $e^{2i\mathfrak h\vartheta(t)}$ comes from the Riemann--von Mangoldt term $N_0$ and the quotient of the two values of $\mathcal P_\theta$ from the primes. Where $\zeta(\frac12+it)\ne0$, the functional equation gives $\zeta(\frac12+it)/\zeta(\frac12-it)=e^{-2i\vartheta(t)}$; with $\zeta$ in place of $\mathcal P_\theta$ the product $e^{2i\mathfrak h\vartheta}(\zeta(\frac12+it)/\zeta(\frac12-it))^{\mathfrak h}$ is identically $1$. Condition \eqref{eq:taucw-harmar-HC} thus asks for cancellation in the mean, modulated by $e^{i\omega t}$, of the deviation of the finite weighted Euler product $\mathcal P_\theta$ from this relation. It is not proved here, for fixed $\omega$ or for $\omega$ growing with $T$.
\end{remark}

The following statement concerns the finite matrix of \cite[\S2.3]{AlpogeFurman2026}. Let $\mathsf L(T)=\log(T/2\pi)$ and let $\widetilde{\mathbf G}_2$ be the matrix $\widetilde G$ of \cite[(2.10)]{AlpogeFurman2026}, built with the indicator window $\psi_{\mathrm{ind}}=\ind_{[-1/2,1/2]}$ in \cite[(2.6)]{AlpogeFurman2026} and with the parameter $L$ of \cite[(2.7)]{AlpogeFurman2026} equal to $2\mathsf L(T)$ instead of $\mathsf L(T)$, so that, in the notation of \cite{AlpogeFurman2026}, $X=e^L=(T/2\pi)^2$, the sample points are $T+2\pi k/L$ and the dimension is $\lfloor LT/2\pi\rfloor$. Write $N(T,2T)$ for the number of zeros with $T<\gamma\le2T$, counted with multiplicity, and $N^{\mathrm{sc}}(T,2T)$ for the number of simple zeros with $\beta=\frac12$ and $T<\gamma\le2T$.

\begin{proposition}[The rank--trace inequality at twice the bandwidth]\label{prop:taucw-AF2}
The following statements hold for the matrix $\widetilde{\mathbf G}_2$ just defined.
\begin{enumerate}[label=(\alph*),leftmargin=*]
\item The archimedean term and the diagonal $n=m$ of the prime term of the double integral of \cite[Prop.~5.2]{AlpogeFurman2026} (the term $D$ of \cite[(5.12)]{AlpogeFurman2026}) add up, after division by $a^2L^2$ with the normalising constant $a$ of \cite[\S2.3]{AlpogeFurman2026}, to $(\frac76+o(1))N(T,2T)$; for a general parameter $L\asymp \mathsf L(T)$ this sum is
$(\mathsf L(T)/L+L/(3\mathsf L(T))+o(1))\,N(T,2T)$.
\item Suppose that, for some real $c_{\mathrm{HS}}$, $\|\widetilde{\mathbf G}_2\|_{\mathrm{HS}}^2\le(\frac76+c_{\mathrm{HS}}+o(1))N(T,2T)$ and $\operatorname{tr}\widetilde{\mathbf G}_2=N(T,2T)+o(N(T,2T))$. Then
\begin{equation}\label{eq:taucw-AF2}
\frac{N^{\mathrm{sc}}(T,2T)}{N(T,2T)}\ge\frac56-c_{\mathrm{HS}}-o(1).
\end{equation}
\item Under the hypotheses of (b), $\liminf_{T\to\infty}N^{\mathrm{sc}}(T,2T)/N(T,2T)\ge\frac56-c_{\mathrm{HS}}$, and the right side exceeds $C_{\mathrm{AF}}=\frac32-\frac1{\sqrt2}\cot\frac1{\sqrt2}=0.6725007\ldots$, the constant obtained in \cite{AlpogeFurman2026} with the Montgomery--Taylor window if and only if $c_{\mathrm{HS}}<\frac56-C_{\mathrm{AF}}=0.1608326\ldots$.
\end{enumerate}
\end{proposition}
\begin{proof}
(a) The archimedean term is evaluated in \cite[Prop.~5.3]{AlpogeFurman2026} and the diagonal of the prime term in the proof of \cite[Prop.~5.4]{AlpogeFurman2026}; neither computation uses $L\le \mathsf L(T)$. For $\psi_{\mathrm{ind}}$ the normalising constants $a,b$ of \cite[\S2.3, \S5.1]{AlpogeFurman2026} are $1+O(1/L)$, the archimedean term is $\frac{LT\mathsf L(T)^2}{2\pi}(1+o(1))$ and the diagonal is $\frac T\pi\sum_{n\le X}\frac{\Lambda(n)^2}n(L-\log n)+O(TL^2)=\frac{TL^3}{6\pi}(1+o(1))$ by \cite[(5.6)]{AlpogeFurman2026}. Dividing by $a^2L^2N(T,2T)=L^2\frac{T\mathsf L(T)}{2\pi}(1+o(1))$ gives $\frac{\mathsf L(T)}L+\frac L{3\mathsf L(T)}$, which is $\frac12+\frac23=\frac76$ at $L=2\mathsf L(T)$. (b) The proof of \cite[Thm.~A]{AlpogeFurman2026} in \cite[\S6]{AlpogeFurman2026} uses \cite[Prop.~4.1, Lemma~2.1, Lemma~3.1, Lemma~3.2]{AlpogeFurman2026} and the count of zeros in short intervals, and none of these depends on $L\le\mathsf L(T)$; it gives $N^{\mathrm{sc}}(T,2T)+o(N)\ge4\operatorname{tr}\widetilde{\mathbf G}_2-2N(T,2T)-\|\widetilde{\mathbf G}_2\|_{\mathrm{HS}}^2$, and the hypotheses give \eqref{eq:taucw-AF2}. (c) is arithmetic.
\end{proof}
\begin{remark}[What is open in \cref{prop:taucw-AF2}]\label{rem:taucw-AF2}
Neither hypothesis of (b) is proved. For $L=2\mathsf L(T)$ the bound $O(X^{1/2}L^2)=O(T\log^2T)$ for $\operatorname{tr}\widetilde{\mathbf G}_2-N(T,2T)$ obtained in the proof of \cite[Prop.~4.2]{AlpogeFurman2026} is larger than $N(T,2T)$; the error terms of \cite[Prop.~5.2]{AlpogeFurman2026}, which after normalisation are of order $X\log L$ in their proof, are $\gg N(T,2T)$; and the bound $O(L^2X)$ of \cite[Prop.~5.4]{AlpogeFurman2026} for the off-diagonal prime terms is $O(T^2\log^2T)$. As \cite[\S7.2(a)]{AlpogeFurman2026} states, for $X\gg T$ the off-diagonal prime sum is no longer dominated by the diagonal; \cite[Remark~6.1]{AlpogeFurman2026} treats $L=\lambda\mathsf L(T)$ only for $0<\lambda<1$. Thus (a) evaluates two terms of the double integral, not the main term of $\|\widetilde{\mathbf G}_2\|_{\mathrm{HS}}^2$. The constant $C_{\mathrm{AF}}$ is also obtained by Lamzouri \cite[Thm.~1.1, (1.1), (3.1), Lemma~3.2]{Lamzouri2026} as $2-C_{\mathrm{MT}}$, $C_{\mathrm{MT}}=\frac12+\frac1{\sqrt2}\cot\frac1{\sqrt2}=1.3274992963205\ldots$, with test functions $\eta_\varepsilon\in C_c^\infty((-\frac12,\frac12))$ in place of a matrix compression. The constant $c_{\mathrm{HS}}$ belongs to the matrix form of \cite{AlpogeFurman2026}, and $\mathscr E^{\mathrm{off}}_T$ of \eqref{eq:taucw-inc1-off} to the second moment of the increment of the taucw phase: they are different quadratic forms, and no transfer between them is proved here.
\end{remark}

\section{Open tasks after merging the research routes}\label{sec:v7-open}
The routes have the following precise mathematical status.
\begin{enumerate}[label=(\arabic*),leftmargin=*]
\item \textbf{Prime-band arithmetic persistence (open).}
 Prove $\limsup_{X\to\infty}\Delta_k(X;W)>0$ for each
 $k$ from the \emph{compensated} prime-power phase, not merely
 from the shrinking single-prime perturbations or a saturating
 auxiliary model. Under the unproved DBSRH hypothesis this
 would imply RH and simplicity.
\item \textbf{MCOM-critical local inequalities (open).}
 Seek a bound in which the actual constants
 the actual constants $A_{\mathrm M}=0.08382$ or $C_{\mathrm H}=0.0748596\ldots$ (the latter satisfies $C_{\mathrm H}<1/(4\pi)$) cross
 a threshold \emph{derived from the ordinal equations}. The
 proximity $2\pi C_{\mathrm H}<1/2$ is not itself such a theorem.
\item \textbf{Compensated excess (open).}
 Prove \eqref{eq:v7-excess-criterion} for an unbounded
 sequence of cutoffs, or a sufficiently strong uniform
 prime-band estimate implying it. By
 \cref{cor:v7-excess}, this would imply RH, without needing
 a separate Cauchy-in-measure statement. It is a one-way
 sufficient criterion, \emph{not} an established bound.
\item \textbf{Parity and local pair residual (open).}
 Exploit the exact Fourier-box identity
 \eqref{eq:v7-fourier-box} or bound the residual in the
 signed-pair model without assuming off-line zero control.
 The current identities alone do not yield the required
 positivity or vanishing.
\item \textbf{Endpoint and analytic inputs (open).}
 Verify the weighted zero-explicit formula including
 branch conventions and the exact IJNT/Fujii ranges (for the Selberg
 weight the formula and the branch of $\log\zeta$ are fixed in
 \cref{lem:taucw-expl}; the
 first-index normalization is settled, under the hypotheses of
 \cref{prop:v8-first}, by \eqref{eq:v8-flat}); no claims about
 $X\to\infty$ at fixed height follow from moment estimates
 established only for $X\le T^{1/4}$.
\item \textbf{Sawtooth coefficient and centroids (open).}
 Prove $\varsigma_T(\omega;F)\to0$ for $F=F_h,F^{\Cm}$ at the
 frequencies $\omega=r\log p$ and $\omega=1$ (via (HC), (NC) of
 \cref{prop:v9-harm}); this decides
 whether the centroids of $\gamma$ and of $\tau^{\mathrm C}$ differ at
 order $\ell/\log\ell$ (\cref{cor:v9-cen}); also the Gaussian law at the
 points $\gamma_k$ and the scale $\omega\asymp\log T$
 (\S\ref{sec:v9-open-centroid}); what is known about the four centroids is in \cref{tab:v11-centroids}.
\end{enumerate}


\section{Research priorities and exclusions}\label{sec:scope-v3}
The current research problem is the \emph{arithmetic structure of the
prime-corrected sublevel constructions}; a stand-alone equivalence with
RH via the Hadamard product, $\xi'/\xi$, a Littlewood integral or
another classical reformulation is \emph{not} proposed as a route (\cref{prop:v8-criterion} is a criterion on the sublevel points, not a stand-alone reformulation of RH).
Traditional analytic techniques remain admissible when they yield new
bounds on the prime phases, local occupation tests, their sublevel
measures, or the evaluation defect \eqref{eq:ek-v3}.

\begin{enumerate}[leftmargin=*,itemsep=4pt]
\item Establish the \emph{actual} weighted explicit formula for
$P_X^W$, retaining the compensation at $u=2$, with an error uniform
in the range needed for $X\to\infty$ at fixed height. For the Selberg weight this is \cref{lem:taucw-expl}, uniformly in $X\ge2$ with remainder $O(1/t)$; it is not established here for the other weights of \eqref{eq:weights}. In particular,
a hypothetical off-line zero cannot be treated as an isolated
maximal-$\beta$ zero unless that property is proved.
\item Prove a prime-only remainder estimate for the spatial-filter inverse in \cref{thm:identifiability-v4}, or derive noncollapse via the local occupation quantities in
\cref{prop:occupation-v3}, or obtain a stronger direct arithmetic
bound on $\Delta_k(X;W)$. The separation of the
$\gamma_k^{\mathrm{RvM}}$ is a structural input but not a substitute
for this estimate.
\item Derive an unconditional second-order bound for
$\sum e_k(F)^2$ on growing index windows for a specified prime
budget. The numerical values of $1-4\eta$ do not prove an improved
asymptotic proportion of zeros on the critical line.
\item Close the explicit $j=1$ budget at $\theta=1/2$ for the actual chosen family, including the $B_T$ term and all endpoint corrections; the reported $82.6\%$ is a numerical conditional estimate, not a theorem.
\item Treat moments of growing order separately from fixed-order
Selberg--Fujii estimates. In particular, numerical successes for
$2j\le14$ at finite heights give no uniform estimate as
$T\to\infty$.
\item Do not substitute the model $C(z,Lg)$ for an unconditional zero decomposition of the Euler phase. When differentiating in $L$, retain the moving-evaluation correction \eqref{eq:commutator-v4}, and use the untruncated phase for derivatives.
\item Study the \emph{consecutive difference} of RvM-relative arithmetic displacements in \eqref{eq:correlated-v5}; do not infer its strict upper bound from individual MCOM proximity or from window-averaged mean-square estimates. The coefficient $C_{\mathrm H}$ belongs to a redefined asymptotic truncation, not to the present fixed $H$.
\item Do not conflate the three different limits: $X\to\infty$
with $k$ fixed (DBSRH), $T\to\infty$ with $X=X(T)$ (mean-value
results), and $j\to\infty$ (strict cell criteria).
\end{enumerate}

\paragraph{Related cell-departure draft.}
The separate source \texttt{teorema\_salida\_celda\_v1.tex} provides
conditional, near-neighbor sufficient conditions for leaving a prescribed
index interval. Its proposed explicit formula is written under RH and
contains steps marked as sketches. It therefore cannot serve as an
unconditional bridge from the finite prime phases to RH. Its local
kernel analysis can nevertheless motivate a direct study of the
occupation sets in \cref{prop:occupation-v3}, with all errors
and quantifiers retained.

\paragraph{Disputed bridges that are not established.}
The implication RH $\Rightarrow$ DBSRH in the separate DBSRH draft
relies on a weighted explicit-formula argument that is not fully
audited here. The unbounded-$X$ uniform fourth moment of the full phase $N_0+P_X^W-N$ (\cref{thm:taucw-online} controls only its on-line part $\mathcal O_L$), the
isolated-maximum-free saturation assertion, and an $L^4$ Cauchy bound
for every $\sqrt T\le X<X'$ are all \emph{open}; no statement that
``only the diagonal term of Montgomery--Vaughan remains'' is justified.
The estimate in \cref{thm:moderate-log} uses the weighted support
$X^2$, so the IJNT prime-index cutoff and $X\asymp\sqrt T$ are not
interchangeable. A classical identity for the evaluation defects
is an exact tool for studying the proposed constructions, not an
independent established route to RH.

\section{Conclusions and theorem-status table}\label{sec:conclusion}
\begingroup\footnotesize
\setlength{\LTpre}{5pt}\setlength{\LTpost}{8pt}
\noindent Labels used in the table: \emph{Exact identity} (derived entirely in this note; the rows marked ``Exact'' or ``Unconditional identity'' are of this kind); \emph{Unconditional} (proved here, using no RH and only identifiable published results); \emph{Conditional on cited input} (depends on an external theorem or manuscript not re-proved here, although it assumes no RH; the input is named in parentheses); \emph{Conditional implication} (follows if an open hypothesis holds; the rows marked ``Proved implication/equivalence; hypothesis open'', ``Conditional'' or ``Conditional on'' an open hypothesis are of this kind); \emph{Numerical observation} (``Numerical only'' is of this kind); \emph{Exploratory/model result} (not a statement about the arithmetic phase; ``Exploratory only'' and ``Exact for the model'' are of this kind). The remaining labels are literal: \emph{Open}, \emph{Heuristic}, and \emph{Derived} (obtained here from cited coefficients).

\begin{longtable}{@{}p{0.34\linewidth}>{\raggedright\arraybackslash}p{0.23\linewidth}p{0.31\linewidth}@{}}
\toprule
Statement & Status & What the proof uses \\
\midrule
\endfirsthead
\toprule
Statement & Status & What the proof uses \\
\midrule
\endhead
MCOM-margin and localized crossing budget & Conditional on cited input (MCOM manuscript) & MCOM effective and asymptotic constants; sublevel comparison, no RH.\\
Quantitative ordinal transport $O(T/(\log T)^3)$ & Conditional on cited input (IJNT/Fujii input; MCOM manuscript) & MCOM interior margin plus cited fourth moment; does not improve Selberg--Fujii moments.\\
Exact finite Euler--pole transform & Unconditional identity & Fourier--Laplace inversion of actual weighted prime phase.\\
Weighted-$L^2$ normal-family criterion for RH & Proved implication; for the Selberg weight the hypothesis is equivalent to RH (\cref{thm:taucw-equiv}); hypothesis open & Montel, Euler--pole transform; no measure convergence assumption.\\
Fourier-box identity and parity test & Exact & Finite-interval integration by parts and endpoint terms; no simplicity claim.\\
Limit phase $P_\infty=S-r_0$ and first-index defect $t_\sharp>0$ & Conditional on weighted-$L^2$ bound, convergence in measure and the MCOM manuscript & Transforms of $N_\vartheta$ and $r_0$, Hadamard product, uniqueness via Cauchy integrals; MCOM bound for inactivity of the truncation.\\
Transform criterion for $\tau_k^\flat$ & Proved equivalence; hypothesis open & Normality from the MCOM barriers, identity theorem, pole comparison; RH together with convergence of all $\tau_k^\flat$.\\
Signed transverse two-source identification & Exploratory only & Analytic continuation of a model; the weighted zero-explicit formula is proved here only for the Selberg weight (\cref{lem:taucw-expl}).\\
Failure of eventual neighboring-zero containment for $h=\log\log t$ & Unconditional & Tsang's $\Omega_\pm$ theorem and an elementary counting argument.\\
Fourth moment and mean-square approximation for taucm & Conditional on cited input (IJNT inputs; MCOM bound) & MCOM bound, projection, IJNT fourth moment, gaps, and sublevel transfer.\\
Structural saturation model & Unconditional & Explicit continuous functions; no Euler phase.\\
Saturation incompatible with a uniform $O(T)$ fourth moment & Unconditional for the model & Weak equidistribution and the identity $\tfrac12(H-G)^4+\tfrac12(-H-G)^4$.\\
Failure for $X=T^{o(1)}$ & Conditional on cited input (Fujii input) & Selberg approximation and Montgomery--Vaughan.\\
Uniform bound for $T^\varepsilon\le X\le T^{1/4}$ & Conditional on cited input (Fujii input; MCOM bound) & Weighted prime-power decomposition and MCOM truncation.\\
Smoothed zero-sum bound for $\log X\ge\varepsilon\log T$ & Conditional on cited input (Fujii increment input) & Zero-count increments and Minkowski.\\
Improved range $X\le T^{1/4}(\log T)^{3/2}$ for two weights & Conditional on cited input (Fujii; no RH) & Endpoint decay of the weights and Montgomery--Vaughan.\\
Exact finite-cutoff margin and Lipschitz spacing & Unconditional & Triangle identity and local Lipschitz regularity; not a uniform lower bound in $X$.\\
Correlated RvM displacement and fixed-$H$ proximity & Exact / unconditional & RvM mean-value theorem, the existing MCOM barrier, and difference of consecutive shifts.\\
Anticorrelation of the increments of a prime band $T^a<p\le T^b$, $\tfrac14<a<b<\tfrac12$, at the RvM points (\cref{thm:v20-anticorr}) & Unconditional; band only, not the complete spacings & Kusmin--Landau on the RvM lattice, summation by parts, Mertens; $b<\tfrac12$ controls the frequencies $\log(pq)$ and the off-diagonal count.\\
$\operatorname{Corr}(s_k,s_{k+1})=-0.3835683$ for the complete points (\cref{rem:v21-two-facts}(i)) & Numerical observation (author's value, quoted, not recomputed here) & Author's computation; cutoff, weight and index window not recorded here.\\
$\mathfrak r^{\mathrm{diag}}_{W,\theta}$ of \eqref{eq:v21-diag} and $\langle(\Phi_X^{W\prime}/L)^2\rangle\approx2\int w^2v\,dv$ ($=\tfrac{11}6$ for $\theta=1$, Selberg weight) (\cref{rem:v21-two-facts}) & Heuristic (diagonal terms); not a theorem & Mertens and the diagonal of the proof of \cref{thm:v20-anticorr}; resonant pairs $nm\approx t/2\pi$, prime powers, compensator and linearization not controlled.\\
Noncollapse criterion under DBSRH & Conditional & Strict one-sided bound on the correlated shifts; requires an additional prime-only proof.\\
Uniform local derivative sufficient condition & Conditional & A uniform $W^{1,p}$ bound; generally incompatible with a discontinuous pointwise phase limit.\\
Two-source kernel: quadratic resolution, global error and full-profile injectivity & Exact for the model & Positive $W_{1/4}$ weight, Taylor remainder and Fourier uniqueness.\\
Spatial filter and cutoff derivative & Exact & Finite Euler sum, compensator, and full chain rule.\\
Quadratic pair inverse from finite samples & Conditional on a model residual & Linear independence, orthogonal projection, and explicit error estimate.\\
Separation of distant indices & Unconditional & $N_0\pm H$ barriers and sublevel measure.\\
Local arithmetic occupation criterion & Exact & Saturated parts of the weighted prime phase.\\
Evaluation defects $e_k$ and repeated heights & Exact & Consecutive index arithmetic and convexity.\\
2,001,052-zero computations & Numerical only & Repository text outputs; no independent interval certification.\\
Stieltjes, second-difference and slip identities for $\sum e^{i\omega\gamma_k}$ vs.\ $\sum e^{i\omega\mathsf x_k}$ & Exact identities (with the exceptional set $E_F$) & Integration by parts; counting function $\lfloor F^*+\frac12\rfloor$ of the sublevel points.\\
Landau cancellation $\frac{2\pi}T\int e^{i\omega t}D_F\to0$ & Derived; uses Landau coefficient & Landau--Gonek (RH) or Ford--Zaharescu ($O(T/\log T)$ as read from their eq.~(3.8), $o(T)$ safely; unconditional); prime Fourier coefficients of $\Phi^\circ$; for $F_h$, $F^{\Cm}$ also the fourth moment \eqref{eq:IJNTinput} (Step 3).\\
Reduction of $\Xi_T^{\mathrm C},\Xi_T^{\Cm}$ to $\varsigma_T$; same coefficient for $\tau^{\mathrm C},\tau^{\Cm}$ & Derived; conditional on cited input (IJNT fourth moment \eqref{eq:IJNTinput}, MCOM bound); endpoint term $O(\log T)$ by MCOM; rearrangement term controlled for $\tau^{\mathrm C}$ (band of width $O(\log\log T/\allowbreak\log T)$) and transferred to $\tau^{\Cm}$ & Level-counting inequality; fourth moments \eqref{eq:IJNTinput}; \cref{prop:v9-rearr}, \cref{prop:v9-CvsCm}.\\
Harmonic criterion: (HC)+(NC) $\Rightarrow\varsigma_T\to0$ & Proved implication; (HC), (NC) open & Fej\'er means of the sawtooth; non-concentration at the half-integers.\\
Exact reconstruction of $\mathsf x_k$ and spacings from $\mathrm{Cen}_\ell(1)$; collision-energy identity & Exact; needs MCOM pointwise bounds (all cases $<\pi$) & $|\mathsf x_k-\gamma_k^{\mathrm{RvM}}|\le2\pi\mathfrak c_{\mathsf x}+o(1)$; no claim on equal limits of centroids.\\
Coherence $\Theta^{\mathrm C,\Cm}_T(0,\omega_T)\to1$ for $|\omega_T|=o(\log T)$ & Conditional on cited input (IJNT/Fujii input) & Mean-square bound \eqref{eq:cm-mean} of \cref{thm:taucm}, \cref{prop:v9-CvsCm}.\\
Gaussian scale $\Theta_T^{\mathrm{RvM}}(0,\omega_T(\psi))\to e^{-\psi^2/2}$ & Conditional on a Gaussian law at the points $\gamma_k$ & Selberg CLT is for uniform $t$; Kargin's law at the zeros not checked against $\varpi_T$.\\
Centroid difference $\mathfrak b\ell/(2\log\ell)$; $\mathrm{Cen}_\ell(1;\gamma)=o(\ell/\log\ell)$ & Conditional on block coefficient; unconditional $o$ via Ford--Zaharescu & $O(\log\ell)$ only under RH.\\
Four centroids ($\gamma$, RvM, $\tau^{\mathrm C}$, $\tau^{\Cm}$): formulas, differences, convergence with and without RH (\cref{tab:v11-centroids}) & RvM centroid converges (unconditional, proved); $\gamma$: RH limit and Tauberian criterion quoted from the author's notes; $\tau^{\mathrm C},\tau^{\Cm}$: reduction to $\varsigma_\infty$, convergence open & Euler--Maclaurin and Kusmin--Landau; \cref{thm:v9-cen}; Cram\'er and Debruyne--Vindas (quoted, not re-proved); repository record quoted.\\
$\varsigma_T\to0$ (i.e.\ $\mathfrak b^{\mathrm C}=0$) & Open & Model numerics only; testable on the points through \cref{cor:v9-test}.\\
Conditional $j=1$ simple-critical-line fraction & Exact conditional inequality & $A_T$, $B_T$, fourth moment, derivative moment and boundary control.\\
Uniform estimate for all $X\ge T^\varepsilon$ for the full phase $N_0+P_X^W-N$ & Open & Requires additional arithmetic cancellation.\\
DBSRH$\Leftrightarrow$RH & Open & Requires a uniform explicit formula and non-cancellation, and careful handling of repeated imaginary parts.\\
Compensator at the cutoff $X=t^\theta$, $\theta\le1$: $\Pi^W_{t^\theta}(t)\ll1/t$ (\cref{lem:taucw-pole}) & Unconditional & Two integrations by parts; the Selberg weight vanishes linearly at $2\log X$.\\
Smoothed explicit formula and decomposition $N_0+P_X^W-N=\mathcal O_L+\mathcal Y_L+O(1/t)$ (\cref{lem:taucw-expl}, \cref{prop:taucw-decomp}) & Unconditional & Selberg weight, residues, integration in $\sigma$; the pole term equals $-\pi\Pi_X^W(t)$ plus a term $\mathscr A(t)$ independent of $X$.\\
Fourth moment of the on-line part $\mathcal O_{L(t)}$, $\theta\ge\theta_0>0$ (\cref{thm:taucw-online}) & Conditional on cited input (Fujii increment input) & Increment estimate \eqref{eq:FujiiB}, Selberg\textquoteright s fourth moment of $S$, Minkowski.\\
Existence, order and coarse localization of the taucw points; $N^{\mathrm w}_\theta=N_0+O_\theta(t^\theta/\log^2t)$; $|\tau_k^{\mathrm w}-\gamma_k^{\mathrm{RvM}}|\ll\gamma_k^{\mathrm{RvM}}/\log^3\gamma_k^{\mathrm{RvM}}$ (\cref{lem:taucw-amp}, \cref{prop:taucw-ord}, \cref{prop:taucw-coarse}) & Unconditional & Partial summation with $\pi(x)\ll x/\log x$; monotone barriers $N_0\pm\mathfrak A_\theta$.\\
Absolute transport $\sum_{k\in\mathcal U}|\tau_k^{\mathrm w}-\gamma_k|\le\int q_F$ and $L^1$ contraction of sublevel distributions (\cref{thm:taucw-L1}, \cref{thm:taucw-contr}) & Exact & Level counting; Fubini--Tonelli; no derivative of $F$ and no simple crossing.\\
Variance $\frac1{2\pi^2}\log\log T+O_\theta(1)$ for $\theta<\frac12$; exact variance of an increment in the cutoff (\cref{thm:taucw-var}, \cref{thm:taucw-inc}) & Unconditional & Montgomery--Vaughan mean value theorem; Mertens; frozen weights (\cref{lem:taucw-freeze}).\\
Signed transport up to $\theta=1$, $O(T^{\theta_+}/\log^2T)$; half-integer sampling identity (\cref{thm:taucw-signed}, \cref{prop:taucw-saw}) & Unconditional & Monotonicity of the moving weights in $t$; one integration by parts; prime number theorem.\\
Fourth moment of the prime-power part; mean of the long-prime part; Chebyshev representation (\cref{thm:taucw-pow}, \cref{lem:taucw-mean}, \cref{cor:taucw-cheb}) & Unconditional & Montgomery--Vaughan; Mertens; prime number theorem; two integrations by parts.\\
Fixed-frequency coefficient of the phase; spectral formula of the points, with error $O_\omega(\log^{-2}T)$ (\cref{thm:taucw-coef}, \cref{prop:taucw-spec}) & Unconditional & Integration by parts against the moving weights; Stieltjes integration by parts.\\
Quadratic estimate for the local reordering defect; low frequencies for $\theta<\frac12$ (\cref{thm:taucw-quad}, \cref{cor:taucw-lowfreq}) & Unconditional & First-moment bound for symmetric differences of sublevel sets; \cref{thm:taucw-var}.\\
Plancherel and band-limited identities; $S-S^{\mathrm w}=N-N^{\mathrm w}$; exact quadratic comparison of two phases (\cref{thm:taucw-plan}, \cref{thm:taucw-fejer}, \cref{prop:taucw-S}, \cref{prop:taucw-polar}) & Exact & Plancherel; matching of sorted lists; algebra.\\
Exact profile of an off-line pair, its spectrum, exact $L^2$ mass with lower bounds and localization; energy identity for finitely many pairs (\cref{prop:taucw-pair}, \cref{cor:taucw-pair}, \cref{prop:taucw-energy}) & Exact & Two-sided Laplace representations of $\mathfrak G_L$; differentiation in $y$; Plancherel for sine transforms.\\
$|P^W_X-\mathcal Y_L|\le C\log(2+t)$ uniformly in $X\ge2$; (no zero off the line) $\Leftrightarrow$ weighted $L^2$ bound of $P^W_{X_j}$ $\Leftrightarrow$ weighted $L^2$ bound of $\mathcal Y_{L_j}$ (\cref{thm:taucw-P}, \cref{thm:taucw-equiv}) & Unconditional; the bound itself is not proved here & Explicit formula of \cref{lem:taucw-expl}; \cref{lem:taucw-kernel}; \cref{thm:v7-normality}.\\
Shared imaginary parts force displacement of the strictly increasing taucw points; $L^1$ cost bounds (\cref{prop:taucw-shared}) & Unconditional & \cref{prop:taucw-ord}; \cref{prop:taucw-S}; triangle inequality; Chebyshev.\\
Mean value with the Selberg endpoint; second moment of the phase with support $n\le t(\log t)^{2\mathsf z}$, $0\le\mathsf z<\frac32$ (\cref{lem:taucw-mvsel}, \cref{prop:taucw-beyond}) & Unconditional & Montgomery--Vaughan; linear vanishing of $W$ at the end of its support; Minkowski's inequality.\\
Increment of the phase from $\theta=\frac12$ to $\theta=1$: diagonal constant $(5\log2-3)/(2\pi^2)$ plus the off-diagonal part (\cref{prop:taucw-inc1}) & Unconditional; the off-diagonal part is not bounded here & Mertens' theorem; Stieltjes integration by parts.\\
Selberg weight as a difference of two triangles; $\mathfrak K_L$ as its Fourier transform; not positive definite (\cref{lem:taucw-fejer}) & Exact & Fourier transforms of triangles.\\
Energy of $\mathcal Y_{\theta\log t}$ on $[T,2T]$ is $O(T\log\log T)$ for $\theta<\frac12$ (\cref{prop:taucw-Yenergy}) & Unconditional & \cref{thm:taucw-var}; Selberg's second moment of $S$; \eqref{eq:taucw-on}.\\
Zeros off the line of bounded height change the taucw error by $O(t^{\theta-2}\log^{-2}t)$ (\cref{prop:taucw-bh}) & Unconditional & \cref{rem:taucw-off}.\\
Tail of the weighted norm, truncation to $t\le\mathsf t_X$, Gram form (\cref{lem:taucw-tail}, \cref{cor:taucw-trunc}, \cref{lem:taucw-gram}) & Unconditional & Montgomery--Vaughan; dyadic decomposition; Fubini's theorem.\\
Harmonics of the phase as quotients of weighted Euler products; equivalent form of (HC) (\cref{lem:taucw-harmar}) & Exact identity; (HC) itself not proved & \eqref{eq:RVM}; \eqref{eq:taucw-phase}.\\
Rank--trace inequality at bandwidth $2\mathsf L(T)$: archimedean and diagonal terms of the double integral give $\frac76$; threshold $0.1608326\ldots$ for $c_{\mathrm{HS}}$ (\cref{prop:taucw-AF2}) & Implication from two matrix bounds that are not proved & \cite{AlpogeFurman2026}, Lemmas~2.1, 3.1, 3.2, Propositions~4.1, 5.2--5.4 and \S6; \cite{Lamzouri2026}.\\
Taucw points, centroids and exponential sums, $\theta=1$, $k\le3000$ (\cref{sec:taucw-data}) & Numerical observation & Points computed from primes only; zeros computed separately for the comparison.\\
\bottomrule
\end{longtable}
\endgroup
The RvM-relative identity \eqref{eq:correlated-v5} expresses adjacent spacings through correlated prime displacements; \cref{prop:rvmbound-v5} shows that the actual fixed truncation gives the coefficient $2\pi\cdot0.08382$ rather than $2\pi C_{\mathrm H}$ for individual RvM deviations. Neither control establishes noncollapse. The author's MCOM estimate removes the truncation-scale defect of the original construction, and the Selberg--Fujii--Montgomery--Vaughan argument supplies a genuine, nontrivial uniform range of prime cutoffs. The structural objection is valid only at the level of the properties it explicitly assumes; it cannot replace analysis of the actual Euler phase. The local quadratic kernel theorem (exact for the model) does not yet bound the residual from the actual Euler phase. The interval $X\ge T^\varepsilon$ without an upper restriction is the appropriate remaining \emph{problem}, not an established theorem. The results here do not prove RH; neither the correlated-shift criterion nor the local inverse bounds have been established for every cutoff and every index.

\section*{Source note}
The IJNT manuscript and the MCOM submission are used as cited inputs and are not modified here. The publicly accessible route notes and the saved ordinal-transport text from commit \texttt{38aac7c} are distinguished from the unpublished manuscripts. The normal-family criterion is a mathematical derivation of this note, not a result represented as already proved by the external notes. The numerical tables taken from the repository were cross-checked against the repository text outputs listed in \cite{CasanaV9,CasanaNotes}; the computations labelled ``this note'' in \S\ref{sec:v9-numbers} and those of \cref{sec:taucw-data} were made for this note. In the section on exponential sums and centroids, $\tau^{\mathrm C}:=u_k(F_h)$ (loglog clip) throughout, as in the current IJNT manuscript (the file with the unclipped definition is discussed in \cref{rem:v9-clip}). That section re-derives the statements of the ChatGPT texts on the joint transform, spectral coherence, second differences, the rearrangement band, the reconstruction at frequency~$1$ and the harmonic criterion from the objects of this note; the Landau and sawtooth coefficients are taken from the author's Fourier notebook, and the cited Landau, Gonek and Ford--Zaharescu statements are quoted, not re-proved. In \S\ref{sec:v11-four} the convergence of the RvM centroid is proved; the Riemann-hypothesis limit of the centroid of $\gamma$, the Tauberian criterion and the window equivalence are quoted from the author's notes on the zeta--RvM paper \cite{CasanaV26}, and $K^0$ and the non-convergence record from the repository README; no numerical computation was made for that subsection. \Cref{sec:taucw} is self-contained: its points, centroids and exponential sums were computed for this note from the primes alone, the zeros being computed separately for the comparison, and the Python code and output files are distributed with this note. In \cref{sec:taucw-beyond} the matrix of Alp\"oge and Furman \cite{AlpogeFurman2026} enters only through the statements cited there.

\begin{thebibliography}{99}

\bibitem{CasanaIJNT}
V.~Casa\~na-Giner, \emph{A prime-corrected Riemann--von Mangoldt approximation of the zeros of the Riemann zeta function}, IJNT manuscript, 2026, source \texttt{casanya\_ordinal\_primes\_ijnt\_v5\_final.tex}; especially the sublevel-transfer theorem, Fujii--Selberg input and fourth moment of the untruncated phase.

\bibitem{CasanaMCOM}
V.~Casa\~na-Giner, \emph{Explicit and asymptotic estimates for the zero-counting function of the Riemann zeta-function}, Mathematics of Computation submission, corrected manuscript (2 November 2026), source \texttt{casanya\_turing\_jensen\_mcom\_CORREGIDO.tex}; explicit $S(t)$ bounds.

\bibitem{CasanaDBSRH}
V.~Casa\~na-Giner, \emph{Deep Blue Sea Riemann Hypothesis: convergence of prime-corrected sublevel points and the Riemann hypothesis}, research draft \texttt{teorema\_DBSRH\_v1.tex}, 2 November 2026; proposed weighted explicit formula and partial converse, with analytic steps still indicated as sketches.

\bibitem{Titchmarsh}
E.~C.~Titchmarsh, \emph{The Theory of the Riemann Zeta-Function}, 2nd edition, revised by D.~R.~Heath-Brown, Oxford University Press (1986).

\bibitem{TitchmarshFI}
E.~C.~Titchmarsh, \emph{Introduction to the Theory of Fourier Integrals}, 2nd edition, Oxford University Press (1948); Theorem~93 (boundary values of the Cauchy integral of an $L^2$ function).

\bibitem{WhittakerWatson}
E.~T.~Whittaker and G.~N.~Watson, \emph{A Course of Modern Analysis}, 4th edition, Cambridge University Press (1927), \S12.31 (Binet's and Malmst\'en's integrals for $\log\Gamma$), \href{https://doi.org/10.1017/CBO9780511608759}{doi:10.1017/CBO9780511608759}.

\bibitem{Tsang1986}
K.-M.~Tsang, \emph{Some $\Omega$-theorems for the Riemann zeta-function}, \emph{Acta Arith.} \textbf{46} (1986), 369--395, \href{https://doi.org/10.4064/aa-46-4-369-395}{doi:10.4064/aa-46-4-369-395}.

\bibitem{Dobner2024}
A.~Dobner, \emph{Large deviations of the argument of the Riemann zeta function}, \emph{Mathematika} \textbf{70} (2024), e12251, \href{https://doi.org/10.1112/mtk.12251}{doi:10.1112/mtk.12251}; arXiv:2101.01747, \href{https://doi.org/10.48550/arXiv.2101.01747}{doi:10.48550/arXiv.2101.01747}.

\bibitem{Selberg1946}
A.~Selberg, \emph{Contributions to the theory of the Riemann zeta-function}, \emph{Arch. Math. Naturvid.} \textbf{48} (1946), no.~5, 89--155.

\bibitem{Fujii1974}
A.~Fujii, \emph{On the zeros of Dirichlet $L$-functions. I}, \emph{Trans. Amer. Math. Soc.} \textbf{196} (1974), 225--235, \href{https://doi.org/10.1090/S0002-9947-1974-0349603-2}{doi:10.1090/S0002-9947-1974-0349603-2}.

\bibitem{Fujii1981}
A.~Fujii, \emph{On the zeros of Dirichlet $L$-functions. II (with corrections to ``On the zeros of Dirichlet $L$-functions. I'' and the subsequent papers)}, \emph{Trans. Amer. Math. Soc.} \textbf{267} (1981), 33--40, \href{https://doi.org/10.1090/S0002-9947-1981-0621970-5}{doi:10.1090/S0002-9947-1981-0621970-5}.

\bibitem{MontgomeryVaughan}
H.~L.~Montgomery and R.~C.~Vaughan, \emph{Hilbert's inequality}, \emph{J. London Math. Soc.} (2) \textbf{8} (1974), 73--82, \href{https://doi.org/10.1112/jlms/s2-8.1.73}{doi:10.1112/jlms/s2-8.1.73}.

\bibitem{Mertens1874}
F.~Mertens, \emph{Ein Beitrag zur analytischen Zahlentheorie}, \emph{J. Reine Angew. Math.} \textbf{78} (1874), 46--62, \href{https://doi.org/10.1515/crll.1874.78.46}{doi:10.1515/crll.1874.78.46}.

\bibitem{Gonek2012}
S.~M.~Gonek, \emph{Finite Euler products and the Riemann hypothesis},
\emph{Trans. Amer. Math. Soc.} \textbf{364} (2012), 2157--2191,
\href{https://doi.org/10.1090/S0002-9947-2011-05546-7}{doi:10.1090/S0002-9947-2011-05546-7}.
The use of Selberg smoothing and Euler products is not claimed as new;
the sublevel-measure constructions and the associated transfer results
are mathematically distinct from the zeros of his functions.

\bibitem{CasanaV9}
V.~Casa\~na-Giner, \emph{Evaluation defects at 2,001,052 zeros},
text outputs and code in
\href{https://github.com/vcasana1/prueba/tree/kc-tauC-2026-11-02/KC_tauC_2026-11-02/v8_identidad_en_ordenadas}{repository directory \texttt{v8\_identidad\_en\_ordenadas}},
2 November 2026; \texttt{draft\_analysis.txt},
\texttt{w14\_th0.5\_analysis.txt}, \texttt{w14\_th0.7\_analysis.txt},
\texttt{gap\_margin\_v9.txt}, \texttt{decomp\_w14\_th0.5.txt},
\texttt{budget\_w14\_th0.5\_j19.txt}, \texttt{mv\_dphi\_th0.5\_j19.txt},
\texttt{tight\_pairs\_v13.txt}, \texttt{margin\_triangle\_v14.txt};
repository commit \texttt{2b24f67dc5348788714bafaf4a4c566aa041a877}.

\bibitem{CasanaNotes}
V.~Casa\~na-Giner, \emph{Notes on the $\tau_k^{\mathrm C}$ research, 1--2 November 2026}, \href{https://github.com/vcasana1/prueba/blob/kc-tauC-2026-11-02/KC_tauC_2026-11-02/NOTAS_investigacion.md}{GitHub branch \texttt{kc-tauC-2026-11-02}}, numerical table and qualifications.

\bibitem{CasanaV10Notes}
V.~Casa\~na-Giner, \emph{Rutas hacia RH desde taucm/taucms --- notas de sesi\'on},
public research notes, 2 November 2026,
\href{https://github.com/vcasana1/prueba/blob/38aac7c/KC_tauC_2026-11-02/v10_rutas_RH/NOTAS_rutas_RH_v1.md}{\texttt{NOTAS\_rutas\_RH\_v1.md}},
commit \texttt{38aac7c}; exploratory statements are not treated as externally verified theorems.

\bibitem{CasanaV10Transport}
V.~Casa\~na-Giner, \emph{Transporte ordinal: texto de la sesi\'on},
2 November 2026,
\href{https://github.com/vcasana1/prueba/blob/38aac7c/KC_tauC_2026-11-02/v10_rutas_RH/transporte_ordinal_texto_externo.txt}{\texttt{transporte\_ordinal\_texto\_externo.txt}},
commit \texttt{38aac7c}. The manuscript above specifies exactly which deductions depend on the IJNT input.

\bibitem{Landau1912}
E.~Landau, \emph{\"Uber die Nullstellen der Zetafunktion}, \emph{Math. Ann.} \textbf{71} (1912), 548--564, \href{https://doi.org/10.1007/BF01456808}{doi:10.1007/BF01456808}.

\bibitem{Gonek1993}
S.~M.~Gonek, \emph{An explicit formula of Landau and its applications to the theory of the zeta-function}, \emph{Contemp. Math.} \textbf{143} (1993), 395--413, \href{https://doi.org/10.1090/conm/143/1210528}{doi:10.1090/conm/143/1210528}.

\bibitem{FordZaharescu2005}
K.~Ford and A.~Zaharescu, \emph{On the distribution of imaginary parts of zeros of the Riemann zeta function}, \emph{J. Reine Angew. Math.} \textbf{579} (2005), 145--158, \href{https://doi.org/10.1515/crll.2005.2005.579.145}{doi:10.1515/crll.2005.2005.579.145}; arXiv:math/0405459, \href{https://doi.org/10.48550/arXiv.math/0405459}{doi:10.48550/arXiv.math/0405459}.

\bibitem{Kargin2013}
V.~Kargin, \emph{On fluctuations of Riemann's zeta zeros}, \emph{Probab. Theory Related Fields} \textbf{157} (2013), 575--604, arXiv:1203.5309, \href{https://doi.org/10.1007/s00440-012-0465-9}{doi:10.1007/s00440-012-0465-9}.

\bibitem{GrahamKolesnik1991}
S.~W.~Graham and G.~Kolesnik, \emph{Van der Corput's Method of Exponential Sums}, London Math. Soc. Lecture Note Ser. \textbf{126}, Cambridge University Press (1991), \href{https://doi.org/10.1017/CBO9780511661976}{doi:10.1017/CBO9780511661976}.


\bibitem{Cramer1919}
H.~Cram\'er, \emph{Studien \"uber die Nullstellen der Riemannschen Zetafunktion}, \emph{Math. Z.} \textbf{4} (1919), 104--130, \href{https://doi.org/10.1007/BF01203390}{doi:10.1007/BF01203390} (volume, year and pages as listed at EuDML, \url{https://eudml.org/doc/167519}). Attribution to this paper of the meromorphic continuation used in \cref{thm:v17-cont} is not confirmed in any source consulted.

\bibitem{DebruyneVindas2019}
G.~Debruyne and J.~Vindas, \emph{Complex Tauberian theorems for Laplace transforms with local pseudofunction boundary behavior}, \emph{J. Anal. Math.} \textbf{138} (2019), 799--833, \href{https://doi.org/10.1007/s11854-019-0045-1}{doi:10.1007/s11854-019-0045-1} (reference and DOI recalled, not verified here).

\bibitem{Hormander1990}
L.~H\"ormander, \emph{The Analysis of Linear Partial Differential Operators I}, 2nd ed., Springer (1990), reprinted in Classics in Mathematics, Springer (2003), \href{https://doi.org/10.1007/978-3-642-61497-2}{doi:10.1007/978-3-642-61497-2}, pp.~63--66 (interval formulation) and 343--344 (distributional formulation); no theorem number found, book not opened.

\bibitem{Rudin1971}
W.~Rudin, \emph{Lectures on the Edge-of-the-Wedge Theorem}, CBMS Regional Conf.\ Ser.\ Math.\ \textbf{6}, Amer.\ Math.\ Soc.\ (1971), listed as a reference, not opened.

\bibitem{BOT1958}
H.~J.~Bremermann, R.~Oehme and J.~G.~Taylor, \emph{Proof of dispersion relations in quantized field theories}, Phys.\ Rev.\ \textbf{109} (1958), 2178--2190, \href{https://doi.org/10.1103/PhysRev.109.2178}{doi:10.1103/PhysRev.109.2178}; listed as a reference, not opened.

\bibitem{CasanaV26}
V.~Casa\~na-Giner, \emph{The midpoint Riemann--von Mangoldt synthetic zeta system}, manuscript \texttt{rvm\_zeta\_pi\_mertens} (version~26), the author's own unpublished work; Section~5: Prop.~5.3, Thm.~5.5, Lemma~5.4, Cor.~5.8, Rem.~5.9, Prop.~5.10, Thm.~5.13, Lemma~5.14, Thm.~5.15.

\bibitem{Davenport2000}
H.~Davenport, \emph{Multiplicative Number Theory}, 3rd ed., Graduate Texts in Mathematics \textbf{74}, Springer (2000), Ch.~17, \href{https://doi.org/10.1007/978-1-4757-5927-3}{doi:10.1007/978-1-4757-5927-3} (DOI recalled, not verified here).

\bibitem{Tenenbaum2015}
G.~Tenenbaum, \emph{Introduction to Analytic and Probabilistic Number Theory}, 3rd ed., Graduate Studies in Mathematics \textbf{163}, Amer.\ Math.\ Soc.\ (2015), Thm.~I.1.10, p.~18, \href{https://doi.org/10.1090/gsm/163}{doi:10.1090/gsm/163} (DOI recalled, not verified here).

\bibitem{AlpogeFurman2026}
L.~Alp\"oge and R.~Furman, \emph{More than two thirds of the zeros of the Riemann zeta function are simple and on the critical line}, preprint, arXiv:2608.13637v2 (19 August 2026), \href{https://doi.org/10.48550/arXiv.2608.13637}{doi:10.48550/arXiv.2608.13637}.

\bibitem{Lamzouri2026}
Y.~Lamzouri, \emph{A new proof that more than $2/3$ of the zeros of the Riemann zeta function are simple and on the critical line}, preprint, arXiv:2609.02882v2 (8 September 2026), \href{https://doi.org/10.48550/arXiv.2609.02882}{doi:10.48550/arXiv.2609.02882}.

\end{thebibliography}
\end{document}
